\documentclass[UTF8, 12pt, fontset=fandol, oneside]{ctexart}
\usepackage{researchnotes}

\graphicspath{{math/数分习题/}}

\title{第六章\quad 多元函数微分学}
\author{}
\date{}

\begin{document}

\maketitle

\section*{§6.1\quad 欧氏空间・多元函数的极限与连续}

\noindent\textbf{导读}\quad 本节第二段是重点，适合本书的各类读者；第一、三段主要针对数学院系的学生.

\subsection*{*一、$m$ 维欧氏空间}

\noindent\textbf{导读}\quad 该段“$m$ 维欧氏空间”，理论性相对较强，主要针对数学院系学生，其他学生可以略，相关考题较少，一般未作重点.

\noindent\textbf{a.\quad 利用模的定义}

\noindent\textbf{例 6.1.1}\quad 设 $x=(x_1,x_2,\cdots,x_m)\in\mathbb{R}^m$，试证模
\[
  |x|=\left(\sum_{i=1}^m x_i^2\right)^{1/2}
\]
有关系：
\[
  1)\quad \frac{\sqrt m}{m}\sum_{i=1}^m |x_i|\le |x|\le \sum_{i=1}^m |x_i|;
\]
\[
  2)\quad \max_{1\le i\le m}|x_i|\le |x|\le
  \sqrt m\max_{1\le i\le m}|x_i|.
\]

\noindent\textbf{证}\quad 1）因为
\[
  \left(\sum_{i=1}^m |x_i|\right)^2
  =\sum_{i=1}^m x_i^2+
  2\sum_{\substack{i=1\\ m\ge j>i}}^{m-1}|x_i||x_j|
  \ge |x|^2,
\]
所以
\[
  |x|=\left(\sum_{i=1}^m |x_i|^2\right)^{1/2}
  \le \sum_{i=1}^m |x_i|.
\]
利用 Cauchy 不等式，
\[
  \sum_{i=1}^m |x_i|
  =\sum_{i=1}^m 1\cdot |x_i|
  \le \left(\sum_{i=1}^m 1^2\cdot\sum_{i=1}^m |x_i|^2\right)^{1/2}
  \le \sqrt m\left(\sum_{i=1}^m x_i^2\right)^{1/2}
  =\sqrt m\,|x|,
\]
故有
\[
  \frac{\sqrt m}{m}\sum_{i=1}^m |x_i|\le |x|.
\]

2）留给读者证明.

\noindent\textbf{b.\quad 利用距离的定义和性质}

\noindent\textbf{例 6.1.2}\quad 设 $E\subseteq\mathbb{R}^m,x,y\in E$，试证 $x$ 到 $E$ 的距离
\[
  \rho(x,E)\le \rho(x,y)+\rho(y,E).
\]

\noindent\textbf{证}\quad 根据距离性质，$\forall x,y,z\in\mathbb{R}^m$，若 $z\in E$，则由 $\rho(x,z)\le \rho(x,y)+\rho(y,z)$ 知
\[
  \rho(x,E)=\inf_{z\in E}\rho(x,z)\le \rho(x,y)+\rho(y,z).
\]
将 $x,y$ 固定，上式对一切 $z\in E$ 成立，因而有
\[
  \rho(x,E)\le \rho(x,y)+\inf_{z\in E}\rho(y,z)
  =\rho(x,y)+\rho(y,E).
\]

\noindent\textbf{例 6.1.3}\quad 设 $A,B\subseteq\mathbb{R}^m,x\in\mathbb{R}^m$，试证：
\[
  \rho(A,B)\le \rho(x,A)+\rho(x,B).
\]

\noindent\textbf{提示}\quad 对 $x\in\mathbb{R}^m,y\in A,z\in B$，应用距离的三角不等式 $\rho(y,z)\le \rho(x,y)+\rho(x,z)$.

\noindent\textbf{c.\quad 利用开集、闭集的定义}

\noindent\textbf{例 6.1.4}\quad 若 $E\subseteq\mathbb{R}^m$ 为闭集，试证 $x\in E$ 的充要条件是 $\rho(x,E)=0$.

\noindent\textbf{证}\quad 必要性明显，只证明充分性：若 $\rho(x,E)=0$，即 $\inf\rho(x,y)=0$，由确界定义知：$\exists y_n\in E$，使得 $|y_n-x|\to0$（当 $n\to\infty$ 时）. 如此，$x$ 为 $E$ 的一个聚点，故 $x\in\overline{E}=E$.

证毕.

\noindent\textbf{**例 6.1.5}\quad 设 $E\subseteq\mathbb{R}^m,r>0$ 为常数，
\[
  A=\{x\mid x\in\mathbb{R}^m,\rho(x,E)<r\},\quad
  B=\{x\mid x\in\mathbb{R}^m,\rho(x,E)\le r\}.
\]
试证：1）$A$ 为开集；2）$B$ 为闭集.

\noindent\textbf{证}\quad 利用例 6.1.2 的结果，易知 $\forall x,y\in\mathbb{R}^m$，有
\[
  |\rho(x,E)-\rho(y,E)|\le \rho(x,y).
\]
从而可知 $f(x)=\rho(x,E)$ 是 $x$ 的连续函数（更确切地说是一致连续函数，见例 6.1.29）. 由此易知，$A$ 为开集，$B$ 为闭集（见例 6.1.25）.

\noindent\textbf{d.\quad 利用边界的定义与聚点性质}

\noindent\textbf{*例 6.1.6}\quad 设 $E\subseteq\mathbb{R}^m$，试证 $E$ 的边界 $\partial E$ 为闭集.

\noindent\textbf{证}\quad 设 $x_0$ 为 $\partial E$ 的任一聚点，我们要证明 $x_0\in\partial E$. 为此我们要证明 $x_0$ 的任一 $\delta$ 邻域
\[
  U(x_0,\delta)=\{x\mid |x-x_0|<\delta\}
\]
里既含有 $E$ 的点，也含有不是 $E$ 的点. 根据聚点定义，至少 $\exists x_1\in U_0(x_0,\delta)\cap\partial E$.（这里 $U_0(x_0,\delta)$ 表示 $x_0$ 的空心邻域
\[
  U_0(x_0,\delta)=\{x\mid 0<|x-x_0|<\delta\}.
\]
）

因为 $x_1\in U_0(x_0,\delta)$，只要取
\[
  \delta_1=\min\{|x_0-x_1|,\delta-|x_0-x_1|\},
\]
则 $x_1$ 的邻域 $U(x_1,\delta_1)\subseteq U(x_0,\delta)$.

又因 $x_1\in\partial E$，所以 $U(x_1,\delta_1)$ 中既含有 $E$ 中的点，又含有不是 $E$ 的点. 由于 $U(x_1,\delta_1)\subseteq U(x_0,\delta)$，所以 $U(x_0,\delta)$ 中既含有 $E$ 中的点，也含有不是 $E$ 中的点. 由于 $\delta>0$ 的任意性，这就证明了 $x_0$ 是 $E$ 的边界点，故 $x_0\in\partial E$. 证毕.

\noindent\textbf{*例 6.1.7}\quad 设 $E\subseteq\mathbb{R}^m$，试证 $\partial\overline{E}\subseteq\partial E$.

\noindent\textbf{证}\quad 设 $x_0\in\partial\overline{E}$，我们要证明 $x_0\in\partial E$，即要证明 $x_0$ 为 $E$ 的边界点，亦即要证明 $\forall\delta>0$，在 $U(x_0,\delta)$ 中既有 $E$ 的点也有不是 $E$ 的点. 因 $x_0\in\partial\overline{E}$，所以 $U(x_0,\delta)$ 中既有 $\overline{E}$ 中的点，又含有不是 $\overline{E}$ 中的点. 设
\[
  x_1\in U(x_0,\delta)\cap\overline{E},\quad
  x_2\in U(x_0,\delta)\setminus\overline{E}.
\]

若 $x_1\in E$，则说明 $U(x_0,\delta)$ 中有 $E$ 的点；若 $x_1\in\overline{E}\setminus E$，则 $x_1$ 为 $E$ 的聚点. 仿上例证法，可知存在邻域 $U(x_1,\delta_1)\subseteq U(x_0,\delta)$，因 $x_1$ 为 $E$ 的聚点，故 $U(x_1,\delta_1)$ 中有 $E$ 中的点，从而知 $U(x_0,\delta)$ 中有 $E$ 的点.

由 $x_2\in U(x_0,\delta)\setminus\overline{E}$ 知，$x_2$ 为 $E$ 的外点，所以存在充分小的邻域在 $E$ 之外，由此知 $U(x_0,\delta)$ 内有不是 $E$ 的点. 证毕.

\noindent\textbf{*例 6.1.8}\quad 设 $F\subseteq E$，其中 $E\subseteq\mathbb{R}^m$ 为有界开区域，$F$ 为闭区域，试证：存在开区域 $V$，使得
\[
  F\subseteq V\subseteq\overline{V}\subseteq E.
\]

\noindent\textbf{证}\quad $\forall x_\lambda\in F\subseteq E$，因 $E$ 为开区域，$\exists r_\lambda>0$（充分小），使得
\[
  \overline{U}(x_\lambda,r_\lambda)
  =\{x\mid x\in\mathbb{R}^m,\ |x-x_\lambda|\le r_\lambda\}\subseteq E.
\]
如此，$\{U(x_\lambda,r_\lambda)\mid x_\lambda\in F\}$ 组成有界闭区域 $F$ 的一个开覆盖. 根据有限覆盖定理，存在有限子覆盖，记作
\[
  \{U(x_i,r_i)\mid i=1,2,\cdots,n\}.
\]
令
\[
  V=\bigcup_{i=1}^n U(x_i,r_i),
\]
则 $F\subseteq V\subseteq\overline{V}\subseteq E$.

\noindent\textbf{※例 6.1.9}\quad 设 $x_0,x_1,\cdots,x_n\in\mathbb{R}^n$，证明 $x_0,x_1,\cdots,x_n$ 在同一超平面上的充要条件是行列式
\[
  \Delta_0=
  \begin{vmatrix}
  x_{01}&x_{02}&\cdots&x_{0n}&1\\
  x_{11}&x_{12}&\cdots&x_{1n}&1\\
  \vdots&\vdots&&\vdots&\vdots\\
  x_{n1}&x_{n2}&\cdots&x_{nn}&1
  \end{vmatrix}=0,
\]
其中 $x_i=(x_{i1},x_{i2},\cdots,x_{in}),i=0,1,2,\cdots,n$.

\noindent\textbf{证}\quad 充分性. 将行列式
\[
  \Delta=
  \begin{vmatrix}
  x_{\lambda1}&x_{\lambda2}&\cdots&x_{\lambda n}&1\\
  x_{11}&x_{12}&\cdots&x_{1n}&1\\
  \vdots&\vdots&&\vdots&\vdots\\
  x_{n1}&x_{n2}&\cdots&x_{nn}&1
  \end{vmatrix}
\]
按第一行展开，可知它是关于动点 $x_\lambda=(x_{\lambda1},x_{\lambda2},\cdots,x_{\lambda n})$ 的坐标的一次式，因此 $\Delta=0$ 代表 $\mathbb{R}^n$ 中一超平面. 因为 $\lambda=i$ 时，行列式 $\Delta$ 变为零，所以点 $x_i$ 皆位于此超平面上（$i=0,1,2,\cdots,n$）.

必要性. 若 $x_0,x_1,\cdots,x_n$ 同位于一超平面上，则存在常数 $a_1,a_2,\cdots,a_n$，使得
\[
\begin{cases}
a_1x_{01}+a_2x_{02}+\cdots+a_nx_{0n}=1,\\
a_1x_{11}+a_2x_{12}+\cdots+a_nx_{1n}=1,\\
\cdots\cdots\cdots\\
a_1x_{n1}+a_2x_{n2}+\cdots+a_nx_{nn}=1.
\end{cases}
\]
此式表明行列式 $\Delta_0$ 中列向量线性相关，因此 $\Delta_0=0$.

当超平面通过原点时，上面方程组右端的 $1$ 应改为 $0$，它表明 $\Delta_0$ 最后一列的余子式恒为 $0$，故 $\Delta_0=0$.

\subsection*{$\star$二、多元函数的极限}

\noindent\textbf{导读}\quad 本段是基础性内容，适合本书各类读者.

\noindent\textbf{a.\quad 多元函数极限的计算}

\noindent\textbf{要点}\quad 计算多元函数的极限常用的方法是：1）利用不等式，使用两边夹法则；2）变量替换化为已知极限，或化为一元函数极限；3）利用极坐标；4）利用初等函数的连续性，利用极限的四则运算性质；5）利用初等变形，特别指数形式常可先求其对数的极限；6）若事先能看出极限值，可用 $\varepsilon$--$\delta$ 方法进行证明.

\noindent\textbf{$\star$例 6.1.10}\quad 1）求
\[
  \lim_{\substack{x\to\infty\\ y\to\infty}}\frac{|x|+|y|}{x^2+y^2};
\]
2）求
\[
  \lim_{\substack{x\to0\\ y\to0}}(x^2+y^2)^{x^2y^2};
\]
3）求
\[
  \lim_{\substack{x\to\infty\\ y\to\infty}}\frac{x^2+y^2}{x^4+y^4};
\]
4）设 $f'(0)=k$，试证明
\[
  \lim_{\substack{a\to0^-\\ b\to0^+}}\frac{f(b)-f(a)}{b-a}=k.
\]
（浙江大学）

\noindent\textbf{解}\quad 1）
\[
\begin{aligned}
0\le \frac{|x|+|y|}{x^2+y^2}
&=\frac{|x|}{x^2+y^2}+\frac{|y|}{x^2+y^2}\\
&\le \frac{|x|}{x^2}+\frac{|y|}{y^2}
=\frac1{|x|}+\frac1{|y|}\to0.
\end{aligned}
\]

2）先求取对数之后的极限：
\[
  \lim_{\substack{x\to0\\ y\to0}}\ln(x^2+y^2)^{x^2y^2}
  =\lim_{\substack{x\to0\\ y\to0}}
  \frac{x^2y^2}{x^2+y^2}(x^2+y^2)\ln(x^2+y^2).
\]
因为
\[
  0\le\frac{x^2y^2}{x^2+y^2}
  \le\frac{(x^2+y^2)^2}{x^2+y^2}
  =x^2+y^2\to0,
\]
\[
  \lim_{\substack{x\to0\\ y\to0}}(x^2+y^2)\ln(x^2+y^2)
  \stackrel{\text{令 }x^2+y^2=t}{=}
  \lim_{t\to0}t\ln t=0,
\]
故原极限 $=e^0=1$.

3）提示\quad 可用极坐标或直接应用不等式.

4）提示\quad （用拟合法.）
\[
  k=\frac{b}{b-a}\cdot k-\frac{a}{b-a}\cdot k,
\]
\[
  \frac{f(b)-f(a)}{b-a}
  =\frac{b}{b-a}\cdot\frac{f(b)-f(0)}{b-0}
  -\frac{a}{b-a}\cdot\frac{f(a)-f(0)}{a-0}.
\]
再提示\quad $a<0<b$，得
\[
  \left|\frac{a}{b-a}\right|<1,\quad
  \left|\frac{b}{b-a}\right|<1.
\]
\[
\begin{aligned}
\left|\frac{f(b)-f(a)}{b-a}-k\right|
&\le \left|\frac{b}{b-a}\right|
\left|\frac{f(b)-f(0)}{b-0}-k\right|\\
&\quad+\left|\frac{a}{b-a}\right|
\left|\frac{f(a)-f(0)}{a-0}-k\right|\\
&\le \left|\frac{f(b)-f(0)}{b-0}-k\right|
+\left|\frac{f(a)-f(0)}{a-0}-k\right|\to0
\end{aligned}
\]
（当 $a\to0^-,b\to0^+$ 时）.

\noindent\textbf{b.\quad 证明二元极限不存在}

\noindent\textbf{要点}\quad 根据全面极限与特殊路径极限的关系，证明二元极限不存在. 通常方法是：1）证明径向路径的极限与辐角（或斜率）有关；2）证明某个特殊路径的极限不存在；3）证明两个特殊极限存在但不相等；4）若二元函数在该点某空心邻域里连续，而两累次极限存在不相等，则该点全面极限不存在.

\noindent\textbf{$\star$例 6.1.11}\quad 证明下列函数在 $(0,0)$ 处全面极限不存在：
\[
  1)\quad f_1(x,y)=\frac{xy}{x^2+y^2};
\]
\[
  2)\quad f_2(x,y)=\frac{xy}{x+y};
\]
\[
  3)\quad f_3(x,y)=\frac{x^6y^8}{(x^2+y^4)^5};
\]
\[
  4)\quad f_4(x,y)=\frac{x^3-y^3}{x^3+y^3}.
\]

\noindent\textbf{提示}\quad 1）令 $y=kx$ 或化为极坐标.

2）分母当 $y=-x$ 时为零，因此可以考虑沿与 $y=-x$ 相切的高次曲线的路径的极限，例如令 $y=mx^2-x$，令 $x\to0$，取极限得 $-\dfrac1m$，与 $m$ 有关.

3）可比较 $x=0$ 与 $x=y^2$ 两路径的极限.

4）$f_4$ 除 $(0,0)$ 点外，处处连续，但在 $(0,0)$ 点两累次极限存在，不相等.

\noindent\textbf{注}\quad 累次极限一般不是特殊路径的极限，但在某空心邻域里若函数连续，则累次极限实为沿坐标轴方向的极限.

\noindent\textbf{例 6.1.12}\quad 函数
\[
  f(x,y)=\frac{x^4y^4}{(x^3+y^6)^2}
\]
在 $(0,0)$ 点的极限 $\displaystyle\lim_{(x,y)\to(0,0)}f(x,y)$ 存在吗？若存在，求其值.（华东师范大学）

\noindent\textbf{提示}\quad 可考虑沿路径 $x=my^2$ 的极限（$m$ 取不同的常数）.

\noindent\textbf{c.\quad 关于全面极限与特殊路径极限的进一步讨论}

\noindent\textbf{例 6.1.13}\quad 证明：

1）$f(x,y)$ 当 $(x,y)$ 沿径向路径趋向 $(x_0,y_0)$ 时极限存在，保持相等，全面极限
\[
  \lim_{(x,y)\to(x_0,y_0)}f(x,y)
\]
仍可以不存在；

2）但若沿径向路径极限存在相等，并关于辐角 $\theta\in[0,2\pi]$ 一致，则全面极限
\[
  \lim_{(x,y)\to(x_0,y_0)}f(x,y)
\]
存在.

\noindent\textbf{提示}\quad 1）可考虑
\[
  \lim_{(x,y)\to(0,0)}\frac{x^2y}{x^4+y^2}.
\]
2）利用极限定义，容易证得.

\noindent\textbf{$\star$例 6.1.14}\quad 设 $f(x,y)$ 是在区域 $D:|x|\le1,|y|\le1$ 上的有界 $k$ 次齐次函数（$k\ge1$），问极限
\[
  \lim_{(x,y)\to(0,0)}[f(x,y)+(x-1)e^y]
\]
是否存在？若存在，试求其值.（南京大学）

\noindent\textbf{解}\quad 因 $f$ 为 $k$ 次齐次函数，故 $\forall t\in\mathbb{R}$，有 $f(tx,ty)=t^kf(x,y)$，因此
\[
  f(r\cos\theta,r\sin\theta)=r^kf(\cos\theta,\sin\theta).
\]
又因 $f(x,y)$ 有界，$\exists M>0$，使得 $|f(x,y)|\le M$（$\forall(x,y)\in D$）. 所以
\[
  |f(r\cos\theta,r\sin\theta)|
  =r^k|f(\cos\theta,\sin\theta)|
  \le r^kM\to0
\]
（当 $r\to0$ 时关于 $\theta\in[0,2\pi]$ 一致）.
于是
\[
  \lim_{(x,y)\to(0,0)}[f(x,y)+(x-1)e^y]=-1.
\]

\noindent\textbf{*例 6.1.15}\quad 设

1）$\forall\theta\in[0,2\pi],f(x,y)=f(r\cos\theta,r\sin\theta)\to0$（当 $r\to0$ 时）；

2）存在 $M>0$，使得对任意两点 $(x,y),(x_0,y_0)$（设它们离原点距离相等，且 $>0$），满足
\[
  |f(x,y)-f(x_0,y_0)|
  \le M\left|\arctan\frac{x_0y-xy_0}{xx_0+yy_0}\right|.
\]
试证：
\[
  \lim_{r\to0}f(x,y)=0\quad (r=\sqrt{x^2+y^2}).
\]

\noindent\textbf{证}\quad 条件 1）表明：$\forall\theta\in[0,2\pi]$，$\forall\varepsilon>0$，$\exists\delta=\delta(\theta,\varepsilon)>0$，当 $0<r<\delta$ 时，有 $|f(x,y)|<\varepsilon$.（注意：这里的 $\delta$ 不仅依赖 $\varepsilon$，还依赖 $\theta$，而 $\theta\in[0,2\pi]$ 有无穷多个选择，因此未必存在最小的 $\delta>0$. 必须求助条件 2）.）

在条件 2）里：记 $(x,y)=(r\cos\alpha,r\sin\alpha)$，$(x_0,y_0)=(r\cos\beta,r\sin\beta)$，$r>0$，则
\[
  \left|\arctan\frac{x_0y-xy_0}{xx_0+yy_0}\right|
  =\left|\arctan\frac{\sin(\alpha-\beta)}{\cos(\alpha-\beta)}\right|
  =|\alpha-\beta|,
\]
这里 $\alpha,\beta$ 分别是点 $(x,y),(x_0,y_0)$ 的辐角. 于是由条件 2）知：$\forall\varepsilon>0$，$\exists\delta_0=\dfrac{\varepsilon}{2M}>0$，当 $|\alpha-\beta|<\delta_0$ 时，有
\[
\begin{aligned}
|f(x,y)-f(x_0,y_0)|
&\le M\left|\arctan\frac{x_0y-xy_0}{xx_0+yy_0}\right|\\
&=M|\alpha-\beta|<M\cdot\frac{\varepsilon}{2M}=\frac{\varepsilon}{2}.
\end{aligned}
\tag{1}
\]

（称原点出发的射线为经线，（原点为圆心的）同心圆为纬线，那么条件 1）表明：“在同一条经线上的任一点 $(x,y)$，当其向径 $r$ 充分小时，对应的函数值就能任意小；条件 2）表明：“在同一纬线上的两点，若它们辐角之差充分小，则对应的函数值相差也能任意小.”）

现从原点出发，画出 $n$ 条射线，将平面 $\mathbb{R}^2$ 均匀分为 $n$ 个相等的扇形. 取 $n$ 充分大，使得每个扇形的辐角 $\gamma$ 足够小，以至于
\[
  \gamma=\frac{2\pi}{n}\le\delta_0=\frac{\varepsilon}{2M},
  \quad\text{即}\quad n\ge N=\left[\frac{4\pi M}{\varepsilon}\right].
\]
此时，平面 $\mathbb{R}^2$ 被等分成 $n$ 个扇形，第 $i$ 号射线 $\ell_i$ 的辐角记为“$i\gamma$”，第 $i$ 个扇形里的任意点 $(x,y)=(r\cos\alpha,r\sin\alpha)$，辐角 $\alpha$ 满足：
\[
  (i-1)\gamma<\alpha\le i\gamma\quad (i=1,2,\cdots,n).
\]
根据条件 1），在 $\ell_i$ 上，$\forall\varepsilon>0$，$\exists\delta_i>0$，当 $r<\delta_i$ 时，有
\[
  |f(x_i,y_i)|=|f(r\cos\alpha_i,r\sin\alpha_i)|<\frac{\varepsilon}{2},
  \quad (x_i,y_i)\in\ell_i\ (i=1,2,\cdots,n).
\tag{2}
\]
令
\[
  \delta=\min\{\delta_i\mid i=0,1,2,\cdots,n\},
\]
则 $(x,y)=(r\cos\alpha,r\sin\alpha)\in\mathbb{R}^2$（必落在某个扇形里，例如在第 $i$ 个（扇形）里）与射线 $\ell_i$ 上向径为 $r$ 的点 $(x_i,y_i)=(r\cos\alpha_i,r\sin\alpha_i)$ 的辐角差：
\[
  |\alpha-\alpha_i|\le\gamma=\frac{2\pi}{n}.
\]
因此当 $r<\delta$ 时，（由式 (1) 和 (2)）有
\[
  |f(x,y)|\le |f(x,y)-f(x_i,y_i)|+|f(x_i,y_i)|
  <\frac{\varepsilon}{2}+\frac{\varepsilon}{2}=\varepsilon.
\]
$\displaystyle\lim_{r\to0}f(x,y)=0$ 获证.

\noindent\textbf{$\star$练习}\quad 设点 $M(x,y)$ 沿任意路径趋向 $M_0(x_0,y_0)$ 时，函数 $f(x,y)$ 的极限恒为 $A$，试证
\[
  \lim_{M\to M_0}f(x,y)=A.
\]

\noindent\textbf{证}\quad （反证法）若 $f(x,y)\nrightarrow A$（当 $x\to x_0,y\to y_0$ 时），则 $\exists\varepsilon_0>0$ 及点列 $\{M_n\}$（$M_n\to M_0$，当 $n\to\infty$ 时）使得
\[
  |f(M_n)-A|\ge\varepsilon_0\quad(n=1,2,\cdots).
\]
如此顺序用直线段将 $M_1,M_2,M_3,\cdots$ 连成折线 $L$，则 $M$ 沿 $L$ 趋向 $M_0$ 时，$f(x,y)\nrightarrow A$，与已知条件矛盾.

\noindent\textbf{new *例 6.1.16}\quad 已知函数 $f(x,y)$ 定义在 $D=(a,b)\times[c,d]$ 上，且

（i）$x$ 固定时，对 $y$ 连续；

（ii）设 $x_0\in(a,b)$ 取定，对任意 $y\in[c,d]$，有
\[
  \lim_{x\to x_0}f(x,y)\stackrel{\text{存在}}{=}g(y).
\]
试证：如下两个断言等价：

（A）重极限
\[
  \lim_{\substack{x\to x_0\\ y\to y_0}} f(x,y)=g(y_0)
\]
对于任意 $y_0\in[c,d]$ 成立；

（B）当 $x\to x_0$ 时，$f(x,y)\to g(y)$ 关于 $y\in[c,d]$ 一致.

当两论断都成立时，$g(y)$ 在 $[c,d]$ 上一致连续.（北京大学）

\noindent\textbf{证}\quad $1^\circ$（（A）$\Rightarrow$（B）.）已知：
\[
  \lim_{\substack{x\to x_0\\ y\to y_0}}f(x,y)=g(y_0)
  \quad(\text{对任意 }y_0\in[c,d]\text{ 成立}).
\]
因此，对每个 $y_0\in[c,d]$：$\forall\varepsilon>0$，$\exists\delta=\delta_{y_0}>0$，当 $|x-x_0|<\delta,|y-y_0|<\delta$（即 $y\in U(y_0,\delta)$）时，有
\[
  |f(x,y)-g(y_0)|<\varepsilon. \tag{1}
\]
在此式中，将 $y\in U(y_0,\delta)$ 固定，令 $x\to x_0$，则得
\[
  |g(y)-g(y_0)|\le\varepsilon. \tag{2}
\]
此式表明：$\forall y_0\in[c,d]$，$g(y)$ 在 $y_0$ 处连续，从而一致连续（Cantor 定理）.

如此，开邻域 $U(y_0,\delta_{y_0})$（$\forall y_0\in[c,d]$）组成 $[c,d]$ 上的开覆盖. 根据有限覆盖定理，从中可选出有限子覆盖，记作 $U(y_k,\delta_{y_k})$（$k=1,2,\cdots,m$）. 这时，$\forall y\in[c,d]\subset\displaystyle\sum_{k=1}^m U(y_k,\delta_{y_k})$，必属于某一个 $U(y_k,\delta_{y_k})$，因此取
\[
  \bar\delta=\min_{1\le k\le m}\delta_k,
\]
则当 $|x-x_0|<\bar\delta$ 时，对任意 $y\in[c,d]$，有式 (1) 成立. 故当 $|x-x_0|<\bar\delta$ 时，利用式 (1) 和 (2)，恒有
\[
  |f(x,y)-g(y_0)|\le |f(x,y)-g(y)|+|g(y)-g(y_0)|<2\varepsilon
\]
（$\forall y\in[c,d]$ 一致成立）.

$2^\circ$（（B）$\Rightarrow$（A）.）已知当 $x\to x_0$ 时，有 $f(x,y)\to g(y)$ 关于 $y\in[c,d]$ 一致. 因此，$\forall\varepsilon>0$，$\exists\delta_{\text{横向}}>0$（横向指 $x$ 的变化范围），当 $x\in(a,b),|x-x_0|<\delta_{\text{横向}}$ 时，就有
\[
  |f(x,y)-g(y)|<\varepsilon\quad(\forall y\in[c,d]). \tag{3}
\]
任意取定 $y_0\in[c,d]$，由已知条件（i），$x$ 固定时，$f(x,y)$ 对 $y$ 连续. 因此在式 (3) 中，若将 $x$ 固定（令 $x=x_i,|x_i-x_0|<\delta$），再让 $y\to y_0$，则对于上述任意 $\varepsilon>0$，$\exists\delta_i>0$，当 $|y-y_0|<\delta_i$ 时，有
\[
  |f(x_i,y)-f(x_i,y_0)|<\varepsilon.
\]
于是
\[
\begin{aligned}
|g(y)-g(y_0)|
&\le |g(y)-f(x_i,y)|+|f(x_i,y)-f(x_i,y_0)|\\
&\quad+|f(x_i,y_0)-g(y_0)|<3\varepsilon.
\end{aligned}
\]
这说明：$\forall y_0\in[c,d]$，$g(y)$ 在 $y_0$ 处连续，（根据 Cantor 定理）从而 $g(y)$ 在 $[c,d]$ 一致连续. 故 $\forall\varepsilon>0$，$\exists\delta_{\text{竖向}}>0$（竖向指 $y$ 的变化范围），$\forall y,y_0\in[c,d]$，当 $|y-y_0|<\delta_{\text{竖向}}$ 时，有
\[
  |g(y)-g(y_0)|<\varepsilon. \tag{4}
\]
取 $\delta=\min\{\delta_{\text{横向}},\delta_{\text{竖向}}\}$，则当 $|x-x_0|<\delta$ 且 $|y-y_0|<\delta$（即方邻域里）时，
\[
  |f(x,y)-g(y_0)|\le |f(x,y)-g(y)|+|g(y)-g(y_0)|<2\varepsilon,
\]
亦即
\[
  \lim_{\substack{x\to x_0\\ y\to y_0}}f(x,y)\stackrel{\text{存在}}{=}g(y_0).
\]
证毕.

\noindent\textbf{d.\quad 累次极限交换次序问题}

\noindent\textbf{$\star$例 6.1.17}\quad $\Omega$ 为 $\mathbb{R}^2$ 中的开集，$(x_0,y_0)\in\Omega$，$f(x,y)$ 为 $\Omega$ 上的函数，且

1）对每个 $(x,y)\in\Omega$ 的 $x$ 存在
\[
  \lim_{y\to y_0}f(x,y)=g(x);
\]
2）
\[
  \lim_{x\to x_0}f(x,y)=h(y)
\]
关于 $(x,y)\in\Omega$ 中的 $y$ 一致.

试证：
\[
  \lim_{x\to x_0}\lim_{y\to y_0}f(x,y)
  =\lim_{y\to y_0}\lim_{x\to x_0}f(x,y).
\tag{1}
\]
（辽宁大学）

\noindent\textbf{方法}\quad 为了证明等式 (1)，只要证明等式左端的累次极限
\[
  \lim_{x\to x_0}\lim_{y\to y_0}f(x,y)=\lim_{x\to x_0}g(x)=A
\]
存在，且右端的函数 $h(y)=\displaystyle\lim_{x\to x_0}f(x,y)$ 当 $y\to y_0$ 时趋向 $A$.

\noindent\textbf{证}\quad $1^\circ$（证明 $\displaystyle\lim_{x\to x_0}g(x)$ 存在.）因 $(x_0,y_0)\in\Omega$（$\Omega$ 为开集），所以 $\exists\delta_1>0$，使得
\[
  \{(x,y)\mid |x-x_0|<\delta_1,\ |y-y_0|<\delta_1\}\subseteq\Omega.
\]
由条件 2），$\forall\varepsilon>0$，$\exists\delta>0$（$\delta<\delta_1$），当 $0<|x'-x_0|<\delta,0<|x''-x_0|<\delta$ 时，有
\[
  |f(x',y)-f(x'',y)|<\varepsilon
  \quad(\forall y\in\{y\mid |y-y_0|<\delta\}).
\]
令 $y\to y_0$ 取极限，（据条件 1））得
\[
  |g(x')-g(x'')|\le\varepsilon.
\]
根据 Cauchy 准则，知 $\displaystyle\lim_{x\to x_0}g(x)$ 存在. 即等式 (1) 左端极限存在. 记之为 $A$.

$2^\circ$（证明 $\displaystyle\lim_{y\to y_0}h(y)=A$.）$\forall\varepsilon>0$，由
\[
  |h(y)-A|\le |h(y)-f(x,y)|+|f(x,y)-g(x)|+|g(x)-A|, \tag{2}
\]
利用条件 2）及 $1^\circ$ 之结论，可取 $x$ 与 $x_0$ 充分接近使得
\[
  |h(y)-f(x,y)|<\frac{\varepsilon}{3},\quad
  |g(x)-A|<\frac{\varepsilon}{3}.
\]
将 $x$ 固定，由条件 1），$\exists\delta>0$，使得 $|y-y_0|<\delta$ 时，$|f(x,y)-g(x)|<\dfrac{\varepsilon}{3}$. 于是由式 (2) 知 $|h(y)-A|<\varepsilon$. 证毕.

\subsection*{*三、多元连续函数}

\noindent\textbf{导读}\quad 本段主要针对数学院系的学生.

\noindent\textbf{a.\quad 连续性的证明}

\noindent\textbf{要点}\quad 要证明 $f(x,y)$ 在 $(x_0,y_0)$ 处连续，即要：$\forall\varepsilon>0$，找 $\delta>0$，使得 $|x-x_0|<\delta,|y-y_0|<\delta$ 时，恒有
\[
  |f(x,y)-f(x_0,y_0)|<\varepsilon;
\]
或等价地，当
\[
  \sqrt{(x-x_0)^2+(y-y_0)^2}<\delta
\]
时，有 $|f(x,y)-f(x_0,y_0)|<\varepsilon$.

\noindent\textbf{$\star$例 6.1.18}\quad 设 $f(x)$ 及 $g(y)$ 分别在区间 $[a,b]$，$[c,d]$ 上连续，定义
\[
  F(x,y)=\int_a^x f(s)\,ds\cdot\int_c^y g(t)\,dt
  \quad(a\le x\le b,\ c\le y\le d).
\]
试用“$\varepsilon$--$\delta$”方法证明 $F(x,y)$ 在
\[
  D=\{(x,y)\mid a\le x\le b,\ c\le y\le d\}
\]
内连续.（大连理工大学）

\noindent\textbf{证}\quad 因 $f(x),g(y)$ 分别在 $[a,b],[c,d]$ 上连续，故 $\exists M>0$，使得 $|f(x)|\le M,|g(y)|\le M$（当 $a\le x\le b,c\le y\le d$ 时）. 于是
\[
\begin{aligned}
|F(x,y)-F(x_0,y_0)|
={}&\left|\int_a^x f(s)\,ds\cdot\int_c^y g(t)\,dt
-\int_a^{x_0} f(s)\,ds\cdot\int_c^{y_0} g(t)\,dt\right|\\
\le{}&\left|\int_a^x f(s)\,ds\cdot\int_c^y g(t)\,dt
-\int_a^{x_0} f(s)\,ds\cdot\int_c^y g(t)\,dt\right|\\
&+\left|\int_a^{x_0} f(s)\,ds\cdot\int_c^y g(t)\,dt
-\int_a^{x_0} f(s)\,ds\cdot\int_c^{y_0} g(t)\,dt\right|\\
={}&\left|\int_{x_0}^x f(s)\,ds\right|
\left|\int_c^y g(t)\,dt\right|
+\left|\int_a^{x_0} f(s)\,ds\right|
\left|\int_{y_0}^y g(t)\,dt\right|\\
\le{}&\int_{x_0}^x |f(s)|\,ds\cdot\int_c^d |g(t)|\,dt
+\int_a^b |f(s)|\,ds\cdot\int_{y_0}^y |g(t)|\,dt\\
\le{}&M^2(d-c)|x-x_0|+M^2(b-a)|y-y_0|.
\end{aligned}
\]
记 $\Delta=\max\{b-a,d-c\}$，于是，$\forall(x_0,y_0)\in D$，$\forall\varepsilon>0$，取
\[
  \delta=\frac{\varepsilon}{2M^2\Delta}>0,
\]
当 $|x-x_0|<\delta,|y-y_0|<\delta,(x,y)\in D$ 时，恒有 $|F(x,y)-F(x_0,y_0)|<\varepsilon$. 证毕.

应当指出的是，如果未限定用 $\varepsilon$--$\delta$ 方法证明，本题用连续函数运算性质做更快. 因为 $\int_a^x f(s)\,ds$ 与 $\int_c^y g(t)\,dt$ 分别为 $x$ 与 $y$ 的一元连续函数，看作二元函数自然也连续. 用连续函数的乘法定理，便知 $F(x,y)$ 连续.

\noindent\textbf{$\star$例 6.1.19}\quad 设 $u=f(x,y,z)$ 在闭立方体 $\overline D[a,b;a,b;a,b]$ 上连续，试证
\[
  g(x,y)=\max_{a\le z\le b} f(x,y,z)
\]
在正方形 $[a,b;a,b]\subset\mathbb{R}^2$ 上连续.

\noindent\textbf{提示}\quad （或作为引理）先证明：若 $f(x,y)$ 在 $D=[a,b]\times[c,d]$ 上连续，则
\[
  g(y)=\max_{x\in[a,b]} f(x,y)
\]
在 $[c,d]$ 上连续.

\noindent\textbf{证}\quad 因 $f(x,y,z)$ 在 $\overline D$ 上连续，故在 $\overline D$ 上一致连续. 于是，$\forall\varepsilon>0$，$\exists\delta>0$，$\overline D$ 上当 $|x-x'|<\delta,|y-y'|<\delta,|z-z'|<\delta$ 时恒有
\[
  |f(x,y,z)-f(x',y',z')|<\varepsilon.
\]
特别，当 $|x-x_0|<\delta,|y-y_0|<\delta$ 时，有
\[
  |f(x,y,z)-f(x_0,y_0,z)|<\varepsilon\quad(\forall z\in[a,b]).
\]
即
\[
  f(x_0,y_0,z)-\varepsilon<f(x,y,z)<f(x_0,y_0,z)+\varepsilon.
\]
固定 $x,y$，让 $z$ 在 $[a,b]$ 上变化，取最大值，可得
\[
  f(x_0,y_0,z)-\varepsilon<f(x,y,z)<g(x_0,y_0)+\varepsilon,
  \quad \forall z\in[a,b].
\]
此不等式中间一项取成最大值 $\displaystyle\max_{a\le z\le b}f(x,y,z)$ 时，上式仍成立，得
\[
  f(x_0,y_0,z)-\varepsilon<g(x,y)<g(x_0,y_0)+\varepsilon,\quad \forall z\in[a,b].
\]
最后令左边第一项取最大值，得
\[
  g(x_0,y_0)-\varepsilon<g(x,y)<g(x_0,y_0)+\varepsilon.
\]
即 $|x-x_0|<\delta,|y-y_0|<\delta$ 时，$|g(x,y)-g(x_0,y_0)|<\varepsilon$.

至此实际已证明了：$g(x,y)$ 不仅在 $[a,b;a,b]$ 上连续，而且一致连续.

\noindent\textbf{向量函数连续性问题}

\noindent\textbf{要点}\quad 设 $(u,v)=F(x,y)=(f(x,y),g(x,y))$ 是 $\mathbb{R}^2\to\mathbb{R}^2$ 的函数. 要证明 $F$ 在 $(x_0,y_0)$ 处连续（即当 $(x,y)\to(x_0,y_0)$ 时，有 $(f(x,y),g(x,y))\to(f(x_0,y_0),g(x_0,y_0))$），等价于要证明 $f(x,y),g(x,y)$ 都在 $(x_0,y_0)$ 处连续，也等价于要证明当
\[
  r=\sqrt{(x-x_0)^2+(y-y_0)^2}\to0
\]
时
\[
  \sqrt{(u-u_0)^2+(v-v_0)^2}\to0.
\]
至于 $\mathbb{R}^n$ 到 $\mathbb{R}^m$ 的映射，情况类似.

\noindent\textbf{※例 6.1.20}\quad 讨论如下向量函数的连续性：设
\[
  (u,v)=F(x,y)=(f(x,y),g(x,y)),
\]
其中
\[
u=f(x,y)=
\begin{cases}
\dfrac{x}{(x^2+y^2)^\alpha}\ln(|x|+|y|), & x^2+y^2\ne0,\\
0, & x^2+y^2=0,
\end{cases}
\]
\[
v=g(x,y)=
\begin{cases}
\dfrac{y}{(x^2+y^2)^\alpha}\ln(|x|+|y|), & x^2+y^2\ne0,\\
0, & x^2+y^2=0.
\end{cases}
\]

\noindent\textbf{解}\quad 显然 $f(x,y),g(x,y)$ 当 $x^2+y^2\ne0$ 时连续，因此 $(x,y)\ne(0,0)$ 时 $F$ 连续. 下面只研究 $(0,0)$ 点的情况. 因为
\begin{equation}
\begin{aligned}
u^2+v^2
&=\frac{x^2+y^2}{(x^2+y^2)^{2\alpha}}\ln^2(|x|+|y|)\\
&=(x^2+y^2)^{1-2\alpha}\ln^2(|x|+|y|)\\
&\to
\begin{cases}
0, & \text{当 }\alpha<\dfrac12\text{ 时},\\
+\infty, & \text{当 }\alpha\ge\dfrac12\text{ 时},
\end{cases}
\quad(\text{当 }r^2=x^2+y^2\to0\text{ 时}).
\end{aligned}
\tag{1}
\end{equation}
故当且仅当 $\alpha<\dfrac12$ 时 $F$ 在 $(0,0)$ 点连续. 下面对式 (1) 中的极限进行补充证明.

当 $\alpha\ge\dfrac12$（即 $1-2\alpha\le0$）时，显然极限为 $+\infty$. 现设 $\alpha<\dfrac12$，记 $\mu=1-2\alpha$，则 $\mu>0$，
\[
  (x^2+y^2)^{1-2\alpha}\ln^2(|x|+|y|)
  =\frac{(x^2+y^2)^\mu}{(|x|+|y|)^{2\mu}}\cdot(|x|+|y|)^{2\mu}\ln^2(|x|+|y|).
\]
这时
\[
  0\le \frac{(x^2+y^2)^\mu}{(|x|+|y|)^{2\mu}}
  =\frac{(x^2+y^2)^\mu}{(x^2+2|x||y|+y^2)^\mu}
  \le \frac{(x^2+y^2)^\mu}{(x^2+y^2)^\mu}=1,
\]
但
\[
  (|x|+|y|)^{2\mu}\ln^2(|x|+|y|)\to0,
\]
所以
\[
  (x^2+y^2)^{1-2\alpha}\ln^2(|x|+|y|)\to0
  \quad(\text{当 }r\to0\text{ 时}).
\]

\noindent\textbf{b.\quad 全面连续与按单变量连续的关系}

\noindent\textbf{要点}\quad 全面连续必按各单变量连续，反之按各单变量连续，不一定全面连续，只有补充某种条件之后，才能保证全面连续.

\noindent\textbf{例 6.1.21}\quad
\[
f(x,y)=
\begin{cases}
\dfrac{xy}{x^2+y^2}, & \text{当 }x^2+y^2\ne0\text{ 时},\\
0, & \text{当 }x^2+y^2=0\text{ 时},
\end{cases}
\]
在 $(0,0)$ 处关于单变量 $x$ 与 $y$ 都是连续的，但在 $(0,0)$ 处不全面连续.

\noindent\textbf{$\star$例 6.1.22}\quad 若 $f(x,y)$ 分别是单变量 $x$ 及 $y$ 的连续函数，又对其中一个变量是单调的，试证 $f(x,y)$ 是二元连续函数.（陕西师范大学）

\noindent\textbf{分析}\quad 假设 $f(x,y)$ 对 $y$ 单调增加，关于 $x,y$ 分别连续. $M_0(x_0,y_0)$ 是任意一点，要证明 $f(x,y)$ 在 $M_0(x_0,y_0)$ 处连续，即要对任意 $\varepsilon>0$，找相应的邻域 $U$，使得 $(x,y)\in U$ 时，有
\[
  |f(x,y)-f(x_0,y_0)|<\varepsilon.
\]
如图 6.1.1，因 $f(x,y)$ 对 $y$ 连续，故 $\delta_1>0$ 充分小时，有
\[
  |f(M_1)-f(M_0)|<\frac{\varepsilon}{2},\quad
  |f(M_2)-f(M_0)|<\frac{\varepsilon}{2},
\]
这里 $M_1(x_0,y_0-\delta_1),M_2(x_0,y_0+\delta_1)$. 又因 $f$ 对 $x$ 连续，所以 $\delta>0$ 充分小时，若 $M\in\overline{M_3M_4}$，有
\[
  |f(M)-f(M_1)|<\frac{\varepsilon}{2};
\]
若 $M\in\overline{M_5M_6}$，有
\[
  |f(M)-f(M_2)|<\frac{\varepsilon}{2}.
\]

\begin{center}
  \includegraphics[width=0.45\textwidth]{6-1-1.png}
\end{center}

至此在 $M_0(x_0,y_0)$ 的方形邻域：矩形 $M_3M_4M_6M_5$ 内恒有
\[
  |f(x,y)-f(x_0,y_0)|<\varepsilon.
\]
原因是 $f(x,y)$ 对 $y$ 单调，$f(x,y)$ 夹于 $f(x,y_0-\delta_1)$ 和 $f(x,y_0+\delta_1)$ 之间. 例如 $f(x,y)$ 单调增加，则
\[
  f(x,y)\le f(x,y_0+\delta_1)<f(x_0,y_0+\delta_1)+\frac{\varepsilon}{2}
  <\left(f(x_0,y_0)+\frac{\varepsilon}{2}\right)+\frac{\varepsilon}{2}
  =f(x_0,y_0)+\varepsilon.
\]
又
\[
  f(x,y)\ge f(x,y_0-\delta_1)>f(x_0,y_0-\delta_1)-\frac{\varepsilon}{2}
  >\left(f(x_0,y_0)-\frac{\varepsilon}{2}\right)-\frac{\varepsilon}{2}
  =f(x_0,y_0)-\varepsilon.
\]
总之
\[
  f(x_0,y_0)-\varepsilon<f(x,y)<f(x_0,y_0)+\varepsilon,
\]
即 $|f(x,y)-f(x_0,y_0)|<\varepsilon$.

\noindent\textbf{$\star$例 6.1.23}\quad 在所讨论区域上设 $f(x,y)$ 分别对 $x$ 和 $y$ 连续，试证在下列条件之一满足时，$f(x,y)$ 全面连续：

1）$f(x,y)$ 对 $x$ 连续关于 $y$ 一致（即 $\forall x_0,\forall\varepsilon>0$，$\exists\delta=\delta(\varepsilon,x_0)>0$（与 $y$ 无关），当 $|x-x_0|<\delta$ 时，对一切 $y$ 恒有 $|f(x,y)-f(x_0,y)|<\varepsilon$）；

2）$f(x,y)$ 对 $y$ 连续关于 $x$ 一致；

3）特别，若对其中一个变量满足 Lipschitz 条件（例如对 $y$ 满足 Lipschitz 条件，即 $\exists L>0$，使得 $\forall y_1,y_2,x$，有 $|f(x,y_1)-f(x,y_2)|\le L|y_1-y_2|$）；

4）设所考虑的范围是某个有界闭区域 $D$，而 $f$ 在包含 $D$ 的某个区域 $G$ 上有意义，且在 $G$ 上对变量 $x$ 或 $y$ 满足局部 Lipschitz 条件（例如对 $y$ 满足局部 Lipschitz 条件，即 $\forall(x_0,y_0)\in G$，存在邻域 $U\subset G$ 及 $L>0$，使得 $\forall(x,y_1),(x,y_2)\in U$，有 $|f(x,y_1)-f(x,y_2)|\le L|y_1-y_2|$）.

\noindent\textbf{证}\quad 1）$\forall(x_0,y_0)$，$\forall\varepsilon>0$，$\exists\delta_1=\delta_1(\varepsilon,x_0)>0$（与 $y$ 无关），当 $|x-x_0|<\delta_1$ 时，对一切 $y$ 有
\[
  |f(x,y)-f(x_0,y)|<\frac{\varepsilon}{2}.
\]
又因 $(x_0,y_0)$ 处 $f(x_0,y)$ 对 $y$ 连续，故对此 $\varepsilon>0$，$\exists\delta_2>0$，使得 $|y-y_0|<\delta_2$ 时，有
\[
  |f(x_0,y)-f(x_0,y_0)|<\frac{\varepsilon}{2}.
\]
取 $\delta=\min\{\delta_1,\delta_2\}$，则 $|x-x_0|<\delta,|y-y_0|<\delta$ 时，有
\[
\begin{aligned}
|f(x,y)-f(x_0,y_0)|
&\le |f(x,y)-f(x_0,y)|+|f(x_0,y)-f(x_0,y_0)|\\
&<\frac{\varepsilon}{2}+\frac{\varepsilon}{2}=\varepsilon.
\end{aligned}
\]
由 $(x_0,y_0)$ 的任意性，就证明了 $f$ 的连续性.

2）类似 1）.

3）可从 Lipschitz 条件导出条件 1）或 2）.

4）可从条件 4）导出条件 1）（用有限覆盖定理）.

\noindent\textbf{例 6.1.24}\quad 设函数 $f(x,y)$ 在原点附近有定义，令
\[
  F(r,\theta)=f(r\cos\theta,r\sin\theta)\quad(r\ge0,\ 0\le\theta<2\pi).
\]
如果 $F(r,\theta)$ 满足如下条件：

1）$\forall\theta\in[0,2\pi]$，$F(r,\theta)$ 对 $r$ 连续；

2）对任意 $\varepsilon>0$，存在 $\delta>0$，当 $|\theta-\theta'|<\delta$ 时，有
\[
  |F(r,\theta)-F(r,\theta')|<\varepsilon,
\]
对于 $r$ 一致成立.

证明：函数 $f(x,y)$ 在原点 $(0,0)$ 处连续.（上海交通大学）

\noindent\textbf{提示}\quad 参见例 6.1.15.

\noindent\textbf{c.\quad 连续性的等价描述}

\noindent\textbf{要点}\quad 连续性除用 $\varepsilon$--$\delta$ 语言描述之外，还可等价地用邻域、序列、开集、闭集等不同的方式进行描述.

\noindent\textbf{例 6.1.25}\quad 设 $f(M)$ 在区域 $D$ 内定义\footnote{“区域”意指开区域，即其内每一点皆为内点.}，则如下诸条件等价：

1）$f(M)$ 在 $D$ 内连续（即：$\forall M_0(x_1^0,\cdots,x_n^0)\in D,\forall\varepsilon>0,\exists\delta>0$，使得当 $|x_i-x_i^0|<\delta(i=1,2,\cdots,n)$ 时，有 $|f(M)-f(M_0)|<\varepsilon$）；

$1'$）$\forall M_0\in D,\forall\varepsilon>0,\exists\delta>0$，使得当 $|M-M_0|<\delta$ 时，有
\[
  |f(M)-f(M_0)|<\varepsilon,
\]
这里
\[
  |M-M_0|=\sqrt{(x_1-x_1^0)^2+\cdots+(x_n-x_n^0)^2};
\]

$1''$）$\forall M_0\in D,\forall U\in\mathcal{N}(f(M_0))$，$\exists V\in\mathcal{N}(M_0)$ 使得 $f(V\cap D)\subseteq U$. 这里 $\mathcal{N}(M_0)$ 表示 $\mathbb{R}^n$ 中点 $M_0$ 的全体邻域组成的集合；$U\in\mathcal{N}(M_0)$ 表示 $U$ 为 $M_0$ 的一个邻域，$\mathcal{N}(f(M_0))$ 表示值域空间 $\mathbb{R}$ 中像点 $f(M_0)$ 的邻域集；

2）$\forall\alpha\in\mathbb{R}$，集合
\[
  E=\{M\mid f(M)>\alpha\}\quad\text{与}\quad
  F=\{M\mid f(M)<\alpha\}
\]
皆为开集；

$2'$）$\forall\alpha<\beta$，
\[
  G=\{M\mid \alpha<f(M)<\beta\}
\]
恒为开集；

3）$\forall\alpha\in\mathbb{R}$，
\[
  E=\{M\mid f(M)\ge\alpha\}\quad\text{与}\quad
  F=\{M\mid f(M)\le\alpha\}
\]
皆为闭集；

4）若 $M_n\in D$，当 $n\to\infty$ 时 $M_n\to M_0\in D$，则必有 $f(M_n)\to f(M_0)$.

\noindent\textbf{证}\quad 与一元情况类似. 下面只就 1），2），$2'$）的等价性进行证明，其余留给读者.

（$1\Rightarrow2$）设 $M_0\in E$，即 $f(M_0)>\alpha$，根据连续函数保号性，知 $\exists\delta>0$，使得 $|M-M_0|<\delta$ 时有 $f(M)>\alpha$. 即 $M_0$ 的邻域 $\{M\mid |M-M_0|<\delta\}\subseteq E$，所以 $E$ 为开集. 同理可证 $F$ 为开集.

（$2\Rightarrow2'$）因 $\forall\alpha<\beta$，
\[
  \{M\mid \alpha<f(M)<\beta\}
  =\{M\mid f(M)>\alpha\}\cap\{M\mid f(M)<\beta\},
\]
故 $\{M\mid \alpha<f(M)<\beta\}$ 亦为开集.

（$2'\Rightarrow1$）$\forall M_0\in D$，令
\[
  \alpha=f(M_0)-\varepsilon,\quad \beta=f(M_0)+\varepsilon,
\]
则 $M_0\in G=\{M\mid \alpha<f(M)<\beta\}$. 因 $G$ 为开集，故存在 $M_0$ 的邻域 $V$，使得 $M_0\in V\subset G$. 即当 $M\in V$ 时，有
\[
  f(M_0)-\varepsilon<f(M)<f(M_0)+\varepsilon,
\]
亦即 $|f(M)-f(M_0)|<\varepsilon$.

\noindent\textbf{注}\quad 条件 $2'$）实际就是：任意开区间 $(\alpha,\beta)\subset\mathbb{R}$ 的逆像
\[
  f^{-1}[(\alpha,\beta)]=\{M\mid \alpha<f(M)<\beta\}
\]
恒为开集.

上述结论很容易推广到向量函数的情况. 由此我们看到拓扑学中连续性概念的渊源.

\noindent\textbf{d.\quad 连续函数性质的应用}

\noindent\textbf{i）有界性的应用}

\noindent\textbf{$\star$例 6.1.26}\quad 设

1）$\varphi(x)$ 在 $(x_0-\delta,x_0)$（$\delta>0$）上具有连续的导数，并存在 $x_n\in(x_0-\delta,x_0)$（$n=1,2,\cdots$），使得 $x_n\to x_0,\varphi(x_n)\to y_0$（当 $n\to\infty$ 时）；

2）$f(x,y)$ 在有界闭区域 $G$ 上连续. 当 $(x_0,y_0)\in G$，且 $x\in(x_0-\delta,x_0)$ 时，有 $(x,\varphi(x))\in G$，$\varphi'(x)=F(x,\varphi(x))$.

试证：
\[
  \lim_{x\to x_0^-}\varphi(x)=y_0.
\]
（山东大学）

\noindent\textbf{分析}\quad 若 $\varphi(x)\nrightarrow y_0$（当 $x\to x_0^-$ 时），那么由条件 1），当 $x\to x_0^-$ 时，$\varphi(x)$ 只能无限振动，如图 6.1.2 所示. $x$ 充分接近 $x_0^-$ 时，$\varphi'(x)$ 无界，从而 $f(x,\varphi(x))=\varphi'(x)$ 亦无界，与 $f(x,y)$ 在 $G$ 上连续相矛盾（因为在有界闭区域上连续必有界）.

\begin{center}
  \includegraphics[width=0.45\textwidth]{6-1-2.png}
\end{center}

\noindent\textbf{证}\quad （反证法）设 $\varphi(x)\nrightarrow y_0$（当 $x\to x_0^-$ 时），则 $\exists\varepsilon_0>0$ 及数列：$x_n'\in G,x_n'<x_0,x_n'\to x_0^-$（当 $n\to\infty$ 时），使得
\[
  |\varphi(x_n')-y_0|\ge\varepsilon_0. \tag{1}
\]
另一方面，已知 $x_n\to x_0,\varphi(x_n)\to y_0$（当 $n\to\infty$ 时），所以 $\forall K\in\mathbb{N}$，$\exists N_1>0$，使得 $n>N_1$ 时，
\[
  |x_n-x_0|<\frac1K,\quad
  |\varphi(x_n)-y_0|<\frac{\varepsilon_0}{2}. \tag{2}
\]
由 $x_n'\to x_0^-$（当 $n\to\infty$ 时）知：$\exists N_2>0$，$n>N_2$ 时，$|x_n'-x_0|<\dfrac1K$. 取 $N=\max\{N_1,N_2\}$，则当 $n>N$ 时有
\[
  |x_n'-x_0|<\frac1K,\quad |x_n-x_0|<\frac1K. \tag{3}
\]
于是由式 (1) 和 (2)，有
\[
  |\varphi(x_n')-\varphi(x_n)|
  \ge ||\varphi(x_n')-y_0|-|y_0-\varphi(x_n)||>\varepsilon_0-\frac{\varepsilon_0}{2}
  =\frac{\varepsilon_0}{2}. \tag{4}
\]
由式 (3) 有
\[
  |x_n-x_n'|\le |x_n-x_0|+|x_0-x_n'|<\frac2K. \tag{5}
\]
对于 $\dfrac1K$，将如此的 $n$ 固定，利用微分中值定理，$\exists\xi_{n_k}$ 于 $x_n,x_n'$ 之间，使得
\[
  \varphi(x_n')-\varphi(x_n)=\varphi'(\xi_{n_k})(x_n'-x_n).
\]
从而由 (4) 和 (5)，
\[
  |\varphi'(\xi_{n_k})|
  =\left|\frac{\varphi(x_n')-\varphi(x_n)}{x_n'-x_n}\right|
  \ge \frac{\varepsilon_0/2}{2/K}
  =\frac{K\varepsilon_0}{4}.
\]
已知 $f(x,\varphi(x))=\varphi'(x)$，根据 $K$ 的任意性，如此我们证明了 $\varphi'(x)=f(x,\varphi(x))$ 无界，与 $f$ 在有界闭区域 $G$ 上连续必有界矛盾.

\noindent\textbf{ii）介值定理的应用}

\noindent\textbf{例 6.1.27}\quad 证明：不存在由闭区间到圆周上的一对一连续对应.

\noindent\textbf{证}\quad （反证法）设
\[
  K=\{(r,\theta)\mid 0\le\theta<2\pi\}
\]
是某个圆，$[a,b]$ 是某个闭区间. 若存在 $[a,b]$ 至 $K$ 的一对一连续对应，则 $[a,b]$ 也就连续一对一地对应于区间 $[0,2\pi)$. 因此 $a,b$ 两点至少有一个对应于 $[0,2\pi)$ 的内点，例如记作 $a$. 记 $f$ 为此对应关系，则有 $0<f(a)<2\pi$. 取 $\theta_1,\theta_2$ 使得 $0<\theta_1<f(a)<\theta_2<2\pi$. 记 $x_1=f^{-1}(\theta_1),x_2=f^{-1}(\theta_2)$，则 $x_1,x_2\in(a,b]$，$0<f(x_1)<f(a)<f(x_2)<2\pi$. 利用介值定理，$\exists\xi$ 在 $x_1$ 与 $x_2$ 之间使得 $f(\xi)=f(a)$，与一对一矛盾. 证毕.

\noindent\textbf{iii）确界技术及有关原理}

\noindent\textbf{要点}\quad 1）若 $f(M)\le B$（$\forall M\in J$），则
\[
  \sup_J f(M)\le B;
\]

2）若 $f(M)\ge A$（$\forall M\in J$），则
\[
  \inf_J f(M)\ge A;
\]

3）若 $f(M)\le g(M)$（$\forall M\in J$），则
\[
  \sup_J f(M)\le \sup_J g(M),\quad
  \inf_J f(M)\le \inf_J g(M);
\]

4）有界闭区域上的连续函数，必达上、下确界.

\noindent\textbf{$\star$例 6.1.28}\quad 设二元函数 $f(x,y)$ 在正方形区域 $[0,1]\times[0,1]$ 上连续，记 $J=[0,1]$.

1）试比较
\[
  \inf_{y\in J}\sup_{x\in J}f(x,y)
\]
与
\[
  \sup_{x\in J}\inf_{y\in J}f(x,y)
\]
的大小并证明之；

2）给出并证明使等式
\[
  \inf_{y\in J}\sup_{x\in J}f(x,y)
  =\sup_{x\in J}\inf_{y\in J}f(x,y)
\tag{1}
\]
成立的（你认为最好的）充分条件.（浙江大学）

\noindent\textbf{证}\quad $1^\circ$ 对任意固定的 $y\in J$，有
\[
  \sup_{x\in J}f(x,y)\ge f(x,y)\ge \inf_{y\in J}f(x,y)
  \quad(\forall x\in J), \tag{2}
\]
所以
\[
  \sup_{x\in J}f(x,y)\ge \sup_{x\in J}\inf_{y\in J}f(x,y).
\]
由 $y$ 的任意性知
\[
  \inf_{y\in J}\sup_{x\in J}f(x,y)
  \ge \sup_{x\in J}\inf_{y\in J}f(x,y).
\]

$2^\circ$ 若 $f(x,y)\equiv$ 常数，则等式 (1) 明显成立. 但这种情况太平凡. 一种简单而有意义的条件是 $f(x,y)$ 关于其中某一变量单调. 下面以 $f(x,y)$ 对变量 $x$ 单增进行证明.

上面已证明了式 (1) 的（左端）$\ge$（右端），现证（左端）$\le$（右端）. 因 $f(x,y)$ 对 $x$ 单调增加，所以固定 $y\in J$，有
\[
  \sup_{x\in J}f(x,y)=f(1,y).
\]
由于 $f(1,y)$ 关于 $y$ 在区间 $J=[0,1]$ 上连续，因此 $\exists y_0\in J$ 使得
\[
  f(1,y_0)=\inf_{y\in J}f(1,y)
  =\inf_{y\in J}\sup_{x\in J}f(x,y)=\text{左端}.
\]
但
\[
  f(1,y_0)=\inf_{y\in J}f(1,y)
  \le \sup_{x\in J}\left[\inf_{y\in J}f(x,y)\right]
  =\text{右端},
\]
故（左端）$\le$（右端）.

\noindent\textbf{e.\quad 一致连续性}

\noindent\textbf{*$\star$例 6.1.29}\quad 设 $\mathbb{R}^n$ 为 $n$ 维欧氏空间，$A$ 是 $\mathbb{R}^n$ 的非空子集，定义 $x$ 到 $A$ 的距离为
\[
  f_A(x)=\inf_{y\in A}\rho(x,y)=\rho(x,A).
\]
证明 $f_A(x)$ 是 $\mathbb{R}^n$ 上的一致连续函数.（南京大学）

\noindent\textbf{证}\quad $\forall x_1,x_2\in\mathbb{R}^n,\forall y\in A$，有
\[
  \rho(x_1,y)\le \rho(x_1,x_2)+\rho(x_2,y).
\]
由此
\[
  \inf_{y\in A}\rho(x_1,y)\le \rho(x_1,x_2)+\rho(x_2,y)
  \quad(\forall y\in A).
\]
从而
\[
  \inf_{y\in A}\rho(x_1,y)\le \rho(x_1,x_2)+\inf_{y\in A}\rho(x_2,y).
\]
所以
\[
  \inf_{y\in A}\rho(x_1,y)-\inf_{y\in A}\rho(x_2,y)\le \rho(x_1,x_2).
\]
$x_1,x_2$ 互换得
\[
  \inf_{y\in A}\rho(x_2,y)-\inf_{y\in A}\rho(x_1,y)\le \rho(x_2,x_1)=\rho(x_1,x_2).
\]
因此
\[
  \left|\inf_{y\in A}\rho(x_2,y)-\inf_{y\in A}\rho(x_1,y)\right|\le \rho(x_1,x_2).
\]
故 $\forall\varepsilon>0$，取 $\delta=\varepsilon$，当 $\rho(x_1,x_2)<\delta$ 时，有
\[
\begin{aligned}
|\rho(x_2,A)-\rho(x_1,A)|
&=\left|\inf_{y\in A}\rho(x_2,y)-\inf_{y\in A}\rho(x_1,y)\right|\\
&\le \rho(x_1,x_2)<\varepsilon,
\end{aligned}
\]
即 $\rho(x,A)$ 在 $\mathbb{R}^n$ 上一致连续.

\noindent\textbf{例 6.1.30}\quad 设 $D\subset A\subset\mathbb{R}^n$，$D$ 为有界闭区域，$A$ 为开集. 试证，存在连续函数 $f:\mathbb{R}^n\to[0,1]$，使得
\[
  f(x)=
  \begin{cases}
  1, & x\in D,\\
  0, & x\notin A.
  \end{cases}
\]
（西南大学）

\noindent\textbf{提示}\quad 利用上例的结果与符号，可取
\[
  f(x)=\frac{\rho(x,A^c)}{\rho(x,A^c)+\rho(x,D)},
\]
其中 $A^c=\mathbb{R}^n-A$ 表示 $A$ 的补集.

一元函数关于一致连续的一些结果，也可以推广到多元. 如

\noindent\textbf{*例 6.1.31}\quad 设 $\varphi(x)$ 在 $\mathbb{R}^n$ 里一致连续，$f(x)$ 在 $\mathbb{R}^n$ 里连续，且
\[
  \lim_{r\to+\infty}[f(x)-\varphi(x)]=0
\]
（这里 $r=|x|=\rho(x,\theta)$，$\theta$ 表示 $\mathbb{R}^n$ 中的原点）. 试证 $f(x)$ 在 $\mathbb{R}^n$ 里一致连续.

\noindent\textbf{提示}\quad 参看例 2.2.8.

\noindent\textbf{*例 6.1.32}\quad 设 $f(x)$ 在 $\mathbb{R}^n$ 中的有界开区域 $D$ 内连续，试证 $f(x)$ 在 $D$ 内一致连续的充要条件是 $\forall x_0\in\partial D$，
\[
  \lim_{\substack{x\to x_0\\ x\in D}}f(x)
\]
存在（这里 $\partial D$ 表示 $D$ 的全体边界点组成之集合）.

\noindent\textbf{证}\quad 必要性. 因为 $f$ 在 $D$ 上一致连续，所以 $\forall\varepsilon>0$，$\exists\delta>0$，当 $x,y\in D,\rho(x,y)<\delta$ 时有 $|f(x)-f(y)|<\varepsilon$. 设 $x_0\in\partial D$，因 $D$ 为开区域，故 $x_0$ 必为 $D$ 的聚点. 设 $x_n\in D,x_n\to x_0$（当 $n\to\infty$ 时）为任一趋向 $x_0$ 的序列. 则对上述 $\delta>0$，$\exists N>0$，当 $n,m>N$ 时，有 $\rho(x_n,x_m)<\delta$，从而 $|f(x_n)-f(x_m)|<\varepsilon$. 根据 Cauchy 准则，知 $\displaystyle\lim_{n\to\infty}f(x_n)$ 存在. 由于 $x_n\to x_0$ 是任取的，由 Heine 定理，
\[
  \lim_{\substack{x\to x_0\\ x\in D}}f(x)
\]
存在.

充分性. 已知 $\forall x_0\in\partial D$，
\[
  \lim_{\substack{x\to x_0\\ x\in D}}f(x)
\]
存在. 今补充定义
\[
  f(x_0)=\lim_{\substack{x\to x_0\\ x\in D}}f(x).
\]
于是 $\forall\varepsilon>0$，$\exists\delta>0$，当 $x\in D,\rho(x,x_0)<\delta$ 时，有 $|f(x)-f(x_0)|<\varepsilon$. 对于 $x_1\in\partial D$，若 $\rho(x_1,x_0)<\delta$，则在上式中令 $x\to x_1$，取极限可得 $|f(x_1)-f(x_0)|\le\varepsilon$.

故 $\forall x\in\overline D=D\cup(\partial D)$，当 $\rho(x,x_0)<\delta$ 时，恒有 $|f(x)-f(x_0)|\le\varepsilon$. 因 $f$ 在 $D$ 内连续，如此证明了 $f$ 在 $\overline D$ 上连续，进而知 $f$ 在 $\overline D$ 上（从而也就在 $D$ 内）一致连续.

\noindent\textbf{new *例 6.1.33}\quad 设

1）$D$ 是 $\mathbb{R}^2$ 里的凸区域（意指：若 $P_1,P_2\in D$ 是 $D$ 内的任意两点，则直线段 $P_1P_2$ 全落在 $D$ 内）；

2）$T$ 是 $D\to\mathbb{R}^2$ 的映射，有连续导数，且 Jacobi 行列式不为零.

试证：$T$ 为单射.（北京大学）

\noindent\textbf{证}\quad （用反证法.）记 $T:(X,Y)=(f(x,y),g(x,y))$（$\forall(x,y)\in D$）. 若 $T$ 不是单射，即至少有两个不同点 $P_1(x_1,y_1),P_2(x_2,y_2)\in D$，使得 $T(P_1)=T(P_2)$，亦即
\[
  f(x_1,y_1)=f(x_2,y_2),\quad g(x_1,y_1)=g(x_2,y_2). \tag{1}
\]
根据条件 1），记
\[
  P_t=(x_1+t\Delta x,y_1+t\Delta y)\in D
\]
（其中 $\Delta x=x_2-x_1,\Delta y=y_2-y_1,t\in[0,1]$），$(X,Y)|_{P_t}=(f(P_t),g(P_t))$，$t\in[0,1]$. 作辅助函数
\[
  W(t)=f(P_t)\Delta x+g(P_t)\Delta y,\quad t\in[0,1]. \tag{2}
\]
根据条件 2），$W(t)$ 有连续导数，由式 (1) 知：$W(0)=W(1)$. 应用 Lagrange 定理，$\exists\xi\in(0,1)$，记 $P_\xi=(x_1+\xi\Delta x,y_1+\xi\Delta y)$，知
\[
\begin{aligned}
0&=W(1)-W(0)=W'(\xi)\cdot1\\
&=(f_x'(P_\xi)\Delta x+f_y'(P_\xi)\Delta y)\Delta x
+(g_x'(P_\xi)\Delta x+g_y'(P_\xi)\Delta y)\Delta y\\
&=f_x'(P_\xi)\Delta x\Delta x+f_y'(P_\xi)\Delta y\Delta x
+g_x'(P_\xi)\Delta x\Delta y+g_y'(P_\xi)\Delta y\Delta y\\
&=(\Delta x,\Delta y)J
\begin{pmatrix}
\Delta x\\
\Delta y
\end{pmatrix},
\end{aligned}
\tag{3}
\]
其中
\[
  J=
  \begin{vmatrix}
  f_x'(P_\xi)&f_y'(P_\xi)\\
  g_x'(P_\xi)&g_y'(P_\xi)
  \end{vmatrix}.
\]
因为偏导数连续，Jacobi 行列式不为零，那么式 (3) 只有零解，得 $P_1=P_2$，矛盾！问题获证.

\noindent\textbf{new 练习（推广到 $\mathbb{R}^n$）}\quad 设 $D\subset\mathbb{R}^n$ 是凸区域，连续映射 $T:D\to\mathbb{R}^n$ 有连续导数，Jacobi 行列式不为零，则 $T$ 必为单射.（北京大学）

\noindent\textbf{提示}\quad 这时，若 $T$ 不是单射，则至少存在不同的两点 $x_k=(x_{k1},x_{k2},\cdots,x_{kn})\in D$（$k=1,2$），使得
\[
  T(x_1)=(f_1(x_1),f_2(x_1),\cdots,f_n(x_1))
  =T(x_2)=(f_1(x_2),f_2(x_2),\cdots,f_n(x_2)).
\]
作辅助函数
\[
  W(t)=\sum_{k=1}^n f_k(x_1+t(x_2-x_1))(x_{2k}-x_{1k}),\quad t\in[0,1].
\]

\noindent\textbf{证}\quad （用 Lagrange 定理.）$\exists\xi\in(0,1)$，使得
\[
\begin{aligned}
0&=W(1)-W(0)=W'(\xi)\cdot1\\
&=\sum_{k=1}^n\left[(x_{2k}-x_{1k})\sum_{j=1}^n
\left(\frac{\partial}{\partial x_j}f_k(x_\xi)\right)(x_{2j}-x_{1j})\right]\\
&=\sum_{k,j=1}^n\left(\frac{\partial}{\partial x_j}f_k(x_\xi)\right)\Delta x_{1k}\Delta x_{1j};
\end{aligned}
\tag{1}
\]
其中
\[
\begin{aligned}
x_\xi&=x_1+\xi(x_2-x_1)\\
&=(x_{11}+\xi(x_{21}-x_{11}),x_{12}+\xi(x_{22}-x_{12}),\cdots,x_{1n}+\xi(x_{2n}-x_{1n}))\\
&=(x_{11}+\xi\Delta x_{11},x_{12}+\xi\Delta x_{12},\cdots,x_{1n}+\xi\Delta x_{1n}),
\end{aligned}
\]
\[
  \Delta x_1=x_2-x_1=(x_{21}-x_{11},x_{22}-x_{12},\cdots,x_{2n}-x_{1n})
  =(\Delta x_{11},\Delta x_{12},\cdots,\Delta x_{1n}).
\]
（二次型）式 (1) 可改写为乘法形式：
\[
  \sum_{k,j=1}^n
  \left(\frac{\partial}{\partial x_j}f_k(x_\xi)\right)\Delta x_{1k}\Delta x_{1j}
  =\det(\Delta x_1J\Delta x_1^T)=0. \tag{2}
\]
其中
\[
  J=\frac{\partial}{\partial x_j}f_k(x_\xi)
  =
  \left.\frac{\partial(f_1,f_2,\cdots,f_n)}{\partial(x_1,x_2,\cdots,x_n)}\right|_{x_\xi}
\]
是 Jacobi 矩阵，$\Delta x_1^T$ 表示向量 $\Delta x_1$ 的转置. 式 (2) 表明，若 $|J|\ne0$，那么例 6.1.33 中的式 (3) 只有零解，得 $x_1=x_2$，矛盾.（证毕.）

\section*{单元练习 6.1}

\subsection*{*$m$ 维欧氏空间}

\noindent\textbf{6.1.1}\quad 设 $l\in\mathbb{R}^m,|l|=1$，$\theta_i$ 表示 $l$ 与坐标向量 $e_i$（$i=1,2,\cdots,m$）的夹角，试证：
\[
  l=(\cos\theta_1,\cos\theta_2,\cdots,\cos\theta_m).
\]

\noindent\textbf{6.1.2}\quad $x,y\in\mathbb{R}^m$，$\theta$ 表示 $x,y$ 的夹角，试证：

1）余弦公式：
\[
  |x-y|^2=|x|^2+|y|^2-2|x||y|\cos\theta;
\]

2）勾股弦定理：$x$ 与 $y$ 正交时，$|x+y|^2=|x|^2+|y|^2$.

\noindent\textbf{6.1.3}\quad $G_1,G_2\subseteq\mathbb{R}^m$ 是任意两开集，$G_1\cap G_2=\varnothing$，试证：
\[
  G_1\cap\overline{G_2}=\varnothing.
\]

\noindent\textbf{提示}\quad 若 $G_1\cap\overline{G_2}\ne\varnothing$，则 $\exists x\in G_1,x\notin G_2$，但 $x\in\overline{G_2}$.

再提示\quad 于是 $x\in\partial G_2$，进而 $\exists\delta>0$，使得 $U_0(x,\delta)\subseteq G_1$，但 $U_0(x,\delta)$ 内含有属于 $G_2$ 中的点：$x_1\in G_1\cap G_2$，矛盾.

\noindent\textbf{6.1.4}\quad 设 $E\subseteq\mathbb{R}^m$ 为任意集合，$E'$ 表示 $E$ 的全体聚点组成的集合，称为 $E$ 的导集，试证 $E'$ 为闭集.

\noindent\textbf{提示}\quad 只需证明 $\forall x\in\overline{E'}$（$E'$ 是导集 $E'$ 的闭包），$\exists x_n^*\in E$（$n=1,2,\cdots$），使得 $x_n^*\to x$（$n\to\infty$）即可.

再提示\quad 因 $x\in\overline{E'}$，故存在彼此不等的 $y_n\in E'$（$n=1,2,\cdots$）使得 $y_n\to x$（$n\to\infty$）. 又由 $y_n\in E'$ 知 $\exists x_n^*\in E$，使得 $|x_n^*-y_n|<\dfrac1{2^n}$. 于是
\[
  |x_n^*-x|\le |x_n^*-y_n|+|y_n-x|\le\frac1{2^n}+|y_n-x|\to0
\]
（当 $n\to\infty$ 时）.

\noindent\textbf{6.1.5}\quad 设 $A,B\subseteq\mathbb{R}^m$ 为开集，$A\cap B=\varnothing$，试证：
\[
  \partial(A\cup B)=\partial A\cup\partial B.
\]

\noindent\textbf{提示}\quad i）$x\in\partial(A\cup B)$，知存在互异的 $x_n\in A\cup B,x_n\to x$，故 $A,B$ 中至少有一个含 $x_n$ 的无穷多项；又存在互异的 $y_n\notin A\cup B,y_n\to x$，而 $y_n\notin A$，且 $y_n\notin B$，所以 $x\in\partial A$ 或 $x\in\partial B$，从而 $x\in\partial A\cup\partial B$.

ii）$x\in\partial A\cup\partial B$，即 $x\in\partial A$ 或 $x\in\partial B$，从而 $x$ 不可能是 $A\cup B$ 之内点（因 $A,B$ 为开集，且 $A\cap B=\varnothing$），$x$ 也不可能是 $A\cup B$ 的外点（否则与 $x\in\partial A$ 或 $x\in\partial B$ 矛盾）. 因此，$x\in\partial(A\cup B)$.

\noindent\textbf{6.1.6}\quad 设 $A,B\subseteq\mathbb{R}^m$ 为有界闭集，$A\cap B=\varnothing$，试证：$\exists$ 开集 $W,V$，使得 $A\subseteq W,B\subseteq V$，且 $W\cap V=\varnothing$.

\noindent\textbf{提示}\quad 记
\[
  d=\rho(A,B)=\inf\{|x-y|\mid x\in A,y\in B\}>0,
\]
\[
  W=\left\{x\mid x\in\mathbb{R}^m,\rho(x,A)<\frac d3\right\},\quad
  V=\left\{x\mid x\in\mathbb{R}^m,\rho(x,B)<\frac d3\right\}
\]
即可（参见例 6.1.4 和 6.1.5）.

\noindent\textbf{$\star$6.1.7}\quad 设 $S\subset\mathbb{R}^2$，$P_0(x_0,y_0)$ 为 $S$ 的内点，$P_1(x_1,y_1)$ 为 $S$ 的外点，证明：直线段 $P_0P_1$ 必与 $S$ 的边界 $\partial S$ 至少有一交点.（华东师范大学）

\noindent\textbf{提示}\quad 可对线段 $P_0P_1$ 进行二等分法，在 $P_0P_1$ 上找出 $S$ 的界点.

再提示\quad $P_0$ 为内点，$P_1$ 为外点，$P_0P_1$ 二等分，若中点为 $S$ 的界点，则问题已解决；否则两半之中，至少有一半其两端点为一内（点）一外（点）. 将此半段再二等分，照此办理，每次二等分后，中点若为界点则问题已解决，否则继续再分. 把 $P_0P_1$ 所在的直线看成是数轴，这就构成了数轴上的区间套 $\{[A_n,B_n]\}$，$A_n,B_n$ 分别为 $S$ 的内点和外点，且 $[A_n,B_n]$ 有唯一的公共点 $\xi\in[A_n,B_n]$（$n=1,2,\cdots$），这里 $A_1=P_0,B_1=P_1$；
\[
  |\xi-A_n|,\ |\xi-B_n|\le |B_n-A_n|
  =\frac1{2^{n-1}}|P_1-P_0|\to0
\]
（当 $n\to\infty$ 时）. $\forall\varepsilon>0$，$\exists N>0$，当 $n>N$ 时 $A_n,B_n\in U(\xi,\varepsilon)$. 即 $\xi$ 的任一 $\varepsilon$ 邻域里既有 $S$ 的内点又有 $S$ 的外点，故 $\xi$ 是 $S$ 的边界点.

\subsection*{$\star$多元极限}

\noindent\textbf{$\star$6.1.8}\quad 求极限

1）
\[
  \lim_{\substack{x\to+\infty\\ y\to+\infty}}(x^2+y^2)e^{-(x+y)}
\]
（北京航空航天大学）；

2）
\[
  \lim_{\substack{x\to0\\ y\to0}}\frac{xy}{\sqrt{x+y+1}-1}
\]
（若极限不存在，说明理由）；（西北轻工业学院）

3）
\[
  \lim_{(x,y)\to(0,0)}\frac{\arctan(x^3+y^3)}{x^2+y^2}.
\]

\noindent\textbf{提示}\quad 1）原式
\[
  =\lim_{\substack{x\to+\infty\\ y\to+\infty}}
  \left(\frac{x^2}{e^x}\cdot\frac1{e^y}
  +\frac{y^2}{e^y}\cdot\frac1{e^x}\right)=0.
\]

2）若原式极限存在，则
\[
  \frac{xy}{x+y}
  =\frac{xy}{\sqrt{x+y+1}-1}\cdot\frac1{\sqrt{x+y+1}+1}
\]
的极限也应存在，但用特殊路径法（例 6.1.11 中 2））已知该极限不存在，矛盾.

3）原式
\[
  =\lim_{\substack{x\to0\\ y\to0}}
  \left[
  \frac{\arctan(x^3+y^3)}{x^3+y^3}\cdot
  \frac{x^3+y^3}{x^2+y^2}
  \right]=1\cdot0=0.
\]

\noindent\textbf{$\star$6.1.9}\quad 设
\[
f(x,y)=
\begin{cases}
x\sin\dfrac1y+y\sin\dfrac1x, & xy\ne0,\\
0, & xy=0.
\end{cases}
\]
试讨论下面三种极限：

1）$\displaystyle\lim_{(x,y)\to(0,0)}f(x,y)$；

2）$\displaystyle\lim_{x\to0}\lim_{y\to0}f(x,y)$；

3）$\displaystyle\lim_{y\to0}\lim_{x\to0}f(x,y)$.（南京工业大学）

\noindent\textbf{提示}\quad 1）$0\le |f(x,y)|\le |x|+|y|\to0$.

2），3）内层极限不存在，从而累次极限不存在.

\noindent\textbf{$\star$6.1.10}\quad 设 $f(x,y)$ 为二元函数，在 $(x_0,y_0)$ 附近有定义，试讨论二重极限
\[
  \lim_{\substack{x\to x_0\\ y\to y_0}}f(x,y)
\]
与累次极限
\[
  \lim_{x\to x_0}\lim_{y\to y_0}f(x,y)
\]
之间的关系.（浙江大学）

\noindent\textbf{提示}\quad （1）上题已证明二重极限存在，两个累次极限可以不存在.

（2）例 6.1.11 中 4）表明：两个累次极限存在可以不相等.

（3）例 6.1.11 中 1）至 4）还表明：二重极限不存在，两个累次极限仍可能存在. 这时两个累次极限可以相等（如 1），2），3））也可以不等（如 4））.

（4）若
\[
  \lim_{\substack{x\to x_0\\ y\to y_0}}f(x,y)\stackrel{\text{存在}}{=}A,
\]
且 $\displaystyle\lim_{x\to x_0}\lim_{y\to y_0}f(x,y)$ 的内层极限 $\displaystyle\lim_{y\to y_0}f(x,y)$ 在 $x_0$ 的某个 $\delta_1$ 邻域里存在（$\delta_1>0$），则
\[
  \lim_{x\to x_0}\lim_{y\to y_0}f(x,y)\stackrel{\text{存在}}{=}A.
\]
（另一累次极限亦然.）

再提示\quad （证 (4)）因 $\forall\varepsilon>0$，$\exists\delta>0$，（取 $\delta<\delta_1$），当 $|x-x_0|<\delta,|y-y_0|<\delta$ 时，
\[
  A-\varepsilon<f(x,y)<A+\varepsilon,
\]
在此不等式里令 $y\to y_0$ 取极限，记 $\displaystyle\lim_{y\to y_0}f(x,y)=g(x)$，得
\[
  A-\varepsilon\le g(x)\le A+\varepsilon\quad(\forall x:|x-x_0|<\delta).
\]
此即表明
\[
  A=\lim_{x\to x_0}g(x)
  =\lim_{x\to x_0}\lim_{y\to y_0}f(x,y).
\]
证毕.

\subsection*{*多元连续函数}

\noindent\textbf{6.1.11}\quad 设 $f(x,y)$ 在 $G=\{(x,y):x^2+y^2<1\}$ 上有定义，若 1）$f(x,0)$ 在点 $x=0$ 处连续；2）$f_y'(x,y)$ 在 $G$ 上有界，证明 $f(x,y)$ 在 $(0,0)$ 处连续.（北京大学）

\noindent\textbf{提示}\quad 可参看例 6.1.23. 并注意由条件 2）可推出 $f(x,y)$ 对 $y$ 满足 Lipschitz 条件，关于 $x$ 一致.

再提示\quad 由条件 1）知：$\forall\varepsilon>0$，$\exists\delta_1>0$，当 $|x|<\delta_1$ 时，$|f(x,0)-f(0,0)|<\dfrac{\varepsilon}{2}$.

由条件 2）知：$\exists M>0$，使 $|f_y'(x,y)|\le M$，$(x,y)\in G$，
\[
  |f(x,y)-f(x,0)|=|f_y'(x,\xi)|\,|y-0|\le M|y|.
\]
取 $\delta=\min\left\{\dfrac{\varepsilon}{2M},\delta_1\right\}$，则当 $|x|<\delta,|y|<\delta$ 时，有
\[
  |f(x,y)-f(0,0)|\le |f(x,y)-f(x,0)|+|f(x,0)-f(0,0)|
  <\frac{\varepsilon}{2}+\frac{\varepsilon}{2}=\varepsilon.
\]

\noindent\textbf{6.1.12}\quad 设 $f(x,y)$ 在有界闭区域 $\overline D\subseteq\mathbb{R}^2$ 上连续，其值域为 $\mathbb{R}$，试证：$\forall\{u_n\}_{n=1}^\infty\subset\mathbb{R}$，$\exists$ 收敛子列 $\{u_{n_k}\}$ 及点 $(x_0,y_0)\in\overline D$，使得
\[
  \lim_{k\to\infty}u_{n_k}=f(x_0,y_0).
\]
（华中师范大学）

\noindent\textbf{提示}\quad $\forall u_n$，$\exists(x_n,y_n)\in\overline D$，使 $f(x_n,y_n)=u_n$，再利用致密性原理.

再提示\quad $\exists\{(x_n,y_n)\}$ 的收敛子列 $\{(x_{n_k},y_{n_k})\}$，记
\[
  \lim_{k\to\infty}(x_{n_k},y_{n_k})=(x_0,y_0)\in\overline D,
\]
则由 $f$ 的连续性知 $u_{n_k}=f(x_{n_k},y_{n_k})\to f(x_0,y_0)$（当 $k\to\infty$ 时）.

\noindent\textbf{$\star$6.1.13}\quad 设 $f(x,y)$ 在矩形 $D:-a\le x\le a,-b\le y\le b$（$a>0,b>0$）上分别是 $x$ 和 $y$ 的连续函数，而且 $f(0,0)=0$. 试证：若当 $x$ 固定时，$f(x,y)$ 是 $y$ 的严格递减函数，则存在 $\delta>0$，使对每个 $x\in(-\delta,\delta)$ 有 $y\in(-b,b)$ 满足 $f(x,y)=0$.（西南大学）

\noindent\textbf{提示}\quad 参看例 6.1.22.

再提示\quad $f(0,0)=0$，$f(x,y)$ 对 $y$ 严格递减 $\Rightarrow f(0,b)<0$，以及 $f(0,-b)>0$，$f(x,y)$ 关于 $x$ 连续 $\Rightarrow\exists\delta>0,\delta<a$，$-\delta<x<\delta$ 时，
\[
  f(x,b)<0,\quad f(x,-b)>0.
\]
$f$ 关于 $y$ 连续 $\Rightarrow \forall x\in(-\delta,\delta)$，$\exists y\in(-b,b)$ 使 $f(x,y)=0$.

\noindent\textbf{$\star$6.1.14}\quad 设 $u=f(x,y,z)$ 在闭立方体 $a\le x\le b,a\le y\le b,a\le z\le b$ 上连续，令
\[
  \varphi(x)=\max_{a\le y\le b}\left\{\min_{a\le z\le b}f(x,y,z)\right\},
\]
试证：$\varphi(x)$ 在 $[a,b]$ 上连续.（辽宁师范大学）

\noindent\textbf{提示}\quad 参看例 6.1.19.

再提示\quad $f(x,y,z)$ 在 $[a,b;a,b;a,b]$ 上连续，可用类似于例 6.1.19 的方法；$g(x,y)=\displaystyle\min_{a\le z\le b}f(x,y,z)$ 在 $[a,b;a,b]$ 上连续，故 $\varphi(x)=\displaystyle\max_{a\le y\le b}g(x,y)$ 在 $[a,b]$ 上连续.

\noindent\textbf{6.1.15}\quad 设 $f(x)$ 在 $\mathbb{R}^n$ 上连续，$x\ne\theta$（$\mathbb{R}^n$ 中的原点）时，$f(x)>0$，且 $\forall x\in\mathbb{R}^n$ 及 $c>0$ 有 $f(cx)=cf(x)$. 试证 $\exists a,b>0$，使得
\[
  a|x|\le f(x)\le b|x|\quad(\forall x\in\mathbb{R}^n).
\]

\noindent\textbf{提示}\quad 利用 $f$ 在单位球面上达到最大、最小值.

再提示\quad $S$ 表示 $\mathbb{R}^n$ 的单位球面，则 $\forall\theta\ne x\in\mathbb{R}^n$，$\dfrac{x}{|x|}\in S$，$f$ 是 $S$ 上的连续函数，在 $S$ 上有最大值 $b(>0)$，最小值 $a(>0)$，故
\[
  a\le f\left(\frac{x}{|x|}\right)\le b,
\]
亦即
\[
  a|x|\le f(x)\le b|x|\quad(\forall x\in\mathbb{R}^n).
\]

\noindent\textbf{6.1.16}\quad 设 $A$ 是 $n\times n$ 矩阵，$\det A\ne0$，试证：$\exists\alpha>0$，使得 $\forall x\in\mathbb{R}^n$，有 $|Ax|\le \alpha|x|$.

\noindent\textbf{提示}\quad 可用 Cauchy 不等式，或用反证法.

再提示\quad 记 $A=(a_{ij}),x=(x_1,x_2,\cdots,x_n)$，则
\[
  |Ax|=\left[\sum_{i=1}^n\left(\sum_{j=1}^n a_{ij}x_j\right)^2\right]^{1/2}
  \le \left(\sum_{i=1}^n\sum_{j=1}^n a_{ij}^2\cdot\sum_{j=1}^n x_j^2\right)^{1/2}
  =\left(\sum_{i=1}^n\sum_{j=1}^n a_{ij}^2\right)^{1/2}|x|\triangleq\alpha|x|.
\]

\noindent\textbf{6.1.17}\quad 设连续函数 $f:\mathbb{R}^n\to\mathbb{R}$ 满足如下三条件：

i）$\forall x\in\mathbb{R}^n,f(x)\ge0$，且 $x=\theta\Leftrightarrow f(x)=0$；

ii）$\forall\lambda\in\mathbb{R}$，有 $f(\lambda x)=|\lambda|f(x)$；

iii）$\forall x,y\in\mathbb{R}^n,f(x+y)\le f(x)+f(y)$.

试证：

1）$\exists M>0$，使得 $\forall x\in\mathbb{R}^n$，有 $f(x)\le M|x|$；

2）$f$ 满足 Lipschitz 条件，即 $\exists L>0$，使得 $|f(x)-f(y)|\le L\rho(x,y)$（$\forall x,y\in\mathbb{R}^n$）；

3）存在常数 $a>0$，使得 $\forall x\in\mathbb{R}^n$ 有 $a|x|\le f(x)$.

\noindent\textbf{提示}\quad 1），2）利用上面习题 6.1.15 可得.

再提示\quad （当 $x\ne y$ 时）$f(x)-f(y)=f(x-y+y)-f(y)\le f(x-y)\le |x-y|f\left(\dfrac{x-y}{|x-y|}\right)\le L|x-y|=L\rho(x,y)$，
其中 $L$ 为 $f$ 在单位球面 $S$ 上的界：$|f(x)|\le L$（$\forall x\in S$）.

同理，有
\[
  f(y)-f(x)\le L\rho(x,y).
\]
故 $|f(x)-f(y)|\le L\rho(x,y)$（$\forall x,y\in\mathbb{R}^n$）（此式是当 $x\ne y$ 时推出的，但 $x=y$ 时明显成立）.

\noindent\textbf{$\ast$6.1.18}\quad 设 $f:\mathbb{R}^n\to\mathbb{R}$ 为连续函数，试证：$E=\{x\mid f(x)=0\}$ 是 $\mathbb{R}^n$ 中的闭集.

\noindent\textbf{提示}\quad 设 $x_0$ 是 $E$ 的任一聚点，证 $x_0\in E$.

\noindent\textbf{$\ast$6.1.19}\quad 试对 $f:\mathbb{R}^n\to\mathbb{R}$ 写出例 2.1.7 对应结果，并给出证明.

\noindent\textbf{$\ast$6.1.20}\quad 试对 $f:\mathbb{R}^n\to\mathbb{R}^m$ 写出例 2.1.6 的相应结果，并给出证明.

\noindent\textbf{$\ast\ast$6.1.21}\quad 设 $f$ 在 $\mathbb{R}^n$ 中点 $x_0$ 的邻域里有界，记
\[
  M_f(x_0,\delta)=\sup\{f(x)\mid \rho(x,x_0)<\delta\},\quad
  m_f(x_0,\delta)=\inf\{f(x)\mid \rho(x,x_0)<\delta\},
\]
则极限 $\omega_f(x_0)=\lim_{\delta\to0+}[M_f(x_0,\delta)-m_f(x_0,\delta)]$ 存在，并称之为 $f$ 在 $x_0$ 处的振幅，试证：$f(x)$ 在 $x_0$ 处连续的充要条件是 $\omega_f(x_0)=0$.

\noindent\textbf{※6.1.22}\quad 设集合 $E\subseteq\mathbb{R}^n$ 是非闭的，函数 $f(x)$ 在 $E$ 上一致连续. 试证：$f$ 只能唯一地连续延拓到 $\overline E$ 上，使在 $\overline E$ 上一致连续.

\noindent\textbf{$\ast$6.1.23}\quad 设 $f(x)$ 在 $\mathbb{R}^n$ 上连续，$\lim_{x\to+\infty}f(x)$ 存在. 试证 $f$ 在 $\mathbb{R}^n$ 上一致连续.

\section*{§6.2\quad 多元函数的偏导数}

\noindent\textbf{导读}\quad 本节是基础性内容，既是教学重点，也是各类考试的热点，一至四段适合本书各类读者；第五段主要针对数学院系学生.

\subsection*{\ensuremath{\star}一、偏导数的计算}

\noindent\textbf{要点}\quad
\[
  f_x'(a,b)=\left[\frac{d}{dx}f(x,b)\right]_{x=a}.
\]

若 $f$ 在 $(a,b)$ 与 $(a,b)$ 的附近用不同的初等函数给出，则需用定义通过极限来计算.

\noindent\textbf{例 6.2.1}\quad 设
\[
  f(x,y)=
  \begin{cases}
  xy\dfrac{x^2-y^2}{x^2+y^2}, & (x,y)\ne(0,0),\\
  0, & (x,y)=(0,0),
  \end{cases}
\]
试求 $f_{xy}''(0,0)$ 与 $f_{yx}''(0,0)$.（北京师范大学）

\noindent\textbf{解}\quad
\[
  f_x'(x,y)=
  \begin{cases}
  y\dfrac{x^4+4x^2y^2-y^4}{(x^2+y^2)^2}, & (x,y)\ne(0,0),\\
  0, & (x,y)=(0,0).
  \end{cases}
\]
由此
\[
  f_x'(0,y)=
  \begin{cases}
  -y, & y\ne0,\\
  0, & y=0
  \end{cases}
  =-y.
\]
所以
\[
  f_{xy}''(0,0)=[f_x'(0,y)]_y'\bigm|_{y=0}=-1.
\]
类似可得 $f_{yx}''(0,0)=1$.

\noindent\textbf{注}\quad 该例说明混合偏导数一般来说与求导次序有关.

\subsection*{\ensuremath{\star}二、复合函数微分法（链式法则）}

\noindent\textbf{例 6.2.2}\quad 设 $u(x,y)$ 的所有二阶偏导数都连续，$\dfrac{\partial^2u}{\partial x^2}-\dfrac{\partial^2u}{\partial y^2}=0$.
\[
  u(x,2x)=x,\quad u_x'(x,2x)=x^2.
\]
试求 $u_{xx}''(x,2x),u_{xy}''(x,2x),u_{yy}''(x,2x)$.（南开大学）

\noindent\textbf{解}\quad 已知 $u(x,2x)=x$，对 $x$ 求导，可得
\[
  u_x'(x,2x)+2u_y'(x,2x)=1.
\]
已知 $u_x'(x,2x)=x^2$，因此有
\[
  u_y'(x,2x)=\frac{1-x^2}{2}.
\]
此式对 $x$ 求导得
\[
  u_{yx}''(x,2x)+2u_{yy}''(x,2x)=-x. \tag{1}
\]
$u_x'(x,2x)=x^2$ 对 $x$ 求导得
\[
  u_{xx}''(x,2x)+2u_{xy}''(x,2x)=2x. \tag{2}
\]
（1），（2）与已知关系 $u_{xx}''-u_{yy}''=0$ 联立，注意到二阶导数连续，故 $u_{xy}''=u_{yx}''$，得
\[
  u_{xx}''(x,2x)=u_{yy}''(x,2x)=-\frac43x,\quad
  u_{xy}''(x,2x)=u_{yx}''(x,2x)=\frac53x.
\]

\noindent\textbf{new 练习}\quad 已知 $f(x)$ 有二阶连续偏导数，且
\[
  f(x,x^2)=x^2,\quad f_x'(x,x^2)=4x,\quad f_{xx}''(x,x^2)=2x^2f_{yy}''(x,x^2),
\]
计算 $4f_y'(x,x^2)-f_{xx}''(x,x^2)+x^2f_{yy}''(x,x^2)$ 之值.

\noindent\textbf{解}\quad （因 $f''$ 连续）$f(x,x^2)=x^2$ 同时对 $x$ 求导，得
\[
  f_x'(x,x^2)+2xf_y'(x,x^2)=2x.
\]
代入已知的 $f_x'(x,x^2)=4x$，得
\[
  f_y'(x,x^2)=-1. \tag{1}
\]
再对 $y$ 求导，得
\[
  f_{xy}''(x,x^2)+2xf_{yy}''(x,x^2)=0. \tag{2}
\]
因 $f''$ 连续，故
\[
  f_{yx}''(x,x^2)+2xf_{yy}''(x,x^2)=0.
\]
由 $f_x'(x,x^2)=4x$ 可知
\[
  f_{xx}''(x,x^2)+2xf_{xy}''(x,x^2)=4. \tag{3}
\]
（3）$-2x\times$（2），得 $f_{xx}''(x,x^2)-4x^2f_{yy}''(x,x^2)=4$. 已知有
\[
  f_{xx}''(x,x^2)=2x^2f_{yy}''(x,x^2), \tag{4}
\]
故
\[
  f_{yy}''(x,x^2)=-\frac2{x^2}. \tag{5}
\]
代入（4）可得
\[
  f_{xx}''(x,x^2)=2x^2f_{yy}''(x,x^2)=2x^2\cdot\left(-\frac2{x^2}\right)=-4. \tag{6}
\]
用式（1），（5），（6）的结果代入，可得
\[
  4f_y'(x,x^2)-f_{xx}''(x,x^2)+x^2f_{yy}''(x,x^2)=-2.
\]

\noindent\textbf{例 6.2.3}\quad 设 $u=f(x+y+z,x^2+y^2+z^2)$，证明：
\[
  \frac{\partial^2u}{\partial x^2}+\frac{\partial^2u}{\partial y^2}+\frac{\partial^2u}{\partial z^2}
  =3f_{11}''+4(x+y+z)f_{12}''+4(x^2+y^2+z^2)f_{22}''+6f_2'.
\]
（注意 $f_1',f_2'$ 仍是中间变量 $x+y+z,x^2+y^2+z^2$ 的函数.）

\noindent\textbf{例 6.2.4}\quad 设 $x^2+xy+y^2=3$，证明：$y''=-\dfrac{18}{(x+2y)^3}$.

\noindent\textbf{注}\quad 求导后，$x,y$ 仍满足原方程.

\noindent\textbf{\ensuremath{\star}例 6.2.5}\quad 设 $u=f(z)$，其中 $z$ 为方程式
\[
  z=x+y\varphi(z) \tag{1}
\]
所定义的变量（为 $x$ 和 $y$ 的隐函数）. 证明 Lagrange 公式
\[
  \frac{\partial^nu}{\partial y^n}
  =\frac{\partial^{n-1}}{\partial x^{n-1}}\left[(\varphi(z))^n\frac{\partial u}{\partial x}\right]. \tag{2}
\]
（广西师范学院）

\noindent\textbf{分析}\quad $u$ 为 $z$ 的函数，$z$ 通过式（1）定义为 $x,y$ 的函数. 因此 $u$ 是 $x,y$ 的复合函数，依赖关系为
\[
  u\longrightarrow z\begin{cases}x\\y.\end{cases}
\]

\noindent\textbf{证}\quad （用数学归纳法）先证明 $n=1$ 的情况. 利用链式法则 $\dfrac{\partial u}{\partial y}=f'(z)z_y'$，由式（1），利用隐函数求导法则，可得 $z_y'=\dfrac{\varphi(z)}{1-y\varphi'(z)}$. 所以
\[
  \frac{\partial u}{\partial y}
  =\frac{f'(z)\varphi(z)}{1-y\varphi'(z)}. \tag{3}
\]
同理可得
\[
  \frac{\partial u}{\partial x}=f'(z)z_x'=\frac{f'(z)}{1-y\varphi'(z)}.
\]
与（3）比较知
\[
  \frac{\partial u}{\partial y}=\varphi(z)\frac{\partial u}{\partial x}. \tag{4}
\]
这就是（2）当 $n=1$ 时的情况.

现假定式（2）对 $n-1$ 的情况成立，求证明对 $n$ 的情况也成立. 事实上，
\[
\begin{aligned}
  \text{式（2）左端}
  &=\frac{\partial}{\partial y}\left(\frac{\partial^{n-1}u}{\partial y^{n-1}}\right)
  =\frac{\partial}{\partial y}\frac{\partial^{n-2}}{\partial x^{n-2}}
  \left\{[\varphi(z)]^{n-1}\frac{\partial u}{\partial x}\right\}\\
  &=\frac{\partial^{n-2}}{\partial x^{n-2}}\frac{\partial}{\partial y}
  \left[(\varphi(z))^{n-1}\frac{\partial u}{\partial x}\right]\\
  &=\frac{\partial^{n-2}}{\partial x^{n-2}}\left[
  (n-1)(\varphi(z))^{n-2}\varphi'(z)\frac{\partial z}{\partial y}\frac{\partial u}{\partial x}
  +(\varphi(z))^{n-1}\frac{\partial^2u}{\partial x\partial y}\right],
\end{aligned} \tag{5}
\]
式（2）右端
\[
\begin{aligned}
  &=\frac{\partial^{n-1}}{\partial x^{n-1}}\left[(\varphi(z))^n\frac{\partial u}{\partial x}\right]
  =\frac{\partial^{n-1}}{\partial x^{n-1}}\left[(\varphi(z))^{n-1}\frac{\partial u}{\partial y}\right]\\
  &=\frac{\partial^{n-2}}{\partial x^{n-2}}\left[
  (n-1)(\varphi(z))^{n-2}\varphi'(z)\frac{\partial z}{\partial x}\frac{\partial u}{\partial y}
  +(\varphi(z))^{n-1}\frac{\partial^2u}{\partial x\partial y}\right].
\end{aligned} \tag{6}
\]
由于 $z_y'=\varphi(z)z_x'$ 及式（4），有
\[
  \frac{\partial z}{\partial y}\frac{\partial u}{\partial x}
  =\frac{\partial z}{\partial x}\frac{\partial u}{\partial y},
\]
故（5），（6）两式右端相等，式（2）获证.

\noindent\textbf{\ensuremath{\star}例 6.2.6}\quad 若 $z=f(x,y)$ 满足
\[
  f(tx,ty)=t^kf(x,y)\quad(t>0), \tag{1}
\]
则称 $f(x,y)$ 为 $k$ 次齐次函数，试证下述关于齐次函数的 Euler 定理：设 $f(x,y)$ 可微，则 $f(x,y)$ 为 $k$ 次齐次函数的充要条件是
\[
  xf_x'(x,y)+yf_y'(x,y)=kf(x,y). \tag{2}
\]

\noindent\textbf{证}\quad 必要性：式（1）对 $t$ 求导，再令 $t=1$ 即得（2）.

充分性：\textbf{方法一}\quad 要证 $t>0$ 时（1）成立，即要证 $\varphi(t)\equiv\dfrac{f(tx,ty)}{t^k}=f(x,y)$. 因 $\varphi(1)=f(x,y)$，因此只要证明 $\varphi'(t)\equiv0$ 即可. 利用复合函数微分法及式（2），这是明显的.

\textbf{方法二}\quad （按必要性的证法倒退回去.）在式（2）中分别用 $tx,ty$ 代替自变量 $x,y$，得
\[
  txf_x'(tx,ty)+tyf_y'(tx,ty)=kf(tx,ty). \tag{3}
\]
记 $\varphi(t)=f(tx,ty)$，则（3）为
\[
  t\varphi'(t)=k\varphi,\quad \text{且}\quad \varphi(1)=f(x,y).
\]
由此解微分方程即得式（1）.

\noindent\textbf{\ensuremath{\star}例 6.2.7}\quad 设 $f(x,y)$ 为 $n$ 次齐次函数，且 $m$ 次可微，证明：
\[
  \left(x\frac{\partial}{\partial x}+y\frac{\partial}{\partial y}\right)^m f
  =n(n-1)\cdots(n-m+1)f. \tag{1}
\]

\textbf{方法一}\quad 直接在 $f(tx,ty)=t^nf(x,y)$ 两端同时对 $t$ 求 $m$ 阶导数，然后令 $t=1$.

\textbf{方法二}\quad 利用数学归纳法.

\noindent\textbf{证 I}\quad 已知：若 $z=f(x,y),x=a+th,y=b+tk$（$a,b,h,k$ 为常数），则
\[
  \frac{d^mz}{dt^m}=\left(h\frac{\partial}{\partial x}+k\frac{\partial}{\partial y}\right)^m f(x,y).
\]
利用此法则，在式 $f(tx,ty)=t^nf(x,y)$ 两端同时对 $t$ 求 $m$ 阶导数，再令 $t=1$，即得式（1）.

\noindent\textbf{证 II}\quad 由上例知：式（1）对 $m=1$ 已成立. 现只要证明：若式（1）对 $m=k$ 成立，则必对 $m=k+1$ 也成立. 事实上，因 $m=k$ 时成立，所以
\[
  \sum_{i=0}^k\left(C_k^i x^i y^{k-i}\frac{\partial^k f}{\partial x^i\partial y^{k-i}}\right)\biggm|_{(tx,ty)}
  =n(n-1)\cdots(n-k+1)f(tx,ty)
\]
\[
  =n(n-1)\cdots(n-k+1)t^nf(x,y).
\]
两端同时除以 $t^k$ 得
\[
  \sum_{i=0}^k C_k^i x^i y^{k-i}
  \left(\frac{\partial^k f}{\partial x^i\partial y^{k-i}}\right)\biggm|_{(tx,ty)}
  =n(n-1)\cdots(n-k+1)t^{n-k}f(x,y). \tag{2}
\]
两边同时对 $t$ 求导，根据复合函数求导法则：
\[
  \frac{d}{dt}\left(\frac{\partial^k f}{\partial x^i\partial y^{k-i}}\biggm|_{(tx,ty)}\right)
  =x\frac{\partial^{k+1}f}{\partial x^{i+1}\partial y^{k-i}}\biggm|_{(tx,ty)}
  +y\frac{\partial^{k+1}f}{\partial x^i\partial y^{k+1-i}}\biggm|_{(tx,ty)}.
\]
故式（2）左端的导数为
\[
\begin{aligned}
  &\sum_{i=0}^k\left(
  C_k^i x^{i+1}y^{k-i}\frac{\partial^{k+1}f}{\partial x^{i+1}\partial y^{k-i}}\biggm|_{(tx,ty)}
  +C_k^i x^iy^{k+1-i}\frac{\partial^{k+1}f}{\partial x^i\partial y^{k+1-i}}\biggm|_{(tx,ty)}
  \right)\\
  &=\sum_{j=1}^{k+1}C_k^{j-1}x^jy^{k+1-j}\frac{\partial^{k+1}f}{\partial x^j\partial y^{k+1-j}}\biggm|_{(tx,ty)}
  +\sum_{i=0}^k C_k^i x^iy^{k+1-i}\frac{\partial^{k+1}f}{\partial x^i\partial y^{k+1-i}}\biggm|_{(tx,ty)}\\
  &=\sum_{i=0}^{k+1}C_{k+1}^i x^iy^{k+1-i}
  \frac{\partial^{k+1}f}{\partial x^i\partial y^{k+1-i}}\biggm|_{(tx,ty)}.
\end{aligned}
\]
故（2）式对 $t$ 求导后得
\[
  \sum_{i=0}^{k+1}C_{k+1}^i x^iy^{k+1-i}
  \frac{\partial^{k+1}f}{\partial x^i\partial y^{k+1-i}}\biggm|_{(tx,ty)}
  =n(n-1)\cdots(n-k)t^{n-k-1}f(x,y).
\]
在此式中令 $t=1$，即得 $m=k+1$ 时的式（1）.

\noindent\textbf{注}\quad 若 $y=f(x_1,x_2,\cdots,x_n)$ 为 $n$ 元 $k$ 次齐次 $m$ 阶可微函数，类似地有
\[
  \left(x_1\frac{\partial}{\partial x_1}+x_2\frac{\partial}{\partial x_2}+\cdots+x_n\frac{\partial}{\partial x_n}\right)^m f
  =k(k-1)\cdots(k-m+1)f,
\]
其中
\[
  f=f(x_1,x_2,\cdots,x_n),
\]
\[
  \left(x_1\frac{\partial}{\partial x_1}+x_2\frac{\partial}{\partial x_2}+\cdots+x_n\frac{\partial}{\partial x_n}\right)^m
  =
  \sum_{r_1+r_2+\cdots+r_n=m}\frac{m!}{r_1!r_2!\cdots r_n!}
  x_1^{r_1}x_2^{r_2}\cdots x_n^{r_n}
  \frac{\partial^m}{\partial x_1^{r_1}\partial x_2^{r_2}\cdots\partial x_n^{r_n}}.
\]

\subsection*{三、偏导数转化为极限}

\noindent\textbf{要点}\quad 利用偏导数的定义，有些关于偏导数的问题，可转化为相应的极限问题.

\noindent\textbf{例 6.2.8}\quad 设 $f_x',f_y',f_{yx}''$ 在 $(x_0,y_0)$ 的某邻域内存在，$f_{yx}''$ 在点 $(x_0,y_0)$ 处连续，证明 $f_{xy}''(x_0,y_0)$ 存在，且 $f_{xy}''(x_0,y_0)=f_{yx}''(x_0,y_0)$.

\noindent\textbf{证}\quad $1^\circ$（将混合偏导数转化为累次极限.）根据偏导数定义，
\[
\begin{aligned}
f_{xy}''(x_0,y_0)
&=\lim_{\Delta y\to0}\frac{f_x'(x_0,y_0+\Delta y)-f_x'(x_0,y_0)}{\Delta y}\\
&=\lim_{\Delta y\to0}\frac1{\Delta y}\left[
\lim_{\Delta x\to0}\frac{f(x_0+\Delta x,y_0+\Delta y)-f(x_0,y_0+\Delta y)}{\Delta x}
-\lim_{\Delta x\to0}\frac{f(x_0+\Delta x,y_0)-f(x_0,y_0)}{\Delta x}\right]\\
&=\lim_{\Delta y\to0}\lim_{\Delta x\to0}\frac1{\Delta x\Delta y}W,
\end{aligned}
\]
其中
\[
  W=f(x_0+\Delta x,y_0+\Delta y)-f(x_0,y_0+\Delta y)-f(x_0+\Delta x,y_0)+f(x_0,y_0).
\]
同理可证
\[
  f_{yx}''(x_0,y_0)=\lim_{\Delta x\to0}\lim_{\Delta y\to0}\frac1{\Delta x\Delta y}W.
\]

$2^\circ$（证明全面极限 $\displaystyle\lim_{\substack{\Delta x\to0\\ \Delta y\to0}}\dfrac{W}{\Delta x\Delta y}$ 存在，且等于 $f_{yx}''(x_0,y_0)$.）令
\[
  \psi(y)=f(x_0+\Delta x,y)-f(x_0,y),
\]
则
\[
\begin{aligned}
\frac{W}{\Delta x\Delta y}
&=\frac1{\Delta x\Delta y}[\psi(y_0+\Delta y)-\psi(y_0)]
=\frac1{\Delta x}\psi_y'(y_0+\theta_1\Delta y)\\
&=\frac1{\Delta x}\bigl[f_y'(x_0+\Delta x,y_0+\theta_1\Delta y)-f_y'(x_0,y_0+\theta_1\Delta y)\bigr]\\
&=f_{yx}''(x_0+\theta\Delta x,y_0+\theta_1\Delta y)\quad(0<\theta,\theta_1<1).
\end{aligned}
\]
因 $f_{yx}''$ 在 $(x_0,y_0)$ 处连续，故
\[
  \lim_{(\Delta x,\Delta y)\to(0,0)}\frac1{\Delta x\Delta y}W
  =\lim_{(\Delta x,\Delta y)\to(0,0)}f_{yx}''(x_0+\theta\Delta x,y_0+\theta_1\Delta y)
  =f_{yx}''(x_0,y_0).
\]

$3^\circ$ 因 $f_x',f_y'$ 在 $(x_0,y_0)$ 的邻域内存在，且 $\Delta y$ 充分小时，$\displaystyle\lim_{\Delta x\to0}\dfrac1{\Delta x}W$ 存在，由累次极限定理（或见前节的习题 6.1.10 的提示（4）及再提示），
\[
  f_{xy}''(x_0,y_0)
  =\lim_{\Delta y\to0}\lim_{\Delta x\to0}\frac1{\Delta x\Delta y}W
  =\lim_{(\Delta x,\Delta y)\to(0,0)}\frac1{\Delta x\Delta y}W
  =f_{yx}''(x_0,y_0).
\]

\subsection*{\ensuremath{\star}四、对微分方程作变量替换}

\noindent\textbf{a.\quad 对自变量作变量替换}

\noindent\textbf{要点}\quad 完成这种替换的关键在于：通过新旧自变量的关系，求出相应导数的关系，再代入方程.

\noindent\textbf{\ensuremath{\star}例 6.2.9}\quad 试将方程
\[
  \frac{\partial^2z}{\partial x^2}+\frac{\partial^2z}{\partial y^2}=0
\]
变换为极坐标的形式.（清华大学，北京师范大学，浙江大学）

\noindent\textbf{解 I}\quad $1^\circ$（由变量之间的关系，求出它们导数之间的关系.）因 $z=z(x,y),x=r\cos\theta,y=r\sin\theta$，所以
\[
  \frac{\partial z}{\partial r}=\frac{\partial z}{\partial x}\cos\theta+\frac{\partial z}{\partial y}\sin\theta,\quad
  \frac{\partial z}{\partial\theta}=-\frac{\partial z}{\partial x}r\sin\theta+\frac{\partial z}{\partial y}r\cos\theta. \tag{1}
\]
解得
\[
  \frac{\partial z}{\partial x}=\frac{\partial z}{\partial r}\cos\theta-\frac1r\frac{\partial z}{\partial\theta}\sin\theta,\quad
  \frac{\partial z}{\partial y}=\frac{\partial z}{\partial r}\sin\theta+\frac1r\frac{\partial z}{\partial\theta}\cos\theta. \tag{2}
\]
即算符有关系：
\[
  \frac{\partial}{\partial x}=\cos\theta\frac{\partial}{\partial r}-\frac1r\sin\theta\frac{\partial}{\partial\theta},\quad
  \frac{\partial}{\partial y}=\sin\theta\frac{\partial}{\partial r}+\frac1r\cos\theta\frac{\partial}{\partial\theta}. \tag{3}
\]
故
\[
\begin{aligned}
\frac{\partial^2 z}{\partial x^2}
&=\frac{\partial}{\partial x}\left(\frac{\partial z}{\partial x}\right)
=\left(\cos\theta\frac{\partial}{\partial r}-\frac1r\sin\theta\frac{\partial}{\partial\theta}\right)\frac{\partial z}{\partial x}\\
&=\frac{\partial^2z}{\partial r^2}\cos^2\theta
-\frac{\partial^2z}{\partial r\partial\theta}\frac1r\sin\theta\cos\theta
+\frac{\partial z}{\partial\theta}\frac1{r^2}\sin\theta\cos\theta\\
&\quad-\frac{\partial^2z}{\partial\theta\partial r}\frac1r\cos\theta\sin\theta
+\frac{\partial z}{\partial r}\frac1r\sin^2\theta
+\frac{\partial^2z}{\partial\theta^2}\frac1{r^2}\sin^2\theta
+\frac{\partial z}{\partial\theta}\frac1{r^2}\sin\theta\cos\theta.
\end{aligned}
\]
同理，有
\[
\begin{aligned}
\frac{\partial^2z}{\partial y^2}
&=\frac{\partial^2z}{\partial r^2}\sin^2\theta
+\frac{\partial^2z}{\partial r\partial\theta}\frac1r\sin\theta\cos\theta
-\frac{\partial z}{\partial\theta}\frac1{r^2}\sin\theta\cos\theta\\
&\quad+\frac{\partial^2z}{\partial\theta\partial r}\frac1r\cos\theta\sin\theta
+\frac{\partial z}{\partial r}\frac{\cos^2\theta}{r}
+\frac{\partial^2z}{\partial\theta^2}\frac{\cos^2\theta}{r^2}
-\frac{\partial z}{\partial\theta}\frac1{r^2}\sin\theta\cos\theta.
\end{aligned}
\]

$2^\circ$（代入化简.）将上述结果代入原方程，化简即得
\[
  \frac{\partial^2z}{\partial r^2}+\frac1r\frac{\partial z}{\partial r}
  +\frac1{r^2}\frac{\partial^2z}{\partial\theta^2}=0. \tag{4}
\]

\noindent\textbf{解 II}\quad 将式（1）分别对 $r,\theta$ 再求一次导数（$x,y$ 作中间变量），得
\[
  z_{rr}''=z_{xx}''\cos^2\theta+z_{xy}''\cos\theta\sin\theta+z_{yx}''\sin\theta\cos\theta+z_{yy}''\sin^2\theta, \tag{5}
\]
\[
  z_{\theta\theta}''=z_{xx}''r^2\sin^2\theta-z_{xy}''r^2\sin\theta\cos\theta
  -z_{yx}''r^2\cos\theta\sin\theta+z_{yy}''r^2\cos^2\theta-z_x'r\cos\theta-z_y'r\sin\theta. \tag{6}
\]
式（5）乘 $r^2$ 与（6）相加，得
\[
  r^2z_{rr}''+z_{\theta\theta}''
  =r^2(z_{xx}''+z_{yy}'')-z_x'r\cos\theta-z_y'r\sin\theta
  \stackrel{\text{式（2）}}{=}r^2(z_{xx}''+z_{yy}'')-rz_r'.
\]
再同除以 $r^2$，移项，即得式（4）.

上例是把新变量理解为自变量，把原来的自变量理解为中间变量；有时也可以反过来.

\noindent\textbf{\ensuremath{\star}例 6.2.10}\quad 设
\[
  \frac{\partial^2W}{\partial x^2}+\frac{\partial^2W}{\partial y^2}=0, \tag{1}
\]
\[
  u=x^2-y^2,\quad v=2xy, \tag{2}
\]
试用关系（2），将（1）变成关于 $u,v$ 的方程.（南京理工大学，北京工业大学）

\noindent\textbf{解}\quad 将依赖关系看成 $W=W(u,v)$，$u,v$ 为 $x,y$ 的函数，于是
\[
  \frac{\partial W}{\partial x}
  =\frac{\partial W}{\partial u}\frac{\partial u}{\partial x}
  +\frac{\partial W}{\partial v}\frac{\partial v}{\partial x}
  =2x\frac{\partial W}{\partial u}+2y\frac{\partial W}{\partial v}.
\]
可见
\[
  \frac{\partial}{\partial x}=2x\frac{\partial}{\partial u}+2y\frac{\partial}{\partial v}.
\]
因此
\[
\begin{aligned}
\frac{\partial^2W}{\partial x^2}
&=2\frac{\partial W}{\partial u}+4x^2\frac{\partial^2W}{\partial u^2}
+4xy\frac{\partial^2W}{\partial v\partial u}
+4xy\frac{\partial^2W}{\partial u\partial v}
+4y^2\frac{\partial^2W}{\partial v^2}.
\end{aligned}
\]
类似可得
\[
  \frac{\partial^2W}{\partial y^2}
  =-2\frac{\partial W}{\partial u}+4y^2\frac{\partial^2W}{\partial u^2}
  -4xy\frac{\partial^2W}{\partial v\partial u}
  -4xy\frac{\partial^2W}{\partial u\partial v}
  +4x^2\frac{\partial^2W}{\partial v^2}.
\]
两式相加得
\[
  \frac{\partial^2W}{\partial x^2}+\frac{\partial^2W}{\partial y^2}
  =(4x^2+4y^2)\left(\frac{\partial^2W}{\partial u^2}+\frac{\partial^2W}{\partial v^2}\right).
\]
故原方程转化为
\[
  \frac{\partial^2W}{\partial u^2}+\frac{\partial^2W}{\partial v^2}=0.
\]

\noindent\textbf{new 练习}\quad 设函数 $z=z(u,v)$ 有连续 2 阶导数，并且变换 $u=x+2y,v=x+ay$ 能将方程 $2z_{xx}+z_{xy}-z_{yy}=0$ 变为 $z_{uv}=0$.

1）求 $a$；\quad 2）求出微分方程的解.（武汉大学，中国科学技术大学）

\noindent\textbf{解}\quad 1）$z=z(u,v)=z(x+2y,x+ay)$，故 $z_x=z_u+z_v$，于是
\[
  z_{xx}=z_{uu}+2z_{uv}+z_{vv},\quad
  z_{xy}=2z_{uu}+(a+2)z_{uv}+az_{vv}.
\]
由 $z_y=2z_u+az_v$，得
\[
  z_{yy}=4z_{uu}+4az_{uv}+a^2z_{vv}.
\]
因此
\[
\begin{aligned}
0&=2z_{xx}+z_{xy}-z_{yy}\\
&=2(z_{uu}+2z_{uv}+z_{vv})+(2z_{uu}+(a+2)z_{uv}+az_{vv})-(4z_{uu}+a^2z_{vv})\\
&=(2+a-a^2)z_{vv}=(1+a)(2-a)z_{vv}.
\end{aligned}
\]
得 $a=-1$ 或 $a=2$. 但 $a=2$ 不合本题，否则 $u=x+2y=v$，整个 $xOy$ 平面变成 $u=v$ 的一条直线，谈不上一对一，谈不上可逆，故 $a=-1$.

2）由 $z_{uv}=0$ 得 $z_u=f(u)$，
\[
  z=\int f(u)\,du+G(v)=F(u)+G(v)=F(x+2y)+G(x-y),
\]
其中 $F(u)$ 和 $G(v)$ 为任意函数（不难验证答案的正确性）.

\noindent\textbf{b.\quad 自变量与因变量都变化的变量替换}

\noindent\textbf{要点}\quad 当自变量与因变量都变化时，要对微分方程作变量替换，可使用如下方法：先将原来的因变量，写成新的因变量的函数；从而原来的函数，可看成是以新变量为中间变量的复合函数；应用复合函数微分法，求出原导数与新导数（新因变量对新自变量的导数）的关系；代入原方程，化为新函数、新变量的方程.

\noindent\textbf{\ensuremath{\star}例 6.2.11}\quad $z$ 为 $x,y$ 的可微函数，试将方程
\[
  x^2\frac{\partial z}{\partial x}+y^2\frac{\partial z}{\partial y}=z^2 \tag{1}
\]
变成 $w=w(u,v)$ 的方程. 假设
\[
  x=u,\quad y=\frac{u}{1+uv},\quad z=\frac{u}{1+uw}. \tag{2}
\]
（河北师范大学，湖南大学）

\noindent\textbf{解}\quad 已知 $z=\dfrac{x}{1+xw}$，$w=w(u,v)$，$u=x$，$v=\dfrac1y-\dfrac1x$，即变量的依赖关系为 $z=z(x,w(u,v))$，$u=x$，$v=v(x,y)$.
利用复合函数微分法
\[
  \frac{\partial z}{\partial x}
  =\frac{1-x^2\dfrac{\partial w}{\partial u}-\dfrac{\partial w}{\partial v}}{(1+xw)^2},\quad
  \frac{\partial z}{\partial y}
  =\frac{(1+xv)^2}{(1+xw)^2}\frac{\partial w}{\partial v}.
\]
将此结果及式（2）代入原方程（1），即得
\[
  \frac{\partial w}{\partial u}=0.
\]

\noindent\textbf{例 6.2.12}\quad 取 $\mu,\nu$ 为新自变量及 $w=w(\mu,\nu)$ 为新函数，变换方程
\[
  \frac{\partial^2z}{\partial x^2}+\frac{\partial^2z}{\partial x\partial y}+\frac{\partial z}{\partial x}=z. \tag{1}
\]
设
\[
  \mu=\frac{x+y}{2},\quad \nu=\frac{x-y}{2},\quad w=ze^x. \tag{2}
\]
（假设出现的导数皆连续.）（河南师范大学）

\noindent\textbf{解}\quad $z$ 看成是 $x,y$ 的复合函数如下：
\[
  z=\frac{w}{e^x},\quad w=w(\mu,\nu),\quad \mu=\frac{x+y}{2},\quad \nu=\frac{x-y}{2}.
\]
求出此复合函数的导数 $\dfrac{\partial^2z}{\partial x^2},\dfrac{\partial^2z}{\partial x\partial y},\dfrac{\partial z}{\partial x}$. 代入原方程（1），并将 $x,y,z$ 变换为 $\mu,\nu,w$，得
\[
  \frac{\partial^2w}{\partial \mu^2}+\frac{\partial^2w}{\partial \mu\partial\nu}=2w.
\]

\noindent\textbf{例 6.2.13}\quad 证明：在变换 $u=\dfrac{x}{y},v=x,w=xz-y$ 之下，方程
\[
  y\frac{\partial^2z}{\partial y^2}+2\frac{\partial z}{\partial y}=\frac2x
\]
可变成
\[
  \frac{\partial^2w}{\partial u^2}=0.
\]
（浙江大学）

\noindent\textbf{提示}\quad $z=\dfrac yx+\dfrac wx$，并将 $w$ 看成 $u,v$ 的函数，$u,v$ 按题设的变换是 $x,y$ 的函数，然后用复合函数求导法则计算 $\dfrac{\partial z}{\partial x},\dfrac{\partial^2z}{\partial x^2}$.

\noindent\textbf{例 6.2.14}\quad 考察变换 $x=a_1u+b_1v+c_1w,y=a_2u+b_2v+c_2w,z=a_3u+b_3v+c_3w$. 问在什么条件下（即 $a_i,b_i,c_i$ 满足什么条件时），对任何二阶连续可微函数 $f$，
\[
  \left(\frac{\partial f}{\partial x}\right)^2+
  \left(\frac{\partial f}{\partial y}\right)^2+
  \left(\frac{\partial f}{\partial z}\right)^2
  =
  \left(\frac{\partial f}{\partial u}\right)^2+
  \left(\frac{\partial f}{\partial v}\right)^2+
  \left(\frac{\partial f}{\partial w}\right)^2,
\]
和
\[
  \frac{\partial^2f}{\partial x^2}+\frac{\partial^2f}{\partial y^2}+\frac{\partial^2f}{\partial z^2}
  =
  \frac{\partial^2f}{\partial u^2}+\frac{\partial^2f}{\partial v^2}+\frac{\partial^2f}{\partial w^2}
\]
在此变换下形式不变，即成立？（复旦大学）

\noindent\textbf{解}\quad 利用复合函数微分法，
\[
  f_u'=f_x'a_1+f_y'a_2+f_z'a_3, \tag{1}
\]
\[
  f_v'=f_x'b_1+f_y'b_2+f_z'b_3, \tag{2}
\]
\[
  f_w'=f_x'c_1+f_y'c_2+f_z'c_3. \tag{3}
\]
以上三式平方后相加，得
\begin{equation}
\begin{aligned}
f_u'^2+f_v'^2+f_w'^2
&=f_x'^2(a_1^2+b_1^2+c_1^2)+f_y'^2(a_2^2+b_2^2+c_2^2)+f_z'^2(a_3^2+b_3^2+c_3^2)\\
&\quad+2f_x'f_y'(a_1a_2+b_1b_2+c_1c_2)
+2f_y'f_z'(a_2a_3+b_2b_3+c_2c_3)\\
&\quad+2f_z'f_x'(a_3a_1+b_3b_1+c_3c_1).
\end{aligned}
\tag{4}
\end{equation}
将式（1），（2），（3）分别对 $u,v,w$ 求导，然后相加得
\begin{equation}
\begin{aligned}
f_{uu}''+f_{vv}''+f_{ww}''
&=f_{xx}''(a_1^2+b_1^2+c_1^2)+f_{yy}''(a_2^2+b_2^2+c_2^2)+f_{zz}''(a_3^2+b_3^2+c_3^2)\\
&\quad+2f_{xy}''(a_1a_2+b_1b_2+c_1c_2)
+2f_{yz}''(a_2a_3+b_2b_3+c_2c_3)\\
&\quad+2f_{zx}''(a_3a_1+b_3b_1+c_3c_1).
\end{aligned}
\tag{5}
\end{equation}
由式（4），（5）可知，

当
\[
  a_1^2+b_1^2+c_1^2=a_2^2+b_2^2+c_2^2=a_3^2+b_3^2+c_3^2=1,
\]
\[
  a_1a_2+b_1b_2+c_1c_2=a_2a_3+b_2b_3+c_2c_3=a_3a_1+b_3b_1+c_3c_1=0
\]
时，恒有
\[
  f_u'^2+f_v'^2+f_w'^2=f_x'^2+f_y'^2+f_z'^2,\quad
  f_{uu}''+f_{vv}''+f_{ww}''=f_{xx}''+f_{yy}''+f_{zz}''.
\]

\subsection*{*五、多元函数的可微性}

\noindent\textbf{要点}\quad $\mathbb{R}^n$ 中，$f$ 在点 $P_0$ 处可微，等价于 $f$ 在 $P_0$ 的某邻域里，满足如下（相互等价的）任一等式：
\[
  \text{i)}\quad
  \Delta f(P_0)=f(P)-f(P_0)
  =\sum_{i=1}^n f_{x_i}'(P_0)\Delta x_i+\sum_{i=1}^n\varepsilon_i\Delta x_i;
\]
其中 $\varepsilon_i\to0$（当 $\Delta x_1,\cdots,\Delta x_n\to0$ 时）. 这里 $P-P_0=(\Delta x_1,\Delta x_2,\cdots,\Delta x_n)$.
\[
  \text{ii)}\quad
  \Delta f(P_0)=f(P)-f(P_0)=\sum_{i=1}^n f_{x_i}'(P_0)\Delta x_i+\varepsilon\cdot\rho,
\]
其中 $\varepsilon\to0$（当 $\rho\to0$ 时），$\rho=\sqrt{\sum_{i=1}^n\Delta x_i^2}$.
\[
  \text{iii)}\quad
  \Delta f(P_0)=f(P)-f(P_0)=\sum_{i=1}^n f_{x_i}'(P_0)\Delta x_i+o(\rho).
\]

\noindent\textbf{\ensuremath{\star}}\quad 以二元函数为例，要证明 $f$ 在 $(x_0,y_0)$ 处可微，通常方法有两种：

其一是根据上述条件 ii），iii），若
\[
  \lim_{\substack{\Delta x\to0\\ \Delta y\to0}}
  \frac{\Delta f-f_x'(x_0,y_0)\Delta x-f_y'(x_0,y_0)\Delta y}
  {\sqrt{\Delta x^2+\Delta y^2}}=0, \tag{A}
\]
则 $f$ 在 $(x_0,y_0)$ 处可微，否则不可微.

方法二是证明 $f$ 在 $(x_0,y_0)$ 的邻域里有
\[
  f(x_0+\Delta x,y_0+\Delta y)-f(x_0,y_0)
  =f_x'(x_0,y_0)\Delta x+f_y'(x_0,y_0)\Delta y+\varepsilon_1\Delta x+\varepsilon_2\Delta y,
\]
其中 $\varepsilon_1,\varepsilon_2\to0$（当 $\Delta x,\Delta y\to0$ 时），否则不可微.

\noindent\textbf{例 6.2.15}\quad 设 $f(x,y)=\varphi(|xy|)$，其中 $\varphi(0)=0$，且在 $u=0$ 附近满足 $|\varphi(u)|\le u^2$，试证 $f(x,y)$ 在 $(0,0)$ 处可微.（中国科学院）

\noindent\textbf{证}\quad
\[
  f_x'(0,0)=\lim_{x\to0}\frac{\varphi(|x\cdot0|)-\varphi(0)}{x}=0.
\]
同理，$f_y'(0,0)=0$. 又
\[
\begin{aligned}
\left|\frac{f(x,y)-f(0,0)-f_x'(0,0)x-f_y'(0,0)y}{\sqrt{x^2+y^2}}\right|
&=\left|\frac{\varphi(|xy|)}{\sqrt{x^2+y^2}}\right|\\
&\le \frac{|xy|^2}{\sqrt{x^2+y^2}}\le |xy|^{3/2}\to0
\end{aligned}
\]
（当 $x\to0,y\to0$ 时）. 因此 $f$ 在 $(0,0)$ 处可微.

\noindent\textbf{例 6.2.16}\quad 证明：函数
\[
  f(x,y)=
  \begin{cases}
  \dfrac{xy}{\sqrt{x^2+y^2}}, & x^2+y^2\ne0,\\
  0, & x^2+y^2=0
  \end{cases}
\]
在 $\mathbb{R}^2$ 上连续，且 $f_x',f_y'$ 有界，但 $f$ 在 $(0,0)$ 处不可微.（山东大学，北京航空航天大学）

\noindent\textbf{提示}\quad $f$ 在 $(0,0)$ 处不可微最后归结为证明：
\[
  \frac{xy}{x^2+y^2}\not\to0\quad(x\to0,y\to0).
\]

\noindent\textbf{注}\quad 按上述方法和步骤证明可微性，原则上没有什么困难. 但因为问题最终归结为求极限，因此进行时，还会碰到许多具体困难；这时应根据具体情况，灵活处理.

\noindent\textbf{\ensuremath{\star}例 6.2.17}\quad 设 $f(x)$ 及 $g(x)$ 分别在区间 $[a,b]$ 及 $[c,d]$ 上连续，定义
\[
  F(x,y)=\int_a^x f(t)\,dt\cdot\int_c^y g(s)\,ds\quad(a\le x\le b,c\le y\le d).
\]
试用全微分的定义证明 $F(x,y)$ 在 $(x_0,y_0)$ 处可微，其中 $a\le x_0\le b,c\le y_0\le d$ 为任意的定点.（大连理工大学）

\noindent\textbf{分析}\quad 要证明 $F(x,y)$ 在 $(x_0,y_0)$ 处可微，即要证明当 $x\to x_0,y\to y_0$ 时，
\[
  \frac{F(x,y)-F(x_0,y_0)-F_x'(x_0,y_0)(x-x_0)-F_y'(x_0,y_0)(y-y_0)}
  {\sqrt{(x-x_0)^2+(y-y_0)^2}}
\]
以 $0$ 为极限. 为此，我们写出式里分子各项，其中
\begin{equation}
\begin{aligned}
F(x,y)-F(x_0,y_0)
&=\int_a^{x_0}f(t)\,dt\int_{y_0}^y g(s)\,ds
+\int_{x_0}^x f(t)\,dt\int_c^{y_0}g(s)\,ds\\
&\quad+\int_{x_0}^x f(t)\,dt\int_{y_0}^y g(s)\,ds.
\end{aligned}
\tag{2}
\end{equation}
式（1）分子的第三项可以化为式（2）右端第二项的形式，
\[
  F_x'(x_0,y_0)(x-x_0)
  =f(x_0)\int_c^{y_0}g(s)\,ds\cdot(x-x_0)
  =\int_{x_0}^x f(x_0)\,dt\int_c^{y_0}g(s)\,ds. \tag{3}
\]
式（1）分子的第四项可化为式（2）右端第一项的形式，
\[
  F_y'(x_0,y_0)(y-y_0)
  =\int_a^{x_0}f(t)\,dt\cdot g(y_0)(y-y_0)
  =\int_a^{x_0}f(t)\,dt\int_{y_0}^y g(y_0)\,ds. \tag{4}
\]
将（2），（3），（4）合并，于是（1）中的分子可以化为三项，即
\begin{equation}
\begin{aligned}
&F(x,y)-F(x_0,y_0)-F_x'(x_0,y_0)(x-x_0)-F_y'(x_0,y_0)(y-y_0)\\
&=\int_a^{x_0}f(t)\,dt\int_{y_0}^y[g(s)-g(y_0)]\,ds
+\int_{x_0}^x[f(t)-f(x_0)]\,dt\int_c^{y_0}g(s)\,ds\\
&\quad+\int_{x_0}^x f(t)\,dt\int_{y_0}^y g(s)\,ds.
\end{aligned}
\tag{5}
\end{equation}
因为 $f(x)$ 在 $[a,b]$ 上连续，即 $\exists M>0$，使得 $|f(x)|\le M$（$\forall x\in[a,b]$），又因 $g(x)$ 在 $y_0$ 处连续，所以 $\forall\varepsilon>0,\exists\delta>0$，当 $y\in[c,d]$，$|y-y_0|<\delta$ 时，$|g(y)-g(y_0)|<\varepsilon$. 于是
\[
\left|
\frac{\int_a^{x_0}f(t)\,dt\int_{y_0}^y[g(s)-g(y_0)]\,ds}
{\sqrt{(x-x_0)^2+(y-y_0)^2}}
\right|
\le M(b-a)\varepsilon.
\]
类似可证明其余两项亦趋向零. 故式（1）趋向零，$F(x,y)$ 可微性获证.

应当指出的是，本例作为试题主要是为了考可微性定义的使用，实际上利用可微函数的乘积定理，本题的结论是明显的.

\noindent\textbf{\ensuremath{\star}例 6.2.18}\quad 若 $f_x'(x,y)$ 在点 $(x_0,y_0)$ 处存在，$f_y'(x,y)$ 在点 $(x_0,y_0)$ 处连续，证明 $f(x,y)$ 在点 $(x_0,y_0)$ 处可微.（吉林大学，辽宁师范大学）

\noindent\textbf{证}\quad
\[
\begin{aligned}
f(x_0+\Delta x,y_0+\Delta y)-f(x_0,y_0)
&=[f(x_0+\Delta x,y_0+\Delta y)-f(x_0+\Delta x,y_0)]\\
&\quad+[f(x_0+\Delta x,y_0)-f(x_0,y_0)]\\
&=f_y'(x_0+\Delta x,y_0+\theta\Delta y)\Delta y
+f_x'(x_0,y_0)\Delta x+\varepsilon_1\Delta x\\
&=[f_y'(x_0,y_0)+\varepsilon_2]\Delta y
+f_x'(x_0,y_0)\Delta x+\varepsilon_1\Delta x\\
&=f_x'(x_0,y_0)\Delta x+f_y'(x_0,y_0)\Delta y+\varepsilon_1\Delta x+\varepsilon_2\Delta y,
\end{aligned}
\]
其中 $\varepsilon_1,\varepsilon_2\to0$（当 $\Delta x,\Delta y\to0$ 时）. 故 $f$ 在点 $(x_0,y_0)$ 处可微.

\noindent\textbf{new 练习}\quad 设函数 $f:\mathbb{R}^n\to\mathbb{R}$，在 $\mathbb{R}^n\backslash\{0\}$ 可微，在 $0$ 处连续，$f(0)=0$ 且
\[
  \lim_{x\to0}\frac{\partial f(x)}{\partial x_i}=0\quad(i=1,2,\cdots,n).
\]
试证 $f$ 在 $0$ 处可微.（北京大学）

\noindent\textbf{提示}\quad $1^\circ$ 利用导数定义易知
\[
  f_{x_k}'(0)=\left.\frac{\partial f(x)}{\partial x_k}\right|_{x=0}=0\quad(k=1,2,\cdots,n).
\]
例如：
\[
  f_{x_1}'(0)=\lim_{\substack{x_k=0(k\ne1)\\x_1\to0}}
  \frac{f(x_1,0,\cdots,0)-f(0)}{x_1}
  =\lim_{\xi_1\to0}f_1'(\xi_1,0,\cdots,0)=0.
\]
$2^\circ$ 仿例 6.2.18，
\[
  |f(x)-f(0)|
  \le |f(x_1,x_2,\cdots,x_n)-f(0,x_2,\cdots,x_n)|+\cdots+
  |f(0,0,\cdots,0,x_n)-f(0,0,\cdots,0)|
\]
\[
  =\sum_{k=1}^n |f_k'(\xi_{(k)})x_k|,
\]
其中 $\xi_{(k)}=(0,0,\cdots,0,\xi_k,x_{k+1},\cdots,x_n)$. 再提示\quad 要证 $f$ 在 $0$ 处可微，即要证：
\[
  \lim_{x\to0}\frac{\Delta f(0)-\sum_{k=1}^n f_k'(0)\Delta x_k}{|x-0|}=0.
\]
事实上，
\[
  \lim_{x\to0}\frac{f(x)-f(0)-\sum_{k=1}^n f_k'(0)(x_k-0)}{|x-0|}
  =\lim_{x\to0}\frac{f(x)-f(0)}{|x|}
  \le\sum_{k=1}^n\lim_{\xi_{(k)}\to0}\left|\frac{\partial f}{\partial x_k}(\xi_{(k)})\right|=0.
\]

\noindent\textbf{$\ast$例 6.2.19}\quad 设 1）$x=\varphi(s,t)$ 及 $y=\psi(s,t)$ 在区域 $D$ 内可微，且 $(s,t)\in D$ 时，$(x,y)=(\varphi(s,t),\psi(s,t))\in E$（区域）；

2）函数 $z=f(x,y)$ 在区域 $E$ 内可微.

试证 $z=f(\varphi(s,t),\psi(s,t))$ 在 $D$ 内可微，且 $df(s,t)=f_x'\,dx+f_y'\,dy$.

\noindent\textbf{证}\quad 记 $P=(s,t)$ 为 $D$ 内任意一点，$Q=(x,y)=(\varphi(s,t),\psi(s,t))$. 由已知条件有
\[
  \Delta z=f_x'(Q)\Delta x+f_y'(Q)\Delta y+\varepsilon\rho, \tag{1}
\]
其中 $\rho=\sqrt{\Delta x^2+\Delta y^2}$，$\varepsilon\to0$（当 $\rho\to0$ 时）.
\[
  \Delta x=\varphi_s'(P)\Delta s+\varphi_t'(P)\Delta t+\varepsilon_1\rho_1, \tag{2}
\]
\[
  \Delta y=\psi_s'(P)\Delta s+\psi_t'(P)\Delta t+\varepsilon_2\rho_1, \tag{3}
\]
其中 $\rho_1=\sqrt{\Delta s^2+\Delta t^2}$，$\varepsilon_1,\varepsilon_2\to0$（当 $\rho_1\to0$ 时）.
将式（2），（3）代入（1）得
\begin{equation}
\begin{aligned}
\Delta z
&=[f_x'(Q)\varphi_s'(P)+f_y'(Q)\psi_s'(P)]\Delta s
+[f_x'(Q)\varphi_t'(P)+f_y'(Q)\psi_t'(P)]\Delta t\\
&\quad+\varepsilon\rho+f_x'(Q)\varepsilon_1\rho_1+f_y'(Q)\varepsilon_2\rho_1.
\end{aligned}
\tag{4}
\end{equation}
要证明 $f(\varphi(s,t),\psi(s,t))$ 在 $P=(s,t)$ 处可微，即要证明式（4）里有
\[
  \frac{\varepsilon\rho+f_x'(Q)\varepsilon_1\rho_1+f_y'(Q)\varepsilon_2\rho_1}{\rho_1}
  =\varepsilon\frac{\rho}{\rho_1}+f_x'(Q)\varepsilon_1+f_y'(Q)\varepsilon_2\to0. \tag{5}
\]
因 $\rho_1\to0$ 时 $\varepsilon_1\to0,\varepsilon_2\to0$，故式（5）中后两项趋向零. 现只需证明 $\varepsilon\dfrac{\rho}{\rho_1}\to0$. 由式（2），（3）知
\[
  \left|\frac{\Delta x}{\rho_1}\right|\le |\varphi_s'(P)|+|\varphi_t'(P)|+1,\quad
  \left|\frac{\Delta y}{\rho_1}\right|\le |\psi_s'(P)|+|\psi_t'(P)|+1
\]
（当 $\rho_1$ 充分小时）. 所以
\[
  \left|\frac{\rho}{\rho_1}\right|
  =\frac{\sqrt{\Delta x^2+\Delta y^2}}{\rho_1}
  =\sqrt{\left(\frac{\Delta x}{\rho_1}\right)^2+\left(\frac{\Delta y}{\rho_1}\right)^2}
\]
有界，式（5）证毕. 即
\[
  dz=[f_x'(Q)\varphi_s'(P)+f_y'(Q)\psi_s'(P)]\Delta s
  +[f_x'(Q)\varphi_t'(P)+f_y'(Q)\psi_t'(P)]\Delta t
  =f_x'(Q)dx+f_y'(Q)dy.
\]
证毕.

$f$ 在 $(x_0,y_0)$ 处可微，意味着 $f$ 在 $(x_0,y_0)$ 附近与一个一次函数近似，只相差一个高阶无穷小量（相对 $\rho=\sqrt{(x-x_0)^2+(y-y_0)^2}$ 而言）：
\[
  f(x,y)=f(x_0,y_0)+f_x'(x_0,y_0)(x-x_0)+f_y'(x_0,y_0)(y-y_0)+o(\rho).
\]
因此，利用此式可将 $f(x,y)$ 变形.

\noindent\textbf{$\ast$例 6.2.20}\quad 设 $f_x',f_y'$ 在 $(x_0,y_0)$ 的某个邻域里存在，且在点 $(x_0,y_0)$ 处可微，证明：$f_{xy}''(x_0,y_0)=f_{yx}''(x_0,y_0)$.

\noindent\textbf{分析}\quad 从例 6.2.8 已看到，$f_{xy}''(x_0,y_0)$ 与 $f_{yx}''(x_0,y_0)$ 是函数
\[
  \frac{W}{\Delta x\Delta y}
  =\frac1{\Delta x\Delta y}[f(x_0+\Delta x,y_0+\Delta y)-f(x_0,y_0+\Delta y)-f(x_0+\Delta x,y_0)+f(x_0,y_0)] \tag{1}
\]
的两个累次极限. 我们利用 $f_x',f_y'$ 在 $(x_0,y_0)$ 处的可微性，将证明 $\dfrac{W}{\Delta x\Delta y}$ 改写成
\[
  \frac{W}{\Delta x\Delta y}=f_{yx}''(x_0,y_0)+\varepsilon_1+\varepsilon_2\theta\frac{\Delta y}{\Delta x}-\varepsilon_3\theta\frac{\Delta y}{\Delta x} \tag{2}
\]
和
\[
  \frac{W}{\Delta x\Delta y}=f_{xy}''(x_0,y_0)+\varepsilon_4+\varepsilon_5\theta_1\frac{\Delta x}{\Delta y}-\varepsilon_6\theta_1\frac{\Delta x}{\Delta y}. \tag{3}
\]
两者对充分小的 $\Delta x,\Delta y$ 同时成立，且当 $\Delta x,\Delta y\to0$ 时，$\varepsilon_i\to0$（$i=1,2,\cdots,6$）. 于是令 $\Delta x=\Delta y\to0$，可得
\[
  f_{yx}''(x_0,y_0)=f_{xy}''(x_0,y_0). \tag{4}
\]

可见，问题归结为证明（2），（3）成立. 为此将 $\Delta x,\Delta y$ 取得足够小，引入辅助函数
\[
  \varphi(y)=f(x_0+\Delta x,y)-f(x_0,y). \tag{5}
\]
式（1）可改写成
\begin{equation}
\begin{aligned}
\frac{W}{\Delta x\Delta y}
&=\frac1{\Delta x}\cdot
\frac{[f(x_0+\Delta x,y_0+\Delta y)-f(x_0,y_0+\Delta y)]
-[f(x_0+\Delta x,y_0)-f(x_0,y_0)]}{\Delta y}\\
&=\frac1{\Delta x}\bigl[f_y'(x_0+\Delta x,y_0+\theta\Delta y)
-f_y'(x_0,y_0+\theta\Delta y)\bigr].
\end{aligned}
\tag{6}
\end{equation}
因 $f_y'$ 在 $(x_0,y_0)$ 处可微，
\[
  f_y'(x_0+\Delta x,y_0+\theta\Delta y)
  =f_y'(x_0,y_0)+f_{yx}''(x_0,y_0)\Delta x
  +f_{yy}''(x_0,y_0)\theta\Delta y+\varepsilon_1\Delta x+\varepsilon_2\theta\Delta y, \tag{7}
\]
\[
  f_y'(x_0,y_0+\theta\Delta y)
  =f_y'(x_0,y_0)+f_{yy}''(x_0,y_0)\theta\Delta y+\varepsilon_3\theta\Delta y, \tag{8}
\]
其中 $\varepsilon_1,\varepsilon_2\to0$（当 $\Delta x,\Delta y\to0$ 时），$\varepsilon_3\to0$（当 $\Delta y\to0$ 时）. 将（7），（8）代入（6）即得（2）. 类似可得（3）.

\subsection*{单元练习 6.2}

\subsection*{\ensuremath{\star}偏导数的计算}

\noindent\textbf{6.2.1}\quad 设 $u=f(r,r\cos\theta)$ 有二阶连续偏导数，求 $\dfrac{\partial u}{\partial r}$，$\dfrac{\partial u}{\partial\theta}$，$\dfrac{\partial^2u}{\partial r\partial\theta}$.（复旦大学）
\[
  \left\langle f_1'+\cos\theta f_2';\ -r\sin\theta f_2';\
  -r\sin\theta f_{21}''-\sin\theta f_2'-r\sin\theta\cos\theta f_{22}''\right\rangle
\]

\noindent\textbf{6.2.2}\quad 设 $u=f(x-y,y-z,z-x)$，$f$ 具有二阶连续偏导数，求 $\dfrac{\partial^2u}{\partial x^2}$ 及 $\dfrac{\partial^2u}{\partial y\partial z}$.（上海交通大学）
\[
  \left\langle f_{11}''-2f_{13}''+f_{33}'';\
  f_{23}''-f_{22}''-f_{13}''+f_{12}''\right\rangle
\]

\noindent\textbf{6.2.3}\quad 设 $u=xyze^{x+y+z}$，求
\[
  \frac{\partial^{p+q+r}u}{\partial x^p\partial y^q\partial z^r}.
\]
（北京航空航天大学）
\[
  \left\langle (x+p)(y+q)(z+r)e^{x+y+z}\right\rangle
\]

\noindent\textbf{6.2.4}\quad 设 $f$ 为可微函数，$u=f(x^2+y^2+z^2)$ 和方程 $3x+2y^2+z^3=6xyz$.

1）若方程确定隐函数 $z=z(x,y)$，求 $\dfrac{\partial u}{\partial x}\bigm|_{P_0}$，其中 $P_0(1,1,1)$.
\[
  \left\langle \frac{\partial u}{\partial x}\biggm|_{P_0}=0\right\rangle
\]

2）若方程确定隐函数 $y=y(x,z)$，求 $\dfrac{\partial u}{\partial x}\bigm|_{P_0}$.（华中师范大学）
\[
  \left\langle \frac{\partial u}{\partial x}\biggm|_{P_0}=-f'(3)\right\rangle
\]

\noindent\textbf{6.2.5}\quad 设 $z=f(x,y)$，$u=x+ay$，$v=x-ay$，$a$ 为常数，$z$ 关于 $u,v$ 有二阶连续偏导数，求 $\dfrac{\partial^2z}{\partial u\partial v}$.（厦门大学）
\[
  \left\langle \frac14f_{xx}''-\frac1{4a^2}f_{yy}''\right\rangle
\]

\noindent\textbf{6.2.6}\quad 设函数 $u(x)$ 由方程组 $u=f(x,y)$，$g(x,y,z)=0$，$h(x,z)=0$ 所确定，且 $\dfrac{\partial h}{\partial z}\ne0$，$\dfrac{\partial g}{\partial y}\ne0$，求 $\dfrac{du}{dx}$.（清华大学）
\[
  \left\langle f_x'-\frac{f_y'g_x'}{g_y'}+\frac{f_y'g_z'h_x'}{g_y'h_z'}\right\rangle
\]

\noindent\textbf{6.2.7}\quad 设 $f,F$ 可微，且 $\dfrac{\partial F}{\partial z}+\dfrac{\partial f}{\partial z}\dfrac{\partial F}{\partial y}\ne0$，求由
\[
  \begin{cases}
  y=f(x,z),\\
  F(x,y,z)=0
  \end{cases}
\]
所确定的函数 $y(x),z(x)$ 的一阶导数.（西安电子科技大学）
\[
  \left\langle
  y_x'=\frac{f_x'F_z'-f_z'F_x'}{F_z'+f_z'F_y'},\quad
  z_x'=\frac{-F_x'-f_x'F_y'}{F_z'+f_z'F_y'}
  \right\rangle
\]

\noindent\textbf{6.2.8}\quad 设函数 $F_i(u)$（$i=1,2,3$）可微，$A=|a_{ij}|$ 是一个三阶的函数行列式，其中 $a_{ij}=F_i(x_j)$，$i,j=1,2,3$，并且 $x_3$ 是由方程 $x_2^2+x_3+\sin(x_2x_3)=1$ 所确定的隐函数，求 $\dfrac{\partial A}{\partial x_1}$ 与 $\dfrac{\partial A}{\partial x_2}$ 在 $x_1=0,x_2=1,x_3=0$ 时的值.（西北大学）
\[
  \left\langle
  \frac{\partial A}{\partial x_1}=
  \begin{vmatrix}
  F_1'(0)&F_1(1)&F_1(0)\\
  F_2'(0)&F_2(1)&F_2(0)\\
  F_3'(0)&F_3(1)&F_3(0)
  \end{vmatrix},\quad
  \frac{\partial A}{\partial x_2}=
  \begin{vmatrix}
  F_1(0)&F_1(1)&-F_1'(0)\\
  F_2(0)&F_2(1)&-F_2'(0)\\
  F_3(0)&F_3(1)&-F_3'(0)
  \end{vmatrix}
  \right\rangle
\]

\noindent\textbf{提示}\quad
\[
  \frac{\partial A}{\partial x_1}=
  \begin{vmatrix}
  F_1'(x_1)&F_1(x_2)&F_1(x_3)\\
  F_2'(x_1)&F_2(x_2)&F_2(x_3)\\
  F_3'(x_1)&F_3(x_2)&F_3(x_3)
  \end{vmatrix},
\]
\[
\begin{aligned}
\frac{\partial A}{\partial x_2}
&=
  \begin{vmatrix}
  F_1(x_1)&F_1(x_2)&F_1'(x_3)\\
  F_2(x_1)&F_2(x_2)&F_2'(x_3)\\
  F_3(x_1)&F_3(x_2)&F_3'(x_3)
  \end{vmatrix}\frac{dx_3}{dx_2}\\
&\quad+
  \begin{vmatrix}
  F_1(x_1)&F_1'(x_2)&F_1(x_3)\\
  F_2(x_1)&F_2'(x_2)&F_2(x_3)\\
  F_3(x_1)&F_3'(x_2)&F_3(x_3)
  \end{vmatrix}.
\end{aligned}
\]

\subsection*{\ensuremath{\star}检验函数满足微分方程}

\noindent\textbf{6.2.9}\quad 设函数 $\varphi(z)$ 和 $\psi(x)$ 具有二阶连续导数，并设 $u=x\varphi(x+y)+y\psi(x+y)$. 试证：
\[
  \frac{\partial^2u}{\partial x^2}-2\frac{\partial^2u}{\partial x\partial y}
  +\frac{\partial^2u}{\partial y^2}=0.
\]
（中国科学院）

\noindent\textbf{6.2.10}\quad 证明：若 $u$ 是 $x,y,z$ 的函数且
\[
  \varphi(u^2-x^2,u^2-y^2,u^2-z^2)=0,
\]
则
\[
  \frac{u_x'}x+\frac{u_y'}y+\frac{u_z'}z=\frac1u.
\]
（东北师范大学）

\noindent\textbf{提示}\quad 原式对 $x$ 求导：
\[
  \varphi_1'(2uu_x'-2x)+\varphi_2'(2uu_x')+\varphi_3'(2uu_x')=0.
\]
整理后，得
\[
  (\varphi_1'+\varphi_2'+\varphi_3')\frac{u_x'}x=\frac1u\varphi_1'.
\]
同理
\[
  (\varphi_1'+\varphi_2'+\varphi_3')\frac{u_y'}y=\frac1u\varphi_2',\quad
  (\varphi_1'+\varphi_2'+\varphi_3')\frac{u_z'}z=\frac1u\varphi_3',
\]
三式相加即得所求.

\noindent\textbf{6.2.11}\quad 设 $u,v,w$ 都是 $x$ 的函数，具有二阶连续导数，试证：
\[
  W(u,v,w)=
  \begin{vmatrix}
  u&v&w\\
  u'&v'&w'\\
  u''&v''&w''
  \end{vmatrix}
\]
满足
\[
  W(u,v,w)=u^3W\left(1,\frac vu,\frac wu\right).
\]
（西北师范大学）

\noindent\textbf{提示}\quad
\[
  W\left(1,\frac vu,\frac wu\right)=
  \begin{vmatrix}
  1&v/u&w/u\\
  0&(v/u)'&(w/u)'\\
  0&(v/u)''&(w/u)''
  \end{vmatrix}
  =\frac1{u^3}W(u,v,w).
\]
（最后等号两端都等于：$\dfrac1{u^3}(A_{31}u''+A_{32}v''+A_{33}w'')$，$A_{ij}$ 是代数余子式.）

\noindent\textbf{6.2.12}\quad 设 $x=f(u,v)$，$y=g(u,v)$，$w=w(x,y)$ 有二阶连续偏导数，满足
\[
  \text{i) }\quad \frac{\partial f}{\partial u}=\frac{\partial g}{\partial v},\quad
  \frac{\partial f}{\partial v}=-\frac{\partial g}{\partial u};
  \qquad
  \text{ii) }\quad
  \frac{\partial^2w}{\partial x^2}+\frac{\partial^2w}{\partial y^2}=0.
\]
试证：
\[
  1)\quad \frac{\partial^2(fg)}{\partial u^2}+\frac{\partial^2(fg)}{\partial v^2}=0;
\]
\[
  2)\quad w(x,y)=w(f(u,v),g(u,v))\text{ 满足 }
  \frac{\partial^2w}{\partial u^2}+\frac{\partial^2w}{\partial v^2}=0.
\]
（北京大学）

\noindent\textbf{提示}\quad 1）$(f\cdot g)_u'=f_u'g+fg_u'=g_v'g-g_u'f,\cdots$ 类似可求 $(f\cdot g)_{vv}''$，并注意 $g_{uu}''=-f_{vu}''$，$g_{vu}''=f_{uu}''$（注 $g_{u^2}''$ 是 $g_{uu}''$ 的简写.）

2）注意使用对称性.

\noindent\textbf{再提示}\quad
\[
  \frac{\partial^2w}{\partial u^2}
  =w_{xx}''x_u'^2+2w_{xy}''x_u'y_u'+w_x'x_{u^2}''
  +w_{yy}''y_u'^2+w_y'y_{u^2}''.
\]
将 $u$ 换成 $v$ 照样成立，故
\[
\begin{aligned}
\frac{\partial^2w}{\partial u^2}+\frac{\partial^2w}{\partial v^2}
&=w_{xx}''(x_u'^2+x_v'^2)+2w_{xy}''[x_u'y_u'+x_v'y_v']\\
&\quad+w_{yy}''(y_u'^2+y_v'^2)+w_x'[x_{u^2}''+x_{v^2}'']
+w_y'[y_{u^2}''+y_{v^2}''].
\end{aligned}
\]
注意由已知条件 i），上式三个方括号为 $0$，两圆括号相等，因此上式
\[
  =(w_{xx}''+w_{yy}'')(x_u'^2+x_v'^2)=0.
\]

\noindent\textbf{注}\quad 事实上结论 2）证明之后，结论 1）自然成立，因为结论 1）中 $fg$ 是 $w=xy$，$x=f(u,v)$，$y=g(u,v)$ 的特殊情况.

\noindent\textbf{6.2.13}\quad 设 $u(x,y)$ 有连续的二阶偏导数，$F(s,t)$ 有连续的一阶偏导数，且满足
\[
  F(u_x',u_y')=0,\qquad (F_s')^2+(F_t')^2\ne0.
\]
证明：
\[
  u_{xx}''u_{yy}''-(u_{xy}'')^2=0.
\]
（华东师范大学）

\noindent\textbf{提示}\quad 题意表明，作为 $x,y$ 的函数，$s=u_x'(x,y)$，$t=u_y'(x,y)$，代入 $F(s,t)$ 后恒为零，即
\[
  F(u_x'(x,y),u_y'(x,y))\equiv0.
\]

\noindent\textbf{再提示}\quad 该式对 $x$ 求导，得
\[
  F_s'u_{xx}''+F_t'u_{yx}''=0,
\]
对 $y$ 求导得
\[
  F_s'u_{xy}''+F_t'u_{yy}''=0.
\]
已知 $(F_s')^2+(F_t')^2\ne0$，表明作为一次齐次线性方程组有非零解，故系数行列式应为零，即
\[
  \begin{vmatrix}
  u_{xx}''&u_{yx}''\\
  u_{xy}''&u_{yy}''
  \end{vmatrix}=0.
\]

\subsection*{\ensuremath{\star}变换微分方程（或微分式）}

\noindent\textbf{6.2.14}\quad 设 $u=x+y$，$v=\dfrac1x+\dfrac1y$，试用 $u,v$ 作新自变量变换方程
\[
  x^2\frac{\partial^2z}{\partial x^2}
  -(x^2+y^2)\frac{\partial^2z}{\partial x\partial y}
  +y^2\frac{\partial^2z}{\partial y^2}=0.
\]
（假设出现的二阶偏导数都连续.）（上海交通大学）

\noindent\textbf{提示}\quad 可用例 6.2.10 中的方法.

\noindent\textbf{解}\quad
\[
  \frac{\partial}{\partial x}=\frac{\partial}{\partial u}-\frac1{x^2}\frac{\partial}{\partial v},\quad
  \frac{\partial}{\partial y}=\frac{\partial}{\partial u}-\frac1{y^2}\frac{\partial}{\partial v}
\]
（其中 $z$ 被省略，成为算子），
\[
  \frac{\partial^2}{\partial x^2}
  =\frac{\partial^2}{\partial u^2}-\frac2{x^2}\frac{\partial^2}{\partial u\partial v}
  +\frac1{x^4}\frac{\partial^2}{\partial v^2}+\frac2{x^3}\frac{\partial}{\partial v},
\]
\[
  \frac{\partial^2}{\partial x\partial y}
  =\frac{\partial^2}{\partial u^2}
  -\left(\frac1{x^2}+\frac1{y^2}\right)\frac{\partial^2}{\partial u\partial v}
  +\frac1{x^2}\frac1{y^2}\frac{\partial^2}{\partial v^2},
\]
\[
  \frac{\partial^2}{\partial y^2}
  =\frac{\partial^2}{\partial u^2}-\frac2{y^2}\frac{\partial^2}{\partial u\partial v}
  +\frac1{y^4}\frac{\partial^2}{\partial v^2}+\frac2{y^3}\frac{\partial}{\partial v}.
\]
因此
\[
\begin{aligned}
&x^2\frac{\partial^2}{\partial x^2}
-(x^2+y^2)\frac{\partial^2}{\partial x\partial y}
+y^2\frac{\partial^2}{\partial y^2}\\
&=uv(uv-4)\frac{\partial^2}{\partial u\partial v}+2v\frac{\partial}{\partial v}.
\end{aligned}
\]
故最后结果是
\[
  uv(uv-4)\frac{\partial^2z}{\partial u\partial v}
  +2v\frac{\partial z}{\partial v}=0.
\]

\noindent\textbf{6.2.15}\quad 通过 $u=x-2\sqrt y$，$v=x+2\sqrt y$，变换方程
\[
  \frac{\partial^2z}{\partial x^2}-y\frac{\partial^2z}{\partial y^2}
  =\frac12\frac{\partial z}{\partial y}\quad (y>0),
\]
假设所出现的偏导数都连续.（复旦大学）
\[
  \left\langle \frac{\partial^2z}{\partial u\partial v}=0\right\rangle
\]

\noindent\textbf{提示}\quad 可用例 6.2.10 中的方法.

\noindent\textbf{注}\quad 作为验证，用逆变换 $x=\dfrac{u+v}{2}$，$y=\dfrac1{16}(v-u)^2$ 很快可将方程 $\dfrac{\partial^2z}{\partial u\partial v}=0$ 变回原方程.

\noindent\textbf{\ensuremath{\star}6.2.16}\quad 设 $z=f(x,y)$ 是二次连续可微函数，又有关系式 $u=x+ay$，$v=x-ay$（$a$ 是不为零的常数），证明：
\[
  a^2\frac{\partial^2z}{\partial x^2}-\frac{\partial^2z}{\partial y^2}
  =4a^2\frac{\partial^2z}{\partial u\partial v}.
\]
（北京大学）

\noindent\textbf{提示}\quad 可用逆变换 $x=\dfrac12(u+v)$，$y=\dfrac1{2a}(u-v)$，将欲证的等式右端变换成左端；也可用 $u=x+ay$，$v=x-ay$ 将左端变换成右端.

\noindent\textbf{6.2.17}\quad 设 $u=f(r)$，$r=\sqrt{x_1^2+x_2^2+\cdots+x_n^2}$，证明：
\[
  \frac{\partial^2u}{\partial x_1^2}+\cdots+\frac{\partial^2u}{\partial x_n^2}
  =\frac{d^2u}{dr^2}+\frac{n-1}{r}\frac{du}{dr}.
\]

\noindent\textbf{提示}\quad
\[
  \frac{\partial u}{\partial x_i}=f_r'\cdot\frac{x_i}{\sqrt{x_1^2+\cdots+x_n^2}},
\]
\[
  \frac{\partial^2u}{\partial x_i^2}
  =f_{rr}''\cdot\frac{x_i^2}{x_1^2+\cdots+x_n^2}
  +f_r'\cdot\left(\frac1r-\frac{x_i^2}{r^3}\right),
\]
代入等式左端可得右端.

\noindent\textbf{\ensuremath{\star}6.2.18}\quad 若 $u(x,y)$ 的二阶导数存在，证明：$u(x,y)=f(x)g(y)$ 的充要条件是
\[
  u\frac{\partial^2u}{\partial x\partial y}
  =\frac{\partial u}{\partial x}\frac{\partial u}{\partial y}\quad(u\ne0).
\]
（清华大学）

\noindent\textbf{提示}\quad 必要性可直接代入验证.（充分性）记 $v=\dfrac{\partial u}{\partial y}$，可将方程变形.

\noindent\textbf{再提示}\quad 这时方程变为
\[
  u\frac{\partial v}{\partial x}=v\frac{\partial u}{\partial x}
  \quad\text{或}\quad
  u\frac{\partial v}{\partial x}-v\frac{\partial u}{\partial x}=0,
\]
亦即
\[
  \frac{u\dfrac{\partial v}{\partial x}-v\dfrac{\partial u}{\partial x}}{u^2}=0,
\]
即
\[
  \frac{\partial}{\partial x}\left(\frac vu\right)=0,
\]
故 $\dfrac vu=h(y)$（与 $x$ 无关）. 于是，$\dfrac{\partial u}{\partial y}=v=uh(y)$，表明 $(\ln u)_y'=h(y)$. 因此
\[
  \ln u=\int h(y)\,dy+c(x),
\]
得
\[
  u=e^{\int h(y)\,dy+c(x)}=f(x)g(y).
\]

\noindent\textbf{6.2.19}\quad 以 $u=\dfrac yx$，$v=xy$ 作自变量，$w=x+y+z$ 作函数，变换方程
\[
  x^2\frac{\partial^2z}{\partial x^2}
  +2xy\frac{\partial^2z}{\partial x\partial y}
  +y^2\frac{\partial^2z}{\partial y^2}=0.
\]
（中南大学）
\[
  \left\langle 2v^2\frac{\partial^2w}{\partial v^2}+v\frac{\partial w}{\partial v}=0\right\rangle
\]

\noindent\textbf{提示}\quad 可参看例 6.2.11--6.2.13.

\subsection*{多元函数可微性}

\noindent\textbf{6.2.20}\quad 设
\[
  f(x,y)=
  \begin{cases}
  \dfrac{x^2y^2}{(x^2+y^2)^{3/2}}, & x^2+y^2\ne0,\\
  0, & x^2+y^2=0.
  \end{cases}
\]
求证：在 $(0,0)$ 处，$f(x,y)$ 连续但不可微.

\noindent\textbf{提示}\quad 连续性：
\[
  |f(x,y)-f(0,0)|\le\frac14\sqrt{x^2+y^2}\to0;
\]
$(0,0)$ 处不可微最后归结为证明
\[
  \lim_{(x,y)\to(0,0)}\frac{x^2y^2}{(x^2+y^2)^2}
\]
不存在.

\noindent\textbf{注}\quad 类似考研题可举出很多，但解法都一样. 参看例 6.2.15--6.2.17.

\noindent\textbf{\ensuremath{\star}6.2.21}\quad 确定 $\alpha$ 的值使得函数
\[
  f(x,y)=
  \begin{cases}
  (x^2+y^2)^\alpha\sin\dfrac1{x^2+y^2}, & (x,y)\ne(0,0),\\
  0, & (x,y)=(0,0)
  \end{cases}
\]
在 $(0,0)$ 处可微.（同济大学）

\noindent\textbf{提示}\quad 证明 $f$ 在 $(0,0)$ 处可微必须 $\alpha>\dfrac12$，而 $\alpha>\dfrac12$，故 $f$ 在 $(0,0)$ 处确实可微.

\noindent\textbf{再提示}\quad $f$ 在 $(0,0)$ 处可微，故 $f_x',f_y'$ 在 $(0,0)$ 处存在：
\[
  f_x'(0,0)=\lim_{x\to0}x^{2\alpha-1}\sin\frac1{x^2}\text{ 要存在，必有 }\alpha>\frac12.
\]

\noindent 反之，若 $\alpha>\dfrac12$，则 $f_x'(0,0)=f_y'(0,0)=0$.
\[
  \left|
  \frac{\Delta f-f_x'(0,0)\Delta x-f_y'(0,0)\Delta y}{\rho}
  \right|
  =
  \left|
  \frac{(x^2+y^2)^\alpha\sin\dfrac1{x^2+y^2}}{\sqrt{x^2+y^2}}
  \right|
  \le |x^2+y^2|^{\alpha-\frac12}\to0\quad(x\to0,y\to0).
\]
故 $f$ 在 $(0,0)$ 处可微.

\noindent\textbf{6.2.22}\quad 设
\[
  f(x,y)=
  \begin{cases}
  g(x,y)\sin\dfrac1{\sqrt{x^2+y^2}}, & (x,y)\ne(0,0),\\
  0, & (x,y)=(0,0),
  \end{cases}
\]
证明：

1）若 $g(0,0)=0,g(x,y)$ 在 $(0,0)$ 处可微，且 $dg(0,0)=0$，则 $f$ 在 $(0,0)$ 处可微，且 $df(0,0)=0$；

2）若 $g$ 在 $(0,0)$ 有偏导数，且 $f$ 在 $(0,0)$ 处可微，则 $df(0,0)=0$.（武汉大学）

\noindent\textbf{提示}\quad 1）先证明 $f_x'(0,0)=f_y'(0,0)=0$，再证
\[
  \frac1\rho[\Delta z-f_x'(0,0)\Delta x-f_y'(0,0)\Delta y]\to0\quad(\rho\to0).
\]

2）任务在于证明 $f_x'(0,0)=f_y'(0,0)=0$.

\noindent\textbf{再提示}\quad 1）由 $dg(0,0)=0$，得 $g_x'(0,0)=g_y'(0,0)=0$. 又因 $g$ 在 $(0,0)$ 处可微，故
\[
  g(x,y)-g(0,0)-g_x'(0,0)\Delta x-g_y'(0,0)\Delta y=g(x,y)=o(\rho)\quad(\rho\to0).
\]
从而
\[
  f_x'(0,0)=\lim_{x\to0}\frac{f(x,0)-0}{x}
  =\lim_{x\to0}\frac{g(x,0)\sin\dfrac1x}{x}=0.
\]
同理，$f_y'(0,0)=0$. 于是
\[
  \frac1\rho[\Delta f-f_x'(0,0)\Delta x-f_y'(0,0)\Delta y]
  =\frac1\rho g(x,y)\sin\frac1{\sqrt{x^2+y^2}}\to0\quad(\rho\to0).
\]

2）$g$ 在 $(0,0)$ 处有偏导数，故 $g$ 在 $(0,0)$ 处对 $x$ 连续，$g(x,0)\to g(0,0)$（当 $x\to0$ 时）. 若
\[
  \lim_{x\to0}g(x,0)=A\ne0,
\]
则
\[
  \sin\frac1{\sqrt{x^2+0^2}}=\frac{f(x,0)}{g(x,0)}\to\frac0A=0
\]
（当 $x\to0$ 时）（因 $f$ 在 $(0,0)$ 可微，故 $f$ 对 $x$ 在 $0$ 处连续，$f(x,0)\to0$），矛盾. 故
\[
  \lim_{x\to0}g(x,0)=0=g(0,0).
\]
由此可推知 $g_x'(0,0)=0$（因已知 $g_x'(0,0)$ 存在，若 $g_x'(0,0)=\alpha\ne0$，则
\[
  \sin\frac1{\sqrt{x^2+0^2}}
  =\frac{f(x,0)}{g(x,0)}
  =\frac{\dfrac{f(x,0)-0}{x}}{\dfrac{g(x,0)-0}{x}}\to\frac{f_x'(0,0)}{\alpha}
\]
当 $x\to0$ 时，矛盾）. 进而得
\[
  f_x'(0,0)=\lim_{x\to0}\frac{g(x,0)\sin\dfrac1{\sqrt{x^2}}}{x}=0.
\]
同理可证 $f_y'(0,0)=0$. 于是 $df(0,0)=0$.

\noindent\textbf{\ensuremath{\star}6.2.23}\quad 设函数 $g(x,y)$ 在 $(x_0,y_0)$ 处可微，$g(x_0,y_0)=0$，且 $\exists M>0$，使得 $|g(x,y)|\le M\rho$（在 $(x_0,y_0)$ 的某个邻域内），其中
\[
  \rho=\sqrt{(x-x_0)^2+(y-y_0)^2}.
\]
试证：任一函数 $f(x,y)$，若
\[
  \lim_{(x,y)\to(x_0,y_0)}f(x,y)=A
\]
存在，则 $z=f(x,y)g(x,y)$ 在 $(x_0,y_0)$ 处可微.（仿武汉大学试题）

\noindent\textbf{提示}\quad 先证 $z_x'(x_0,y_0)=Ag_x'(x_0,y_0)$，$z_y'(x_0,y_0)=Ag_y'(x_0,y_0)$，然后由 $g(x,y)$ 在 $(x_0,y_0)$ 处可微，导出 $f(x,y)g(x,y)$ 在 $(x_0,y_0)$ 处可微.

\noindent\textbf{再提示}\quad
\[
  z_x'(x_0,y_0)=\lim_{x\to x_0}
  \frac{f(x,y_0)g(x,y_0)-0}{x-x_0}
  =\lim_{x\to x_0}f(x,y_0)
  \left(\frac{g(x,y_0)-0}{x-x_0}\right)
  =Ag_x'(x_0,y_0).
\]
同理 $z_y'(x_0,y_0)=Ag_y'(x_0,y_0)$. 这时有 $\Delta z=f(x,y)g(x,y)=(A+\alpha)g(x,y)$，其中 $\alpha\to0$（当 $\rho\to0$ 时）. 于是
\[
  \frac1\rho[\Delta z-z_x'(x_0,y_0)\Delta x-z_y'(x_0,y_0)\Delta y]
\]
\[
  =\frac1\rho[(A+\alpha)g(x,y)-Ag_x'(x_0,y_0)\Delta x-Ag_y'(x_0,y_0)\Delta y]
\]
\[
  =\frac A\rho[g(x,y)-g_x'(x_0,y_0)\Delta x-g_y'(x_0,y_0)\Delta y]
  +\frac\alpha\rho g(x,y)\to0\quad(\rho\to0).
\]
由 $g$ 在 $(x_0,y_0)$ 处可微，第一项 $\to0$；由 $|g(x,y)|\le M\rho$ 知第二项 $\to0$（当 $\rho\to0$ 时）. 故 $z=f(x,y)$ 在 $(x_0,y_0)$ 处可微.

\noindent\textbf{\ensuremath{\star}6.2.24}\quad 设 $f_x'',f_y'$ 在 $(x_0,y_0)$ 的某个邻域里存在，在 $(x_0,y_0)$ 的某个空心邻域里 $f_{xy}''$ 存在，且
\[
  \lim_{(x,y)\to(x_0,y_0)}f_{xy}''(x,y)
\]
存在，试证：$f_{xy}''$ 在 $(x_0,y_0)$ 处连续，$f_{yx}''(x_0,y_0)$ 存在，且 $f_{yx}''(x_0,y_0)=f_{xy}''(x_0,y_0)$.

\noindent\textbf{提示}\quad 根据例 6.2.8，只需两混合偏导数 $f_{xy}'',f_{yx}''$ 其中一个在 $(x_0,y_0)$ 处连续即可. 因已知
\[
  \lim_{(x,y)\to(x_0,y_0)}f_{xy}''(x,y)=A
\]
存在，故只需证明 $f_{xy}''(x_0,y_0)$ 存在，且等于 $A$.

\noindent\textbf{再提示}\quad
\[
  f_{xy}''(x_0,y_0)=\lim_{y\to y_0}
  \frac{f_x'(x_0,y)-f_x'(x_0,y_0)}{y-y_0}.
\]
利用 Lagrange 微分中值定理，$\exists\xi$ 在 $y$ 与 $y_0$ 之间，使得
\[
  f_x'(x_0,y)-f_x'(x_0,y_0)=f_{xy}''(x_0,\xi)(y-y_0),
\]
因此
\[
  \text{上式}=\lim_{\xi\to y_0}f_{xy}''(x_0,\xi)=A.
\]

\section*{§6.3\quad 多元 Taylor 公式・凸函数・几何应用・极值}

\noindent\textbf{导读}\quad 本节中的几何应用与极值两部分是各类考试热点，适合各类读者. 多元 Taylor 公式及凸函数主要适合数学院系学生，非数学院系学生不作太多要求.

\subsection*{*一、多元 Taylor 公式}

这里只讨论 Taylor 公式的唯一性及 Taylor 公式的某些应用. 求初等函数的展开，一般不感觉困难，此处从略.

\noindent\textbf{例 6.3.1（Taylor 公式的唯一性）}\quad 假设 $f(x,y)$ 具有 $n+1$ 阶连续偏导数，若用某种方法得到展开式
\[
  f(x,y)=\sum_{i+j=0}^{n}A_{ij}(x-x_0)^i(y-y_0)^j+o(\rho^n), \tag{1}
\]
其中 $\rho=\sqrt{(x-x_0)^2+(y-y_0)^2}$，则必有
\[
  A_{ij}=\frac{C_{i+j}^i}{(i+j)!}
  \left.
  \frac{\partial^{i+j}}{\partial x^i\partial y^j}f(x,y)
  \right|_{(x_0,y_0)}
  =\frac1{i!\,j!}
  \frac{\partial^{i+j}}{\partial x^i\partial y^j}f(x_0,y_0).
\]

（多元的情况有类似结论.）

\noindent\textbf{证}\quad 已知 $f(x,y)$ 有 $n+1$ 阶连续偏导数，故 $f(x,y)$ 的 Taylor 公式成立：
\[
  f(x,y)=\sum_{i+j=0}^{n}
  \frac{C_{i+j}^i}{(i+j)!}
  \left(\frac{\partial^{i+j}}{\partial x^i\partial y^j}f(x_0,y_0)\right)
  (x-x_0)^i(y-y_0)^j+o(\rho^n). \tag{2}
\]
式（1）减式（2），便得 0 函数的展开式：
\[
  0=\sum_{i+j=0}^{n}B_{ij}(x-x_0)^i(y-y_0)^j+o(\rho^n), \tag{3}
\]
其中
\[
  B_{ij}=A_{ij}-\frac{C_{i+j}^i}{(i+j)!}
  \frac{\partial^{i+j}}{\partial x^i\partial y^j}f(x_0,y_0).
\]
因此，我们只要由式（3）推出 $B_{ij}=0\ (i+j=0,1,\cdots,n)$ 即可. 作变量替换 $\xi=x-x_0,\eta=y-y_0$. 对于新变量 $(\xi,\eta)$，式（3）变成
\[
  0=\sum_{i+j=0}^{n}B_{ij}\xi^i\eta^j+o(\rho^n)
  \quad(\rho=\sqrt{\xi^2+\eta^2}).
\]
为了照顾习惯，仍把 $(\xi,\eta)$ 记作 $(x,y)$. 于是问题化为由式
\[
  0=\sum_{i+j=0}^{n}B_{ij}x^iy^j+o(\rho^n)
  \quad(\rho=\sqrt{x^2+y^2}) \tag{4}
\]
证明
\[
  B_{ij}=0\quad(i+j=0,1,2,\cdots,n;\ i,j\text{ 为非负整数}).
\]

首先，在式（4）中，令 $\rho\to0$，便得 $B_{00}=0$. 然后令 $y=\alpha x$，则式（4）变成
\[
  \sum_{i+j=1}^{n}\alpha^jB_{ij}x^{i+j}+o(x^n)=0. \tag{5}
\]
设 $x\ne0$，用 $x$ 除此式，令 $x\to0$，得 $B_{10}+\alpha B_{01}=0$. 因 $\alpha$ 为任意实数，故知 $B_{10}=B_{01}=0$.

式（5）成为
\[
  \sum_{i+j=2}^{n}\alpha^jB_{ij}x^{i+j}+o(x^n)=0. \tag{6}
\]
同样，式（6）除以 $x^2$，令 $x\to0$，得 $B_{20}+\alpha B_{11}+\alpha^2B_{02}=0$. 由 $\alpha$ 的任意性，可知 $B_{20}=B_{11}=B_{02}=0$. 从而式（6）变成
\[
  \sum_{i+j=3}^{n}\alpha^jB_{ij}x^{i+j}+o(x^n)=0.
\]
如此继续下去，可得一切 $B_{ij}=0\ (i+j=0,1,2,\cdots,n)$. 证毕.

\noindent\textbf{注}\quad 有了唯一性，求 Taylor 公式展开式，不一定要用求导数的方法，只要余项是 $\rho^n$ 的高级无穷小量，所得的展开式必是 Taylor 公式的展开式（见本节后面的有关练习）.

下面两例说明 Taylor 公式的某些应用.

\noindent\textbf{例 6.3.2}\quad 设 $D\subseteq\mathbb{R}^n$ 为凸的有界闭区域，$f(P)$ 在 $D$ 上有连续的一阶偏导数. 试证：$f(P)$ 在 $D$ 上满足 Lipschitz 条件. 即：$\exists L>0,\forall P,P_1\in D$，有
\[
  |f(P)-f(P_1)|\le L|P-P_1|. \tag{1}
\]

\noindent\textbf{证}\quad 根据已知条件可知：$\exists M>0$，使得 $|f_{x_i}'(P)|\le M,\forall P\in D,i=1,2,\cdots,n$. 因 $D$ 为凸区域，由 Taylor 公式，$\forall P,P_1\in D,\exists P^*\in\overline{PP_1}\subset D$，使得
\[
  |f(P)-f(P_1)|
  =\left|\sum_{i=1}^{n}\frac{\partial f}{\partial x_i}(P^*)(x_i-x_{1i})\right|
\]
\[
  \le\sum_{i=1}^{n}\left|\frac{\partial f}{\partial x_i}(P^*)\right||x_i-x_{1i}|
  \le Mn\rho(P,P_1).
\]
这里 $P$ 和 $P_1$ 分别为 $P(x_1,x_2,\cdots,x_n)$ 和 $P_1(x_{11},x_{12},\cdots,x_{1n})$. 令 $L=Mn$，则得式（1）. 由本例可知，在有界凸区域上函数 $f$ 的一阶偏导数连续有界，则 $f$ 在此区域上一致连续.

\noindent\textbf{\ensuremath{\star}例 6.3.3}\quad 设 $F(x,y,z)$ 在 $\mathbb{R}^3$ 中有连续的一阶偏导数 $\dfrac{\partial F}{\partial x},\dfrac{\partial F}{\partial y},\dfrac{\partial F}{\partial z}$，并满足
\[
  y\frac{\partial F}{\partial x}-x\frac{\partial F}{\partial y}
  +\frac{\partial F}{\partial z}\ge\alpha>0,\quad \forall(x,y,z)\in\mathbb{R}^3,
\]
其中 $\alpha$ 为常数. 试证明：当 $(x,y,z)$ 沿着曲线 $\Gamma:x=-\cos t,y=\sin t,z=t,t\ge0$ 趋向无穷远时，$F(x,y,z)\to+\infty$.（北京大学）

\noindent\textbf{方法}\quad 利用推导多元 Taylor 公式的方法，对函数 $\Phi(t)=F(-\cos t,\sin t,t)$ 应用一元 Taylor 公式：
\[
  \Phi(t)=\Phi(0)+\Phi'(\tau)t.
\]

\noindent\textbf{证}\quad 对曲线 $\Gamma$ 上的点 $(x,y,z)=(-\cos t,\sin t,t)\in\Gamma$，有
\[
  F(x,y,z)=F(-\cos t,\sin t,t)
  =F(-1,0,0)+[F(-\cos t,\sin t,t)]_{t=\tau}'\cdot t.
\]
记 $\beta=F(-1,0,0)$，$\tau$ 对应的点为 $Q:Q=(-\cos\tau,\sin\tau,\tau)=(\xi,\eta,\zeta)$，则
\[
  [F(-\cos t,\sin t,t)]_{t=\tau}'
  =\sin\tau\left.\frac{\partial F}{\partial x}\right|_Q
  +\cos\tau\left.\frac{\partial F}{\partial y}\right|_Q
  +\left.\frac{\partial F}{\partial z}\right|_Q
\]
\[
  =\eta\left.\frac{\partial F}{\partial x}\right|_Q
  -\xi\left.\frac{\partial F}{\partial y}\right|_Q
  +\left.\frac{\partial F}{\partial z}\right|_Q
  \ge\alpha>0.
\]
于是有
\[
  F(x,y,z)\ge\beta+\alpha t\to+\infty\quad(\text{当 }t\to+\infty\text{ 时}).
\]
证毕.

\subsection*{二、凸函数}

作为 Taylor 公式的一个应用，我们来研究凸函数.

\noindent\textbf{定义}\quad 区域 $D\subseteq\mathbb{R}^n$ 称为凸的，当且仅当 $\forall x,y\in D,\forall\lambda\in[0,1]$，有
\[
  \lambda x+(1-\lambda)y\in D.
\]
$y=f(x)$ 称为凸区域 $D$ 上的凸函数，当且仅当 $\forall x,y\in D,\forall\lambda\in[0,1]$，有
\[
  f[\lambda x+(1-\lambda)y]\le \lambda f(x)+(1-\lambda)f(y).
\]
（若“$\le$”换成“$<$”，则 $f$ 称为严格凸的）.

其几何意义如图 6.3.1 所示.
\begin{center}
  \includegraphics[width=0.42\textwidth]{6-3-1.png}

  图 6.3.1
\end{center}

\noindent\textbf{定理 1}\quad 设 $f$ 在凸区域 $D$ 上定义并且有连续的一阶偏导数，则 $f$ 在 $D$ 内为凸函数的充要条件是：$\forall x,y\in D$，有
\[
  f(y)\ge f(x)+(y-x)\nabla f(x), \tag{1}
\]
其中 $x=(x_1,x_2,\cdots,x_n),y=(y_1,y_2,\cdots,y_n)$，
\[
  \nabla f(x)=\left(\frac{\partial f}{\partial x_1},\frac{\partial f}{\partial x_2},\cdots,\frac{\partial f}{\partial x_n}\right),
\]
\[
  (y-x)\nabla f(x)=\frac{\partial f}{\partial x_1}(y_1-x_1)
  +\frac{\partial f}{\partial x_2}(y_2-x_2)+\cdots+
  \frac{\partial f}{\partial x_n}(y_n-x_n). \tag{2}
\]

\noindent\textbf{证}\quad 1$^\circ$（必要性）由于 $f$ 在 $D$ 上为凸函数，故 $\forall x,y\in D,\forall\lambda\in[0,1]$，有
\[
  f[\lambda y+(1-\lambda)x]\le \lambda f(y)+(1-\lambda)f(x),
\]
即
\[
  f[x+\lambda(y-x)]-f(x)\le\lambda f(y)-\lambda f(x). \tag{3}
\]
因 $f$ 有连续的一阶导数，故 $f$ 可微：
\[
  f[x+\lambda(y-x)]-f(x)
  =\left(\frac{\partial}{\partial x_1}f(x)\right)\lambda(y_1-x_1)+\cdots+
  \left(\frac{\partial}{\partial x_n}f(x)\right)\lambda(y_n-x_n)
\]
\[
  +\varepsilon_1\lambda(y_1-x_1)+\cdots+\varepsilon_n\lambda(y_n-x_n), \tag{4}
\]
其中 $\varepsilon_1,\cdots,\varepsilon_n\to0$（当 $\lambda\to0$ 时）. 将式（4）代入式（3），令 $\lambda\ne0$，以 $\lambda$ 同除式（3）两端，再令 $\lambda\to0$，得
\[
  \left(\frac{\partial}{\partial x_1}f(x)\right)(y_1-x_1)+\cdots+
  \left(\frac{\partial}{\partial x_n}f(x)\right)(y_n-x_n)
  \le f(y)-f(x).
\]
注意到式（2），此式左端即为 $(y-x)\nabla f(x)$. 故式（1）得证.

2$^\circ$（充分性）$\forall x,y\in D,\forall\lambda\in[0,1]$，记 $z=\lambda x+(1-\lambda)y\in D$，按已知条件，
\[
  f(x)\ge f(z)+(x-z)\nabla f(z), \tag{5}
\]
\[
  f(y)\ge f(z)+(y-z)\nabla f(z). \tag{6}
\]
将（5）、（6）分别乘 $\lambda$ 与 $(1-\lambda)$，然后相加，得
\[
  \lambda f(x)+(1-\lambda)f(y)
  \ge f(z)+[\lambda(x-z)+(1-\lambda)(y-z)]\nabla f(z)
  =f(z)=f[\lambda x+(1-\lambda)y],
\]
即
\[
  f[\lambda x+(1-\lambda)y]\le\lambda f(x)+(1-\lambda)f(y).
\]
这就证明了 $f$ 为凸函数.

\noindent\textbf{定理 2}\quad 设 $D\subseteq\mathbb{R}^n$ 为凸区域，$f(x)=f(x_1,\cdots,x_n)$ 在 $D$ 上定义，有连续的二阶偏导数，证明 $f(x)$ 在 $D$ 上为凸函数的充要条件是 Hesse 矩阵 $\left(\dfrac{\partial^2 f}{\partial x_i\partial x_j}\right)_{i,j=1}^{n}$ 在 $D$ 上为半正定的.（浙江大学）

\noindent\textbf{证}\quad 1$^\circ$（充分性）$\forall x,y\in D$，根据 Taylor 公式，$\exists\xi=x+\theta(y-x)\ (0<\theta<1)$，使得
\[
  f(y)=f(x)+(y-x)\nabla f(x)+\frac12
  \left[(y_1-x_1)\frac{\partial}{\partial x_1}+\cdots+
  (y_n-x_n)\frac{\partial}{\partial x_n}\right]^2 f(\xi). \tag{1}
\]
注意到
\[
  \left[(y_1-x_1)\frac{\partial}{\partial x_1}+\cdots+
  (y_n-x_n)\frac{\partial}{\partial x_n}\right]^2f(\xi)
  =\sum_{i,j=1}^{n}(y_i-x_i)(y_j-x_j)
  \frac{\partial^2f(\xi)}{\partial x_i\partial x_j}
\]
\[
  =(y_1-x_1,y_2-x_2,\cdots,y_n-x_n)
  \left(\frac{\partial^2f(\xi)}{\partial x_i\partial x_j}\right)
  \begin{pmatrix}
  y_1-x_1\\ y_2-x_2\\ \vdots\\ y_n-x_n
  \end{pmatrix}. \tag{2}
\]
若矩阵 $\left(\dfrac{\partial^2f}{\partial x_i\partial x_j}\right)(i,j=1,2,\cdots,n)$ 在 $D$ 上为半正定的，则式（2）非负，式（1）成为
\[
  f(y)\ge f(x)+(y-x)\nabla f(x).
\]
从而 $f$ 在 $D$ 上为凸函数（定理 1）.

2$^\circ$（必要性）用反证法. 假设 $\left(\dfrac{\partial^2f}{\partial x_i\partial x_j}\right)$ 为非半正定的，则 $\exists x\in D$ 及 $h=(h_1,\cdots,h_n)$，使得
\[
  (h_1,\cdots,h_n)
  \left(\frac{\partial^2f(x)}{\partial x_i\partial x_j}\right)
  \begin{pmatrix}h_1\\ \vdots\\ h_n\end{pmatrix}<0. \tag{3}
\]
另一方面，由 Taylor 公式，当 $\lambda\to0$ 时，
\[
  f(x+\lambda h)=f(x)+\lambda h\nabla f(x)
  +\lambda^2\left[
  \frac12(h_1,\cdots,h_n)
  \left(\frac{\partial^2f(x)}{\partial x_i\partial x_j}\right)
  \begin{pmatrix}h_1\\ \vdots\\ h_n\end{pmatrix}
  +o(1)\right].
\]
由式（3），当 $\lambda$ 充分小时，此式右端第三项为负，于是
\[
  f(x+\lambda h)\le f(x)+\lambda h\nabla f(x),
\]
与 $f$ 的凸性矛盾（见定理 1）.

\subsection*{三、几何应用}

\noindent\textbf{要点}\quad 空间曲线 $x=x(t),y=y(t),z=z(t)$ 的切向量为 $(x'(t),y'(t),z'(t))$；曲面 $F(x,y,z)=0$ 的法向量为 $(F_x',F_y',F_z')$. 让流动向量与之平行或垂直，就可写出空间曲线切线与法平面、空间曲面的法线与切平面的方程，和解决与之有关的问题.

\noindent\textbf{例 6.3.4}\quad 求 $x=r\cos\varphi,y=r\sin\varphi,z=r\cot\alpha$ 在点 $M_0(\varphi_0,r_0)$ 处的切面与法线（其中 $\alpha$ 为常数）.

\noindent\textbf{解}\quad $r=r_0$ 对应的曲线为 $x=r_0\cos\varphi,y=r_0\sin\varphi,z=r_0\cot\alpha$. 它在 $M_0$ 的切向量为
\[
  \tau_1=r_0(-\sin\varphi_0,\cos\varphi_0,0).
\]
类似可得 $\varphi=\varphi_0$ 曲线在 $M_0$ 的切向量
\[
  \tau_2=(\cos\varphi_0,\sin\varphi_0,\cot\alpha).
\]
从而曲面在 $M_0$ 的法向量为
\[
  n=\tau_1\times\tau_2
  =r_0(\cos\varphi_0\cot\alpha,\sin\varphi_0\cot\alpha,-1).
\]
由此可得切平面为
\[
  x\cos\varphi_0+y\sin\varphi_0-z\tan\alpha=0.
\]
法线为
\[
  \frac{x-r_0\cos\varphi_0}{\cos\varphi_0}
  =\frac{y-r_0\sin\varphi_0}{\sin\varphi_0}
  =\frac{z-r_0\cot\alpha}{-\tan\alpha}.
\]

\noindent\textbf{\ensuremath{\star}例 6.3.5}\quad 证明：若 $F(u,v)$ 有连续偏导数，则曲面 $S:F(nx-lz,ny-mz)=0$ 上任意一点的切平面都平行于直线
\[
  L:\frac xl=\frac ym=\frac zn.
\]
（东北师范大学）

\noindent\textbf{证}\quad 曲面 $S$ 上任意一点 $(x_0,y_0,z_0)$ 的法向量与切平面分别为
\[
  n=(nF_u',nF_v',-lF_u'-mF_v'),
\]
\[
  nF_u'(x-x_0)+nF_v'(y-y_0)-(lF_u'+mF_v')(z-z_0)=0.
\]
直线 $L$ 的方向数为 $(l,m,n)$，
\[
  n\cdot(l,m,n)=l\cdot nF_u'+m\cdot nF_v'
  +n(-lF_u'-mF_v')=0.
\]
因此该直线与 $n$ 垂直，故 $L$ 与任意一点的切平面平行.

\noindent\textbf{*例 6.3.6}\quad 设 $D$ 为凸的有界闭区域，曲面的方程为 $z=f(x,y),(x,y)\in D$. $f(x,y)$ 在 $D$ 上有有界的二阶导函数. 令 $\varphi=\varphi(P_1,P)$ 表示曲面在 $P(x,y)\in D,P_1(x_1,y_1)\in D$ 两点法线之间夹角. 试证：当 $P$ 与 $P_1$ 充分接近时，$\varphi(P_1,P)$ 满足 A. M. Lyapunov 不等式：
\[
  \varphi(P_1,P)\le c\rho(P_1,P),\quad P_1,P\in D,
\]
其中 $c$ 为常数，$\rho(P_1,P)$ 表示 $P_1$ 与 $P$ 之间的距离.

\noindent\textbf{分析}\quad 我们只要证明 $\varphi\le\dfrac\pi2\sin\varphi\le c\rho(P_1,P)$ 即可. 已知 $0\le\varphi\le\dfrac\pi2$ 时，第一个不等式成立（$\left(0,\dfrac\pi2\right)$ 内 $\dfrac{\sin x}x\ge\dfrac2\pi$）. 现只需证明第二个不等式.

\noindent\textbf{证}\quad 法向量为
\[
  n(P)=(f_x'(P),f_y'(P),-1),\quad
  n(P_1)=(f_x'(P_1),f_y'(P_1),-1).
\]
可见 $|n(P)|\ge1,\ |n(P_1)|\ge1$. 故
\[
  \sin^2\varphi
  =\frac{|n(P_1)\times n(P)|^2}{|n(P_1)|^2|n(P)|^2}
  \le |n(P_1)\times n(P)|^2
\]
\[
  =
  \bigl(f_y'(P_1)-f_y'(P)\bigr)^2+
  \bigl(f_x'(P_1)-f_x'(P)\bigr)^2+
  \bigl(f_x'(P_1)f_y'(P)-f_x'(P)f_y'(P_1)\bigr)^2. \tag{1}
\]
由已知条件知，$\exists M>0$ 使得 $D$ 内有 $|f_x'|,|f_y'|,|f_{xx}''|,|f_{xy}''|,|f_{yy}''|\le M$. 利用中值定理，$\exists P^*\in\overline{PP_1}$，使得
\[
  |f_x'(P_1)-f_x'(P)|
  \le2M\rho(P_1,P),
\]
同理 $|f_y'(P_1)-f_y'(P)|\le2M\rho(P_1,P)$. 故
\[
  |f_y'(P_1)f_x'(P)-f_x'(P)f_y'(P_1)|
  \le4M^2\rho(P_1,P).
\]
由式（1），
\[
  \sin^2\varphi\le(1+2M^2)8M^2\rho^2(P_1,P).
\]
可见当 $\rho(P_1,P)$ 充分小时，$0\le\varphi\le\dfrac\pi2$，从而
\[
  0\le\varphi(P_1,P)\le\frac\pi2\sin\varphi\le c\rho(P_1,P),
\]
其中 $c=\sqrt2\pi M\sqrt{1+2M^2}$ 为常数.

\noindent\textbf{*例 6.3.7}\quad 从原点向单叶双曲面
\[
  \frac{x^2}{a^2}+\frac{y^2}{b^2}-\frac{z^2}{c^2}=1
\]
的切平面引垂线，求垂足的轨迹.

\noindent\textbf{解}\quad 所谓垂足，即切平面与垂线的交点. 曲面上任意一点 $(x_1,y_1,z_1)$ 的切平面为
\[
  \frac{x_1x}{a^2}+\frac{y_1y}{b^2}-\frac{z_1z}{c^2}=1. \tag{1}
\]
过原点向此切平面引的垂线为
\[
  \frac{a^2x}{x_1}=\frac{b^2y}{y_1}=-\frac{c^2z}{z_1}. \tag{2}
\]
我们的问题是：当 $(x_1,y_1,z_1)$ 沿单叶双曲面
\[
  \frac{x_1^2}{a^2}+\frac{y_1^2}{b^2}-\frac{z_1^2}{c^2}=1 \tag{3}
\]
移动时，求方程式（1）、（2）所决定的垂足 $(x,y,z)$ 的轨迹. 因此，只要在（1）、（2）、（3）中消去 $x_1,y_1,z_1$，求出 $(x,y,z)$ 满足的方程即可. 令式（2）等于 $\dfrac1k$，则得
\[
  x_1=ka^2x,\quad y_1=kb^2y,\quad z_1=-kc^2z.
\]
代入（1）、（3），得
\[
  k(x^2+y^2+z^2)=1,\quad
  k^2(a^2x^2+b^2y^2-c^2z^2)=1.
\]
从而
\[
  (x^2+y^2+z^2)^2=a^2x^2+b^2y^2-c^2z^2,
\]
即为所求.

\noindent\textbf{*例 6.3.8}\quad 设 $a>b>c>0$ 为三个正数，试证 $\mathbb{R}^3$ 中任意一点 $M(x,y,z)$ 处有三个二次曲面
\[
  \frac{x^2}{a^2-\lambda_i^2}
  +\frac{y^2}{b^2-\lambda_i^2}
  +\frac{z^2}{c^2-\lambda_i^2}=-1
  \quad(i=1,2,3) \tag{1}
\]
（其中 $\lambda_1,\lambda_2,\lambda_3$ 为三个彼此不同的实数），它们通过点 $M$，并在点 $M$ 相互正交.

\noindent\textbf{提示}\quad 要求 $\lambda_i(i=1,2,3)$ 使式（1）成立，即要求函数
\[
  F(\lambda^2)=x^2(b^2-\lambda^2)(c^2-\lambda^2)
  +y^2(a^2-\lambda^2)(c^2-\lambda^2)
\]
\[
  +z^2(a^2-\lambda^2)(b^2-\lambda^2)
  +(a^2-\lambda^2)(b^2-\lambda^2)(c^2-\lambda^2)
\]
的根. $F(\lambda^2)$ 为 $\lambda^2$ 的三次多项式，且在区间 $[c,b],[b,a],[a,+\infty)$ 端点上异号，因此有且仅有三个不同实根.

在 $(x,y,z)$ 处，三曲面的法向量
\[
  n_i=\left(\frac{2x}{a^2-\lambda_i^2},
  \frac{2y}{b^2-\lambda_i^2},
  \frac{2z}{c^2-\lambda_i^2}\right)\quad(i=1,2,3)
\]
相互正交，因为 $i\ne j$ 时，
\[
  n_i\cdot n_j=
  \frac4{\lambda_i^2-\lambda_j^2}
  \left[
  \left(\frac{x^2}{a^2-\lambda_i^2}
  +\frac{y^2}{b^2-\lambda_i^2}
  +\frac{z^2}{c^2-\lambda_i^2}\right)
  -\left(\frac{x^2}{a^2-\lambda_j^2}
  +\frac{y^2}{b^2-\lambda_j^2}
  +\frac{z^2}{c^2-\lambda_j^2}\right)\right]=0.
\]

\subsection*{\ensuremath{\star}四、极值}

\noindent\textbf{a.\quad 自由极值}

\noindent\textbf{要点}\quad 自由极值又称局部极值. $f$ 在点 $P_0$ 有极大（小）值，指函数 $f$ 在 $P_0$ 的某邻域里，恒有 $f(P_0)\ge f(P)$（或 $f(P_0)\le f(P)$）（将 $\ge(\le)$ 改为 $>(<)$，则称为严格极值）.

求自由极值的方法步骤：

1）求可疑点. 可疑点包括：i）稳定点（即一阶偏导数同时等于零的点）；ii）使至少某一阶偏导数不存在的点.

2）对可疑点进行判断. 基本方法是：i）用定义判断；ii）利用实际背景进行判断；iii）利用二阶导数：设 $P_0$ 为稳定点，若在 $P_0$ 处 Hesse 矩阵
\[
  H(P_0)=
  \begin{pmatrix}
  f_{x_1x_1}''&f_{x_1x_2}''&\cdots&f_{x_1x_n}''\\
  f_{x_2x_1}''&f_{x_2x_2}''&\cdots&f_{x_2x_n}''\\
  \vdots&\vdots&&\vdots\\
  f_{x_nx_1}''&f_{x_nx_2}''&\cdots&f_{x_nx_n}''
  \end{pmatrix}
\]
为正定的，则 $f$ 在 $P_0$ 处取极小值；若 $H(P_0)$ 为负定的，则 $f$ 在 $P_0$ 处取极大值；若 $H(P_0)$ 为不定的，则 $f$ 在 $P_0$ 处无极值. 具体到二元函数即是：若 $f_x'(x_0,y_0)=f_y'(x_0,y_0)=0$，记
\[
  \Delta=(f_{xx}''f_{yy}''-f_{xy}''{}^2)\bigm|_{(x_0,y_0)},
\]
则当 $\Delta>0$ 且 $f_{xx}''\bigm|_{(x_0,y_0)}>0$ 时取（严格）极小值；当 $\Delta>0$ 且 $f_{xx}''\bigm|_{(x_0,y_0)}<0$ 时取（严格）极大值；当 $\Delta<0$ 时无极值；当 $\Delta=0$ 时，情况待定.

多元极值有两个观念值得澄清，其一是：一元函数的极大值与极小值总是交替地出现，多元函数谈不上交替，甚至只有一种极值（无穷多个）. 如

\noindent\textbf{\ensuremath{\star}例 6.3.9}\quad 证明：函数
\[
  z=f(x,y)=(1+e^y)\cos x-ye^y
\]
有无穷多个极大值，但无极小值.（大连海事大学，中国人民大学）

\noindent\textbf{证}\quad $f_x'=-(1+e^y)\sin x,\quad f_y'=(\cos x-1-y)e^y$.

令 $f_x'=0,f_y'=0$，解方程，可得无穷多个稳定点
\[
  (x_n,y_n)=(n\pi,\cos n\pi-1)\quad(n=0,\pm1,\pm2,\cdots).
\]
当 $n$ 为偶数时，在 $(x_n,y_n)$ 内，
\[
  \Delta=f_{xx}''f_{yy}''-f_{xy}''{}^2=2>0,\quad f_{xx}''=-2<0.
\]
故 $f$ 在 $(2k\pi,0)$ 上取极大值 $(k=0,\pm1,\pm2,\cdots)$. 当 $n$ 为奇数时，在 $(x_n,y_n)$ 内，
\[
  \Delta=f_{xx}''f_{yy}''-f_{xy}''{}^2=-(1+e^{-2})e^{-2}<0,
\]
此处无极值. 总之，$f$ 有无穷多个极大值而无极小值.

另一个值得澄清的问题是：（以极小值为例）$f$ 在某点 $P_0$ 取极小值，是指 $f$ 在 $P_0$ 点的值比某邻域里其他点的值小. 假设在过点 $P_0$ 的每一直线上，$f$ 在 $P_0$ 取极小值，问是否能断言 $f$ 在 $P_0$ 处取极小值？回答是否定的. 如

\noindent\textbf{例 6.3.10}\quad $f(x,y)=(y-x^2)(y-2x^2)$，当限定 $(x,y)$ 在过 $(0,0)$ 的直线上时，$f$ 在 $(0,0)$ 处为极小，但作为二元函数，$f$ 在 $(0,0)$ 处无极值.

\noindent\textbf{证明}\quad 读者用定义很容易证明.

关于自由极值的求法，下面还会讲到.

\noindent\textbf{b.\quad 条件极值与 Lagrange 乘数法}

\noindent\textbf{要点}\quad 若 $y=f(x_1,\cdots,x_n)$ 及 $\varphi_i(x_1,\cdots,x_n)(i=1,2,\cdots,m;m<n)$ 有连续偏导数，且 Jacobi 矩阵
\[
  \frac{\partial(\varphi_1,\cdots,\varphi_m)}{\partial(x_1,\cdots,x_n)}
\]
的秩为 $r=m$（不妨设行列式 $\dfrac{\partial(\varphi_1,\cdots,\varphi_m)}{\partial(x_1,\cdots,x_m)}\ne0$），那么函数 $y=f(x_1,\cdots,x_n)$ 在条件 $\varphi_i(x_1,\cdots,x_n)=0(i=1,2,\cdots,m)$ 限制之下的极值点，可用 Lagrange 乘数法寻找.

具体做法是：首先作 Lagrange 函数：
\[
  L(x_1,\cdots,x_n)=f(x_1,\cdots,x_n)+\sum_{i=1}^{m}\lambda_i\varphi_i(x_1,\cdots,x_n),
\]
然后解方程组
\[
  L_{x_1}'=0,\quad L_{x_2}'=0,\quad\cdots,\quad L_{x_n}'=0,\quad
  \varphi_i(x_1,\cdots,x_n)=0\quad(i=1,2,\cdots,m),
\]
求出稳定点，这里 $\lambda_i$ 为待定常数，有时不一定要求出.

最后，对稳定点进行判别，常用的方法是

i）利用极值点的定义进行判别；

ii）利用实际背景进行判别；

iii）利用 Lagrange 函数的二阶微分进行判别. 若在某稳定点 $P_0$ 处取条件极小（大）值，$d^2L(P_0)>0(<0)$，则 $f$ 在此点 $P_0$ 取条件极小（大）值，其中
\[
  d^2L(P_0)=
  \left(\frac{\partial}{\partial x_1}dx_1+\cdots+
  \frac{\partial}{\partial x_n}dx_n\right)^2L(P_0)
  =\sum_{k,j=1}^{n}L_{x_kx_j}''(P_0)dx_kdx_j,
\]
$dx_i(i=1,2,\cdots,n)$ 应满足方程
\[
  \sum_{i=1}^{n}\frac{\partial\varphi_j}{\partial x_i}(P_0)dx_i=0
  \quad(j=1,2,\cdots,m).
\]

\noindent\textbf{\ensuremath{\star}例 6.3.11}\quad 求 $f(x,y,z)=x^4+y^4+z^4$ 在条件 $xyz=1$ 下的极值；最大值与最小值各为多少？为什么？

\noindent\textbf{解}\quad 记
\[
  L=x^4+y^4+z^4+\lambda(xyz-1).
\]
则
\[
  \begin{cases}
    L_x'=4x^3+\lambda yz=0,\\
    L_y'=4y^3+\lambda xz=0,\\
    L_z'=4z^3+\lambda xy=0,\\
    xyz=1.
  \end{cases}
\]
解得四个点
\[
  (1,1,1),\quad(-1,-1,1),\quad(-1,1,-1),\quad(1,-1,-1).
\]
在这四点处 $f=3$，它们是极小值点. 由对称性，只证 $P_1=(1,1,1)$ 是极小值点.

在曲面 $xyz=1$ 的第一卦限内，有
\[
  f=x^4+y^4+\frac{1}{x^4y^4}.
\]
在 $xOy$ 平面上，直线 $x=\frac14,x=2,y=\frac14,y=2$ 围成矩形 $ABCD$（图 6.3.2），在其边界上至少有一项 $\ge16$，故
\[
  f\left(x,y,\frac1{xy}\right)\ge16>3=f(P_1).
\]
因此，若 $f$ 取极小值，只能在矩形 $ABCD$ 内部取得. 而内部唯一的稳定点是 $(1,1)$，所以 $(1,1,1)$ 为极小值点.

\begin{center}
  \includegraphics[width=0.42\textwidth]{6-3-2.png}

  图 6.3.2
\end{center}

\noindent\textbf{解法 II}\quad 利用二阶微分判别. 在 $P_1=(1,1,1)$ 处，由上方程组得 $\lambda=-4$，且
\[
  L_{xx}''=L_{yy}''=L_{zz}''=12,\quad
  L_{xy}''=L_{yz}''=L_{zx}''=\lambda=-4.
\]
所以
\[
  d^2L(P_1)=12(dx^2+dy^2+dz^2)-8(dxdy+dydz+dzdx).
\]
又由 $xyz=1$ 得 $dz=-dx-dy$，代入上式得
\[
\begin{aligned}
  d^2L(P_1)
  &=12(dx^2+dy^2+dz^2)+4(dx^2+2dxdy+dy^2)+4dx^2+4dy^2\\
  &=12(dx^2+dy^2+dz^2)+4dz^2+4dx^2+4dy^2>0.
\end{aligned}
\]
于是 $P_1$ 为极小值点. 其余各点由对称性可得.

\noindent\textbf{练习 1}\quad 利用 Lagrange 乘数法证明：若 $a_k\ge0(k=1,2,\cdots,n)$，则
\[
  \sqrt[n]{a_1a_2\cdots a_n}\le
  \frac{a_1+a_2+\cdots+a_n}{n},
\]
等号成立当且仅当 $a_1=a_2=\cdots=a_n$.

几何上，不等式表明：周长一定的矩形中，正方形面积最大；表面积一定的长方体中，正方体体积最大；等等.

\noindent\textbf{证}\quad 记 $a=a_1+\cdots+a_n$，在条件 $a_1+\cdots+a_n=a$ 下，求 $V=a_1a_2\cdots a_n$ 的最大值. 令
\[
  L=a_1a_2\cdots a_n+\lambda\left(\sum_{k=1}^{n}a_k-a\right).
\]
则
\[
  L_{a_k}'=\prod_{i\ne k}a_i+\lambda=0\quad(k=1,2,\cdots,n).
\]
将第 $k$ 个方程乘以 $a_k$ 后相加，并利用约束条件，得
\[
  n\prod_{i=1}^{n}a_i+\lambda a=0,
\]
即
\[
  \lambda=-\frac{n}{a}\prod_{i=1}^{n}a_i.
\]
代回方程得
\[
  a_k=\frac{a}{n}=\frac{a_1+\cdots+a_n}{n}\quad(k=1,2,\cdots,n).
\]
故
\[
  a_1a_2\cdots a_n\le\left(\frac{a_1+\cdots+a_n}{n}\right)^n,
\]
且等号成立当且仅当 $a_1=a_2=\cdots=a_n$.

\noindent\textbf{练习 2}\quad 若 $u_1,u_2,\cdots,u_n\ge0,\ u_1u_2\cdots u_n=1$，证明
\[
  u_1+u_2+\cdots+u_n\ge n,
\]
且等号成立当且仅当 $u_1=u_2=\cdots=u_n=1$. 此不等式与上题等价.

\noindent\textbf{c.\quad 求函数在闭区域上的最大最小值}

\noindent\textbf{要点}\quad 求函数在闭区域上的最大最小值，可先求内部的极值，再求边界上的极值，最后比较各候选值.

\noindent\textbf{例 6.3.12}\quad 求函数 $f(x,y)=ax^2+2bxy+cy^2$ 在圆 $x^2+y^2\le1$ 上的最大值和最小值，其中 $b^2-ac>0,\ a,b,c>0$.

\noindent\textbf{解}\quad 在 $x^2+y^2<1$ 内，令
\[
  f_x'=2ax+2by=0,\quad f_y'=2bx+2cy=0.
\]
因 $ac-b^2\ne0$，只得 $(x,y)=(0,0)$，此时 $f=0$.

在边界 $x^2+y^2=1$ 上，令
\[
  L=ax^2+2bxy+cy^2-\lambda(x^2+y^2-1).
\]
则
\[
  \begin{cases}
    (a-\lambda)x+by=0,\\
    bx+(c-\lambda)y=0.
  \end{cases}
\]
因为 $(x,y)\ne(0,0)$，故
\[
  \begin{vmatrix}
    a-\lambda&b\\
    b&c-\lambda
  \end{vmatrix}=0,
\]
从而
\[
  \lambda_{1,2}=\frac{a+c\pm\sqrt{(a-c)^2+4b^2}}{2}.
\]
把上面两个方程分别乘以 $x,y$ 后相加，并利用 $x^2+y^2=1$，得 $f=\lambda$. 因此
\[
  \max f=\max\{0,\lambda_1,\lambda_2\}
  =\frac{a+c+\sqrt{(a-c)^2+4b^2}}{2},
\]
\[
  \min f=\min\{0,\lambda_1,\lambda_2\}
  =\frac{a+c-\sqrt{(a-c)^2+4b^2}}{2}.
\]

\noindent\textbf{例 6.3.13}\quad 求函数 $f(x,y)=4x+xy^2+y^2$ 在圆 $x^2+y^2\le1$ 上的最大值和最小值.（四川大学）

\noindent\textbf{解}\quad 因为 $f_x'=4+y^2>0$，故无内部极值，极值只能在圆周上取得. 当 $x^2+y^2=1$ 时，
\[
  f=1+5x-x^2-x^3\equiv\varphi(x),\quad -1\le x\le1.
\]
又
\[
  \varphi'(x)=5-2x-3x^2=(5+3x)(1-x).
\]
于是 $\varphi(x)$ 在 $[-1,1]$ 上于 $x=-1$ 取最小值，于 $x=1$ 取最大值. 因此
\[
  \max f=f(1,0)=4,\quad \min f=f(-1,0)=-4.
\]

\noindent\textbf{d.\quad 用极值证明不等式}

\noindent\textbf{i）用自由极值证明不等式}

\noindent\textbf{要点}\quad 若 $f$ 在区域 $D$ 上的最大值为 $B$，最小值为 $A$，则 $A\le f(P)\le B(P\in D)$. 对于不等式 $f\le g$，可证明 $\psi=f-g$ 的最大值不超过 $0$；对于 $f\ge g$，可证明 $\psi=f-g$ 的最小值不小于 $0$.

\noindent\textbf{例 6.3.14}\quad 证明：当 $t\ge1,s\ge0$ 时，
\[
  ts\le t\ln t-t+e^s.
\]
（武汉竞赛题）

\noindent\textbf{证}\quad 如图 6.3.3，要证
\[
  \varphi(s,t)=t\ln t-t+e^s-ts
\]
在 $D=\{(s,t)\mid s\ge0,t\ge1\}$ 上的最小值为 $0$.

当 $t$ 固定时，
\[
  \varphi_s'=-t+e^s.
\]
若 $0\le s<\ln t$，则 $\varphi_s'<0$；若 $s>\ln t$，则 $\varphi_s'>0$. 因此 $\varphi$ 的最小值在曲线 $t=e^s$ 上取得，而
\[
  \varphi(s,e^s)=e^ss-e^s+e^s-e^ss=0.
\]
故 $\varphi(s,t)\ge0$，即所证不等式成立.

\begin{center}
  \includegraphics[width=0.42\textwidth]{6-3-3.png}

  图 6.3.3
\end{center}

\noindent\textbf{例 6.3.15}\quad 证明：当 $0<x<1,0<y<+\infty$ 时，
\[
  f(x,y)=yx^y(1-x)<e^{-1}.
\]
（吉林大学）

\noindent\textbf{证}\quad 只需证明 $f$ 在区域 $0<x<1,0<y<+\infty$ 上的最大值小于 $e^{-1}$. 在边界上 $f$ 趋于 $0$，而在内部 $f>0$，故最大值只能在内部取得.

令
\[
  f_x'=y^2x^{y-1}(1-x)-yx^y=yx^{y-1}(y-xy-x),
\]
\[
  f_y'=x^y(1-x)+yx^y(1-x)\ln x=x^y(1-x)(1+y\ln x).
\]
由 $f_x'=f_y'=0$ 得
\[
  y(1-x)=x,\quad x^y=e^{-1}.
\]
于是，在稳定点处
\[
  f=yx^y(1-x)=e^{-1}x<e^{-1}.
\]
故原不等式成立.

\noindent\textbf{ii）用条件极值证明不等式}

\noindent\textbf{要点}\quad 若在条件 $\varphi(P)=a$ 下，$u=f(P)$ 的最大值为 $B(a)$，则对任意 $P$ 有
\[
  f(P)\le B(\varphi(P)).
\]
同理也可处理最小值情形.

\noindent\textbf{\ensuremath{\star}例 6.3.16}\quad 求 $x>0,y>0,z>0$ 时函数
\[
  f(x,y,z)=\ln x+2\ln y+3\ln z
\]
在球面 $x^2+y^2+z^2=6r^2$ 上的最大值，并证明当 $a,b,c$ 为正实数时，
\[
  ab^2c^3\le108\left(\frac{a+b+c}{6}\right)^6. \tag{1}
\]
（清华大学）

\noindent\textbf{解}\quad 令
\[
  L=\ln x+2\ln y+3\ln z+\lambda(x^2+y^2+z^2-6r^2).
\]
解得
\[
  x=r,\quad y=\sqrt2\,r,\quad z=\sqrt3\,r.
\]
由于 $f$ 在第一卦限内的球面上连续，而在边界处 $f\to-\infty$，最大值只能在内部取得；又内部仅有上述稳定点，故
\[
  \max f=\ln(6\sqrt3\,r^6).
\]
从而
\[
  \ln(xy^2z^3)\le
  \ln\left[6\sqrt3\left(\frac{x^2+y^2+z^2}{6}\right)^3\right],
\]
即
\[
  xy^2z^3\le6\sqrt3\left(\frac{x^2+y^2+z^2}{6}\right)^3.
\]
两边平方，并令 $a=x^2,b=y^2,c=z^2$，便得不等式 (1).

\noindent\textbf{注 1}\quad 上述方法还可以证明一些著名不等式. 例如：

1）在条件 $\sum a_ix_i=A$ 下求函数
\[
  \left(\sum a_i^k\right)^{1/k}\left(\sum x_i^{k'}\right)^{1/k'}
\]
的最小值，可证明 Holder 不等式
\[
  \sum a_ix_i\le
  \left(\sum a_i^k\right)^{1/k}
  \left(\sum x_i^{k'}\right)^{1/k'}
\]
其中 $a_i\ge0,x_i\ge0(i=1,2,\cdots,n),\ k>1,\ \frac1k+\frac1{k'}=1$.

2）在条件 $\sum x_{ij}^2=s_i(i=1,2,\cdots,n)$ 下，求行列式
\[
  \begin{vmatrix}
    x_{11}&x_{12}&\cdots&x_{1n}\\
    x_{21}&x_{22}&\cdots&x_{2n}\\
    \vdots&\vdots&&\vdots\\
    x_{n1}&x_{n2}&\cdots&x_{nn}
  \end{vmatrix}
\]
绝对值的最大值，可证明 Hadamard 不等式
\[
  |\det(x_{ij})|^2\le
  \prod_{i=1}^{n}s_i=
  \prod_{i=1}^{n}\sum_{j=1}^{n}x_{ij}^2.
\]

\noindent\textbf{注 2}\quad 不等式 (1) 也可直接用 6 个正数的平均值不等式证明：
\[
  ab^2c^3=108\left(a\cdot\frac b2\cdot\frac b2\cdot
  \frac c3\cdot\frac c3\cdot\frac c3\right)
  \le108\left(\frac{a+b+c}{6}\right)^6,
\]
等号当且仅当 6 个正数彼此相等时成立，亦即当且仅当 $6a=3b=2c$ 时，等号才成立：
\[
  ab^2c^3=108\left(\frac{a+b+c}{6}\right)^6.
\]

\noindent\textbf{注 3}\quad （将此不等式推广到 $n$ 元的情况）对任意 $n$ 个正数 $a_k>0(k=1,2,\cdots,n)$，有如下不等式成立：
\[
  a_1a_2^2a_3^3\cdots a_n^n
  \le2^23^3\cdots n^n
  \left[\frac{a_1+a_2+\cdots+a_n}{n(n+1)/2}\right]^{n(n+1)/2},
\]
且等号成立的充分必要条件是
\[
  a_1=\frac{a_2}{2}=\frac{a_3}{3}=\cdots=\frac{a_n}{n}.
\]

\noindent\textbf{练习}\quad 设
\[
  S=\{(x,y,z)\in\mathbb{R}^3\mid xy^2z^3=1\}.
\]
1）证明 $S$ 在 $\mathbb{R}^3$ 中确定一张隐式的曲面，并求出一个在点 $(1,1,1)$ 附近的参数方程；

2）$S$ 是否连通，是否紧致？

3）点 $q\in S$，$\|q\|$ 表示 $q$ 到原点的距离，点 $p$ 满足 $\|p\|=\inf_{q\in S}\|q\|$，求 $p$ 组成的集合.（中国科学技术大学）

\noindent\textbf{提示}\quad $xy^2z^3=1\Rightarrow x$ 和 $z$ 必须同号，故 $S$ 只在 1，4，6，7 卦限，第 2，3，5，8 卦限无 $S$ 的图像.

$xy^2z^3=1\Rightarrow x,y,z\ne0$，曲面 $S$ 与坐标面无交点，故 $S$ 被坐标面隔成 4 叶（不连通）.

$xy^2z^3=1\Rightarrow$ 不论 $y$ 变号，或者 $x,y,z$ 同时变号，$xy^2z^3=1$ 都不会变，故 $S$ 关于平面 $xOz$（镜面）对称，而且关于原点有点对称，只需研究第一卦限里的情况，其余由对称性可得.

$S$ 第 1 卦限里的图像如下：$\forall y_0\in\mathbb{R}$，用平面 $y=y_0$ 截 $S$，交线是
\[
  x=\frac{1}{y_0^2z^3},
\]
截线类似于双曲线，以坐标轴为渐近线.

同理，用平面 $z=z_0$ 或 $x=x_0$ 截 $S$，有类似结果.

可见，曲面 $S$ 分别以三坐标面为渐近切面. 形象地说：$S$ 像 4 口大锅，位于 1，4，6，7 卦限，锅底朝原点，锅与三坐标面越来越贴近.

用求极值的方法，可求出离原点最近的 4 个顶点 $p_k$：
\[
  \|p_k\|=\inf_{q\in S}\|q\|\quad(k=1,2,3,4).
\]

\noindent\textbf{解}\quad 1）$S$ 的参数方程可写为
\[
  x=\frac{1}{s^2t^3},\quad y=s,\quad z=t\quad(s\ne0,t\ne0).
\]

2）图形不连通，不紧致. 例如点列
\[
  \left\{M_n=\left(n,n,\frac1n\right)\ (n=1,2,\cdots)\right\}
\]
就无聚点.

3）\noindent\textbf{解法 I}\quad 设
\[
  L=x^2+y^2+z^2-2\lambda(xy^2z^3-1).
\]
令 $xL_x'=0$（注意 $xy^2z^3=1$），得 $x^2=\lambda$. 类似可得
\[
  y^2=2\lambda,\quad z^2=3\lambda.
\]
代入 $xy^2z^3=1$，得
\[
  6\sqrt3\,\lambda^{1/2+1+3/2}=1,\quad
  \lambda=\frac1{\sqrt[6]{108}}.
\]
于是
\[
  (x,y,z)=(\sqrt\lambda,\sqrt{2\lambda},\sqrt{3\lambda})
  =\left(\frac1{\sqrt[12]{108}},\frac{\sqrt2}{\sqrt[12]{108}},
  \frac{\sqrt3}{\sqrt[12]{108}}\right).
\]
利用对称性，全部 4 点是
\[
  \left(\frac1{\sqrt[12]{108}},\pm\frac{\sqrt2}{\sqrt[12]{108}},
  \frac{\sqrt3}{\sqrt[12]{108}}\right),\quad
  \left(-\frac1{\sqrt[12]{108}},\pm\frac{\sqrt2}{\sqrt[12]{108}},
  -\frac{\sqrt3}{\sqrt[12]{108}}\right).
\]

\noindent\textbf{解法 II}\quad 例 6.3.16 中的不等式，当且仅当 $6a=3b=2c$ 时，等号才成立. 应用到这里，由 $1=xy^2z^3$ 得
\[
  1=x^2y^4z^6
  =108\left(x^2\cdot\frac{y^2}{2}\cdot\frac{y^2}{2}\cdot
  \frac{z^2}{3}\cdot\frac{z^2}{3}\cdot\frac{z^2}{3}\right)
  \le108\left(\frac{x^2+y^2+z^2}{6}\right)^6. \tag{1}
\]
当“$\le$”变为“$=$”号时，右端的值取最小，等于 1. 因此得
\[
  \inf_{(x,y,z)\in S}\sqrt{x^2+y^2+z^2}
  =\sqrt[12]{\frac{6^6}{108}}
  =\sqrt[12]{2\times6^3}
  =\sqrt[12]{432}.
\]
不等式 (1) 变成等式的充分必要条件是
\[
  x^2=\frac{y^2}{2}=\frac{z^2}{3}.
\]
因 $xy^2z^3=1$，故 $x^2y^4z^6=1$，即 $x^2(2x^2)^2(3x^2)^3=1$，解得
\[
  x=\pm\frac1{\sqrt[12]{108}}.
\]
从而 $y=\pm\sqrt2\,x,\ z=\sqrt3\,x$. 所以，$S$ 距离原点最近之四点是
\[
  \left(\frac1{\sqrt[12]{108}},\pm\frac{\sqrt2}{\sqrt[12]{108}},
  \frac{\sqrt3}{\sqrt[12]{108}}\right),\quad
  \left(-\frac1{\sqrt[12]{108}},\pm\frac{\sqrt2}{\sqrt[12]{108}},
  -\frac{\sqrt3}{\sqrt[12]{108}}\right).
\]
（不难验证：它们满足 $S$ 的方程，且至原点的距离均为 $\sqrt[12]{432}=\inf_{(x,y,z)\in S}\sqrt{x^2+y^2+z^2}$.）

\noindent\textbf{e.\quad 极值应用问题}

\noindent\textbf{要点}\quad 求解极值应用题，关键在于选取适当的变量，使之能方便地表示目标函数以及约束条件.

\noindent\textbf{\ensuremath{\star}例 6.3.17}\quad 将长度为 $l$ 的铁丝分为三段，用此三段分别作成圆、正方形、等边三角形，问采用何种分法，才能使这三个图形的面积之和最小.（可以利用以下结论：二元二次函数 $F(x,y)=ax^2+bxy+cy^2+dx+ey+f$ 的极大（小）值，必为其最大（小）值.）（南开大学）

\noindent\textbf{解}\quad $1^\circ$ 为了便于表达周长之和与总面积，我们不妨取 $x,y,z$ 分别表示圆之半径、正方形的边长、等边三角形的边长. 于是总面积
\[
  S=\pi x^2+y^2+\frac{\sqrt3}{4}z^2\quad(\text{目标函数}) \tag{1}
\]
应满足方程
\[
  2\pi x+4y+3z=l\quad(\text{约束条件}). \tag{2}
\]
Lagrange 函数
\[
  L=\pi x^2+y^2+\frac{\sqrt3}{4}z^2-\lambda(2\pi x+4y+3z-l).
\]
令
\[
  \begin{cases}
    L_x'=2\pi x-2\pi\lambda=0,\\
    L_y'=2y-4\lambda=0,\\
    L_z'=\frac{\sqrt3}{2}z-3\lambda=0,
  \end{cases}
\]
得
\[
  x=\lambda,\quad y=2\lambda,\quad z=\frac6{\sqrt3}\lambda.
\]
这时铁丝三段长度的比为
\[
  2\pi x:4y:3z=2\pi\lambda:8\lambda:6\sqrt3\,\lambda=\pi:4:3\sqrt3.
\]

$2^\circ$（判断）将 (2) 代入 (1)，得
\[
  S=\pi x^2+y^2+\frac{\sqrt3}{36}(l-2\pi x-4y)^2
\]
\[
  =\pi x^2+y^2+\frac{\sqrt3}{36}(4\pi^2x^2+16y^2+16\pi xy-4\pi lx-8ly+l^2).
\]
由此可知
\[
  \Delta=S_{xx}''S_{yy}''-S_{xy}''{}^2
  =4\pi\left[1+\frac{\sqrt3}{9}(\pi+4)\right]>0,\quad S_{xx}''>0,
\]
从而按上述比例分割铁丝，所围的面积极小，也是最小.

\noindent\textbf{例 6.3.18}\quad 设 $\triangle ABC$ 为正三角形，边长为 $a$. $P$ 为 $\triangle ABC$ 内任意一点，由 $P$ 向三边引垂线（如图 6.3.4），其与三边的交点分别为 $D,E,F$. 试求 $\triangle DEF$ 的面积最大值.（武汉大学赛题）

\begin{center}
  \includegraphics[width=0.35\textwidth]{6-3-4.png}

  图 6.3.4
\end{center}

\noindent\textbf{解}\quad 记点 $P$ 至三边的距离分别为 $x,y,z$. 注意到
\[
  \angle DPF=\angle DPE=\angle EPF=\frac23\pi,
\]
所以 $\triangle DEF$ 的面积
\begin{equation}
\begin{aligned}
  S_{\triangle DEF}
  &=S_{\triangle DPF}+S_{\triangle DPE}+S_{\triangle EPF}\\
  &=\frac12\sin\frac{\pi}{3}(xz+xy+yz)\\
  &=\frac{\sqrt3}{4}(xz+xy+yz)\quad(\text{目标函数}).
\end{aligned}
\tag{1}
\end{equation}
由
\[
  S_{\triangle PBC}+S_{\triangle CPA}+S_{\triangle APB}=S_{\triangle ABC}
\]
得约束方程为
\[
  \frac12ax+\frac12ay+\frac12az=\frac12a\cdot\frac{\sqrt3}{2}a,
\]
即
\[
  x+y+z=\frac{\sqrt3}{2}a\quad(\text{约束条件}). \tag{2}
\]
由此可得
\[
  x=y=z=\frac{\sqrt3}{6}a,\quad
  \max S_{\triangle DEF}=\frac{\sqrt3}{16}a^2.
\]
（计算和判断过程从略. 本题求解的方法很多.）

\noindent\textbf{例 6.3.19}\quad 在平面上给一边长分别为 $a,b,c$ 的三角形，在它上面作无数个定高 $h$ 的锥体，求侧面积最小的锥体.（大连理工大学）

\begin{center}
  \includegraphics[width=0.42\textwidth]{6-3-5.png}

  图 6.3.5
\end{center}

\noindent\textbf{解}\quad 锥顶 $H$ 在底面的投影记为 $O$，从 $O$ 到三边 $BC,CA,AB$ 的距离分别记为 $x,y,z$（如图 6.3.5），则锥的侧面积（目标函数）为
\[
  S=\frac12a\sqrt{h^2+x^2}+\frac12b\sqrt{h^2+y^2}
  +\frac12c\sqrt{h^2+z^2}. \tag{1}
\]
规定：若点 $O$ 与三角形 $ABC$ 的内心在 $BC$ 同侧，则 $x$ 算作正；在异侧，则算作负. 对 $y,z$ 也作类似的规定. 此时，不论点 $O$ 在 $\triangle ABC$ 的内部还是外部，恒有
\[
  \frac12ax+\frac12by+\frac12cz=S_{\triangle ABC}. \tag{2}
\]
即
\[
  ax+by+cz=2S_{\triangle ABC}\overset{\text{记}}{=}m,
\]
其中
\[
  S_{\triangle ABC}=\sqrt{p(p-a)(p-b)(p-c)},\quad p=\frac{a+b+c}{2}.
\]
记
\[
  L=a\sqrt{h^2+x^2}+b\sqrt{h^2+y^2}
  +c\sqrt{h^2+z^2}-\lambda(ax+by+cz-m).
\]
令 $L_x'=L_y'=L_z'=0$，得
\[
  \lambda=\frac{x}{\sqrt{h^2+x^2}}
  =\frac{y}{\sqrt{h^2+y^2}}
  =\frac{z}{\sqrt{h^2+z^2}}. \tag{3}
\]
从实际背景看，问题有最小值，无最大值. 现在只有一个可疑点，故它对应最小值. 式 (3) 表明最小值发生在三侧面与底面成等角的时候. 因此，当 $x=y=z$，即 $O$ 与三角形内心重合时，侧面积最小. 此时
\[
  S=\frac12\sqrt{h^2+r^2}(a+b+c),
\]
其中（内接圆半径）
\[
  r=\frac{2\sqrt{p(p-a)(p-b)(p-c)}}{a+b+c},\quad
  p=\frac{a+b+c}{2}.
\]

下面这道命题在中学里已学会用几何方法进行推断，现在我们可用分析方法进行证明了. 例 6.3.20 还说明有时可疑点可不必求出.

\noindent\textbf{\ensuremath{\star}例 6.3.20}\quad 试由费马原理“光线总是按费时最短的路径传播”来证明镜面反射时，入射角等于反射角.

\begin{center}
  \includegraphics[width=0.42\textwidth]{6-3-6.png}

  图 6.3.6
\end{center}

\noindent\textbf{提示}\quad 可令
\[
  L=\sqrt{x^2+h_1^2}+\sqrt{y^2+h_2^2}-\lambda(x+y-a)
\]
（意义见图 6.3.6）. 由 $L_x'=L_y'=0$ 得
\[
  \frac{x}{\sqrt{x^2+h_1^2}}=
  \frac{y}{\sqrt{y^2+h_2^2}}.
\]
这表明
\[
  \sin\theta_1=\sin\theta_2,\quad \theta_1=\theta_2.
\]

\noindent\textbf{解}\quad 光程
\[
  s(x)=AP+PB=\sqrt{x^2+h_1^2}+\sqrt{(a-x)^2+h_2^2}.
\]
令
\[
  s_x'=\frac{x}{\sqrt{x^2+h_1^2}}-
  \frac{a-x}{\sqrt{(a-x)^2+h_2^2}}=0,
\]
得
\[
  \frac{x}{\sqrt{x^2+h_1^2}}=
  \frac{a-x}{\sqrt{(a-x)^2+h_2^2}},
\]
即 $\sin\theta_1=\sin\theta_2$.

\noindent\textbf{例 6.3.21}\quad 设
\[
  ax^2+by^2+cz^2+2exy+2fyz+2gzx=1 \tag{1}
\]
为一椭球面，求证其三个半轴之长恰为矩阵
\[
  \begin{pmatrix}
    a&e&g\\
    e&b&f\\
    g&f&c
  \end{pmatrix} \tag{2}
\]
的三个特征值的平方根的倒数.（北京航空航天大学）

\noindent\textbf{方法}\quad 求 $u=x^2+y^2+z^2$ 在条件 (1) 下的稳定点，以求出
\[
  \rho=\sqrt{x^2+y^2+z^2}
\]
的极值. 技巧是：回避求出稳定点，而从稳定点满足的方程中直接解出 $\rho=\sqrt{x^2+y^2+z^2}$.

\noindent\textbf{解}\quad 因为 $\rho=\sqrt{x^2+y^2+z^2}$ 与 $u=x^2+y^2+z^2$ 有相同的极值点. 记
\[
  L=x^2+y^2+z^2-\frac1{\lambda}
  (ax^2+by^2+cz^2+2exy+2fyz+2gzx-1).
\]
令 $L_x'=L_y'=L_z'=0$，得方程组
\[
  \begin{cases}
    (a-\lambda)x+ey+gz=0,\\
    ex+(b-\lambda)y+fz=0,\\
    gx+fy+(c-\lambda)z=0.
  \end{cases} \tag{3--5}
\]
此方程组有非零解必须系数行列式为零，故
\[
  \begin{vmatrix}
    a-\lambda&e&g\\
    e&b-\lambda&f\\
    g&f&c-\lambda
  \end{vmatrix}=0. \tag{6}
\]
这是 $\lambda$ 的三次方程. 因为椭球面是有心的非退化二次曲面，故
\[
  \begin{vmatrix}
    a&e&g\\
    e&b&f\\
    g&f&c
  \end{vmatrix}\ne0.
\]
因而方程 (6) 有三个实根 $\lambda_1,\lambda_2,\lambda_3$. 根据线性代数的知识，它们是矩阵 (2) 的特征值. 将 $\lambda_i$ 代回方程 (3), (4), (5)，对方程 (3), (4), (5) 分别乘 $x,y,z$ 再相加，注意到 $(x,y,z)$ 满足 (1)，则
\[
  \lambda_i(x^2+y^2+z^2)=
  ax^2+by^2+cz^2+2exy+2fyz+2gzx=1.
\]
从而
\[
  x^2+y^2+z^2=\frac1{\lambda_i},
\]
\[
  \rho=\sqrt{x^2+y^2+z^2}=\frac1{\sqrt{\lambda_i}}
  \quad(i=1,2,3).
\]
这就证明了半轴之长等于矩阵 (2) 的特征值的平方根的倒数.

\noindent\textbf{练习}\quad 设 $A$ 为三阶实对称方阵，定义函数：
\[
  h(x,y,z)=(x,y,z)A
  \begin{pmatrix}
    x\\y\\z
  \end{pmatrix},
\]
求证：$h(x,y,z)$ 在条件 $x^2+y^2+z^2=1$ 下的最大值是矩阵 $A$ 的最大特征值.（南开大学）

\noindent\textbf{证}\quad $1^\circ$ 记
\[
  A=(a_{ij})=
  \begin{pmatrix}
    a_{11}&a_{12}&a_{13}\\
    a_{21}&a_{22}&a_{23}\\
    a_{31}&a_{32}&a_{33}
  \end{pmatrix},
\]
则
\[
  h(x,y,z)=a_{11}x^2+a_{22}y^2+a_{33}z^2
  +2a_{12}xy+2a_{13}xz+2a_{23}yz
\]
（是连续函数），$h(x,y,z)$ 在有限闭集
\[
  D=\{(x,y,z)\mid x^2+y^2+z^2=1\}
\]
上必有最大值.

$2^\circ$（用 Lagrange 乘数法，求函数 $h$ 的条件极值.）设
\[
  L=a_{11}x^2+a_{22}y^2+a_{33}z^2+2a_{12}xy+2a_{13}xz+2a_{23}yz
  -\lambda(x^2+y^2+z^2-1).
\]
求 $L$ 对 $x,y,z$ 的偏导数，并令之为零，得
\[
  (a_{11}-\lambda)x+a_{12}y+a_{13}z=0, \tag{1}
\]
\[
  a_{12}x+(a_{22}-\lambda)y+a_{23}z=0, \tag{2}
\]
\[
  a_{13}x+a_{23}y+(a_{33}-\lambda)z=0. \tag{3}
\]
式 (1), (2), (3) 分别乘 $x,y,z$，再将三式相加，可得（注意 $x^2+y^2+z^2=1$，且 $A$ 为对称方阵）
\[
  \lambda=a_{11}x^2+a_{22}y^2+a_{33}z^2
  +2a_{12}xy+2a_{13}xz+2a_{23}yz=h(x,y,z). \tag{4}
\]
另一方面，既然 $h$ 在附球面 $D$ 上有最大值点，其坐标 $x,y,z$ 就不会同时为零，那么齐次方程组 (1), (2), (3) 应有非零解，故（方程组的）系数行列式应为零，即
\[
  \begin{vmatrix}
    a_{11}-\lambda&a_{12}&a_{13}\\
    a_{21}&a_{22}-\lambda&a_{23}\\
    a_{31}&a_{32}&a_{33}-\lambda
  \end{vmatrix}=0. \tag{5}
\]
此式左端是矩阵 $A$ 的特征多项式，而 $\lambda$ 是矩阵 $A$ 的特征值. 式 (4) 表明极值点上 $h=\lambda$，所以
\[
  \max_{(x,y,z)\in D}h(x,y,z)=\max_{k=1,2,3}\lambda_k
\]
（是矩阵 $A$ 的最大特征值）. 证毕.

\noindent\textbf{\ensuremath{\star}例 6.3.22}\quad 试求平面 $\alpha x+\beta y+\gamma z=0$ 与圆柱面
\[
  \frac{x^2}{A^2}+\frac{y^2}{B^2}=1\quad(A,B>0)
\]
相交所成椭圆的面积.

\noindent\textbf{解 I}\quad （用极值方法.）要计算椭圆的面积 $S=\pi ab$，只需算出椭圆的半轴 $a$ 与 $b$.

现在来说，也就是要求 $\rho=\sqrt{x^2+y^2+z^2}$ 在条件 $\alpha x+\beta y+\gamma z=0$ 及
\[
  \frac{x^2}{A^2}+\frac{y^2}{B^2}=1
\]
之下的极值. 以 $\rho^2$ 替换 $\rho$，取
\[
  L=x^2+y^2+z^2+2\lambda(\alpha x+\beta y+\gamma z)
  -\mu\left(\frac{x^2}{A^2}+\frac{y^2}{B^2}-1\right).
\]
令 $L_x'=L_y'=L_z'=0$，得
\[
  x\left(1-\mu\frac1{A^2}\right)+\lambda\alpha=0, \tag{1}
\]
\[
  y\left(1-\mu\frac1{B^2}\right)+\lambda\beta=0, \tag{2}
\]
\[
  z+\lambda\gamma=0, \tag{3}
\]
\[
  \alpha x+\beta y+\gamma z=0, \tag{4}
\]
\[
  \frac{x^2}{A^2}+\frac{y^2}{B^2}=1. \tag{5}
\]
(1), (2), (3) 中消去 $\lambda$ 得 $\mu$ 的方程：
\[
  \mu^2-\left(\frac{\beta^2}{\gamma^2}B^2+
  \frac{\alpha^2}{\gamma^2}A^2+B^2+A^2\right)\mu
  +\frac{A^2B^2}{\gamma^2}(\alpha^2+\beta^2+\gamma^2)=0. \tag{6}
\]
(1), (2), (3) 分别乘 $x,y,z$，相加得
\[
  x^2+y^2+z^2=\mu\left(\frac{x^2}{A^2}+\frac{y^2}{B^2}\right)
  -\lambda(\alpha x+\beta y+\gamma z)=\mu. \tag{7}
\]
$a,b$ 为所讨论椭圆的两个半轴，则它们应是 $\rho$ 的极值. 而 $\rho$ 与 $\rho^2$ 的极值点相同，(7) 表明 $\rho^2$ 的极值等于 $\mu$. $\mu$ 满足方程 (6)，设 (6) 的两根为 $\mu_1,\mu_2$，则 $\mu_1,\mu_2$ 应分别为 $a^2,b^2$. 因而 $S=\pi ab=\pi\sqrt{\mu_1\mu_2}$. 根据 Viete 定理，
\[
  \mu_1\mu_2=\frac{A^2B^2}{\gamma^2}(\alpha^2+\beta^2+\gamma^2).
\]
所以
\[
  S=\pi\frac{AB}{|\gamma|}\sqrt{\alpha^2+\beta^2+\gamma^2}.
\]

\noindent\textbf{解 II}\quad （利用面积的投影.）所讨论椭圆在 $xOy$ 平面上的投影为椭圆
\[
  \frac{x^2}{A^2}+\frac{y^2}{B^2}=1,
\]
其面积为 $\pi AB$. 所讨论椭圆位于平面 $\alpha x+\beta y+\gamma z=0$ 上，该法线与 $z$ 轴夹角的余弦为
\[
  \frac{|\gamma|}{\sqrt{\alpha^2+\beta^2+\gamma^2}}.
\]
根据面积投影关系，所讨论椭圆的面积
\[
  S=\pi AB\frac1{|\gamma|}\sqrt{\alpha^2+\beta^2+\gamma^2}.
\]

\noindent\textbf{\ensuremath{\star}例 6.3.23}\quad 设函数 $F(x,y)$ 在平面区域 $G$ 内有定义，其一阶偏导数连续，又设方程 $F(x,y)=0$ 的图形是一条自身不相交的封闭曲线 $\Gamma$，$\Gamma\subset G$，并且在 $\Gamma$ 的每一点上 $F_x'(x,y)$ 和 $F_y'(x,y)$ 不同时为零. 试证：若 $AB$ 是 $\Gamma$ 的一条极大弦（即 $A,B$ 是 $\Gamma$ 上的两点，且存在点 $A$ 的邻域 $U(A)$ 和点 $B$ 的邻域 $U(B)$，使得当点 $C\in U(A)\cap\Gamma$，点 $D\in U(B)\cap\Gamma$ 时总有 $\overline{CD}\le\overline{AB}$），则 $\Gamma$ 在 $A,B$ 两点的切线必互相平行.（武汉大学）

\noindent\textbf{证}\quad 记 $A(x_A,y_A),B(x_B,y_B)$. 根据题意，函数
\[
  u=(x-x_B)^2+(y-y_B)^2
\]
应在条件 $F(x,y)=0$ 限制下，在点 $A(x_A,y_A)$ 处达到极大. 因此，
\[
  u_x'\big|_A=\left[2(x-x_B)-2(y-y_B)y_x'\right]\big|_A=0. \tag{1}
\]
由 $F(x,y)\equiv0$ 得 $F_x'+F_y'\cdot y_x'=0$，$y_x'=-\dfrac{F_x'}{F_y'}$. 代入式 (1) 得
\[
  (x_A-x_B)+(y_A-y_B)\left(-\frac{F_x'(x_A,y_A)}{F_y'(x_A,y_A)}\right)=0. \tag{2}
\]
故曲线 $F(x,y)=0$ 在点 $A(x_A,y_A)$ 的切线斜率
\[
  k_A=-\frac{F_x'(x_A,y_A)}{F_y'(x_A,y_A)}
  =-\frac{x_A-x_B}{y_A-y_B}. \tag{3}
\]
同理，点 $B$ 的切线斜率
\[
  k_B=-\frac{x_B-x_A}{y_B-y_A}. \tag{4}
\]
比较 (3), (4)，知 $A,B$ 两点的切线相互平行.

\noindent\textbf{\ensuremath{\star\star}例 6.3.24}\quad 若 $x^0=(x_1^0,x_2^0,\cdots,x_n^0)$ 是函数
\[
  f(x_1,x_2,\cdots,x_n)=\sum_{i,j=1}^n a_{ij}x_ix_j\quad(a_{ij}=a_{ji})
\]
在 $x_1^2+x_2^2+\cdots+x_n^2=1$ 上的极值点，
\[
  A=\begin{pmatrix}
    a_{11}&\cdots&a_{1n}\\
    \vdots&&\vdots\\
    a_{n1}&\cdots&a_{nn}
  \end{pmatrix},
\]
试证 $Ax^0=0$ 或者 $x^0$ 是 $A$ 的特征向量.（兰州大学）

\noindent\textbf{证}\quad 设
\[
  L(x_1,x_2,\cdots,x_n)=\sum_{i,j=1}^n a_{ij}x_ix_j-\lambda(x_1^2+\cdots+x_n^2-1),
\]
则 $x^0=(x_1^0,x_2^0,\cdots,x_n^0)$ 应满足方程 $\dfrac{\partial L}{\partial x_i}=0$，即
\[
  a_{i1}x_1+a_{i2}x_2+\cdots+a_{in}x_n-\lambda x_i=0
  \quad(i=1,2,\cdots,n),
\]
亦即
\[
  A\begin{pmatrix}
    x_1\\x_2\\ \vdots\\x_n
  \end{pmatrix}
  =
  \lambda\begin{pmatrix}
    x_1\\x_2\\ \vdots\\x_n
  \end{pmatrix}.
\]
故
\[
  Ax^0=\lambda x^0.
\]
这说明 $x^0$ 或为方程 $Ax^0=0$ 的根（当 $\lambda=0$ 时），或为矩阵 $A$ 的特征向量（当 $\lambda\ne0$ 时）.

下面考虑 $\mathbb{R}^n$ 中二次函数在 $m$ 个超平面交线上的极值问题.

\noindent\textbf{※例 6.3.25}\quad 函数为 $f(x)=\dfrac12x^{\mathrm T}Ax+B^{\mathrm T}x,x\in\mathbb{R}^n$，约束方程为
\[
  Cx=D. \tag{1}
\]
求 $f$ 的极值点应满足的条件，其中 $A$ 为 $n\times n$ 方阵，$B$ 为 $n$ 维向量，$C$ 为 $m\times n$ 矩阵 $(m<n)$，$D$ 为 $m$ 维向量.

\noindent\textbf{解}\quad 设
\[
  L=\frac12x^{\mathrm T}Ax+B^{\mathrm T}x+\lambda^{\mathrm T}(Cx-D),
  \quad \lambda=(\lambda_1,\cdots,\lambda_m)^{\mathrm T}.
\]
令 $\nabla L=0$，其中 $\nabla L=(L_{x_1}',\cdots,L_{x_n}')$. 由此
\[
  Ax+C^{\mathrm T}\lambda=-B.
\]
与式 (1) 联立，即为
\[
  \begin{pmatrix}
    A&C^{\mathrm T}\\
    C&O
  \end{pmatrix}
  \begin{pmatrix}
    x\\ \lambda
  \end{pmatrix}
  =
  \begin{pmatrix}
    -B\\D
  \end{pmatrix}.
\]
此即为极值点应满足的条件. 记
\[
  H=\begin{pmatrix}
    A&C^{\mathrm T}\\
    C&O
  \end{pmatrix},
\]
若 $H^{-1}$ 存在，则
\[
  \begin{pmatrix}
    x\\ \lambda
  \end{pmatrix}
  =
  H^{-1}
  \begin{pmatrix}
    -B\\D
  \end{pmatrix}.
\]

\noindent\textbf{\ensuremath{\star}例 6.3.26}\quad 设 $P_0$ 是椭球面
\[
  S:\frac{x^2}{a^2}+\frac{y^2}{b^2}+\frac{z^2}{c^2}=1
\]
外一点，而 $P_1\in S$. 若
\[
  |\overline{P_1P_0}|=\max_{P\in S}|\overline{PP_0}|,
\]
求证：直线 $P_1P_0$ 是 $S$ 在点 $P_1$ 处的法线.（华东师范大学）

\noindent\textbf{证}\quad $1^\circ$ 令
\[
  L=(x-x_0)^2+(y-y_0)^2+(z-z_0)^2
  +\lambda\left(\frac{x^2}{a^2}+\frac{y^2}{b^2}+\frac{z^2}{c^2}-1\right).
\]
由 $\dfrac{\partial L}{\partial x}=\dfrac{\partial L}{\partial y}=\dfrac{\partial L}{\partial z}=0$，可得极值的必要条件：
\[
  (x-x_0,y-y_0,z-z_0)=-\lambda\left(\frac{x}{a^2},\frac{y}{b^2},\frac{z}{c^2}\right). \tag{1}
\]
因 $|\overline{PP_0}|=\sqrt{(x-x_0)^2+(y-y_0)^2+(z-z_0)^2}$ 在 $S$ 上连续，$S$ 是有界闭集，故 $|\overline{PP_0}|$ 在 $S$ 上有最值（$|\overline{PP_0}|^2$ 亦然）. 因此，最远点 $P_1(x_1,y_1,z_1)$ 应满足式 (1)，即有
\[
  (x_1-x_0,y_1-y_0,z_1-z_0)
  =-\lambda\left(\frac{x_1}{a^2},\frac{y_1}{b^2},\frac{z_1}{c^2}\right). \tag{2}
\]
$2^\circ$
\[
  S:F(x,y,z)=\frac{x^2}{a^2}+\frac{y^2}{b^2}+\frac{z^2}{c^2}-1=0,
\]
$F_x',F_y',F_z'$ 是 $S$ 的法线方向数. 故法向量：
\[
  (F_x',F_y',F_z')\big|_{P_1}
  =2\left(\frac{x_1}{a^2},\frac{y_1}{b^2},\frac{z_1}{c^2}\right),
\]
$S$ 在 $P_1$ 处的法线
\[
  \frac{x-x_1}{\dfrac{x_1}{a^2}}
  =\frac{y-y_1}{\dfrac{y_1}{b^2}}
  =\frac{z-z_1}{\dfrac{z_1}{c^2}}. \tag{3}
\]
式 (2) 表明：$P_0(x_0,y_0,z_0)$ 满足法线方程 (3). 可见直线 $P_0P_1$ 是 $S$ 在 $P_1$ 点的法线. 证毕.

\noindent\textbf{练习}\quad 设 $F(x,y,z)$ 是定义在开区域 $D$ 上有连续偏导数的三元函数，且
\[
  F_x^2(x,y,z)+F_y^2(x,y,z)+F_z^2(x,y,z)\ne0,\quad
  \forall(x,y,z)\in D. \tag{1}
\]
$S$ 是由 $F(x,y,z)=0$ 定义的封闭光滑曲面，$P_1,P_2\in S$ 是 $S$ 上相隔距离最远的两点. 试证：$S$ 在点 $P_1,P_2$ 处的两切平面相互平行且与 $P_1,P_2$ 连线垂直.（华东师范大学）

\noindent\textbf{提示}\quad 设 $P_1$ 与 $P_2$ 的坐标为 $(x_1,y_1,z_1),(x_2,y_2,z_2)$. 取
\[
  L=(x_1-x_2)^2+(y_1-y_2)^2+(z_1-z_2)^2
  +\lambda F(x_1,y_1,z_1)-\mu F(x_2,y_2,z_2).
\]
令 $L_{x_k}'=L_{y_k}'=L_{z_k}'=0\ (k=1,2)$，得最远两点 $P_1,P_2$ 必满足方程：
\[
  (x_1-x_2,y_1-y_2,z_1-z_2)
  =\lambda(F_x,F_y,F_z)\big|_{P_1}
  =\mu(F_x,F_y,F_z)\big|_{P_2}. \tag{2}
\]
因 $F(x,y,z)=0$ 是曲面 $S$ 的方程，故 $F_x\big|_{P_k},F_y\big|_{P_k},F_z\big|_{P_k}$ 是曲面 $S$ 在点 $P_k(k=1,2)$ 法线的方向数. 式 (2) 表明：$P_1,P_2$ 两点处的法线与 $P_1,P_2$ 两点的连线重合，故在 $P_1,P_2$ 处的切平面都与 $P_1,P_2$ 连线垂直，两者相互平行.

\noindent\textbf{\ensuremath{\star}例 6.3.27}\quad 设椭球面
\[
  S:\frac{x^2}{a^2}+\frac{y^2}{b^2}+\frac{z^2}{c^2}=1,
\]
求：$S$ 在 $x>0,y>0,z>0$ 的切平面与三个坐标平面所围成的几何体的最小体积.（华东师范大学）

\noindent\textbf{提示}\quad $S:\dfrac{x^2}{a^2}+\dfrac{y^2}{b^2}+\dfrac{z^2}{c^2}=1$（在 $x>0,y>0,z>0$ 的部分）在点 $P_0(x_0,y_0,z_0)\in S$ 处的切面为
\[
  \frac{xx_0}{a^2}+\frac{yy_0}{b^2}+\frac{zz_0}{c^2}=1.
\]
可见，此切面（在三坐标轴上）的截距分别为 $\dfrac{a^2}{x_0},\dfrac{b^2}{y_0},\dfrac{c^2}{z_0}$，因此欲求的体积为
\[
  V=\frac{a^2b^2c^2}{6x_0y_0z_0}.
\]

再提示\quad 当切点 $P_0(x_0,y_0,z_0)$ 靠近 $S$ 的边界时，体积 $V\to+\infty$，故最小值点必在内部.

设 $P_0(x_0,y_0,z_0)\in S$ 使得 $\dfrac{a^2b^2c^2}{6x_0y_0z_0}$ 取最小值，亦即使得
\[
  \frac{x_0^2}{a^2}\frac{y_0^2}{b^2}\frac{z_0^2}{c^2}
\]
取最大值. 但
\[
  \sqrt[3]{\frac{x_0^2}{a^2}\frac{y_0^2}{b^2}\frac{z_0^2}{c^2}}
  \le \frac13\left(\frac{x_0^2}{a^2}+\frac{y_0^2}{b^2}+\frac{z_0^2}{c^2}\right)
\]
（均值不等式），当且仅当三正量相等时等号才成立. 所以，
\[
  \frac{x_0^2}{a^2}=\frac{y_0^2}{b^2}=\frac{z_0^2}{c^2},
  \quad
  \frac{x_0^2}{a^2}+\frac{y_0^2}{b^2}+\frac{z_0^2}{c^2}=1.
\]
因此，
\[
  \frac{x_0^2}{a^2}=\frac{y_0^2}{b^2}=\frac{z_0^2}{c^2}=\frac13,
  \quad
  (x_0,y_0,z_0)=\left(\frac{a}{\sqrt3},\frac{b}{\sqrt3},\frac{c}{\sqrt3}\right),
\]
此时
\[
  V_{\min}=\min_{(x,y,z)\in S}\frac{a^2b^2c^2}{6xyz}
  =\frac{a^2b^2c^2}{6xyz}\bigg|_{P_0}
  =\frac{\sqrt3}{2}abc.
\]

\noindent\textbf{\ensuremath{\star}例 6.3.28}\quad 设 $x,y,z\ge0,x+y+z=\pi$，求 $2\cos x+3\cos y+4\cos z$ 的最大值和最小值.（北京大学）

\noindent\textbf{提示}\quad 问题是：求连续函数 $u=2\cos x+3\cos y+4\cos z$ 在平面 $\Sigma:x+y+z=\pi$ 位于第一卦限部分（记作 $D$）的最大值和最小值.（因 $D$ 是有界闭集，故 $u$ 在 $D$ 上有最大、最小值.）因此，只需分别求出 $u$ 在 $D$ 内和 $D$ 边界上的可疑点，比较 $u$ 值的大小即可得到.

\noindent\textbf{解}\quad 设
\[
  L=2\cos x+3\cos y+4\cos z+2\lambda(x+y+z-\pi).
\]
令 $L_x'=L_y'=L_z'=0$，得
\[
  \sin x=\lambda,\quad 3\sin y=2\lambda,\quad 2\sin z=\lambda,
\]
\[
  \sin x:\sin y:\sin z=1:\frac23:\frac12=6:4:3.
\]
根据三角形的正弦定理，知三角形三边比 $a:b:c=\sin x:\sin y:\sin z=6:4:3$.

根据三角形的余弦定理：
\[
  a^2=b^2+c^2-2bc\cos x,
\]
得 $36=16+9-24\cos x$，解得
\[
  \cos x=-\frac{11}{24}.
\]
同理可得
\[
  \cos y=\frac{29}{36},\quad \cos z=\frac{43}{48}.
\]
于是
\[
  u=2\cos x+3\cos y+4\cos z
  =2\times\left(-\frac{11}{24}\right)
  +3\times\frac{29}{36}+4\times\frac{43}{48}
  =\frac{61}{12}. \tag{1}
\]
（这是 $D$ 内唯一可疑值）.

平面 $\Sigma$ 在三坐标轴上的截距都为 $\pi$. $\Sigma$ 与 $xOy$ 的交线（$z=0,x+y=\pi$）上，
\[
  u=2\cos x+3\cos(\pi-x)+4=4-\cos x\quad(0\le x\le\pi),
\]
因此最小值为 $3$，最大值为 $5$. 同样可求出 $u$ 在另外两条边上的最小、最大值分别为 $1,5$ 和 $1,3$. 因此，$u$ 在边界上的最小、最大值分别为 $1$ 和 $5$.

而在 $D$ 内只有唯一可疑值 $\dfrac{61}{12}$（见式 (1)），知函数 $u=2\cos x+3\cos y+4\cos z$ 在 $D$ 上最大值为 $\dfrac{61}{12}\approx5.083$，最小值为 $1$.

\noindent\textbf{\ensuremath{\star}例 6.3.29}\quad 设 $D$ 是 $\mathbb{R}^3$ 中的有界闭区域，$f$ 在 $D$ 上连续且有偏导数. 证明：如果在 $D$ 上有
\[
  f_x'+f_y'+f_z'=f,\quad f\big|_{\partial D}=0
\]
（$\partial D$ 表示 $D$ 的边界），则 $f$ 在 $D$ 上恒等于 $0$.（中国科学技术大学）

\noindent\textbf{提示}\quad 在有界闭区域上连续，必有最大、最小值点（统称为最值点，统一记作 $P_k$）. 若最值点 $P_k$ 是内点，则 $P_k$ 也是 $f$ 的极值点，故
\[
  f_x'(P_k)=f_y'(P_k)=f_z'(P_k)=0,
\]
从而
\[
  f(P_k)=f_x'(P_k)+f_y'(P_k)+f_z'(P_k)=0.
\]
若 $P_k\in\partial D$，因已知 $f|_{\partial D}=0$，故也有 $f(P_k)=0$.

因此，$f$ 在 $D$ 上，既有最大值点，也有最小值点；但在所有最值点上 $f(P_k)=0$，故
\[
  f\equiv0\quad(D\text{ 上}).
\]

\begin{center}
  单元练习 6.3
\end{center}

\noindent\textbf{多元 Taylor 公式}\quad （此类考研试题相对较少）

\noindent\textbf{6.3.1}\quad 写出函数 $f(x,y)=y^x$ 在点 $(1,1)$ 附近的 Taylor 公式（写出二阶项，余项形式可不具体写出）.（兰州大学）
\[
  \left\langle
  1+2(y-1)+2\ln2\,(x-1)(y-1)+(y-1)^2+o(\rho^2)
  \right\rangle
\]

\noindent\textbf{提示}\quad （可直接使用公式计算.）
\[
  f(x,y)=f(x_0,y_0)+
  \left(h\frac{\partial}{\partial x}+k\frac{\partial}{\partial y}\right)f(x_0,y_0)
  +\frac1{2!}\left(h\frac{\partial}{\partial x}+k\frac{\partial}{\partial y}\right)^2f(x_0,y_0)
  +o(\rho^2),
\]
其中 $h=x-x_0,k=y-y_0,\rho^2=h^2+k^2$，这里 $f(x,y)\equiv y^x,(x_0,y_0)=(1,1)$（本题目的是考公式与计算能力）.

\noindent\textbf{6.3.2}\quad 求 $f(x,y)=e^x\cos y$ 在 $(0,0)$ 点带 Peano 余项的 Taylor 展开式至四阶项.（北京大学）
\[
  \left\langle
  1+x+\frac12(x^2-y^2)+\frac16x^3-\frac12xy^2
  +\frac1{24}x^4-\frac14x^2y^2+\frac1{24}y^4+o(\rho^4)
  \right\rangle
\]

\noindent\textbf{提示}\quad 利用 $e^x$ 和 $\cos y$ 的展开式相乘，或直接使用公式计算.

\noindent\textbf{再提示}\quad
\[
  f(x,y)=\left(1+x+\frac12x^2+\frac1{3!}x^3+\frac1{4!}x^4+o(x^4)\right)
  \left(1-\frac1{2!}y^2+\frac1{4!}y^4+o(y^4)\right).
\]

\noindent\textbf{6.3.3}\quad 求函数 $f(x,y)=2x^2-xy+y^2$ 在 $(1,-2)$ 处的 Taylor 展开式.
\[
  \left\langle
  8+6(x-1)-5(y+2)+2(x-1)^2-(x-1)(y+2)+(y+2)^2
  \right\rangle
\]

\noindent\textbf{提示}\quad 可用代换法，令 $x-1=u,y+2=v$.

\noindent\textbf{6.3.4}\quad $|x|,|y|$ 很小时，求
\[
  \arctan\frac{1+x+y}{1-x+y}
\]
的近似多项式，准确到 $x,y$ 的二次项.
\[
  \left\langle\frac{\pi}{4}+x-xy\right\rangle
\]

\noindent\textbf{提示}\quad 可用
\[
  \arctan\frac{1+u}{1-u}=\arctan1+\arctan u,\quad
  u=\frac{x}{1+y}.
\]

\noindent\textbf{再提示}\quad 原式
\[
  =\frac{\pi}{4}+\arctan\frac{x}{1+y}
  =\frac{\pi}{4}+\arctan x(1-y+y^2+o(y^2))
\]
\[
  =\frac{\pi}{4}+x(1-y+y^2+o(y^2))+o(x^2)
  =\frac{\pi}{4}+x-xy+o(xy^2)+o(x^2)
  \approx\frac{\pi}{4}+x-xy.
\]

\noindent\textbf{6.3.5}\quad 写出
\[
  f(x,y)=\int_0^1(1+x)^{t^2y}\,dt
\]
的 Maclaurin 级数的前面不为零的三项.
\[
  \left\langle1+\frac13xy-\frac16x^2y\right\rangle
\]

\noindent\textbf{提示}\quad
\[
  (1+x)^{t^2y}=e^{t^2y\ln(1+x)}
  =1+t^2y\left(x-\frac{x^2}{2}+\cdots\right)+\cdots.
\]

\noindent\textbf{6.3.6}\quad 设 $z$ 为由方程
\[
  z^3-2xz+y=0
\]
所定义的 $x$ 和 $y$ 的隐函数，当 $x=1$ 和 $y=1$ 时它的值为 $z=1$. 试写出函数 $z$ 按二项式 $x-1$ 和 $y-1$ 的升幂排列的展开式中的若干项.
\[
  \left\langle
  1+2(x-1)-(y-1)+[8(x-1)^2-10(x-1)(y-1)+3(y-1)^2]+\cdots
  \right\rangle
\]

\noindent\textbf{提示}\quad 用隐函数求导方法求 $z$ 在 $(1,1,1)$ 处对 $x,y$ 的低阶（如一至二阶）的偏导数，代入 Taylor 级数.

\noindent\textbf{偏导数的几何应用}

\noindent\textbf{注}\quad 曲面 $F(x,y,z)=0$ 的法向量为 $(F_x',F_y',F_z')$，空间曲线 $x=x(t),y=y(t),z=z(t)$ 在 $t=t_0$ 处的切向量为 $(x'(t_0),y'(t_0),z'(t_0))$.

\noindent\textbf{6.3.7}\quad 求曲面 $e^x-z+xy=3$ 在点 $(2,1,0)$ 处的切平面方程.（四川大学）
\[
  \left\langle x-2+2(y-1)=0\right\rangle
\]

\noindent\textbf{提示}\quad 记 $F(x,y,z)=e^x-z+xy-3$，则有
\[
  F_x'(2,1,0)(x-2)+F_y'(2,1,0)(y-1)+F_z'(2,1,0)z=0.
\]

\noindent\textbf{\ensuremath{\star}6.3.8}\quad 过直线
\[
  l:\begin{cases}
    10x+2y-2z=27,\\
    x+y-z=0
  \end{cases} \tag{1}
\]
作曲面
\[
  3x^2+y^2-z^2=27 \tag{2}
\]
的切面，求此切平面方程.（中南大学）

\noindent\textbf{提示}\quad 过直线 $l$ 之平面束为
\[
  10x+2y-2z-27+k(x+y-z)=0\quad(\forall k\in\mathbb{R}). \tag{3}
\]
求 $k$ 使 (3) 与 (2) 相切.

\noindent\textbf{再提示}\quad 让 (3) 与 (2) 的法向量相互平行，则方向数成比例：
\[
  \frac{6x}{10+k}=\frac{2y}{2+k}=\frac{-2z}{-2-k}=2t,
\]
得
\[
  x=\frac13(10+k)t,\quad y=(2+k)t,\quad z=(2+k)t.
\]
代入 (2), (3)，得
\[
  \begin{cases}
    \left[\dfrac13(10+k)^2+(2+k)^2-(2+k)^2\right]t^2=27,\\
    100t+20kt+k^2t-81=0,
  \end{cases}
\]
解得 $t=1,k=-1$ 或 $k=-19$. 再代入 (3) 得
\[
  9x+y-z=27
  \quad\text{和}\quad
  9x+17y-17z+27=0.
\]

\noindent\textbf{6.3.9}\quad 已知平面方程为
\[
  lx+my+nz=p, \tag{1}
\]
与椭球面方程
\[
  \frac{x^2}{a^2}+\frac{y^2}{b^2}+\frac{z^2}{c^2}=1 \tag{2}
\]
相切，证明系数满足方程
\[
  a^2l^2+b^2m^2+c^2n^2=p^2. \tag{3}
\]
（武汉大学）

\noindent\textbf{提示}\quad 由 (1) 与 (2) 相切可知，在切点处两法向量平行.

\noindent\textbf{再提示}\quad 设在 $(x,y,z)$ 处相切，法向量
\[
  n_1=(l,m,n),\quad
  \frac12n_2=\left(\frac{x}{a^2},\frac{y}{b^2},\frac{z}{c^2}\right).
\]
由 $n_1\parallel n_2$ 知
\[
  \frac{\dfrac{x}{a^2}}{l}
  =\frac{\dfrac{y}{b^2}}{m}
  =\frac{\dfrac{z}{c^2}}{n}=k,
\]
即
\[
  x=kla^2,\quad y=kmb^2,\quad z=knc^2. \tag{4}
\]
(4) 代入 (2) 得
\[
  k=(a^2l^2+b^2m^2+c^2n^2)^{-1/2}.
\]
代回 (4)，再将 (4) 代入 (1) 即得 (3).

\noindent\textbf{6.3.10}\quad 试证曲面 $S:xyz=a^2$ 在任何一点处切平面与三坐标平面所围成的立体体积为定值.（合肥工业大学）

\noindent\textbf{提示}\quad 切平面 $F_x'\cdot(X-x)+F_y'\cdot(Y-y)+F_z'\cdot(Z-z)=0$ 上，流动点 $(X,Y,Z)$ 在坐标轴上可分别求出三个截距
\[
  \frac{3a^2}{yz},\quad \frac{3a^2}{xz},\quad \frac{3a^2}{xy},
\]
可验证四面体体积为常数. 事实上，
\[
  V=\frac13\left[\frac12\left(\frac{3a^2}{yz}\cdot\frac{3a^2}{xz}\right)\right]:
  \frac{3a^2}{xy}
  =\frac92a^2\quad(\forall(x,y,z)\in S).
\]

\noindent\textbf{6.3.11}\quad 证明：曲面
\[
  \sqrt{x}+\sqrt{y}+\sqrt{z}=\sqrt{a}\quad(a>0)
\]
的切平面在坐标轴上割下的诸线段，其和为常量.（华东理工大学）

\noindent\textbf{提示}\quad $(x,y,z)$ 处之切平面
\[
  \frac{X}{\sqrt{ax}}+\frac{Y}{\sqrt{ay}}+\frac{Z}{\sqrt{az}}=1,
\]
截距之和
\[
  \sqrt{ax}+\sqrt{ay}+\sqrt{az}
  =\sqrt a\cdot\sqrt a=a\quad(\forall(x,y,z)\in\text{该曲面}).
\]

\noindent\textbf{\ensuremath{\star}6.3.12}\quad 求曲线 $C:x=t,y=-t^2,z=t^3$ 上与平面 $\pi:x+2y+z=4$ 平行的切线方程.（大连理工大学）

\noindent\textbf{提示}\quad 可求出曲线 $C$ 的切向量 $\tau$ 和平面 $\pi$ 的法向量 $n$，$\tau\parallel\pi\Leftrightarrow \tau\perp n\Leftrightarrow \tau\cdot n=0$.

\noindent\textbf{再提示}\quad
\[
  \tau=(x_t',y_t',z_t')=(1,-2t,3t^2),\quad n=(F_x',F_y',F_z')=(1,2,1).
\]
\[
  \tau\cdot n=1\cdot1+2\cdot(-2t)+1\cdot3t^2=0,
\]
解得 $t=1,\dfrac13$.

当 $t=1$ 时，切线为
\[
  \frac{x-1}{1}=\frac{y+1}{-2}=\frac{z-1}{3};
\]
当 $t=\dfrac13$ 时，切线为
\[
  \frac{x-\dfrac13}{1}=\frac{y+\dfrac19}{-\dfrac23}
  =\frac{z-\dfrac1{27}}{\dfrac13},
\]
即
\[
  \frac{3x-1}{3}=\frac{9y+1}{-6}=\frac{27z-1}{9}.
\]

\noindent\textbf{\ensuremath{\star}6.3.13}\quad 求曲线
\[
  C:\begin{cases}
    x^2+y^2+z^2=6,\\
    x+y+z=0
  \end{cases}
\]
在点 $M(1,-2,1)$ 的切线及法平面方程.（北京科技大学）
\[
  \left\langle
  \frac{x-1}{-6}=\frac{z-1}{6},\ y+2=0;\ x-z=0
  \right\rangle
\]

\noindent\textbf{提示}\quad 曲线 $C$ 是椭球面 ① 与平面 ② 的交线，可分别求出 ①, ② 在点 $M$ 处的法向量 $n_1,n_2$，这时 $n_1\times n_2$ 就是曲线 $C$ 在点 $M$ 的切向量. 由此可写出 $C$ 在 $M$ 处的切线与法平面方程.

\noindent\textbf{再提示}\quad
\[
  n_1=(2x,2y,2z),\quad n_2=(1,1,1),
\]
\[
  n_1\times n_2\big|_M=
  \begin{vmatrix}
    \mathbf{i}&\mathbf{j}&\mathbf{k}\\
    2x&2y&2z\\
    1&1&1
  \end{vmatrix}_M
  =(-6,0,6).
\]

\noindent\textbf{\ensuremath{\star}6.3.14}\quad 求椭球面
\[
  S:3x^2+y^2+z^2=16, \tag{1}
\]
与球面
\[
  x^2+y^2+z^2=14 \tag{2}
\]
在点 $P_0(-1,2,3)$ 处的交角.（武汉大学）
\[
  \left\langle \theta=\arccos\frac{8}{\sqrt{77}}\right\rangle
\]

\noindent\textbf{提示}\quad
\[
  n_1\cdot n_2=|n_1||n_2|\cos\theta,\quad
  \cos\theta\big|_{P_0}=
  \frac{12+16+36}{\sqrt{88}\cdot\sqrt{56}}
  =\frac8{\sqrt{77}}.
\]

\noindent\textbf{\ensuremath{\star}6.3.15}\quad 在曲面
\[
  x^2+2y^2+3z^2+2xy+2xz+4yz=8
\]
上求出切平面平行于坐标平面的诸切点.

\noindent\textbf{提示}\quad 利用法向量与坐标轴平行.

\noindent\textbf{再提示}\quad 曲面上某点 $(x,y,z)$ 处法向量 $n$，有
\[
  \frac12n=(x+y+z,x+2y+2z,x+2y+3z)\parallel k=(0,0,1),
\]
可得
\[
  \begin{cases}
    x+y+z=0,\\
    x+2y+2z=0,\\
    x+2y+3z=\alpha,
  \end{cases}
\]
即 $x=0,-y=z=\alpha$. 代入曲面方程得 $\alpha=\pm2\sqrt2$，故平行 $xOy$ 平面的切平面之切点为
\[
  (0,\pm2\sqrt2,\mp2\sqrt2).
\]
类似可求出平行另外两个坐标平面之切面切点.

\noindent\textbf{6.3.16}\quad 在附球面
\[
  \frac{x^2}{a^2}+\frac{y^2}{b^2}+\frac{z^2}{c^2}=1
\]
上怎样的点，使椭球面的法线与三坐标轴成等角？
\[
  \left\langle
  \text{共有两点：}(\pm a^2|k|,\pm b^2|k|,\pm c^2|k|),
  \ |k|=(a^2+b^2+c^2)^{-1/2}
  \right\rangle
\]

\noindent\textbf{提示}\quad 法向量与三坐标轴成等角，则法向量的三分量应相等.

\noindent\textbf{再提示}\quad
\[
  \frac12n=\left(\frac{x}{a^2},\frac{y}{b^2},\frac{z}{c^2}\right)
\]
与三坐标轴成等角 $\alpha$，则
\[
  \frac{x}{a^2}=\frac{y}{b^2}=\frac{z}{c^2}=|n|\cos\alpha\triangleq k,
\]
代入曲面方程得
\[
  k=\frac{\pm1}{\sqrt{a^2+b^2+c^2}}.
\]

\noindent\textbf{6.3.17}\quad 证明：锥面
\[
  z=xf\left(\frac{y}{x}\right)
\]
的切平面经过其顶点.

\noindent\textbf{说明}\quad 设从原点出发的射线 $l$，与 $x,y,z$ 轴夹角记为 $\alpha,\beta,\gamma$，则该曲面正是由此类射线组成，当且仅当 $\alpha,\beta,\gamma$ 满足
\[
  \cos\gamma=\cos\alpha\,f\left(\frac{\cos\beta}{\cos\alpha}\right)
\]
时，整个射线 $l$ 就在该锥面上，与动点的向径无关. 可见该曲面是锥面，且原点为锥顶.

\noindent\textbf{提示}\quad $\forall(x_0,y_0,z_0)\in S$，此点之切平面
\[
  z=\left[f\left(\frac{y_0}{x_0}\right)
  -\frac{y_0}{x_0}f'\left(\frac{y_0}{x_0}\right)\right]x
  +f'\left(\frac{y_0}{x_0}\right)y
\]
恒过原点 $(0,0,0)$.

\noindent\textbf{6.3.18}\quad 求椭球面
\[
  x^2+y^2+z^2-xy=1
\]
在坐标平面上的投影.

\noindent\textbf{提示}\quad 例如可通过椭球面上法线与 $z$ 轴垂直的点在 $xOy$ 平面上的投影，以得到该椭球面在 $xOy$ 平面上的投影区域.

\noindent\textbf{再提示}\quad 椭球面的法向量
\[
  n=(2x-y,2y-x,2z),
\]
\[
  n\perp(0,0,1)\Leftrightarrow n\cdot(0,0,1)=2z=0.
\]
代入椭球面方程得
\[
  x^2+y^2-xy=1.
\]
在空间这是平行 $z$ 轴的柱面，在 $xOy$ 平面上它是包围原点的封闭曲线. 椭球面在 $xOy$ 平面上的投影区域就是它的内部：
\[
  r^2\le\frac1{1-\frac12\sin2\theta}.
\]
类似可求出椭球面在另外两个坐标平面上的投影区域.

\noindent\textbf{\ensuremath{\star}多元极值及其应用}

\noindent\textbf{注}\quad 本类考研试题甚多，适合各类读者，须倍加注意.

\noindent\textbf{\ensuremath{\star}6.3.19}\quad 若 $M_0$ 是 $f(x,y)$ 的极小值点，且在 $M_0$ 点 $f_{xx}'',f_{yy}''$ 存在，试证在 $M_0$ 点，$f_{xx}''+f_{yy}''\ge0$.（江西师范大学）

\noindent\textbf{提示}\quad 可用 Taylor 公式.

\noindent\textbf{再提示}\quad 因 $f_x'\big|_{M_0}=f_y'\big|_{M_0}=0$，
\[
  0\le f(x_0+h,y_0)-f(x_0,y_0)
  =f_{xx}''\big|_{M_0}\cdot\frac{h^2}{2}+o(h^2),
\]
\[
  0\le f(x_0,y_0+h)-f(x_0,y_0)
  =f_{yy}''\big|_{M_0}\cdot\frac{h^2}{2}+o(h^2).
\]
两式相加，同除以 $h^2$，有
\[
  \frac12(f_{xx}''+f_{yy}'')\big|_{M_0}+o(1)\ge0,
\]
再令 $h\to0$ 即得.

\noindent\textbf{注}\quad 本题给出了极小值点的一个必要条件.

\noindent\textbf{6.3.20}\quad 设 $z=f(x,y)$ 在有界闭区域 $D$ 内有二阶连续的偏导数，且 $f_{xx}''+f_{yy}''=0,f_{xx}''\ne0$. 证明：$z=f(x,y)$ 的最大值和最小值只能在区域的边界上取得.（华中师范大学）

\noindent\textbf{提示}\quad 只需证 $D$ 的内部
\[
  f_{xx}''f_{yy}''-(f_{xy}'')^2<0.
\]

\noindent\textbf{再提示}\quad 在 $D$ 之内部，
\[
  2f_{xx}''f_{yy}''\le(f_{xx}''+f_{yy}'')^2=0,\quad f_{xx}''\ne0,
\]
$f_{xy}''{}^2>0$，故
\[
  f_{xx}''f_{yy}''-(f_{xy}'')^2<0.
\]
可知 $D$ 内
部无极值. 但 $z=f(x,y)$ 在有界闭区域上连续，必有最大、最小值. 得证.

\noindent\textbf{\ensuremath{\star}6.3.21}\quad 求曲面 $z=xy-1$ 上与原点最近的点的坐标.（中山大学）

\noindent\textbf{解 I}\quad （直接法）曲面上一点 $(x,y,xy-1)$ 到原点距离为
\[
  \rho=\sqrt{x^2+y^2+(xy-1)^2}
  \ge\sqrt{(x-y)^2+1}\ge1=\rho(0,0,-1).
\]
故点 $(0,0,-1)$ 是该曲面上与原点最近之点，其距离为 $1$.

\noindent\textbf{解 II}\quad （化为自由极值.）设
\[
  s=x^2+y^2+(xy-1)^2,
\]
令 $\dfrac{\partial s}{\partial x}=\dfrac{\partial s}{\partial y}=0$，得
\[
  \begin{cases}
    x-y+xy^2=0,\\
    y-x+x^2y=0.
  \end{cases}
\]
解之，有唯一可疑点 $(0,0)$. 因
\[
  \Delta\big|_{(0,0)}
  =\left[\frac{\partial^2s}{\partial x^2}\cdot
  \frac{\partial^2s}{\partial y^2}
  -\left(\frac{\partial^2s}{\partial x\partial y}\right)^2\right]\bigg|_{(0,0)}
  =4+4>0,\quad
  \frac{\partial^2s}{\partial x^2}\bigg|_{(0,0)}=2>0,
\]
故 $(0,0)$ 是极小值点. 点 $(0,0,-1)$ 至原点最近.

\noindent\textbf{解 III}\quad （用 Lagrange 乘数法.）$\rho=\sqrt{x^2+y^2+z^2}$ 与 $\rho^2$ 的极值点相同，设
\[
  L=x^2+y^2+z^2+\lambda(z-xy+1).
\]
令 $L_x'=L_y'=L_z'=L_\lambda'=0$，解方程得 $x=y=0,z=-1$. $z=xy-1$ 是鞍点在 $(0,0,-1)$ 的马鞍面，站在 $xOy$ 平面观察，曲面在一、三象限无限向上延伸，二、四象限无限向下延伸，$\rho=\sqrt{x^2+y^2+z^2}$ 无最大值. $(0,0,-1)$ 只能是最小值点. 若以 $G_r$ 表示原点为中心、半径为 $r$ 的闭圆域，则
\[
  \rho=\sqrt{x^2+y^2+(xy-1)^2}\ (>r)
\]
在 $G_r$ 上连续，有最大、最小值. $r\to+\infty$ 时 $\rho\to+\infty$，可见 $r$ 充分大时，$\rho$ 的最小值只能在内部达到. 而内部只有唯一的可疑点，故它必是最小值点（这种判断方法时常用到）.

\noindent\textbf{\ensuremath{\star}6.3.22}\quad 求两曲面
\[
  x+2y=1\quad\text{和}\quad x^2+2y^2+z^2=1
\]
的交线上距原点最近的点.（中国科学院）
\[
  \left\langle\left(-\frac13,\frac23,0\right)\right\rangle
\]

\noindent\textbf{提示}\quad 可用 Lagrange 乘数法，还可转化成一元最值问题.

\noindent\textbf{解 I}\quad 设
\[
  L=x^2+y^2+z^2+\lambda(x+2y-1)+\mu(x^2+2y^2+z^2-1).
\]
\[
  L_x'=2x+\lambda+2x\mu=0, \tag{1}
\]
\[
  L_y'=2y+2\lambda+4y\mu=0, \tag{2}
\]
\[
  L_z'=2z+2z\mu=0, \tag{3}
\]
\[
  x+2y=1, \tag{4}
\]
\[
  x^2+2y^2+z^2=1. \tag{5}
\]
由 (3) 得 $z=0$. 代入 (5)，并与 (4) 联立，可得
\[
  (x,y,z)=(1,0,0)\quad\text{和}\quad
  (x,y,z)=\left(-\frac13,\frac23,0\right).
\]
$\rho=\sqrt{x^2+y^2+z^2}$ 在交线（有界闭集）上连续，必有最大、最小值，现只有两个可疑点
\[
  \rho\big|_{(1,0,0)}=1>\rho\bigg|_{\left(-\frac13,\frac23,0\right)}
  =\frac13\sqrt5,
\]
故距原点最近的点为 $\left(-\dfrac13,\dfrac23,0\right)$.

\noindent\textbf{解 II}\quad （直接法）（化为一元函数最值.）将 $x=1-2y$ 代入椭球面方程得
\[
  1-4y+6y^2+z^2=1,
\]
即
\[
  z=\pm\sqrt2\sqrt{y(2-3y)}\quad
  \left(0\le y\le\frac23\right),\quad y_{\max}=\frac23.
\]
\[
  \rho=\sqrt{x^2+y^2+z^2}
  =\sqrt{1-y^2}\ge\sqrt{1-y_{\max}^2}
  =\sqrt{1-\left(\frac23\right)^2}=\frac{\sqrt5}{3}.
\]
由 $y=y_{\max}$，知 $x=-\dfrac13,z=0$，故交线上点 $\left(-\dfrac13,\dfrac23,0\right)$ 距原点最近.

\noindent\textbf{6.3.23}\quad 在平面上求一点，使它到 $n$ 个定点 $(x_1,y_1),(x_2,y_2),\cdots,(x_n,y_n)$ 的距离之平方和最小.（西北工业大学）
\[
  \left\langle
  \frac1n\sum_{i=1}^n x_i,\ \frac1n\sum_{i=1}^n y_i
  \right\rangle
\]

\noindent\textbf{提示}\quad 可以
\[
  u=\sum_{i=1}^n[(x-x_i)^2+(y-y_i)^2]
\]
作目标函数，令 $u_x'=u_y'=0$，得可疑点 $M$ 的坐标
\[
  (x_0,y_0)=\left(\frac1n\sum_{i=1}^n x_i,\frac1n\sum_{i=1}^n y_i\right).
\]
而
\[
  (u_{xx}''u_{yy}''-u_{xy}''{}^2)\big|_M=2n\cdot2n-0>0,\quad
  u_{xx}''\big|_M>0.
\]

\noindent\textbf{6.3.24}\quad 抛物面 $z=x^2+y^2$ 被平面 $x+y+z=1$ 截成一椭圆，求原点至该椭圆最近、最远距离.（北京航空航天大学）

\noindent\textbf{提示}\quad 以 $u=x^2+y^2+z^2$ 作目标函数，约束条件为 $z-x^2-y^2=0$ 和 $x+y+z-1=0$，用 Lagrange 乘数法可得可疑点为
\[
  P=\left(\frac{-1-\sqrt3}{2},\frac{-1-\sqrt3}{2},2+\sqrt3\right),
  \quad
  M=\left(\frac{-1+\sqrt3}{2},\frac{-1+\sqrt3}{2},2-\sqrt3\right).
\]
根据实际背景，原点上方的一条椭圆曲线必有一个最近点和一个最远点（理论上讲：连续函数在有界闭集上必有最大、最小值）. 而
\[
  u(P)=9+5\sqrt3>u(M)=9-5\sqrt3.
\]
可见离原点最近距离为
\[
  \sqrt{u(M)}=\sqrt{9-5\sqrt3},
\]
最远距离为
\[
  \sqrt{u(P)}=\sqrt{9+5\sqrt3}.
\]

\noindent\textbf{\ensuremath{\star}6.3.25}\quad 求 $u=ky^3+zx$ 在条件 $x^2+y^2+z^2=1,z\ge0$ 下的最大值和最小值.（清华大学）

\noindent\textbf{解}\quad $1^\circ$ 若 $k$ 为变数，$u$ 作为四元函数，在所给条件下明显无最大、最小值，因为：如 $k\to+\infty$，当 $x=z=0,y=1$ 时，$u=k\to+\infty$；当 $x=z=0,y=-1$ 时，$u=-k\to-\infty$.

$2^\circ$ 若 $k$ 为常数.

i）求 $z>0,x^2+y^2+z^2=1$ 上 $u$ 的极值点. 作
\[
  L=ky^3+zx+\lambda(x^2+y^2+z^2-1).
\]
令 $L_x'=L_y'=L_z'=L_\lambda'=0$，得
\[
  \begin{cases}
    z+2\lambda x=0,\\
    3ky^2+2\lambda y=0,\\
    x+2\lambda z=0,\\
    x^2+y^2+z^2=1.
  \end{cases}
\]
由②可知，$y=0$ 或 $y=-\dfrac{2\lambda}{3k}$. 由①, ③解得 $\lambda=\pm\dfrac12,x=\mp z$.

当 $y=0$ 时，得 $x=\mp\dfrac1{\sqrt2}$；

当 $y=-\dfrac{2\lambda}{3k}=\mp\dfrac1{3k}$ 时，
\[
  x=\mp\frac{\sqrt{9k^2-1}}{3\sqrt2|k|},\quad
  z=\pm\frac{\sqrt{9k^2-1}}{3\sqrt2|k|}
  \quad(\text{要求 }9k^2-1\ge0,\ |k|\ge\frac13).
\]
记
\[
  P_{1,2}=\left(\mp\frac1{\sqrt2},0,\frac1{\sqrt2}\right),
\]
\[
  M_{1,2}=\left(\mp\frac{\sqrt{9k^2-1}}{3\sqrt2|k|},
  \mp\frac1{3k},
  \frac{\sqrt{9k^2-1}}{3\sqrt2|k|}\right),
\]
\[
  u\big|_{P_1}=-\frac12,\quad
  u\big|_{P_2}=\frac12,\quad
  u\big|_{M_1}=-\frac12+\frac1{54k^2}>-\frac12=u\big|_{P_1},
\]
\[
  u\big|_{M_2}=\frac12-\frac1{54k^2}<\frac12=u\big|_{P_2}.
\]
可见 $u(P_2)>u(M_2)>u(M_1)>u(P_1)$，$M_1,M_2$ 不是最值点.

$u$ 通过 $z$ 依赖于 $x,y$，由隐函数求导不难求出
\[
  u_{xx}''\big|_{P_2}=-4,\quad
  u_{yy}''\big|_{P_2}=-1,\quad
  u_{xy}''\big|_{P_2}=0,\quad
  \Delta\big|_{P_2}>0,\quad
  u_{xx}''\big|_{P_2}<0,
\]
故 $P_2$ 为极大值点. 类似可得 $P_1$ 为极小值点. 故在上半椭球面上（不包括边界）：
\[
  u_{\max}=u(P_2)=\frac12,\quad
  u_{\min}=u(P_1)=-\frac12.
\]

ii）在半椭球面的边界上（$z=0,x^2+y^2=1$），$u=ky^3+zx=ky^3$，
\[
  u_{\max}=|k|,\quad u_{\min}=-|k|.
\]
总之，在上半椭球面包括边界上，
\[
  u_{\max}=
  \begin{cases}
    |k|, & |k|>\dfrac12,\\
    \dfrac12, & |k|\le\dfrac12,
  \end{cases}
  \quad
  u_{\min}=
  \begin{cases}
    -|k|, & |k|>\dfrac12,\\
    -\dfrac12, & |k|\le\dfrac12.
  \end{cases}
\]

\noindent\textbf{注}\quad 解法第 $2^\circ$ 部分，也可不用 Lagrange 乘数法，而把 $u=ky^3+zx$ 看成通过 $z(z=\sqrt{1-x^2-y^2})$ 依赖于 $x,y$ 的二元函数（$x^2+y^2\le1$）. 用隐函数微分法可求出 $u_x',u_y'$，令之为零，也可找出同样的可疑点.

\noindent\textbf{\ensuremath{\star}6.3.26}\quad 求函数 $u=x^2-y^2+2xy$ 在单位圆 $x^2+y^2\le1$ 上的最大、最小值.（北京科技大学）
\[
  \left\langle\sqrt2,-\sqrt2\right\rangle
\]

\noindent\textbf{提示}\quad 在圆内为自由极值，边界上 $x^2+y^2=1$ 为条件极值.

\noindent\textbf{再提示}\quad 在圆内令 $u_x'=0,u_y'=0$ 得
\[
  \begin{cases}
    2x+2y=0,\\
    -2y+2x=0,
  \end{cases}
\]
只有唯一解 $(0,0)$，这里 $u(0,0)=0$.

在边界上 $r^2=x^2+y^2=1$，
\[
  u=\cos^2\theta-\sin^2\theta+2\sin\theta\cos\theta
  =\cos2\theta+\sin2\theta
\]
\[
  =\sqrt2\left(\sin\frac\pi4\cos2\theta+\cos\frac\pi4\sin2\theta\right)
  =\sqrt2\sin\left(2\theta+\frac\pi4\right).
\]
当 $\theta=\dfrac\pi8,\dfrac58\pi$ 时分别达到最大值 $\sqrt2$ 和最小值 $-\sqrt2$.

因为连续函数在有界闭区域上必有最大、最小值，现只有三个可疑点，故比较三点之值即得.

\noindent\textbf{6.3.27}\quad 在直线 $x+y=\dfrac\pi2$ 位于第一象限的那一段上求一点，使该点横坐标的余弦与纵坐标的余弦之乘积最大并求出最大值.（华中理工大学，西北工业大学）
\[
  \left\langle\left(\frac\pi4,\frac\pi4\right),\frac12\right\rangle
\]

\noindent\textbf{提示}\quad 除用 Lagrange 乘数法之外，还可用直接法（用定义）.

\noindent\textbf{再提示}\quad 目标函数
\[
  f(x,y)=\cos x\cos y\big|_{y=\frac\pi2-x}
  =\cos x\cos\left(\frac\pi2-x\right)
  =\frac12\sin2x
  \le f\left(\frac\pi4,\frac\pi4\right)=\frac12
  \quad\left(\forall x\in\left[0,\frac\pi2\right]\right).
\]

\noindent\textbf{\ensuremath{\star}6.3.28}\quad 求直线
\[
  4x+3y=16 \tag{1}
\]
与椭圆
\[
  18x^2+5y^2=45 \tag{2}
\]
之间的最短距离.（华中科技大学）

\noindent\textbf{提示}\quad $1^\circ$ 可用 Lagrange 乘数法，任一直线 $ax+by=c$，法线方程为
\[
  \frac{ax+by-c}{\sqrt{a^2+b^2}}=0,
\]
这时线外一点 $(x_0,y_0)$ 到直线的距离为
\[
  \rho(x_0,y_0)=\frac{|ax_0+by_0-c|}{\sqrt{a^2+b^2}}.
\]
因此本题是求目标函数
\[
  \rho=\frac{|4x+3y-16|}{5}
\]
在条件 $18x^2+5y^2=45$ 下的最小值.

$2^\circ$ 还可用几何、代数方法. 平行直线 (1) 的直线 $4x+3y=c$，如果是椭圆 (2) 的切线，则
\[
  \left(x,\frac13(c-4x)\right)
\]
应能满足方程 (2) 且仅有唯一根. 关于 $x$ 的二次方程
\[
  18x^2+5\left[\frac13(c-4x)\right]^2=45
\]
之判别式为零，可求出待定常数 $c=\pm11$. 再在切线 $4x+3y=\pm11$ 上任取一点，如 $\left(0,\pm\dfrac{11}{3}\right)$，代入
\[
  \frac{|4x+3y-16|}{5},
\]
即可得相应距离 $1$ 和 $\dfrac{27}{5}$，故 $\rho_{\min}=1$.

\noindent\textbf{\ensuremath{\star}6.3.29}\quad 证明：在光滑曲面 $F(x,y,z)=0$ 上离原点最近的点处的法线必过原点.（武汉理工大学）

\noindent\textbf{提示}\quad 求曲面上距离原点最近的点，作出此点的法线，验证原点在法线上即可.

\noindent\textbf{再提示}\quad
\[
  L=x^2+y^2+z^2+\lambda F(x,y,z).
\]
令
\[
  L_x'=2x+\lambda F_x'(x,y,z)=0,
\]
\[
  L_y'=2y+\lambda F_y'(x,y,z)=0,
\]
\[
  L_z'=2z+\lambda F_z'(x,y,z)=0.
\]
可见曲面上某点 $(x_0,y_0,z_0)$ 若离原点最近，则
\[
  x_0:y_0:z_0=F_x'(x_0,y_0,z_0):F_y'(x_0,y_0,z_0):F_z'(x_0,y_0,z_0).
\]
此处的法线可写成
\[
  \frac{x-x_0}{x_0}=\frac{y-y_0}{y_0}=\frac{z-z_0}{z_0},
\]
显然 $(0,0,0)$ 满足方程.

\noindent\textbf{\ensuremath{\star}6.3.30}\quad 利用导数证明周长一定的三角形中以等边三角形的面积最大.（清华大学）

\noindent\textbf{提示}\quad 设三角形三边长为 $x,y,z$，三角形的周长 $l=2p$，则面积
\[
  S=\sqrt{p(p-x)(p-y)(p-z)}.
\]
因此可取 Lagrange 函数为
\[
  L=p(p-x)(p-y)(p-z)+\lambda(x+y+z-2p),
\]
约束条件为 $x+y+z=2p$ 和 $0<x,y,z<p$.

\noindent\textbf{6.3.31}\quad 在曲面
\[
  x^2+y^2+\frac{z^2}{4}=1\quad(x>0,y>0,z>0)
\]
上求一点，使过该点的切平面在三个坐标轴上的截距平方和最小.（复旦大学）
\[
  \left\langle\left(\frac12,\frac12,\sqrt2\right)\right\rangle
\]

\noindent\textbf{提示}\quad 椭球上点 $(x,y,z)$ 处的切平面为
\[
  xX+yY+\frac{z}{4}Z=1\quad((X,Y,Z)\text{ 为切面上流动点}),
\]
在坐标上的截距分别为
\[
  \frac1x,\quad \frac1y,\quad \frac4z.
\]
可取
\[
  L=\frac1{x^2}+\frac1{y^2}+\frac{16}{z^2}
  +\lambda\left(x^2+y^2+\frac14z^2-1\right).
\]

\noindent\textbf{利用极值证明不等式}

\noindent\textbf{\ensuremath{\star}6.3.32}\quad 证明
\[
  \sin x\sin y\sin(x+y)\le\frac{3\sqrt3}{8}\quad(0<x,y<\pi),
\]
并确定何时等号成立.（中国科学院）

\noindent\textbf{提示}\quad 可考虑 $f(x,y)=\sin x\sin y\sin(x+y)$，证明
\[
  f_{\max}\le\frac{3\sqrt3}{8}\quad(0<x,y<\pi).
\]

\noindent\textbf{再提示}\quad 由 $f_x'=f_y'=0$ 可知，
\[
  x=y=\frac\pi3,\quad f\left(\frac\pi3,\frac\pi3\right)=\frac{3\sqrt3}{8}.
\]
而在有界闭区域 $[0,\pi]\times[0,\pi]$ 的边界上 $f(x,y)=0$.

\noindent\textbf{6.3.33}\quad 若 $n\ge1$ 及 $x\ge0,y\ge0$，试用求极值的方法证明不等式：
\[
  \frac{x^n+y^n}{2}\ge\left(\frac{x+y}{2}\right)^n.
\]

\noindent\textbf{提示}\quad 可考虑
\[
  z=\frac{x^n+y^n}{2}
\]
在条件 $x+y=a(a>0,x\ge0,y\ge0)$ 下的极值问题，得
\[
  z=\frac{x^n+y^n}{2}\ge\left(\frac a2\right)^n.
\]
或当 $x\ne0$ 时令 $u=\dfrac yx$，化为一元问题.

\noindent\textbf{6.3.34}\quad 用条件极值法证明不等式：
\[
  \frac{x_1^2+x_2^2+\cdots+x_n^2}{n}
  \ge
  \left(\frac{x_1+x_2+\cdots+x_n}{n}\right)^2
  \quad(x_i>0,k=1,2,\cdots,n).
\]

\noindent\textbf{提示}\quad $n$ 个正数和为定数，其平方和以 $n$ 数彼此相等时为最小. 即 $f(x_1,\cdots,x_n)=x_1^2+\cdots+x_n^2$ 在条件 $x_1+\cdots+x_n=a$ 下的极值.

\noindent\textbf{6.3.35}\quad 设 $a_i>0(i=1,2,\cdots,n)$，证明：
\[
  n\left(\frac1{a_1}+\frac1{a_2}+\cdots+\frac1{a_n}\right)^{-1}
  \le(a_1a_2\cdots a_n)^{1/n}.
\]

\noindent\textbf{提示}\quad 可考虑 $f(x_1,x_2,\cdots,x_n)=x_1x_2\cdots x_n$ 在条件
\[
  \frac1{x_1}+\frac1{x_2}+\cdots+\frac1{x_n}=\frac1a\quad(x_i>0,a>0)
\]
下的极值；或将此式看作 $\dfrac1{a_1},\dfrac1{a_2},\cdots,\dfrac1{a_n}$ 的几何平均与算术平均的关系.

\noindent\textbf{6.3.36}\quad 证明不等式：
\[
  e^y+x\ln x-x-xy\ge0\quad(x\ge1,y\ge0).
\]
（厦门大学）

\noindent\textbf{提示}\quad 记
\[
  F(x,y)=e^y+x\ln x-x-xy.
\]
令 $F_y'(x,y)=e^y-x=0$，在 $xOy$ 平面上得一曲线 $C:x=e^y$. 在 $C$ 上：
\[
  F_y'(x,\ln x)\equiv0\quad(\forall x\ge1).
\]

\noindent\textbf{再提示}\quad $\forall x_0\ge1$，$z=F(x_0,y)$ 是关于 $y$ 的一元函数 $(y\ge0)$. 当 $y=y_0=\ln x_0$ 时，
\[
  F_y'(x_0,y_0)=0,\quad F_{yy}''(x_0,y_0)=e^{y_0}=x_0\ge1>0.
\]
所以 $z=F(x_0,y)$ 在 $y=y_0$ 处取最小值. 而 $F(x_0,y_0)=0$，所以
\[
  F(x,y)=e^y+x\ln x-x-xy\ge F_{\min}(x,y)
  =F(x_0,y_0)=0\quad(\text{当 }x\ge1,y\ge0\text{ 时}).
\]

\noindent\textbf{6.3.37}\quad 费马原理指出，从 $A$ 射出到达 $B$ 的光线，是沿费时最短的路线传播. 假设点 $A$ 和点 $B$ 位于以平面分开的不同的光介质中，并且光的传播速度在两个介质中分别为 $v_1$ 与 $v_2$，试由费马原理推出光的折射定律.

\begin{center}
  \includegraphics[width=0.42\textwidth]{6-3-7.png}

  图 6.3.7
\end{center}

\noindent\textbf{提示}\quad （转化为数学问题.）（如图 6.3.7）设光线以速度 $v_1$ 从界面上方点 $A$ 出发，照射到界面的 $O$ 点，穿进新介质（例如水）. 界面处被折射，入射角记为 $\alpha$，折射角记为 $\beta$. 在新介质里，速度变为 $v_2$，最后到达界面下方的点 $B$. 假设 $A,B$ 与界面的距离分别为 $a$ 和 $b$，两点之间的水平距离为 $l$，问：$x$ 为多少时，光线从 $A$ 至 $B$ 所花的时间才最节省. 亦即：求 $x_0$，使得
\begin{equation}
  T(x_0)=\left(\frac{AO}{v_1}+\frac{OB}{v_2}\right)\bigg|_{x_0}
  =
  \left(
    \frac{\sqrt{x^2+a^2}}{v_1}
    +\frac{\sqrt{(l-x)^2+b^2}}{v_2}
  \right)\bigg|_{x=x_0}
  =\min_{0\le x\le l}T(x).
\end{equation}
进而证明：（折射）费时最短的路径，正好是折射定律指示的路径：
\begin{equation}
  \frac{\sin\alpha}{\sin\beta}=\frac{v_1}{v_2}=\text{常数}.
\end{equation}

\noindent\textbf{解}\quad 令 $T'(x)=0$，有
\[
  \frac{x}{v_1\sqrt{x^2+a^2}}
  -
  \frac{l-x}{v_2\sqrt{(l-x)^2+b^2}}=0,
\]
于是
\[
  \frac{\sin\alpha}{v_1}=\frac{\sin\beta}{v_2},
\]
或
\[
  \frac{\sin\alpha}{\sin\beta}=\frac{v_1}{v_2}.
\]
由式 (1) 知 $T(x)$ 连续，当 $x$ 在 $[0,l]$ 外时，光线不满足费时最短的要求，故费时最短的点在 $[0,l]$ 内部. 由 $T'(x)=0$ 得到的点即为费时最短点，折射定律 (2) 获证.

\section*{** §6.4\quad 隐函数存在定理及函数相关}

\noindent\textbf{导读}\quad 该节理论性较强，难度较大. 在数学分析中应用不像一致收敛那样广泛，考题相对较少. 此次修订，函数相关内容从略，需要可查原书第一、二版.

\subsection*{*一、隐函数存在定理}

\noindent\textbf{a.\quad 一个方程的情况}\quad 对方程 $F(x,y)=0$ 而言（多元的情况 $x=(x_1,\cdots,x_n)$），隐函数存在定理告诉我们只要验证了条件：

1) $F(P_0)=0,F_y'(P_0)\ne0$，其中 $P_0=(x_0,y_0)$；多元情形下 $x_0=(x_{01},\cdots,x_{0n})$，$F(P)=F(x_1,\cdots,x_n,y)$；

2) $F(x,y)$ 与 $F_y'(x,y)$ 在 $P_0$ 的某邻域内连续，

\noindent 就可断言：存在 $\delta,\eta>0$ 及函数 $y=y(x)$ 满足：

i) $y_0=y(x_0)$；

ii) $F(x,y(x))\equiv0,|y(x)-y_0|<\delta,x\in U(x_0,\eta)$，其中
\[
  U(x_0,\eta)=\{x:|x-x_0|<\eta\};
\]

iii) 满足 i), ii) 的函数是唯一的；

iv) $y=y(x)$ 在 $U(x_0,\eta)$ 内连续.

此外，若 $F_x'(x,y)$ 在 $P_0$ 的某邻域内也连续，则 $y'(x)$ 存在，且
\[
  y_x'=-\frac{F_x'(x,y)}{F_y'(x,y)}
\]
（多元情形下，$y_{x_i}'=-\dfrac{F_{x_i}'(x,y)}{F_y'(x,y)}$）.

\noindent 值得注意的是，上述条件，只是充分条件而不是必要条件. 条件不满足时，隐函数是否存在，有待讨论.

\noindent\textbf{例 6.4.1}\quad 设方程
\begin{equation}
  x^2+y+\sin(xy)=0.
\end{equation}
试证：

1) 在 $(0,0)$ 的某小邻域内，方程 (1) 确定唯一的连续函数 $y=y(x),y(0)=0$；

2) 在 $x=0$ 的某小邻域内讨论函数 $y=y(x)$ 的可微性；

3) 在 $x=0$ 的某小邻域内讨论 $y=y(x)$ 的单调性；

4) 在 $(0,0)$ 的某小邻域内方程 (1) 能否确定唯一的单值函数 $x=x(y)$，满足 $x(0)=0$，为什么？（武汉大学）

\noindent\textbf{分析}\quad 1) 令
\[
  F(x,y)=x^2+y+\sin(xy),
\]
则 $F(0,0)=0,F_y'(0,0)=1\ne0$，且 $F,F_y'$ 在 $(0,0)$ 的某邻域内连续. 由隐函数存在定理可知，方程 (1) 在 $(0,0)$ 的某小邻域内确定唯一的连续函数 $y=y(x),y(0)=0$.

2) 因 $F_x'(x,y)=2x+y\cos(xy)$ 在 $(0,0)$ 的某邻域内连续，故 $y=y(x)$ 在 $x=0$ 的某邻域内可微，且
\begin{equation}
  y'(x)=-\frac{2x+y\cos(xy)}{1+x\cos(xy)}.
\end{equation}

3) 为讨论 $y=y(x)$ 的单调性，需讨论 $y'(x)$ 的符号. 由式 (2)，当 $x,y$ 充分小时，$y'(x)$ 的符号取决于 $-[2x+y\cos(xy)]$ 的符号. 由 $y(0)=0$ 及式 (2)，可知 $y'(0)=0$，从而当 $x\to0$ 时，$y=o(x)$. 于是
\[
  |y\cos(xy)|\le |y|=o(x).
\]
故 $y'(x)$ 的符号与 $-2x$ 的符号相同. 即当 $x>0$ 时，$y'<0,y(x)\searrow$；当 $x<0$ 时，$y'>0,y(x)\nearrow$. 因而 $y(x)$ 在 $x=0$ 处有严格极大值.

4) 由于 $F_x'(0,0)=0$，由隐函数存在定理不能判定方程 (1) 能否确定唯一函数 $x=x(y)$. 但由 3) 可知，这样的函数不存在. 因为 $y(x)$ 在 $x=0$ 处有严格极大值，在充分小的邻域内，当 $y<0$ 时，至少有两个 $x$ 与之对应，而当 $y>0$ 时无 $x$ 与之对应，使得 $F(x,y)=0$.

\noindent\textbf{练习}\quad 已知 $F(x,y)=2-\sin x+y^3e^{-y}$，证明 $F(x,y)=0$ 在整个平面上确定唯一的隐函数 $y=y(x)$，且 $y=y(x)$ 连续可微. （北京大学）

\noindent\textbf{提示}\quad 将方程写成
\begin{equation}
  \sin x-2=y^3e^{-y}.
\end{equation}
式 (1) 左端 $z=\sin x-2$ 的取值范围为 $[-3,-1]$. 而右端 $z=y^3e^{-y}$ 当 $y\ge0$ 时非负；当 $y<0$ 时为负. 且当 $y<3$ 时，
\[
  z_y'=(y^3e^{-y})'=y^2(3-y)e^{-y}>0,
\]
所以 $z=y^3e^{-y}$ 在 $(-\infty,0)$ 上单调，反函数存在. 因此，对每一个 $x$，$z=\sin x-2\in[-3,-1]$，并通过反函数（记为 $y=f(z)$）唯一确定一个 $y\in(-\infty,0)$. 故 $F(x,y)=0$ 在整个平面上存在唯一的隐函数
\[
  y=f(\sin x-2).
\]
又因 $F,F_y'$ 连续，且当 $(x,y)=(x,f(\sin x-2))$ 时，$F_y'=y^2(3-y)e^{-y}>0$，根据隐函数存在定理，$y=y(x)$ 连续且可微.

\noindent\textbf{例 6.4.2}\quad 设
\[
  F(x,y)=\sum_{n=1}^\infty nye^{-n(x+y)}.
\]
是否存在一个 $a>0$ 与唯一的函数 $h(x)$，使得 $h(x)$ 在 $(1-a,1+a)$ 上可微，$h(1)=0$，且 $F(x,h(x))=0$. （北京大学）

\noindent\textbf{分析}\quad 题意是：用级数定义的函数 $F$ 在点 $P(1,0)$ 的邻域里，方程 $F=0$ 是否能唯一确定可导的隐函数 $y=h(x)$. 因此任务是检验隐函数存在定理的条件是否能满足.

\noindent\textbf{解}\quad 1) 零函数满足该方程. 当 $y=0$ 时，级数各项均为零，故 $P(1,0)$ 满足方程，且零函数 $y=h(x)\equiv0$ 满足方程及 $h(1)=0$. 存在性得证，下面证明唯一性.

2) 证明 $F(x,y)$ 连续. 因各项函数连续，只需证明级数在 $P(1,0)$ 的某邻域内一致收敛. 任取 $0<\alpha<\dfrac12$，在邻域
\[
  U(P,\alpha)=\{(x,y):|x-1|<\alpha,|y|<\alpha\}
\]
内，有 $x+y>1-2\alpha>0$，故
\[
  |nye^{-n(x+y)}|
  \le n\alpha e^{-n(1-2\alpha)}
  =
  \frac{n\alpha}{\sum_{k=0}^\infty\dfrac{n^k(1-2\alpha)^k}{k!}}
  \le
  \frac{6\alpha}{n^2(1-2\alpha)^3}.
\]
（分母中，$k\ne3$ 的项删除）因 $\sum\dfrac{6\alpha}{n^2(1-2\alpha)^3}$ 收敛，故由 Weierstrass 判别法可知级数一致收敛，从而 $F(x,y)$ 连续.

3) 同理可知，由偏导函数构成的级数也一致收敛，故可逐项求偏导，且
\[
  F_y'(1,0)
  =
  \sum_{n=1}^\infty
  \left[
    ne^{-n(x+y)}-n^2ye^{-n(x+y)}
  \right]\bigg|_{x=1,y=0}
  =
  \sum_{n=1}^\infty ne^{-n}>0.
\]
因而隐函数存在定理的条件全都满足. 故在点 $P(1,0)$ 的某邻域内，方程 $F(x,y)=0$ 唯一确定可微隐函数，且该函数即为零函数.

\noindent\textbf{练习}\quad 设 $f(x,y)$ 及一阶偏导数在点 $(0,1)$ 的某邻域内存在且连续，并且 $f_y'(0,1)\ne0,f(0,1)=0$. 证明：
\[
  f\left(x,\int_0^t\sin\tau\,d\tau\right)=0
\]
在 $(0,\pi/2)$ 的某邻域内确定单值函数 $t=\varphi(x)$，并求 $\varphi'(0)$. （南京大学）

\noindent\textbf{提示}\quad 令
\[
  F(x,t)=f\left(x,\int_0^t\sin\tau\,d\tau\right)
  =f(x,1-\cos t).
\]
验证 $F(x,t)$ 在 $(0,\pi/2)$ 处满足隐函数存在定理条件即可.

\noindent\textbf{例 6.4.3}\quad 设函数 $f(x,y)$ 在 $(x_0,y_0)$ 的某邻域内二次连续可微，且
\[
  f_x'(x_0,y_0)=0,\quad f_{xx}''(x_0,y_0)>0.
\]
试证：

1) 存在 $y_0$ 的邻域 $U(y_0,\delta)$，使得对任意 $y\in U(y_0,\delta)$，函数 $f(x,y)$ 作为 $x$ 的函数有一极小值 $g(y)$；

2) $g'(y_0)=f_y'(x_0,y_0)$.

\noindent\textbf{证明}\quad 对给定的 $y$，要研究 $f(x,y)$ 作为 $x$ 的函数的极值，自然要解方程
\[
  f_x'(x,y)=0,
\]
即找出隐函数 $x=x(y)$，使得
\[
  f_x'(x(y),y)=0.
\]
由于 $f$ 二次连续可微，且
\[
  f_x'(x_0,y_0)=0,\quad f_{xx}''(x_0,y_0)>0,
\]
故上述方程满足隐函数存在定理的条件. 因而在 $(x_0,y_0)$ 的某邻域内，方程唯一确定可微函数 $x=x(y)$，满足
\[
  x_0=x(y_0),\quad f_x'(x(y),y)\equiv0.
\]
又因 $f_{xx}''(x_0,y_0)>0$，当邻域取得充分小时，有
\[
  f_{xx}''(x(y),y)>0.
\]
故 $f(x,y)$ 作为 $x$ 的函数在 $x=x(y)$ 处取极小值. 记
\[
  g(y)=f(x(y),y).
\]
最后来证 $g'(y_0)=f_y'(x_0,y_0)$. 由定义，
\[
\begin{aligned}
  g'(y_0)
  &=
  \lim_{\Delta y\to0}
  \frac{g(y_0+\Delta y)-g(y_0)}{\Delta y}       \\
  &=
  \lim_{\Delta y\to0}
  \frac{
    f(x(y_0+\Delta y),y_0+\Delta y)-f(x(y_0),y_0)
  }{\Delta y}                                  \\
  &=
  \lim_{\Delta y\to0}
  \frac{
    f_x'(x_0,y_0)[x(y_0+\Delta y)-x(y_0)]
    +f_y'(x_0,y_0)\Delta y
    +\varepsilon_1[x(y_0+\Delta y)-x(y_0)]
    +\varepsilon_2\Delta y
  }{\Delta y}.
\end{aligned}
\]
其中 $\varepsilon_1,\varepsilon_2\to0\;(\Delta y\to0)$. 因 $f_x'(x_0,y_0)=0$，且
\[
  \lim_{\Delta y\to0}
  \frac{x(y_0+\Delta y)-x(y_0)}{\Delta y}
  =x_y'(y_0),
\]
故
\[
  g'(y_0)=f_y'(x_0,y_0).
\]

\noindent\textbf{例 6.4.4}\quad 证明：对于方程
\[
  \psi(x)-y-\varphi(y)=0,
\]
若满足：

1) $\psi(x_0)-y_0-\varphi(y_0)=0$；

2) $\psi(x)$ 在 $x_0$ 的某邻域 $U(x_0,r)$ 内连续；

3) $\varphi(y)$ 在 $J=\{y:|y-y_0|\le\delta\}$ 上满足 Lipschitz 条件：存在 $\alpha,0<\alpha<1$，使得对任意 $y_1,y_2\in J$，均有
\begin{equation}
  |\varphi(y_2)-\varphi(y_1)|\le\alpha|y_2-y_1|.
\end{equation}

\noindent 则存在隐函数 $y=y(x)$ 及 $\eta>0$，使得：

i) $y_0=y(x_0)$；

ii) $|y(x)-y_0|\le\delta$，$\psi(x)-y(x)-\varphi(y(x))\equiv0$，当 $|x-x_0|<\eta$ 时；

iii) 满足 i), ii) 的函数唯一；

iv) $y=y(x)$ 连续.

\noindent\textbf{证明}\quad （迭代法）因 $\psi$ 在 $x_0$ 处连续，所以对 $(1-\alpha)\delta>0$，$\exists\eta>0(\eta<r)$，使得 $|x-x_0|<\eta$ 时，有
\begin{equation}
  |\psi(x)-\psi(x_0)|<(1-\alpha)\delta.
\end{equation}
令
\begin{equation}
  y_n=\psi(x)-\varphi(y_{n-1})\quad(n=1,2,\cdots).
\end{equation}
我们证明：对任意 $n\in\{0,1,2,\cdots\}$，有
\begin{equation}
  |y_{n+1}-y_n|\le\alpha^n(1-\alpha)\delta,\quad
  |y_{n+1}-y_0|<(1-\alpha^{n+1})\delta<\delta.
\end{equation}
首先，
\[
  |y_1-y_0|
  =
  |\psi(x)-\varphi(y_0)-[\psi(x_0)-\varphi(y_0)]|
  =
  |\psi(x)-\psi(x_0)|
  <
  (1-\alpha)\delta,
\]
且
\[
  |y_2-y_1|
  =
  |\psi(x)-\varphi(y_1)-[\psi(x)-\varphi(y_0)]|
  \le
  \alpha|y_1-y_0|
  \le
  \alpha(1-\alpha)\delta<\delta.
\]
故
\[
  |y_2-y_0|\le |y_2-y_1|+|y_1-y_0|
  <
  \alpha(1-\alpha)\delta+(1-\alpha)\delta
  =
  (1-\alpha^2)\delta<\delta.
\]
若式 (4) 对 $n=k-1$ 成立，则
\[
\begin{aligned}
  |y_{k+1}-y_k|
  &=
  |\psi(x)-\varphi(y_k)-[\psi(x)-\varphi(y_{k-1})]| \\
  &\le |\varphi(y_k)-\varphi(y_{k-1})|              \\
  &\le \alpha|y_k-y_{k-1}|                          \\
  &\le \alpha^k(1-\alpha)\delta,
\end{aligned}
\]
并且
\[
  |y_{k+1}-y_0|
  \le |y_{k+1}-y_k|+|y_k-y_0|
  \le \alpha^k(1-\alpha)\delta+(1-\alpha^k)\delta
  =
  (1-\alpha^{k+1})\delta.
\]
式 (4) 获证.

下面证 $\{y_n(x)\}$ 一致收敛. 当 $|x-x_0|<\eta$ 时，由式 (3)，
\[
  y_n=\sum_{k=1}^n(y_k-y_{k-1})+y_0,
\]
又 $|y_k-y_{k-1}|\le\alpha^{k-1}(1-\alpha)\delta$，故级数收敛，从而 $y_n(x)$ 一致收敛. 记其极限为 $y(x)$.

由递推式可知，当 $x=x_0$ 时，各 $y_n(x_0)=y_0$，故 $y(x_0)=y_0$. 又由 $|y_n(x)-y_0|<\delta$，令 $n\to\infty$，得 $|y(x)-y_0|\le\delta$. 由条件 3)，
\[
  |\varphi(y)-\varphi(y_{n-1})|\le\alpha|y-y_{n-1}|\to0.
\]
在
\[
  y_n=\psi(x)-\varphi(y_{n-1})
\]
中取极限，得
\[
  y=\psi(x)-\varphi(y),
\]
即方程成立.

若另一函数 $\bar y(x)$ 也满足 i), ii)，则
\[
  \bar y=\psi(x)-\varphi(\bar y),\quad
  y=\psi(x)-\varphi(y).
\]
于是
\[
  |\bar y-y|=|\varphi(y)-\varphi(\bar y)|
  \le\alpha|\bar y-y|.
\]
因 $0<\alpha<1$，故 $\bar y-y=0$，唯一性得证.

最后，由于各 $y_n(x)$ 均连续，且 $y_n(x)\to y(x)$ 在 $|x-x_0|<\eta$ 内一致收敛，故 $y(x)$ 连续. 证毕.

\noindent\textbf{练习}\quad 证明下面的方程在点 $(0,0,0)$ 附近唯一确定了隐函数 $z=z(x,y)$：
\[
  x+\frac12 y^2+\frac12 z+\sin z=0,
\]
并将 $z(x,y)$ 在点 $(0,0)$ 展开为带 Peano 型余项的 Taylor 公式，展开到二阶.（北京大学）

\noindent\textbf{提示}\quad 记
\[
  F(x,y,z)=x+\frac12 y^2+\frac12 z+\sin z,
\]
则 $F$ 连续，$F(0,0,0)=0$，$F'_z(0,0,0)=\frac32$. 因此在 $(0,0,0)$ 的邻域里存在唯一隐函数 $z=z(x,y)$，使得
\[
  F(x,y,z(x,y))\equiv
  x+\frac12 y^2+\frac12 z(x,y)+\sin z(x,y)=0.
\]
显然是恒等式，故可两边同时求导. 固定 $y$ 对 $x$ 求导得 $z'_x$；同样，固定 $x$ 对 $y$ 求导得 $z'_y$；$z'_x(0,0)=-\frac23$ 和 $z'_y(0,0)=0$. 注意 $z'_x$ 和 $z'_y$ 仍是 $(x,y)$ 的函数，继续求导，可得 $z''_{xx}(0,0)=0$，$z''_{xy}(0,0)=0$，$z''_{yy}(0,0)=-\frac23$. 代入二元 Taylor 公式，可得隐函数 $z=z(x,y)$ 在 $(0,0)$ 的 Taylor 展开式：
\[
\begin{aligned}
  z={}&z(0,0)+z'_x(0,0)x+z'_y(0,0)y  \\
  &+\frac1{2!}\left[z''_{xx}(0,0)x^2+2z''_{xy}(0,0)xy+z''_{yy}(0,0)y^2\right]+o(\rho) \\
  ={}&-\frac23x-\frac13y^2+o(x^2+y^2).
\end{aligned}
\]

\noindent\textbf{b.\quad 多个方程的情况}

\noindent\textbf{要点}\quad 方程组
\[
\begin{cases}
  F_1(x_1,x_2,\cdots,x_n,y_1,y_2,\cdots,y_m)=0,\\
  F_2(x_1,x_2,\cdots,x_n,y_1,y_2,\cdots,y_m)=0,\\
  \cdots\cdots\cdots\\
  F_m(x_1,x_2,\cdots,x_n,y_1,y_2,\cdots,y_m)=0,
\end{cases}\qquad\text{(A)}
\]
简记为
\[
  \mathbf F(\mathbf x,\mathbf y)=0,\qquad\text{($A'$)}
\]
其中 $\mathbf F=(F_1,F_2,\cdots,F_m)$ 为向量函数：
$\mathbb R^{n+m}\to\mathbb R^m$，$(\mathbf x,\mathbf y)\mapsto
\mathbf F(\mathbf x,\mathbf y)$. 这里
$\mathbf x=(x_1,x_2,\cdots,x_n)$，$\mathbf y=(y_1,y_2,\cdots,y_m)$，
$(\mathbf x,\mathbf y)=(x_1,x_2,\cdots,x_n,y_1,y_2,\cdots,y_m)$.

\noindent\textbf{定理 1}\quad 假若

1) $\mathbf F(P_0)=0$，
$\det\left(\dfrac{\partial F_i}{\partial y_j}\right)_{P_0}\ne0$；

2) $F_i,\dfrac{\partial F_i}{\partial y_j}$ 在 $P_0$ 的邻域里连续，

\noindent 这里
\[
  P_0=(\mathbf x_0,\mathbf y_0)
  =(x_{01},x_{02},\cdots,x_{0n},y_{01},y_{02},\cdots,y_{0m}).
\]
则方程 $(A')$ 在 $P_0$ 的邻域里确定了 $\mathbf x$ 的唯一隐函数
$\mathbf y=\mathbf y(\mathbf x)$，并且连续. 具体来说，即存在
$\delta,\eta>0$ 及函数
\[
  \mathbf y=\mathbf y(\mathbf x)
  =(y_1(\mathbf x),y_2(\mathbf x),\cdots,y_m(\mathbf x)),
  \quad \mathbf x=(x_1,\cdots,x_n),
\]
使得

i) $\mathbf y(\mathbf x_0)=\mathbf y_0$（这里
$\mathbf x_0=(x_{01},x_{02},\cdots,x_{0n})$，
$\mathbf y_0=(y_{01},y_{02},\cdots,y_{0m})$）；

ii) $|y_j(\mathbf x)-y_{0j}|\le\delta$，
$\mathbf F(\mathbf x,\mathbf y(\mathbf x))=0$（当 $\mathbf x\in I$ 时），这里
\[
  I=\{\mathbf x=(x_1,x_2,\cdots,x_n): |x_i-x_{0i}|<\eta,\ i=1,2,\cdots,n\};
\]

iii) 满足条件 i)，ii) 的这种函数是唯一的；

iv) $\mathbf y=\mathbf y(\mathbf x)$ 在 $I$ 上连续（等价地每个 $y_j(\mathbf x)$ 在 $I$ 上连续，$j=1,2,\cdots,m$）.

\noindent 隐函数微分法：若 $F_i$ 有连续偏导数，则隐函数也有连续（偏）导数；并且它的导数可按如下的方法求出：将隐式方程中的 $y_j$ 看成由方程所确定的隐函数，从而隐式方程成为恒等式，在等式两端同时求导，便可求得隐函数导数的线性方程组，解之即可求得隐函数的导数.

\noindent\textbf{例 6.4.5}\quad 证明：
\[
\begin{cases}
  e^{xu}\cos yv=\dfrac{x}{\sqrt2},\\[4pt]
  e^{xu}\sin yv=\dfrac{y}{\sqrt2}
\end{cases}
\]
在点 $P_0(x_0,y_0,u_0,v_0)=\left(1,1,0,\dfrac\pi4\right)$ 的邻域里确定了唯一的隐函数
$u=u(x,y)$，$v=v(x,y)$，并求 $du,dv,d^2u,d^2v$ 在点 $P_0$ 处的值.

\noindent\textbf{解}\quad （隐函数存在定理的条件容易验证，实际上由方程也容易解出 $u,v$. 作为例题，这里只对微分用隐函数微分法进行计算，为了训练计算能力，请读者先自己计算，再与这里比较.）将 $u,v$ 看成是 $x,y$ 的函数，对原式求微分得
\[
  e^{xu}\cos yv\,d(xu)-e^{xu}\sin yv\,d(yv)-\frac1{\sqrt2}dx=0, \qquad\text{(1)}
\]
\[
  e^{xu}\sin yv\,d(xu)+e^{xu}\cos yv\,d(yv)-\frac1{\sqrt2}dy=0. \qquad\text{(2)}
\]
用点 $P_0\left(1,1,0,\dfrac\pi4\right)$ 代入，可得
\[
  du=\frac12(dx+dy),\qquad
  dv=\left(\frac12-\frac\pi4\right)dy-\frac12dx.
\]
对 (1)，(2) 再微分，再用 $P_0$ 代入可得
\[
  d^2u=-dx^2-2dxdy,\qquad
  d^2v=\frac12dx^2+dxdy+\left(\frac\pi2-\frac32\right)dy^2.
\]

\noindent\textbf{例 6.4.6}\quad 设空间曲线 $C$ 的方程是：
\[
  x=f(t),\quad y=\varphi(t),\quad
  z=\frac{f'(t)}{\varphi'(t)}\quad(-1<t<1),
\]
其中函数 $f(t),\varphi(t)$ 在 $(-1,1)$ 内都具有二阶连续导数，且一阶导数处处不等于零；而点集
\[
  E=\left\{(x,y,z)\ \middle|\ 
  x=s^2+\frac{f'(t)}{\varphi'(t)}s+f(t),\
  y=2s+\varphi(t),\
  z=\frac{f'(t)}{\varphi'(t)},\
  s,t\in(-1,1)\right\}.
\]
试证明：$E$ 中与曲线 $C$ 充分接近（即 $|s|$ 充分小）的一些点，组成一张连续的曲面 $z=z(x,y)$.（武汉大学）

\noindent\textbf{分析}\quad 我们看到在 $E$ 的定义中，
\[
  x=s^2+\frac{f'(t)}{\varphi'(t)}s+f(t),\quad s,t\in(-1,1), \qquad\text{(1)}
\]
\[
  y=2s+\varphi(t),\quad s,t\in(-1,1), \qquad\text{(2)}
\]
\[
  z=\frac{f'(t)}{\varphi'(t)},\quad t\in(-1,1). \qquad\text{(3)}
\]
令 $s=0$，恰变成曲线
\[
  C:\ x=f(t),\quad y=\varphi(t),\quad z=\frac{f'(t)}{\varphi'(t)},\quad t\in(-1,1),
\]
故 $C\subset E$. 为了证明 $E$ 中与 $C$ 充分邻近的点组成一张连续曲面 $z=z(x,y)$，我们只要证明，当 $|s|$ 充分小时，方程 (1)，(2) 即
\[
  F\equiv x-s^2-\frac{f'(t)}{\varphi'(t)}s-f(t)=0, \qquad\text{(4)}
\]
\[
  G\equiv y-2s-\varphi(t)=0 \qquad\text{(5)}
\]
确定了唯一的单值连续函数 $t=t(x,y)$，代入 (3) 即得 $z$ 作为 $x,y$ 的函数. 可见问题归结为对 (4)，(5) 验证隐函数存在定理的条件.

事实上，任取 $t_0\in(-1,1)$，对应曲线 $C$ 上一点 $(x_0,y_0,z_0)$：
\[
  x_0=f(t_0),\quad y_0=\varphi(t_0),\quad
  z_0=\frac{f'(t_0)}{\varphi'(t_0)}.
\]
当 $(x,y,s,t)$ 取点 $P_0(x_0,y_0,0,t_0)$ 时，方程 (4)，(5) 满足，且
\[
\left.\begin{vmatrix}
  F'_s & F'_t\\
  G'_s & G'_t
\end{vmatrix}\right|_{P_0}
=
\left.\begin{vmatrix}
  -2s-\dfrac{f'(t)}{\varphi'(t)} &
  -\dfrac{f''(t)\varphi'(t)-\varphi''(t)f'(t)}{(\varphi'(t))^2}s-f'(t)\\[6pt]
  -2 & -\varphi'(t)
\end{vmatrix}\right|_{P_0}
\]
\[
  =
\begin{vmatrix}
  -\dfrac{f'(t_0)}{\varphi'(t_0)} & -f'(t_0)\\[6pt]
  -2 & -\varphi'(t_0)
\end{vmatrix}
=f'(t_0)-2f'(t_0)=-f'(t_0)\ne0.
\]
于是由 $F,G$ 连续并有连续偏导数，便知 (4)，(5) 在 $P_0$ 的邻域里确定了唯一且连续的隐函数 $t=t(x,y)$. 再由 $t_0$ 的任意性，这样就证明了 $E$ 在 $C$ 的邻近点组成了一个连续曲面 $z=z(x,y)$.

\noindent\textbf{例 6.4.7}\quad 设 $f(x,y)$ 存在二阶连续偏导数，且
$f''_{xx}f''_{yy}-f''_{xy}{}^2\ne0$，证明变换
\[
  u=f'_x(x,y), \qquad\text{(1)}
\]
\[
  v=f'_y(x,y), \qquad\text{(2)}
\]
\[
  w=-z+xf'_x(x,y)+yf'_y(x,y) \qquad\text{(3)}
\]
存在唯一的逆变换：
\[
  x=g'_u(u,v), \qquad\text{(4)}
\]
\[
  y=g'_v(u,v), \qquad\text{(5)}
\]
\[
  z=-w+ug'_u(u,v)+vg'_v(u,v). \qquad\text{(6)}
\]
（华中师范大学）

\noindent\textbf{分析}\quad 关键在于证明有式 (4)，(5)，因为将 (4)，(5) 代入 (3)，即可得式 (6). 欲证有 (4)，(5)，等价于求证存在函数 $g$ 使得
\[
  dg=xdu+ydv. \qquad\text{(7)}
\]
根据式 (1)，(2)，有
\[
\begin{aligned}
  xdu+ydv
  &=
  x f''_{xx}dx+xf''_{xy}dy+yf''_{yx}dx+yf''_{yy}dy \\
  &=(xf''_{xx}+yf''_{yx})dx+(xf''_{xy}+yf''_{yy})dy \\
  &=(xf'_x+yf'_y-f)'_x dx+(xf'_x+yf'_y-f)'_y dy.
\end{aligned}
\]
可见若令 $g=xf'_x+yf'_y-f$，则得式 (7)，从而式 (4)，(5) 成立. 另一方面，对于方程
\[
  F\equiv u-f'_x(x,y)=0,\qquad
  G\equiv v-f'_y(x,y)=0,
\]
作为 $u,v,x,y$ 的函数，因为 $F,G$ 连续，有连续偏导数，
\[
  \begin{vmatrix}
    F'_x & F'_y\\
    G'_x & G'_y
  \end{vmatrix}
  =
  f''_{xx}f''_{yy}-f''_{xy}{}^2\ne0.
\]
故逆变换 (4)，(5) 存在且唯一，从而 (6) 存在且唯一.

\noindent\textbf{例 6.4.8}\quad 设

1) $F(x,y,z)$ 在 $\mathbb R^3$ 上连续，$F$ 是三次齐次函数，即 $\forall t\in\mathbb R$，
\[
  F(tx,ty,tz)=t^3F(x,y,z)\quad(\forall (x,y,z)\in\mathbb R^3); \qquad\text{(1)}
\]

2) $F$ 有连续偏导数，且
\[
  F'_z(x,y,z)\ne0; \qquad\text{(2)}
\]

3) $z=f(x,y)$ 是方程 $F(x,y,z)=0$ 所确定的隐函数，且 $f$ 可微.

\noindent 试证：$z=f(x,y)$ 是一次齐次函数.（华中师范大学）

\noindent\textbf{证}\quad $1^\circ$ 因 $z=f(x,y)$ 是方程 $F(x,y,z)=0$ 确定的隐函数，故 $(x,y,z)=(x,y,f(x,y))$ 必满足方程，即有
\[
  F(x,y,f(x,y))=0. \qquad\text{(3)}
\]

$2^\circ$ 因 $F'_z$ 的介值性，如果 $F'_z$ 变号，就会出现零点，与 $F'_z\ne0$ 矛盾，故 $F'_z(x,y,z)\ne0$ 从而 $F'_z$ 不会变号，可知 $F(x,y,z)$ 对 $z$ 单调. 因此，$\mathbb R^2$ 上有隐函数 $z=f(x,y)$ 且必唯一.

$3^\circ$（证明 $z=f(x,y)$ 是一次齐次函数.）即 $\forall t\in\mathbb R$，
\[
  f(tx,ty)=tf(x,y)\quad(\forall(x,y)\in\mathbb R^2).
\]
根据式 (1) 和 (3)：
\[
  F(tx,ty,tf(x,y))\overset{\text{式(1)}}{=}t^3F(x,y,f(x,y))
  \overset{\text{式(3)}}{=}0. \qquad\text{(4)}
\]
在式 (3) 里，$\forall t\in\mathbb R$，用 $(tx,ty)$ 替代 $(x,y)$，则得
\[
  F(tx,ty,f(tx,ty))=0. \qquad\text{(5)}
\]
(4)--(5) 得
\[
  F(tx,ty,tf(x,y))-F(tx,ty,f(tx,ty))=0.
\]
此式左端应用 Lagrange 定理：
\[
  F'_z(tx,ty,\xi)[tf(x,y)-f(tx,ty)]=0, \qquad\text{(6)}
\]
其中 $\xi$ 在 $z_1=tf(x,y)$，$z_2=f(tx,ty)$ 之间，且（据式 (2)）$F'_z(tx,ty,\xi)\ne0$，故
\[
  tf(x,y)=f(tx,ty)\quad(\forall t\in\mathbb R).
\]
即 $z=f(x,y)$ 是一次齐次函数.（证毕.）

\noindent\textbf{例 6.4.9}\quad $T:D\to\mathbb R^2$ 是给定的映射，其中 $D\subset\mathbb R^2$ 是一个凸区域. 当 $x\in D$ 时，$T(x)$ 在 $D$ 上有连续的二阶偏导数，且 $T(x)$ 的 Jacobi 矩阵 $JT(x)$ 是正定的. 试证 $T:D\to\mathbb R^2$ 为单射.（北京大学）

\noindent\textbf{提示}\quad 记 $T$ 为：
\[
  (x,y)\mapsto(u,v)=(f(x,y),g(x,y)),
\]
其中 $(x,y)\in D\subset\mathbb R^2$，$(u,v)\in\mathbb R^2$. $T$ 为单射，意即：“若 $P_0\ne P_1$，则 $T(P_0)\ne T(P_1)$.”或等价地：“若 $T(P_0)=T(P_1)$，则 $P_0=P_1$.”

\noindent\textbf{证}\quad （反证法）假设 $D$ 中有两个不同的点 $P_0(x_0,y_0)$，$P_1(x_1,y_1)$，但 $T(P_0)=T(P_1)$. 即
\[
  u_0=f(x_0,y_0)=f(x_1,y_1)=u_1,\quad
  v_0=g(x_0,y_0)=g(x_1,y_1)=v_1. \qquad\text{(1)}
\]
$1^\circ$ 因 $D$ 为凸域，所连线段 $P_0P_1\subset D$. $P_0P_1$ 的方程可写为
\[
\begin{cases}
  x=x_0+\lambda(x_1-x_0),\\
  y=y_0+\lambda(y_1-y_0),
\end{cases}\quad \lambda\in[0,1].
\]
$\lambda=0$ 对应 $P_0(x_0,y_0)\mapsto T(P_0)=(u_0,v_0)$.

$\lambda=1$ 对应 $P_1(x_1,y_1)\mapsto T(P_1)=(u_1,v_1)\overset{\text{式(1)}}{=}(u_0,v_0)=T(P_0)$.

$\lambda\in(0,1)$ 对应
\[
  P_\lambda(x_0+\lambda(x_1-x_0),y_0+\lambda(y_1-y_0))\mapsto T(P_\lambda)\overset{\text{记}}{=}(u_\lambda,v_\lambda).
\]
作辅助函数：
\begin{equation}
\begin{aligned}
  F(\lambda)
  &=(x_1-x_0)f(P_\lambda)+(y_1-y_0)g(P_\lambda)\\
  &=(x_1-x_0)f[x_0+\lambda(x_1-x_0),y_0+\lambda(y_1-y_0)]\\
  &\quad +(y_1-y_0)g[x_0+\lambda(x_1-x_0),y_0+\lambda(y_1-y_0)].
\end{aligned}
\end{equation}
因为 $D$ 为凸域，故线段 $P_0P_1\subset D$，$P_\lambda\in P_0P_1\subset D$（$\lambda\in[0,1]$）. 式 (2) 表明：$F(\lambda)$ 是 $[0,1]$ 上一元函数，且
\[
  F(0)=(x_1-x_0)f(x_0,y_0)+(y_1-y_0)g(x_0,y_0)
  =(x_1-x_0)u_0+(y_1-y_0)v_0,
\]
\[
  F(1)=(x_1-x_0)f(x_1,y_1)+(y_1-y_0)g(x_1,y_1)
  =(x_1-x_0)u_1+(y_1-y_0)v_1.
\]
又从式 (1) 看出：$F(0)=F(1)$，可见 $F(\lambda)$ 在 $[0,1]$ 上满足 Lagrange 定理条件. 故 $\exists\xi\in(0,1)$，使得 $F'(\xi)=0$. 利用式 (2)，求导得
\begin{equation}
\begin{aligned}
  &(x_1-x_0)\bigl[f'_x(P_\xi)(x_1-x_0)+f'_y(P_\xi)(y_1-y_0)\bigr]\\
  &\quad +(y_1-y_0)\bigl[g'_x(P_\xi)(x_1-x_0)+g'_y(P_\xi)(y_1-y_0)\bigr]=0.
\end{aligned}
\end{equation}
其中 $P_\xi=(x_0+\xi(x_1-x_0),y_0+\xi(y_1-y_0))\in P_0P_1\subset D$. 借助矩阵乘法，式 (3) 可写为如下二次型：
\[
  \bigl((x_1-x_0),(y_1-y_0)\bigr)
  \begin{pmatrix}
    f'_x(P_\xi) & f'_y(P_\xi)\\
    g'_x(P_\xi) & g'_y(P_\xi)
  \end{pmatrix}
  \binom{x_1-x_0}{y_1-y_0}=0.
\]
其中 $P_0,P_1\in D$ 是两不同之点，故 $(x_1-x_0)$ 和 $(y_1-y_0)$ 不同时为 $0$. 因此，Jacobi 行列式
\[
  \begin{vmatrix}
    f'_x(P_\xi) & f'_y(P_\xi)\\
    g'_x(P_\xi) & g'_y(P_\xi)
  \end{vmatrix}
  =0,
\]
跟正定条件矛盾.（证毕.）

\noindent\textbf{例 6.4.10}\quad 假设 $D\subset\mathbb R^2$ 是凸区域，函数 $u=u(x,y),v=v(x,y)$ 在区域 $D$ 上连续，有连续一阶偏导数，$W(P)=(u(x,y),v(x,y))$ 是 $\mathbb R^2\to\mathbb R^2$ 的映射. 求证：对 $D$ 上任意两个不同的点 $P_1(x_1,y_1)$ 和 $P_2(x_2,y_2)$，若 $W(P_1)=W(P_2)$，其中 $W(P_i)=(u(x_i,y_i),v(x_i,y_i))$（$i=1,2$），则在 $P_1,P_2$ 之间必存在点 $P_\xi$，使得 Jacobi 行列式：
\[
  \left.\frac{\partial(u,v)}{\partial(x,y)}\right|_{P_\xi}=0.
\]
（北京大学）

\noindent\textbf{证}\quad 设 $P_k(x_k,y_k)$（$k=1,2$），因 $D$ 是凸区域，$\forall P_1,P_2\in D$，
\[
  P(t)=P_1+t(P_2-P_1)=(x_1+t(x_2-x_1),y_1+t(y_2-y_1))\in D
  \quad(t\in[0,1]).
\]
则
\[
  P(0)=P_1=(x_1,y_1),\quad P(1)=P_2=(x_2,y_2),
\]
\[
  W(t)=(U(t),V(t))=(u(P(t)),v(P(t))).
\]
作辅助函数
\[
  F(t)=t(W(1)-W(0))-W(t),
\]
则 $F(t)$ 在 $[0,1]$ 连续，有连续导数，且 $F(0)=F(1)$. 可应用微分中值定理：$\exists\xi\in(0,1)$，使得
\[
  F'(\xi)=(W(1)-W(0))-W'(\xi)=0.
\]
当 $W(P_1)=W(P_2)$（亦即 $W(1)=W(0)$）时有 $W'(\xi)=0$，亦即
\[
  \left.u'_t\right|_\xi
  =
  \left.\frac d{dt}u(x_1+t(x_2-x_1),y_1+t(y_2-y_1))\right|_\xi
\]
\[
  =u'_x(P(\xi))(x_2-x_1)+u'_y(P(\xi))(y_2-y_1)=0.
\]
类似地，
\[
  \left.v'_t\right|_\xi
  =v'_x(P(\xi))(x_2-x_1)+v'_y(P(\xi))(y_2-y_1)=0.
\]
因为 $P_1\ne P_2$，说明此一次齐次联立方程组
\[
\begin{cases}
  \left.u'_t\right|_\xi=0,\\
  \left.v'_t\right|_\xi=0
\end{cases}
\]
有非零解. 故系数行列式（即 Jacobi 行列式）
\[
  \begin{vmatrix}
    u'_x(P(\xi)) & u'_y(P(\xi))\\
    v'_x(P(\xi)) & v'_y(P(\xi))
  \end{vmatrix}
  =
  \left.\frac{\partial(u,v)}{\partial(x,y)}\right|_{P(\xi)}
  =0.
\]

\subsection*{※二、函数相关}

这里只列出定义和两个重要定理，如欲了解更多具体内容，请查看本书第 2 版.

\noindent\textbf{定义}\quad 设函数
\[
\begin{cases}
  y_1=y_1(x_1,\cdots,x_n),\\
  \cdots\cdots\cdots\\
  y_m=y_m(x_1,\cdots,x_n)
\end{cases}
\quad(\text{简记作 } \mathbf y=\mathbf y(\mathbf x))
\]
在 $\mathbf x_0=(x_{01},\cdots,x_{0n})$ 的某邻域内有定义，又设
$\mathbf y(\mathbf x_0)=\mathbf y_0=(y_{01},\cdots,y_{0m})$，则当且仅当存在函数
$F(y_1,\cdots,y_m)$ 在 $\mathbf y_0$ 的任意邻域里不恒为零，使得
\[
  F(y_1(x_1,\cdots,x_n),\cdots,y_m(x_1,\cdots,x_n))=0
\]
在 $\mathbf x_0$ 的邻域里成立时，$y_1,\cdots,y_m$ 称为在 $\mathbf x_0$ 处函数相关；否则称为函数无关. 当且仅当在区域 $D\subset\mathbb R^n$ 内处处函数相关时，称为在 $D$ 内函数相关；当且仅当在 $D$ 内处处函数无关时，才称在 $D$ 内函数无关.

\noindent\textbf{定理 2}\quad 若 Jacobi 矩阵
\[
  \frac{\partial(y_1,\cdots,y_m)}{\partial(x_1,\cdots,x_n)}
  =
  \left(\frac{\partial y_i}{\partial x_j}\right)_
  {\substack{i=1,\cdots,m\\ j=1,\cdots,n}}
\]
在点 $\mathbf x_0=(x_{01},\cdots,x_{0n})$ 处的秩 $r=m$，则
$y_1,y_2,\cdots,y_m$ 在 $\mathbf x_0$ 处函数无关.

\noindent\textbf{定理 3}\quad 设（齐次线性函数）
\[
\begin{aligned}
  y_1&=a_{11}x_1+\cdots+a_{1n}x_n,\\
  &\cdots\cdots\cdots\\
  y_m&=a_{m1}x_1+\cdots+a_{mn}x_n,
\end{aligned}
\]
则以下三条等价：

1) $y_1,\cdots,y_m$ 在某点 $\mathbf x_0$ 处函数相关；

2) $y_1,\cdots,y_m$ 线性相关；

3) 在 $\mathbb R^n$ 中处处函数相关.

\noindent\textbf{定理 4}\quad 设 $y_i=y_i(x_1,\cdots,x_n)$（$i=1,2,\cdots,m$）在点
$\mathbf x_0=(x_{01},x_{02},\cdots,x_{0n})$ 的某邻域里有连续的一阶偏导数，Jacobi 矩阵
\[
  \left(\frac{\partial y_i}{\partial x_j}\right)_
  {\substack{i=1,\cdots,m\\ j=1,\cdots,n}}
\]
在 $\mathbf x_0$ 附近的秩 $r$，$0<r<m$，
\[
\left.
\begin{vmatrix}
  \dfrac{\partial y_1}{\partial x_1} & \cdots & \dfrac{\partial y_1}{\partial x_r}\\
  \vdots & & \vdots\\
  \dfrac{\partial y_r}{\partial x_1} & \cdots & \dfrac{\partial y_r}{\partial x_r}
\end{vmatrix}
\right|_{\mathbf x_0}
\ne0,
\]
则 1) $y_1,\cdots,y_r$ 在 $\mathbf x_0$ 处函数无关；2) $y_1,\cdots,y_m$ 在 $\mathbf x_0$ 处函数相关.

\subsection*{单元练习 6.4}

\noindent\textbf{6.4.1}\quad $P_0(x_0,y_0)$ 是右半平面 $(x>0)$ 内任意一点，试证方程组
\[
\begin{cases}
  u=\phi(x,y)=(e^x+1)\sin y,\\
  v=\psi(x,y)=(e^x-1)\cos y
\end{cases}
\]
能在 $P_0$ 的（充分小的）邻域内确定连续可微的反函数.（北京师范大学）

\noindent\textbf{6.4.2}\quad 设
\[
  F(u,v,w,x,y)=uy+vx+w+x^2,\quad
  G(u,v,w,x,y)=uvw+x+y+1,
\]
$P_0$ 的坐标为 $(2,1,0,-1,0)$，又 $F(P_0)=0$，$G(P_0)=0$.

1) 证明：在 $(2,1,0)$ 的某一邻域内能由方程组 $F=0,G=0$ 定义唯一的一对函数
\[
  x=f(u,v,w),\quad y=g(u,v,w);
\]

2) 求 Jacobi 矩阵
\[
  \left.\frac{\partial(f,g)}{\partial(u,v,w)}\right|_{P_0}.
\]
（上海师范大学）

\noindent\textbf{6.4.3}\quad 设函数 $f(x,y),g(x,y)$ 是定义在平面开区域 $G$ 内的两个函数，在 $G$ 内均有连续的一阶偏导数，且在 $G$ 内任意点处，均有
\[
  \frac{\partial f}{\partial x}\frac{\partial g}{\partial y}
  -
  \frac{\partial f}{\partial y}\frac{\partial g}{\partial x}
  \ne0.
\]
又设有界闭区域 $D\subset G$. 试证：在 $D$ 中满足方程组
\[
\begin{cases}
  f(x,y)=0,\\
  g(x,y)=0
\end{cases}
\]
的点至多有有限个.（武汉大学）

\noindent\textbf{提示}\quad 可用反证法、聚点原理和隐函数存在定理.

\noindent\textbf{6.4.4}\quad 已知方程 $F(x,y,z)=0$ 和 $G(x,y,z)=0$.

1) 在什么条件下，由此两个方程能确定一条通过点 $P(x_0,y_0,z_0)$ 的曲线？

2) 在什么条件下，上述曲线在点 $P$ 处有切线？

3) 在什么条件下，上述切线平行于 $z$ 轴？

4) 导出上述曲线从点 $(x_0,y_0,z_0)$ 到点 $(x_1,y_1,z_1)$ 之间的一个弧长公式（用函数 $F,G$ 及其偏导数来表示）；

5) 上述弧长公式成立的条件是什么？（华东师范大学）

\noindent\textbf{6.4.5}\quad 设函数 $F(x,y)$ 在点 $(x_0,y_0)$ 的某邻域内有连续的二阶偏导数，且
$F(x_0,y_0)=0$，$F'_y(x_0,y_0)=0$，$F'_x(x_0,y_0)>0$，
$F''_{yy}(x_0,y_0)<0$. 试证：由方程 $F(x,y)=0$ 确定的定义于点 $x_0$ 邻近的隐函数
$y=y(x)$ 在点 $x_0$ 达到（局部）极小.（武汉大学）

\noindent\textbf{提示}\quad $y'(x_0)=0,y''(x_0)>0$.

\noindent\textbf{6.4.6}\quad 设 $\varphi(x,y),\varphi'_y(x,y)$ 在 $(x_0,y_0)$ 的邻域
$D: |x-x_0|, |y-y_0|\le\Delta$ 上连续，$|\varphi'_y(x,y)|<\lambda<1$，
$|\varphi(x,y)|<(1-\lambda)\Delta$. 令
$y_n=y_0+\varphi(x,y_{n-1})$，则由 $y_0$ 可依次得到
$y_1(x),y_2(x),\cdots,y_n(x),\cdots$. 试证此序列收敛，且极限函数为隐式方程
$y=y_0+\varphi(x,y)$ 的唯一连续解.（大连理工大学）

\noindent\textbf{6.4.7}\quad 设 $f(x)$ 是完备距离空间 $(X,d)$ 上将 $X$ 映为自身的连续映射，若存在正实数列 $a_n\to0$，使
$\sum_{n=1}^\infty a_n$ 收敛，且
\[
  d(f^n(x),f^n(y))\le a_n d(x,y)\quad(n\ge1,\ x,y\in X),
\]
其中 $f^{n+1}(x)=f(f^n(x))$，$f^1(x)=f(x)$，$d(x,y)$ 表示空间中 $x,y$ 两点的距离. 证明：$f(x)$ 在 $X$ 内有唯一的一点 $\xi$，使得 $f(\xi)=\xi$.（吉林工业大学）

\noindent\textbf{提示}\quad 取 $x_1=f(x)$，$x_{n+1}=f(x_n)$（$n=1,2,\cdots$），证明 $\{x_n\}$ 是 Cauchy 序列.

\noindent\textbf{6.4.8}\quad 设函数 $f(x)$ 当 $a<x<b$ 时连续，并且函数 $\varphi(y)$ 当 $c<y<d$ 时单调增加而且连续. 问在怎样的条件下，方程 $\varphi(y)=f(x)$ 定义出单值的函数
$y=\varphi^{-1}[f(x)]$？研究例子：

1) $\sin y+\operatorname{sh} y=x$；

2) $e^{-y}=-\sin^2x$.

\noindent\textbf{6.4.9}\quad 设
\[
  x=y+\varphi(y), \qquad\text{(1)}
\]
其中 $\varphi(0)=0$，且当 $-a<y<a$ 时 $\varphi'(y)$ 连续并满足
$|\varphi'(y)|\le k<1$. 证明：存在 $\delta>0$，当 $-\delta<x<\delta$ 时存在唯一的可微函数 $y(x)$ 满足方程 (1)，且 $y(0)=0$.

\noindent\textbf{6.4.10}\quad 方程 $xy+z\ln y+e^{xz}=1$ 在点 $(0,1,1)$ 的邻域内能否确定出某一变量为另两个变量的函数？

\noindent\textbf{6.4.11}\quad 设 $y=y(x)$ 是方程 $x=ky+\varphi(y)$ 所定义的隐函数，其中常数 $k\ne0$，$\varphi(y)$ 为以 $\omega$ 为周期的周期函数，且 $|\varphi'(y)|<|k|$. 证明：
$y=\dfrac{x}{k}+\psi(x)$，其中 $\psi(x)$ 是以 $|k|\omega$ 为周期的周期函数.

\noindent\textbf{证}\quad $1^\circ$ 原始方程改写为
\[
  F(x,y)=ky+\varphi(y)-x=0. \qquad\text{(1)}
\]
i) 按题意，$y=y(x)$ 是该方程确定的隐函数，因此 $(x,y(x))$ 满足方程 (1).

ii) 又（已知）$\varphi(y)$ 有导数，当然 $\varphi(y)$ 连续；由式 (1) 知 $F(x,y)$ 是二元连续函数.

iii) 对式 (1) 求导，得
\[
  |F'_y(x,y)|=|k+\varphi'(y)|\ge \bigl||k|-|\varphi'(y)|\bigr|>0
\]
（因为 $|\varphi'(y)|<|k|$）.

故隐函数存在定理的条件满足，说明 $y=y(x)$ 不仅是隐函数，而且此方程的隐函数是唯一的.

$2^\circ$ 将隐函数 $y=y(x)$ 写为 $y=\dfrac{x}{k}+\psi(x)$，即令
\[
  \psi(x)=y-\frac{x}{k}. \qquad\text{(2)}
\]
将式 (2) 代入方程 (1)，得
\[
  k\psi(x)+\varphi\left(\frac{x}{k}+\psi(x)\right)=0\quad(\forall x\in\mathbb R). \qquad\text{(3)}
\]
用 $x+|k|\omega$ 替代 (3) 中的 $x$，得
\[
  k\psi(x+|k|\omega)
  +\varphi\left(\frac{x+|k|\omega}{k}+\psi(x+|k|\omega)\right)=0. \qquad\text{(4)}
\]
因 $\varphi(y)$ 以 $\omega$ 为周期，故
\[
  \varphi\left(\frac{x+|k|\omega}{k}+\psi(x+|k|\omega)\right)
  =
  \varphi\left(\frac{x}{k}\pm\omega+\psi(x+|k|\omega)\right)
  =
  \varphi\left(\frac{x}{k}+\psi(x+|k|\omega)\right).
\]
因而式 (4) 可写为
\[
  k\psi(x+|k|\omega)
  +\varphi\left(\frac{x}{k}+\psi(x+|k|\omega)\right)=0. \qquad\text{(5)}
\]
此式表明 $\psi(x+|k|\omega)$ 也满足方程 (3)，从而
\[
  y_1(x)\overset{\text{记}}{=}\frac{x}{k}+\psi(x+|k|\omega)
\]
也满足方程 (1).（因为将 $y_1(x)=\frac{x}{k}+\psi(x+|k|\omega)$ 代入方程 (1)，即得式 (5)，说明 $y_1(x)$ 也是方程 (1) 的解.）但上面已证 $y=y(x)$ 是方程 (1) 的解且是唯一解，故
$y_1(x)=y(x)$（$\forall x\in\mathbb R$）. 因此
\[
  \psi(x+|k|\omega)=\psi(x)\quad(\forall x\in\mathbb R).
\]
亦即表明 $|k|\omega$ 是 $\psi(x)$ 的周期. 证毕.

\noindent\textbf{6.4.12}\quad 设 $\mathbf x\in\mathbb R^n,\mathbf y\in\mathbb R^m$，
$\mathbf F(\mathbf x,\mathbf y)\in\mathbb R^m$，$\mathbf F$ 有连续偏导数，
$\mathbf F=(F_1,F_2,\cdots,F_m)$，Jacobi 行列式
\[
  \frac{\partial(F_1,F_2,\cdots,F_m)}
       {\partial(y_1,y_2,\cdots,y_m)}\ne0,
\]
$\mathbf y=\mathbf y(\mathbf x)=(y_1(\mathbf x),y_2(\mathbf x),\cdots,y_m(\mathbf x))$
是方程 $\mathbf F(\mathbf x,\mathbf y)=0$ 的隐函数.

1) 证明：
\[
  D\mathbf y(\mathbf x)=-[D_y\mathbf F(\mathbf x,\mathbf y)]^{-1}D_x\mathbf F(\mathbf x,\mathbf y),
\]
其中 $D\mathbf y(\mathbf x),D_x\mathbf F(\mathbf x,\mathbf y),D_y\mathbf F(\mathbf x,\mathbf y)$ 为矩阵：
\[
D\mathbf y(\mathbf x)=
\begin{pmatrix}
  (y_1)'_{x_1} & (y_1)'_{x_2} & \cdots & (y_1)'_{x_n}\\
  (y_2)'_{x_1} & (y_2)'_{x_2} & \cdots & (y_2)'_{x_n}\\
  \vdots & \vdots & & \vdots\\
  (y_m)'_{x_1} & (y_m)'_{x_2} & \cdots & (y_m)'_{x_n}
\end{pmatrix},
\]
\[
D_x\mathbf F(\mathbf x,\mathbf y)=
\begin{pmatrix}
  (F_1)'_{x_1} & (F_1)'_{x_2} & \cdots & (F_1)'_{x_n}\\
  (F_2)'_{x_1} & (F_2)'_{x_2} & \cdots & (F_2)'_{x_n}\\
  \vdots & \vdots & & \vdots\\
  (F_m)'_{x_1} & (F_m)'_{x_2} & \cdots & (F_m)'_{x_n}
\end{pmatrix},
\]
\[
D_y\mathbf F(\mathbf x,\mathbf y)=
\begin{pmatrix}
  (F_1)'_{y_1} & (F_1)'_{y_2} & \cdots & (F_1)'_{y_m}\\
  (F_2)'_{y_1} & (F_2)'_{y_2} & \cdots & (F_2)'_{y_m}\\
  \vdots & \vdots & & \vdots\\
  (F_m)'_{y_1} & (F_m)'_{y_2} & \cdots & (F_m)'_{y_m}
\end{pmatrix};
\]

2) 证明：当 $n=m$ 时，Jacobi 行列式
\[
  \frac{\partial(y_1,y_2,\cdots,y_n)}
       {\partial(x_1,x_2,\cdots,x_n)}
  =
  (-1)^n
  \frac{\partial(F_1,F_2,\cdots,F_n)}
       {\partial(x_1,x_2,\cdots,x_n)}
  \bigg/
  \frac{\partial(F_1,F_2,\cdots,F_n)}
       {\partial(y_1,y_2,\cdots,y_n)}.
\]

\noindent\textbf{6.4.13}\quad 设 $(x_1,x_2,\cdots,x_n)$ 与
$(r,\theta_1,\cdots,\theta_{n-1})$ 为 $\mathbb R^n$ 中的直角坐标与球坐标：
\[
\begin{aligned}
  x_1&=r\cos\theta_1,\\
  x_2&=r\sin\theta_1\cos\theta_2,\\
  x_3&=r\sin\theta_1\sin\theta_2\cos\theta_3,\\
  &\cdots\cdots\cdots\\
  x_{n-1}&=r\sin\theta_1\sin\theta_2\cdots\sin\theta_{n-2}\cos\theta_{n-1},\\
  x_n&=r\sin\theta_1\cdots\sin\theta_{n-1}.
\end{aligned}
\]

1) 证明：
\[
  F_1=r^2-(x_1^2+x_2^2+\cdots+x_n^2)=0,
\]
\[
  F_2=r^2\sin^2\theta_1-(x_2^2+\cdots+x_n^2)=0,
\]
\[
  \cdots\cdots\cdots,
\]
\[
  F_n=r^2\sin^2\theta_1\cdots\sin^2\theta_{n-1}-x_n^2=0;
\]

2) 利用上题最后结果计算 Jacobi 行列式
\[
  \frac{\partial(x_1,x_2,\cdots,x_n)}
       {\partial(r,\theta_1,\cdots,\theta_{n-1})}.
\]

\noindent\textbf{提示}\quad 1) $F_n=0$ 明显；再顺次证明 $F_{n-1}=0,F_{n-2}=0,\cdots$，直至最后 $F_1=0$.

2) 应用上题结论 2)：
\[
  \frac{\partial(x_1,x_2,\cdots,x_n)}
       {\partial(r,\theta_1,\cdots,\theta_{n-1})}
  =
  (-1)^n
  \frac{\partial(F_1,F_2,\cdots,F_n)}
       {\partial(r,\theta_1,\cdots,\theta_{n-1})}
  \bigg/
  \frac{\partial(F_1,F_2,\cdots,F_n)}
       {\partial(x_1,x_2,\cdots,x_n)},
\]
而
\[
  \frac{\partial(F_1,F_2,\cdots,F_n)}
       {\partial(r,\theta_1,\cdots,\theta_{n-1})}
  =
  2^n r^{2n-1}\sin^{2n-3}\theta_1\sin^{2n-5}\theta_2
  \cdots\sin^3\theta_{n-2}\sin\theta_{n-1}
  \cos\theta_1\cdots\cos\theta_{n-1}.
\]
（注意：这里 $F_k(k=1,2,\cdots,n)$ 对 $r,\theta_1,\cdots,\theta_{n-1}$ 求导时，$F_k$ 中的 $x_1,x_2,\cdots,x_n$ 被固定；同样地，下面 $F_i$ 对 $x_1,x_2,\cdots,x_n$ 求导时，$r,\theta_1,\cdots,\theta_{n-1}$ 也被固定.）
\[
  \frac{\partial(F_1,F_2,\cdots,F_n)}
       {\partial(x_1,x_2,\cdots,x_n)}
  =
  (-1)^n2^n x_1x_2\cdots x_n
\]
\[
  =
  (-1)^n2^n r^n\sin^{n-1}\theta_1\sin^{n-2}\theta_2
  \cdots\sin\theta_{n-1}\cos\theta_1\cos\theta_2\cdots\cos\theta_{n-1}.
\]
所以
\[
  \frac{\partial(x_1,x_2,\cdots,x_n)}
       {\partial(r,\theta_1,\cdots,\theta_{n-1})}
  =
  (-1)^n r^{n-1}\sin^{n-2}\theta_1\sin^{n-3}\theta_2\cdots\sin\theta_{n-2}.
\]

\noindent\textbf{6.4.14}\quad 设 $\varphi(x,y)$ 为 $\mathbb R^2$ 中有二阶连续偏导数的二次齐次函数：
\[
  \varphi(tx,ty)=t^2\varphi(x,y).
\]

1) 证明：$\varphi'_x=x\varphi''_{xx}+y\varphi''_{xy}$，
$\varphi'_y=x\varphi''_{xy}+y\varphi''_{yy}$；

2) 设
\[
  \begin{vmatrix}
    \varphi''_{xx} & \varphi''_{xy}\\
    \varphi''_{yx} & \varphi''_{yy}
  \end{vmatrix}
  \ne0,
\]
证明：令 $u=\varphi'_x(x,y)$，$v=\varphi'_y(x,y)$ 时，函数 $\varphi(x,y)$ 可变成 $\psi(u,v)$ 的形式；

3) 证明：$\psi'_u=x,\psi'_v=y$；

4) 试将此结果推广到 $\mathbb R^n$ 空间.

\noindent\textbf{提示}\quad 1) 恒等式 $\varphi(tx,ty)=t^2\varphi(x,y)$ 对 $t$ 求导得
\[
  x\varphi'_x(tx,ty)+y\varphi'_y(tx,ty)=2t\varphi(x,y).
\]
令 $t=1$，仍是 $(x,y)$ 的恒等式，再对 $x$ 求偏导，即得
\[
  \varphi'_x=x\varphi''_{xx}+y\varphi''_{xy}. \qquad\text{(1)}
\]
类似地，有
\[
  \varphi'_y=x\varphi''_{xy}+y\varphi''_{yy}. \qquad\text{(2)}
\]

2) $u=\varphi'_x(x,y)\overset{\text{记}}{=}f(x,y)$，
$v=\varphi'_y(x,y)\overset{\text{记}}{=}g(x,y)$，则
\[
  \left|\frac{\partial(f,g)}{\partial(x,y)}\right|\ne0.
\]
根据方程隐函数存在定理，存在反函数：
\[
  x=x(u,v),\quad y=y(u,v).
\]
因此
\[
  \varphi(x,y)=\varphi(x(u,v),y(u,v))\overset{\text{记}}{=}\psi(u,v).
\]

3) 将反函数代入原式
\[
\begin{cases}
  u=\varphi'_x(x(u,v),y(u,v)),\\
  v=\varphi'_y(x(u,v),y(u,v)),
\end{cases}
\]
是关于 $u,v$ 的（两个）恒等式. 固定 $v$，对 $u$ 求偏导，得到关于 $x'_u$ 和 $y'_u$ 的方程组：
\[
\begin{cases}
  1=\varphi''_{xx}x'_u+\varphi''_{xy}y'_u,\\
  0=\varphi''_{yx}x'_u+\varphi''_{yy}y'_u.
\end{cases}\qquad\text{(3--4)}
\]
$x\times(3)+y\times(4)$ 得
\[
  x=(x\varphi''_{xx}+y\varphi''_{yx})x'_u
  +(x\varphi''_{xy}+y\varphi''_{yy})y'_u
  \overset{\text{式(1),(2)}}{=}
  \varphi'_x x'_u+\varphi'_y y'_u.
\]
于是
\[
  \psi'_u(u,v)=\frac{\partial}{\partial u}\varphi(x(u,v),y(u,v))
  =\varphi'_x x'_u+\varphi'_y y'_u=x.
\]
类似可证：
\[
  \psi'_v(u,v)=y.
\]

4)（推广到多元）设 $\varphi(x_1,x_2,\cdots,x_n)$ 是 $\mathbb R^n$ 中有直到二阶连续偏导数的二次齐次函数：
\[
  \varphi(tx_1,tx_2,\cdots,tx_n)=t^2\varphi(x_1,x_2,\cdots,x_n),
\]
则
\[
  \varphi'_{x_i}=\sum_{j=1}^n x_j(\varphi'_{x_i})'_{x_j}
  \quad(i=1,2,\cdots,n).
\]
若
\[
  \left|
  \frac{\partial(\varphi'_{x_1},\varphi'_{x_2},\cdots,\varphi'_{x_n})}
       {\partial(x_1,x_2,\cdots,x_n)}
  \right|\ne0,
\]
则 $u_i=\varphi'_{x_i}(x_1,x_2,\cdots,x_n)$（$i=1,2,\cdots,n$）有反函数：
\[
  x_i=x_i(u_1,u_2,\cdots,u_n)\quad(i=1,2,\cdots,n).
\]
将
\[
  \varphi(x_1,x_2,\cdots,x_n)
  =
  \varphi(x_1(u_1,u_2,\cdots,u_n),\cdots,x_n(u_1,u_2,\cdots,u_n))
\]
记为 $\psi(u_1,u_2,\cdots,u_n)$，则有
\[
  \psi'_{u_i}(u_1,u_2,\cdots,u_n)=x_i\quad(i=1,2,\cdots,n).
\]

\noindent\textbf{6.4.15}\quad 设
\[
  \Delta=
  \begin{vmatrix}
    a&b&d\\
    b&c&e\\
    d&e&f
  \end{vmatrix}\ne0.
\]
证明：对于二次曲线
\[
  ax^2+2bxy+cy^2+2dx+2ey+f=0,
\]
有如下等式成立：
\[
  \frac{d^3}{dx^3}\left[(y'')^{-2/3}\right]=0.
\]

\noindent\textbf{提示}\quad 已知 $\Delta\ne0$，用隐函数微分法：将 $y$ 看成原方程定义的（关于 $x$ 的）函数（代回方程，则方程变成恒等式），对 $x$ 求导，再除以 $2$（以消掉公因子 $2$），得
\[
  ax+by+bxy'+cyy'+d+ey'=0, \qquad\text{(1)}
\]
得
\[
  y'=-\frac{ax+by+d}{bx+cy+e}. \qquad\text{(2)}
\]
在式 (1) 中再对 $x$ 求一次导数，得
\[
  a+2by'+c(y')^2+(bx+cy+e)y''=0. \qquad\text{(3)}
\]
解出 $y''$，并用 (2) 代入 $y'$：
\[
\begin{aligned}
  y''
  &=-\frac1{(bx+cy+e)^3}
  \bigl[a(bx+cy+e)^2\\
  &\quad -2b(ax+by+d)(bx+cy+e)
  +c(ax+by+d)^2\bigr]\\
  &=-\frac1{(bx+cy+e)^3}
  \bigl[(ax^2+2bxy+cy^2+2dx+2ey)(ac-b^2)
  +ae^2-2bed+cd^2\bigr].
\end{aligned}
\]
利用原始方程，得
\[
  y''=
  -\frac1{(bx+cy+e)^3}
  \bigl[(-f)(ac-b^2)+ae^2-2bed+cd^2\bigr]
  =
  \frac{\Delta}{(bx+cy+e)^3}.
\]
原始方程可写为
\[
  cy^2+2bxy+2ey=-(ax^2+2dx+f). \qquad\text{(4)}
\]
利用式 (4)，可将 $(y'')^{-2/3}$ 写成 $x$ 的二次函数（如下）：
\[
\begin{aligned}
  (y'')^{-2/3}
  &=\Delta^{-2/3}(bx+cy+e)^2\\
  &=\Delta^{-2/3}\bigl[b^2x^2+c^2y^2+e^2+2(bx\cdot cy+cy\cdot e+ebx)\bigr]\\
  &=\Delta^{-2/3}\bigl[b^2x^2+e^2+2ebx+c(cy^2+2bxy+2ey)\bigr]\\
  &\overset{\text{式(4)}}{=}
  \Delta^{-2/3}\bigl[b^2x^2+e^2+2ebx-c(ax^2+2dx+f)\bigr].
\end{aligned}
\]
该式是 $x$ 的二次函数. 可见：$(y'')^{-2/3}$ 对 $x$ 的三阶导数必为 $0$. 证毕.

\noindent\textbf{注}\quad 每次求导之后，其变量仍然满足原方程，常常可用原方程化简计算结果. 若忘记这一点，有些题目很难得到欲求的结果，本题很典型.

\noindent\textbf{6.4.16}\quad 设
\[
  e^{u/x}\cos\frac vy=\frac{x}{\sqrt2},\qquad
  e^{u/x}\sin\frac vy=\frac{y}{\sqrt2}.
\]
求 $du,dv,d^2u$ 和 $d^2v$ 在 $x=1,y=1,u=0,v=0$ 时的表达式.
\[
  \left(du=\frac{dx+dy}{2},\quad dv=dy-dx,\quad
  d^2u=(dx)^2,\quad d^2v=2(dy-dx)^2\right)
\]

\noindent\textbf{提示}\quad （小窍门：动手之前，先想想算式能不能简化；微分之后不忙整理，代入数字变简单后再整理；若分母能去掉，尽量先去分母.）这里先将算式化简：
\[
  e^{2u/x}=\frac{x^2+y^2}{2}, \qquad\text{(1)}
\]
\[
  \tan\frac vy=\frac yx. \qquad\text{(2)}
\]

\noindent\textbf{再提示}\quad $1^\circ$ 对式 (1) 微分：
\[
  e^{2u/x}\left(2\frac{xdu-udx}{x^2}\right)=xdx+ydy. \qquad\text{(3)}
\]
在式 (3) 中令 $x=y=1,u=0$，得
\[
  \left.du\right|_{x=y=1,u=0}=\frac{dx+dy}{2}. \qquad\text{(4)}
\]
（下面求二阶微分 $d^2u$.）将式 (3) 写为
\[
  2e^{2u/x}(xdu-udx)=x^3dx+x^2ydy,
\]
对此式两端同时求微分，得
\[
  2e^{2u/x}\left[\frac{2(xdu-udx)^2}{x^2}+xd^2u\right]
  =
  3x^2(dx)^2+2xydxdy+x^2(dy)^2.
\]
代入 $x=y=1,u=0$，并注意式 (4)，得
\[
  \left.d^2u\right|_{x=y=1,u=0}=(dx)^2.
\]

$2^\circ$ 由式 (2)：
\[
  \tan\frac vy=\frac yx,
\]
得
\[
  x^2(ydv-vdy)=y^2\cos^2\frac vy(xdy-ydx). \qquad\text{(5)}
\]
在式 (5) 中令 $x=y=1,v=0$，得
\[
  \left.dv\right|_{x=y=1,v=0}=dy-dx. \qquad\text{(6)}
\]
（为了求在给定点的二阶微分，考虑式 (5) 比较复杂，先分别对式 (5) 两端求微分，并用 $x=y=1,v=0$ 代入，然后将两结果列为等式，解出 $d^2v|_{x=y=1,v=0}$.）
\[
\left.d[x^2(ydv-vdy)]\right|_{x=y=1,v=0}
=
\left[2xdx(ydv-vdy)+x^2(dydv+yd^2v-dvdy)\right]_{x=y=1,v=0}
=2dxdv+d^2v. \qquad\text{(7)}
\]
\[
\begin{aligned}
&\left.d\left[y^2\cos^2\frac vy(xdy-ydx)\right]\right|_{x=y=1,v=0}\\
&=\left\{2ydy\cos^2\frac vy(xdy-ydx)
y^2\cdot 2\cos\frac vy\left(-\sin\frac vy\right)
\frac{ydv-vdy}{y^2}(xdy-ydx)\right.\\
&\quad\left.
y^2\cos^2\frac vy(dxdy-dydx)\right\}_{x=y=1,v=0}\\
&=2dy\cdot1\cdot(dy-dx)+0+0
=2[(dy)^2-dxdy]. \qquad\text{(8)}
\end{aligned}
\]
再令式 (7) 与式 (8) 相等，即
\[
  2dxdv+d^2v=2[(dy)^2-dxdy],
\]
即亦
\[
  d^2v=2(dy)^2-2dxdy-2dxdv
  \overset{\text{式(6)}}{=}
  2(dy)^2-2dxdy-2dx(dy-dx).
\]
故
\[
  \left.d^2v\right|_{x=y=1,v=0}=2(dy-dx)^2.
\]

\noindent\textbf{6.4.17}\quad 函数 $u=u(x)$ 由方程组
\[
  u=f(x,y,z),\quad g(x,y,z)=0,\quad h(x,y,z)=0
\]
定义，求 $\dfrac{du}{dx}$.
\[
  \left(f'_x+f'_y\frac{dy}{dx}+f'_z\frac{dz}{dx};\quad
  y'_x=\frac KJ,\ z'_x=\frac LJ,\ 
  J=\frac{\partial(g,h)}{\partial(y,z)},\
  K=\frac{\partial(g,h)}{\partial(z,x)},\
  L=\frac{\partial(g,h)}{\partial(x,y)}
  \right)
\]

\noindent\textbf{提示}\quad 因 $u=u(x)$ 是复合函数：
$u=f(x,y(x),z(x))$，$y=y(x)$，$z=z(x)$，其中出现的
$\dfrac{dy}{dx}$ 和 $\dfrac{dz}{dx}$ 可从后两个方程对 $x$ 求导解得.

\noindent\textbf{再提示}\quad $1^\circ$ 对 $u=f(x,y(x),z(x))$ 求导，得
\[
  \frac{du}{dx}=f'_x+f'_y y'_x+f'_z z'_x. \qquad\text{(1)}
\]
$2^\circ$（对后两式求导.）
\[
  g'_x+g'_y y'_x+g'_z z'_x=0,\qquad
  h'_x+h'_y y'_x+h'_z z'_x=0. \qquad\text{(2)}
\]
记
\[
  J=
  \begin{vmatrix}
    g'_y&g'_z\\
    h'_y&h'_z
  \end{vmatrix}
  =
  \frac{\partial(g,h)}{\partial(y,z)},\quad
  K=
  \begin{vmatrix}
    g'_z&g'_x\\
    h'_z&h'_x
  \end{vmatrix}
  =
  \frac{\partial(g,h)}{\partial(z,x)},\quad
  L=
  \begin{vmatrix}
    g'_x&g'_y\\
    h'_x&h'_y
  \end{vmatrix}
  =
  \frac{\partial(g,h)}{\partial(x,y)}.
\]
解方程组 (2)，可得
\[
  y'_x=\frac KJ,\qquad z'_x=\frac LJ. \qquad\text{(3)}
\]
以式 (3) 代入 (1)：
\[
  \frac{du}{dx}=f'_x+\frac KJ f'_y+\frac LJ f'_z.
\]

\noindent\textbf{6.4.18}\quad 设一对一变换
\[
T:
\begin{cases}
  x=x(u,v),\\
  y=y(u,v)
\end{cases}
\]
在 $D$ 上具有连续的偏导数 $x'_u,x'_v,y'_u,y'_v$，且行列式
\[
  \frac{\partial(x,y)}{\partial(u,v)}\ne0,
\]
则 $T$ 将 $uv$ 平面上由分段光滑闭曲线围成的闭区域 $D$ 变为 $xy$ 平面上相应闭区域 $D'$，且其边界也是分段光滑的闭曲线.（陕西师范大学）

\noindent\textbf{6.4.19}\quad $f_1(x),\cdots,f_n(x)$ 是定义在 $[a,b]$ 上的连续函数，证明：$f_1(x),\cdots,f_n(x)$ 线性无关的充分条件是行列式
$\det(\alpha_{ij})\ne0$，其中
\[
  \alpha_{ij}=\int_a^b f_i(x)f_j(x)\,dx.
\]

\noindent\textbf{提示}\quad 要证明：“若 $\det(\alpha_{ij})\ne0$，则 $\{f_i(x)\}_{i=1}^n$ 在 $[a,b]$ 上线性无关.”

\noindent\textbf{反证法}\quad “假设 $\{f_i(x)\}_{i=1}^n$ 在 $[a,b]$ 上线性相关，则必导致 $\det(\alpha_{ij})=0$.”

\noindent\textbf{再提示}\quad 若 $\{f_i(x)\}_{i=1}^n$ 线性相关，意味着存在不全为 $0$ 的实数 $\{b_i\}_{i=1}^n$ 使得
\[
  \sum_{i=1}^n b_i f_i(x)=0\quad(\forall x\in[a,b]). \qquad\text{(1)}
\]
于是
\[
  \sum_{i=1}^n b_i f_i(x)f_j(x)=0\quad(\forall x\in[a,b])\quad(j=1,2,\cdots,n).
\]
进而有
\[
  \sum_{i=1}^n b_i\int_a^b f_i(x)f_j(x)\,dx=0
  \quad(j=1,2,\cdots,n). \qquad\text{(2)}
\]
此即有不全为 $0$ 的 $\{b_i\}_{i=1}^n$ 满足式 (2)，说明齐次线性方程组 (2) 有非零解，故系数行列式
\[
  \det\left(\int_a^b f_i(x)f_j(x)\,dx\right)_{n\times n}=0,
\]
矛盾.

\noindent\textbf{6.4.20}\quad 证明：若一元函数组
$\varphi_1(x),\varphi_2(x),\cdots,\varphi_n(x)$ 在区间 $(a,b)$ 线性相关，则（出现的各阶导数都存在）
\[
\begin{vmatrix}
  \varphi_1(x)&\varphi_2(x)&\cdots&\varphi_n(x)\\
  \varphi'_1(x)&\varphi'_2(x)&\cdots&\varphi'_n(x)\\
  \vdots&\vdots&&\vdots\\
  \varphi_1^{(n-1)}(x)&\varphi_2^{(n-1)}(x)&\cdots&\varphi_n^{(n-1)}(x)
\end{vmatrix}
\equiv0\quad(x\in(a,b)).
\]

\noindent\textbf{提示}\quad 可用上题类似的方法.

\noindent\textbf{再提示}\quad 如果 $\{\varphi_k(x)\}_{k=1}^n$ 线性相关，则存在不全为零的常数组
$\{a_k\}_{k=1}^n$，使得
\[
  \sum_{k=1}^n a_k\varphi_k(x)=0.
\]
两端同时对 $x$ 求 $1$ 至 $(n-1)$ 阶导数，得 $n$ 个关于 $\{a_k\}_{k=1}^n$ 的一次齐次线性方程组：
\[
  \sum_{k=1}^n a_k\varphi_k(x)=0,\quad
  \sum_{k=1}^n a_k\varphi'_k(x)=0,\quad
  \sum_{k=1}^n a_k\varphi''_k(x)=0,\quad\cdots,\quad
  \sum_{k=1}^n a_k\varphi_k^{(n-1)}(x)=0.
\]
因 $\{a_k\}_{k=1}^n$（不全为 $0$）满足方程组，说明该方程组有非零解，（根据代数方程的理论）其系数行列式应为零. 即欲求的等式成立.

\noindent\textbf{6.4.21}\quad 证明：
\[
  x=r\cos\theta\cos\varphi,\quad
  y=r\cos\theta\sin\varphi,\quad
  z=r\sin\theta
\]
函数独立（即函数无关）.

\noindent\textbf{提示}\quad （利用定理 2.）因此时 Jacobi 行列式：
\[
  |J|=\left|\frac{\partial(x,y,z)}{\partial(r,\theta,\varphi)}\right|
  =|r^2\cos\theta|\ne0,
\]
故函数独立.

\noindent\textbf{6.4.22}\quad 讨论下列函数的相关性：

1)
\[
  \frac{x-y}{x-z},\quad \frac{y-z}{y-x},\quad \frac{z-x}{z-y};
\]

2)
\[
  \frac{x}{1-x-y-z},\quad
  \frac{y}{1-x-y-z},\quad
  \frac{z}{1-x-y-z}.
\]

\noindent\textbf{提示}\quad （利用定理 2.）

1) 函数无关，因为
\[
J=
\begin{vmatrix}
  \dfrac{y-z}{(x-z)^2}&\dfrac{z-x}{(x-z)^2}&\dfrac{x-y}{(x-z)^2}\\[6pt]
  \dfrac{y-z}{(y-x)^2}&\dfrac{z-x}{(y-x)^2}&\dfrac{x-y}{(y-x)^2}\\[6pt]
  \dfrac{y-z}{(z-y)^2}&\dfrac{z-x}{(z-y)^2}&\dfrac{x-y}{(z-y)^2}
\end{vmatrix}
=
\frac{(x-y)(y-z)(z-x)}
{(x-y)^2(y-z)^2(z-x)^2}
\]
\[
  =
  \frac1{(x-y)(y-z)(z-x)}\ne0.
\]

2) 函数无关. 因为
\[
  J=\frac1{(1-x-y-z)^6}
  \begin{vmatrix}
    1-y-z&x&x\\
    y&1-x-z&y\\
    z&z&1-x-y
  \end{vmatrix}
\]
\[
  =
  \frac1{(1-x-y-z)^6}
  \begin{vmatrix}
    1-x-y-z&0&x\\
    0&1-x-y-z&y\\
    -(1-x-y-z)&-(1-x-y-z)&1-x-y
  \end{vmatrix}
\]
\[
  =
  \frac1{(1-x-y-z)^6}
  \begin{vmatrix}
    1-x-y-z&0&x\\
    0&1-x-y-z&y\\
    0&0&1
  \end{vmatrix}
  =
  \frac1{(1-x-y-z)^4}\ne0.
\]

\noindent\textbf{6.4.23}\quad 设函数组
\[
\begin{cases}
  u=u(x,y),\\
  v=v(x,y),
\end{cases}
\quad (x,y)\in\mathbb R^2
\]
中的函数 $u,v$ 有处处连续的一阶偏导数，记 $W=(u,v)$，$P=(x,y)$. 当
$|W|=\sqrt{u^2+v^2}$，$|P|=\sqrt{x^2+y^2}$ 时，存在数 $C>0$，使得对于任意的
$P_1\in\mathbb R^2,P_2\in\mathbb R^2$，成立不等式
\[
  |W_2-W_1|\ge C|P_2-P_1|,
\]
这里 $W_i$ 为与 $P_i$ 相对应的点（$i=1,2$）. 试证：
Jacobi 行列式
\[
  \frac{\partial(u,v)}{\partial(x,y)}\ne0,\quad \forall(x,y)\in\mathbb R^2.
\]
（武汉大学）

\noindent\textbf{提示}\quad 如果在某个点 $P_1$ 处，Jacobi 行列式
\[
  \left.\frac{\partial(u,v)}{\partial(x,y)}\right|_{P_1}=0,
\]
那么当 $W_2=W_1$ 时，二元一次联立方程组
\[
  u'_x(P_1)(x_2-x_1)+u'_y(P_1)(y_2-y_1)=u(P_2)-u(P_1)=0,
\]
\[
  v'_x(P_1)(x_2-x_1)+v'_y(P_1)(y_2-y_1)=v(P_2)-v(P_1)=0
\]
应有非零解. 换句话说：存在 $P_2\ne P_1$ 使得 $W_2=W_1$. 这与已知条件
$|W_2-W_1|\ge C|P_2-P_1|$ 矛盾. 故不可能！

\noindent\textbf{再提示}\quad 设 $P_k=(x_k,y_k)$（$k=1,2$），
\[
  P(t)=P_1+t(P_2-P_1)=(x_1+t(x_2-x_1),y_1+t(y_2-y_1)).
\]
则
\[
  P(0)=P_1=(x_1,y_1),\quad P(1)=P_2=(x_2,y_2).
\]
令 $W(t)=(U(t),V(t))=(u(P(t)),v(P(t)))$，作辅助函数
\[
  F(t)=t(W(1)-W(0))-W(t),
\]
则 $F(0)=F(1)$，
\[
  W^2(t)=U^2(x_1+t(x_2-x_1),y_1+t(y_2-y_1))
  +V^2(x_1+t(x_2-x_1),y_1+t(y_2-y_1)).
\]
若记 $W_k=(U_k,V_k)=(U(P_k),V(P_k))$（$k=1,2$），那么 $W(0)=W_1$，
$W(1)=W_2$. 根据条件 $|W_2-W_1|\ge C|P_2-P_1|$ 知：当 $W_2=W_1$ 时，有 $P_2=P_1$. 应用 Lagrange 定理，
\[
  \exists\xi\in(0,1),\quad
  W^2(1)-W^2(0)=W^{2\,\prime}(\xi)(1-0)=2WW'(\xi)(1-0)=0,
\]
\[
  W'(\xi)=(u(P(t)),v(P(t)))'_t|_{t=\xi}=0,
\]
\[
  0=U(1)-U(0)=u'(P(\xi))(1-0)
  =u'_x(P(\xi))(x_2-x_1)+u'_y(P(\xi))(y_2-y_1), \qquad\text{(1)}
\]
\[
  0=V(1)-V(0)=v'(P(\xi))(1-0)
  =v'_x(P(\xi))(x_2-x_1)+v'_y(P(\xi))(y_2-y_1), \qquad\text{(2)}
\]
其中 $P(\xi)$ 表示
\[
  (\xi x_2+(1-\xi)x_1,\xi y_2+(1-\xi)y_1)\quad(0<\xi<1).
\]
把方程 (1)，(2) 看成是关于 $(x_2-x_1)$ 和 $(y_2-y_1)$ 的一次线性方程组. 亦即：方程 (1)，(2) 左端为 $0$ 时，右端 $(x_2-x_1)$ 和 $(y_2-y_1)$ 必为零，也就等于说，一次齐次线性方程组只有零解，故系数行列式
\[
  \begin{vmatrix}
    u'_x(P(\xi)) & u'_y(P(\xi))\\
    v'_x(P(\xi)) & v'_y(P(\xi))
  \end{vmatrix}
  =
  \left.\frac{\partial(u,v)}{\partial(x,y)}\right|_{P(\xi)}
  \ne0
\]
（题设 Jacobi 行列式 $\dfrac{\partial(u,v)}{\partial(x,y)}\ne0$）.

\noindent\textbf{6.4.24}\quad 设 $u=u(x,y,z),v=v(x,y,z),x=x(s,t),y=y(s,t),z=z(s,t)$ 都有连续一阶偏导数，证明行列式
\[
  \frac{\partial(u,v)}{\partial(s,t)}
  =
  \frac{\partial(u,v)}{\partial(x,y)}
  \frac{\partial(x,y)}{\partial(s,t)}
  +
  \frac{\partial(u,v)}{\partial(y,z)}
  \frac{\partial(y,z)}{\partial(s,t)}
  +
  \frac{\partial(u,v)}{\partial(z,x)}
  \frac{\partial(z,x)}{\partial(s,t)}.
\]

\noindent\textbf{提示}\quad 展开后，右边多三项，但此三项为零.

\noindent\textbf{6.4.25}\quad 设
\[
  F(x,y)=x^2y^3+|x|y+y-5.
\]

1) 证明方程 $F(x,y)=0$ 在 $(-\infty,+\infty)$ 上确定唯一的隐函数 $y=f(x)$；

2) 求 $y=f(x)$ 的极值.（北京大学）

\noindent\textbf{提示}\quad 对每个固定的 $x$，$x^2y^3+|x|y+y-5=0$ 是 $y$ 的三次方程，有实根，而 $F'_y(x,y)>0$，所以实根是唯一的.（注意，由原式可看出：$y\ne0$.）利用隐函数微分法，
\[
  f'(x)=-\frac{F'_x}{F'_y}
  =
  -\frac{2xy^3\pm y}{3x^2y^2+|x|+1}
  =
  \begin{cases}
    \dfrac{-y}{3x^2y^2+|x|+1}(2xy^2+1)<0,\quad x>0,\\[6pt]
    \dfrac{-y}{3x^2y^2+|x|+1}(2xy^2-1)>0,\quad x<0.
  \end{cases}
\]
所以 $f(0)=5$ 为最大值（亦极大值）.

\section*{§6.5\quad 方向导数与梯度}

方向导数对理工科专业是重要内容.

\subsection*{一、方向导数的计算}

\noindent\textbf{要点}\quad 计算方向导数的基本方法如下：

1) 利用定义. 函数 $y=f(\mathbf x)$（$\mathbf x\in\mathbb R^n$）在点
$P(x_1,x_2,\cdots,x_n)$ 处沿单位向量
$\mathbf l=(l_1,l_2,\cdots,l_n)$ 方向的方向导数定义为
\[
\left.\frac{\partial f}{\partial l}\right|_P
=\lim_{t\to0}
\frac{f(x_1+tl_1,x_2+tl_2,\cdots,x_n+tl_n)-f(x_1,x_2,\cdots,x_n)}{t}
\]
\[
=
\left.
\frac{d f(x_1+tl_1,x_2+tl_2,\cdots,x_n+tl_n)}{dt}
\right|_{t=0}. \qquad\text{(A)}
\]

2) 利用偏导数与方向导数的关系. 若 $f$ 在 $P$ 处可微，则 $f$ 在点 $P$ 沿任意方向
\[
  \mathbf l=(\cos\alpha_1,\cos\alpha_2,\cdots,\cos\alpha_n)
\]
的方向导数存在，并且
\[
\left.\frac{\partial f}{\partial l}\right|_P
=f'_{x_1}(P)\cos\alpha_1+f'_{x_2}(P)\cos\alpha_2+\cdots+f'_{x_n}(P)\cos\alpha_n. \qquad\text{(B)}
\]

3)
\[
  \operatorname{grad} f(P)\overset{\text{记}}{=}
  (f'_{x_1}(P),f'_{x_2}(P),\cdots,f'_{x_n}(P)),
\]
称为 $f(P)$ 在点 $P$ 处的梯度. 若 $f(P)$ 在点 $P$ 可微，则 $f(P)$ 在点 $P$ 沿任意方向
$\mathbf l=(\cos\alpha_1,\cos\alpha_2,\cdots,\cos\alpha_n)$ 的方向导数
\[
\left.\frac{\partial f}{\partial l}\right|_P
=\operatorname{grad} f(P)\cdot(\cos\alpha_1,\cos\alpha_2,\cdots,\cos\alpha_n)
=\operatorname{grad}_{\mathbf l}f(P)
\]
\[
  =|\operatorname{grad} f(P)|\cos\theta, \qquad\text{(C)}
\]
其中 $\theta$ 表示 $\operatorname{grad}f(P)$ 与 $\mathbf l$ 的夹角.

4) $f(P)$ 在点 $P$ 的梯度 $\operatorname{grad}f(P)$ 是一个向量. 它的方向和大小如下确定：
若 $f(P)$ 在点 $P$ 可微，则对于从点 $P$ 出发的不同方向，$f(P)$ 有大小不同的方向导数（值），方向导数最大的那个方向就定为 $\operatorname{grad}f(P)$ 的方向；$f(P)$ 在点 $P$ 处最大方向导数（之值），就定为向量 $\operatorname{grad}f(P)$ 的值.

\noindent\textbf{例 6.5.1}\quad 设
\[
f(x,y)=
\begin{cases}
  \dfrac{xy}{\sqrt{x^2+y^2}},&x^2+y^2\ne0,\\[6pt]
  0,&x^2+y^2=0,
\end{cases}
\]
试证 $f(x,y)$ 在 $(0,0)$ 沿任意方向的方向导数存在，但在 $(0,0)$ 处不可微.

\noindent\textbf{证}\quad 任取方向 $\mathbf l=(\cos\alpha,\sin\alpha)$，则
\[
f(t\cos\alpha,t\sin\alpha)
=
\begin{cases}
  t\cos\alpha\sin\alpha,&t\ne0,\\
  0,&t=0
\end{cases}
=t\cos\alpha\sin\alpha.
\]
于是
\[
\left.\frac{\partial f}{\partial l}\right|_{(0,0)}
=
\left.\frac d{dt}f(t\cos\alpha,t\sin\alpha)\right|_{t=0}
=(t\cos\alpha\sin\alpha)'|_{t=0}
=\cos\alpha\sin\alpha.
\]
可见在 $(0,0)$ 处沿任意方向导数存在.（不可微性留给读者证明.）

\noindent\textbf{例 6.5.2}\quad 证明：
\[
f(x,y)=
\begin{cases}
  \dfrac{x^2y}{x^4+y^2},&x^2+y^2\ne0,\\[6pt]
  0,&x^2+y^2=0
\end{cases}
\]
沿任意方向 $\mathbf l=(\cos\alpha,\sin\alpha)$ 的方向导数为
\[
\frac{\partial f(0,0)}{\partial l}
=
\begin{cases}
  \dfrac{\cos^2\alpha}{\sin\alpha},&\sin\alpha\ne0,\\[6pt]
  0,&\sin\alpha=0.
\end{cases}
\]

\noindent\textbf{证}\quad
\[
\begin{aligned}
\frac{\partial f(0,0)}{\partial l}
&=\lim_{t\to0}
\left[
\frac{(t^2\cos^2\alpha)(t\sin\alpha)}
     {t^4\cos^4\alpha+t^2\sin^2\alpha}
-0
\right]\bigg/t\\
&=\lim_{t\to0}\frac{\cos^2\alpha\sin\alpha}{t^2\cos^4\alpha+\sin^2\alpha}
=
\begin{cases}
  \dfrac{\cos^2\alpha}{\sin\alpha},&\sin\alpha\ne0,\\[6pt]
  0,&\sin\alpha=0.
\end{cases}
\end{aligned}
\]

\noindent\textbf{练习}\quad 若
\[
f(x,y)=
\begin{cases}
  \dfrac{x^2y}{x^4+y^2},&x^2+y^2\ne0,\\[6pt]
  0,&x^2+y^2=0,
\end{cases}
\]
在点 $(x_0,y_0)$ 沿任意方向的方向导数都存在，问 $f(x,y)$ 在点 $(x_0,y_0)$ 连续吗？回答并说明理由.（华东师范大学）

\noindent\textbf{提示}\quad $f(x,y)$ 在原点 $(0,0)$ 沿任意方向的方向导数都存在，但 $f$ 在此点并不连续. 如令 $y=kx^2$，当 $x\to0$ 时，
\[
  f(x,y)=\left.\frac{x^2y}{x^4+y^2}\right|_{y=kx^2}
  \to \frac{k}{1+k^2}
\]
（不同路径极限不同）.

\noindent\textbf{例 6.5.3}\quad 求
\[
  f(x,y,z)=x^2+y^2+z^2
\]
在椭球面
\[
  \frac{x^2}{a^2}+\frac{y^2}{b^2}+\frac{z^2}{c^2}=1
\]
上的点 $P(x_0,y_0,z_0)$ 的外法线方向的导数.

\noindent\textbf{解}\quad 法向量为
\[
  \mathbf n=\left(\frac{2x_0}{a^2},\frac{2y_0}{b^2},\frac{2z_0}{c^2}\right),
\]
单位法向量为
\[
  \mathbf n=\frac1{|\mathbf n|}
  \left(\frac{2x_0}{a^2},\frac{2y_0}{b^2},\frac{2z_0}{c^2}\right),
\]
朝外，其中
\[
  |\mathbf n|
  =2\sqrt{\frac{x_0^2}{a^4}+\frac{y_0^2}{b^4}+\frac{z_0^2}{c^4}}
  \overset{\text{记}}{=}2\mu.
\]
因此
\[
\left.\frac{\partial f}{\partial n}\right|_P
=\operatorname{grad}_{\mathbf n}f
=(2x_0,2y_0,2z_0)\cdot
\left(\frac{x_0}{a^2\mu},\frac{y_0}{b^2\mu},\frac{z_0}{c^2\mu}\right)
=
\frac{2}{\sqrt{\dfrac{x_0^2}{a^4}+\dfrac{y_0^2}{b^4}+\dfrac{z_0^2}{c^4}}}.
\]

\noindent\textbf{例 6.5.4}\quad 设 $f(x,y)$ 在点 $P(x_0,y_0)$ 处可微，
$l_1,l_2,\cdots,l_n$ 为 $P_0$ 处给定的 $n$ 个单位向量，相邻两向量夹角为
$\dfrac{2\pi}{n}$，证明：
\[
  \sum_{i=1}^n\frac{\partial f(x_0,y_0)}{\partial l_i}=0. \qquad\text{(1)}
\]

\noindent\textbf{证}\quad 在 $\mathbb R^2$ 中利用公式 (B)，
\[
\begin{aligned}
  \sum_{i=1}^n\frac{\partial f(x_0,y_0)}{\partial l_i}
  &=
  \sum_{i=1}^n\left[f'_x(x_0,y_0)\cos(l_i,x)
  +f'_y(x_0,y_0)\cos(l_i,y)\right]\\
  &=f'_x(x_0,y_0)\sum_{i=1}^n\cos(l_i,x)
  +f'_y(x_0,y_0)\sum_{i=1}^n\cos(l_i,y). \qquad\text{(2)}
\end{aligned}
\]
不妨设在点 $P(x_0,y_0)$ 处 $x$ 轴沿逆时针方向转动遇到的第一个向量为 $l_1$，记 $l_1$ 与 $x$ 轴的夹角为 $\alpha$，则
$l_1,l_2,\cdots,l_n$ 与 $x$ 轴的夹角顺次为
\[
  \alpha,\alpha+\frac{2\pi}{n},\cdots,\alpha+(n-1)\frac{2\pi}{n}.
\]
因此
\[
\begin{aligned}
  \sum_{i=1}^n\cos(l_i,x)
  &=
  \cos\alpha+\cos\left(\alpha+\frac{2\pi}{n}\right)+\cdots+
  \cos\left[\alpha+(n-1)\frac{2\pi}{n}\right]\\
  &=
  \frac1{2\sin\frac{2\pi}{n}}
  \sum_{i=0}^{n-1}
  2\cos\left(\alpha+i\frac{2\pi}{n}\right)\sin\frac{2\pi}{n}\\
  &=
  \frac1{2\sin\frac{2\pi}{n}}
  \sum_{i=0}^{n-1}
  \left\{
  \sin\left[\alpha+(i+1)\frac{2\pi}{n}\right]
  -
  \sin\left[\alpha+(i-1)\frac{2\pi}{n}\right]
  \right\}\\
  &=0. \qquad\text{(3)}
\end{aligned}
\]
同理
\[
  \sum_{i=1}^n\cos(l_i,y)=0. \qquad\text{(4)}
\]
将 (3)，(4) 代入 (2) 即得 (1).

\noindent\textbf{例 6.5.5}\quad 设 $y=\varphi(x)$ 是区间 $[a,b]$ 上的可微函数，在 $xOy$ 直角坐标平面内其图像为曲线 $\Gamma$. 若二元函数 $f(x,y)$ 在包含曲线 $\Gamma$ 的某区域上连续可微（即具有连续的偏导数），且在曲线 $\Gamma$ 上恒为 $0$，求证：$f(x,y)$ 在曲线 $\Gamma$ 上任一给定点处沿该曲线切线方向的方向导数等于 $0$.（湘潭大学）

\noindent\textbf{分析}\quad 设 $\mathbf l=(\cos\alpha,\cos\beta)$ 是曲线 $\Gamma$ 上 $P$ 处的单位切向量. 利用公式 (B)，
\[
  \frac{\partial f}{\partial l}=f'_x\cos\alpha+f'_y\cos\beta. \qquad\text{(1)}
\]
可见，要计算 $\dfrac{\partial f}{\partial l}$，关键在于求出 $\cos\alpha,\cos\beta$. 按已知条件，
$f(x,\varphi(x))=0$，因此
\[
  f'_x(P)+f'_y(P)\varphi'(x)=0\quad(P\in\Gamma).
\]
故
\[
  \tan\alpha=\varphi'(x)=-\frac{f'_x(P)}{f'_y(P)}.
\]
从而
\[
  \cos\alpha=\frac1{\pm\sqrt{1+\tan^2\alpha}}
  =
  \frac{|f'_y|}{\pm\sqrt{(f'_x)^2+(f'_y)^2}},
\]
\[
  \cos\beta=\sin\alpha=\tan\alpha\cos\alpha
  =
  \frac{-f'_x}{f'_y}\cdot
  \frac{|f'_y|}{\pm\sqrt{(f'_x)^2+(f'_y)^2}}.
\]
代入 (1) 得
\[
\left.\frac{\partial f}{\partial l}\right|_P
=f'_x\cos\alpha+f'_y\cos\beta
=
\frac{f'_x|f'_y|-f'_x|f'_y|}
{\pm\sqrt{(f'_x)^2+(f'_y)^2}}
=0.
\]

\noindent\textbf{例 6.5.6}\quad 设 $l_1,l_2,\cdots,l_n$ 为 $\mathbb R^n$ 中 $n$ 个线性无关的单位向量，函数 $f(x)$ 在 $\mathbb R^n$ 中可微，方向导数
$\dfrac{\partial f}{\partial l_i}=0$（$i=1,2,\cdots,n$），试证：$f(x)\equiv$ 常数.

\noindent\textbf{证}\quad 记 $l_i=(a_{i1},a_{i2},\cdots,a_{in})$（$i=1,2,\cdots,n$）. 因
$\dfrac{\partial f}{\partial l_i}=0$，应用公式 (B)，得
\[
  0=\frac{\partial f}{\partial l_i}
  =
  f'_{x_1}a_{i1}+f'_{x_2}a_{i2}+\cdots+f'_{x_n}a_{in}
  \quad(i=1,2,\cdots,n). \qquad\text{(1)}
\]
因为 $l_i$ 线性无关，故 $\det(a_{ij})_{i,j=1}^n\ne0$，从而式 (1) 只有零解：
$f'_{x_i}=0$（$i=1,2,\cdots,n$）.

记点 $P$ 和 $P_0$ 的坐标分别为 $(x_1,x_2,\cdots,x_n)$ 和
$(x_{01},x_{02},\cdots,x_{0n})$，根据微分中值公式，
\[
  f(P)=f(P_0)+\sum_{i=1}^n f'_{x_i}(P^*)(x_i-x_{0i})=f(P_0)
  \quad(P^*\in\mathbb R^n).
\]
此即表明
\[
  f(P)\equiv\text{常数}.
\]

\subsection*{二、梯度的计算}

\noindent\textbf{要点}\quad 梯度的计算（以 $\mathbb R^3$ 为例）主要使用如下公式：
\[
  \nabla f=\operatorname{grad}f
  =
  \left(\frac{\partial f}{\partial x},\frac{\partial f}{\partial y},\frac{\partial f}{\partial z}\right)
  =
  \frac{\partial f}{\partial x}\mathbf i
  +\frac{\partial f}{\partial y}\mathbf j
  +\frac{\partial f}{\partial z}\mathbf k,
\]
其中 $\nabla$ 为 Hamilton 算符，$\mathbf i,\mathbf j,\mathbf k$ 分别表示 $x,y,z$ 轴上的单位向量.

梯度是向量，因此关于它的运算，要遵从向量的运算法则.

\noindent\textbf{例 6.5.7}\quad 设 $u=f(x,y)$，$x=r\cos\theta$，$y=r\sin\theta$，求证：
\[
  \nabla u=\frac{\partial f}{\partial r}\mathbf r_0+\frac1r\frac{\partial f}{\partial\theta}\boldsymbol\theta_0,
\]
其中 $\mathbf r_0$ 和 $\boldsymbol\theta_0$ 分别是径向与圆周方向的单位向量（如图 6.5.1）.

\begin{center}
  \includegraphics[width=0.35\textwidth]{6-5-1.png}

  图 6.5.1
\end{center}

\noindent\textbf{证}\quad 按向量的分解原理，
\[
  \nabla u=(\nabla u\cdot\mathbf r_0)\mathbf r_0+
  (\nabla u\cdot\boldsymbol\theta_0)\boldsymbol\theta_0.
\]
因
\[
  \mathbf r_0=(\cos\theta,\sin\theta),\quad
  \boldsymbol\theta_0=\left(\cos\left(\theta+\frac\pi2\right),\cos\theta\right),
\]
$\nabla u=\left(\dfrac{\partial u}{\partial x},\dfrac{\partial u}{\partial y}\right)$，故
\[
  \nabla u\cdot\mathbf r_0
  =
  \frac{\partial u}{\partial x}\cos\theta+
  \frac{\partial u}{\partial y}\sin\theta
  =
  \frac{\partial f}{\partial r},
\]
\[
  \nabla u\cdot\boldsymbol\theta_0
  =
  \frac{\partial u}{\partial x}\cos\left(\theta+\frac\pi2\right)
  +
  \frac{\partial u}{\partial y}\cos\theta
  =
  \frac1r\left[
  \frac{\partial u}{\partial x}(-r\sin\theta)
  +
  \frac{\partial u}{\partial y}r\cos\theta
  \right]
  =
  \frac1r\frac{\partial f}{\partial\theta}.
\]
从而
\[
  \nabla u=\frac{\partial f}{\partial r}\mathbf r_0+
  \frac1r\frac{\partial f}{\partial\theta}\boldsymbol\theta_0.
\]

梯度的等式是向量的等式，请看

\noindent\textbf{例 6.5.8}\quad 在直角坐标系 $xOy$ 中引入变换
\[
  x=x(u,v),\quad y=y(u,v), \qquad\text{(1)}
\]
并将坐标系中任一点的位置向量记为 $\mathbf r=\mathbf r(u,v)$. 若变换式中函数 $x,y$ 连续可微，Jacobi 行列式
\[
  \frac{\partial(x,y)}{\partial(u,v)}\ne0,
\]
且 $\dfrac{\partial\mathbf r}{\partial u}$ 与 $\dfrac{\partial\mathbf r}{\partial v}$ 垂直，试证：对任何可微函数 $F(u,v)$，其梯度可表示为
\[
  \operatorname{grad}F
  =
  \frac1{H_u^2}\frac{\partial F}{\partial u}\frac{\partial\mathbf r}{\partial u}
  +
  \frac1{H_v^2}\frac{\partial F}{\partial v}\frac{\partial\mathbf r}{\partial v}, \qquad\text{(2)}
\]
其中
\[
  H_u=\left|\frac{\partial\mathbf r}{\partial u}\right|,\quad
  H_v=\left|\frac{\partial\mathbf r}{\partial v}\right|.
\]
（复旦大学）

\noindent\textbf{分析}\quad 这是一个兼有梯度计算与变量替换的问题. 式 (2) 为向量等式. 因
\[
  \mathbf r=x(u,v)\mathbf i+y(u,v)\mathbf j,
\]
得
\[
  \frac{\partial\mathbf r}{\partial u}=x'_u\mathbf i+y'_u\mathbf j,\quad
  \frac{\partial\mathbf r}{\partial v}=x'_v\mathbf i+y'_v\mathbf j,\quad
  H_u^2=x_u'^2+y_u'^2,\quad H_v^2=x_v'^2+y_v'^2.
\]
故要证明式 (2)，等价于要证明
\[
  \operatorname{grad}F
  =
  F'_u\frac{x'_u\mathbf i+y'_u\mathbf j}{x_u'^2+y_u'^2}
  +
  F'_v\frac{x'_v\mathbf i+y'_v\mathbf j}{x_v'^2+y_v'^2}. \qquad\text{(3)}
\]
由 $F(u,v)$ 的可微性：
\[
\begin{aligned}
  \operatorname{grad}F
  &=F'_x\mathbf i+F'_y\mathbf j
  =(F'_u u'_x+F'_v v'_x)\mathbf i+(F'_u u'_y+F'_v v'_y)\mathbf j\\
  &=F'_u(u'_x\mathbf i+u'_y\mathbf j)+F'_v(v'_x\mathbf i+v'_y\mathbf j). \qquad\text{(4)}
\end{aligned}
\]
比较 (3)，(4) 可知，要从 (4) 推出 (3)，只要证明：
\[
  u'_x=\frac{x'_u}{x_u'^2+y_u'^2},\quad
  v'_x=\frac{x'_v}{x_v'^2+y_v'^2},\quad
  u'_y=\frac{y'_u}{x_u'^2+y_u'^2},\quad
  v'_y=\frac{y'_v}{x_v'^2+y_v'^2}. \qquad\text{(5)}
\]
事实上，
\[
  \frac{\partial(x,y)}{\partial(u,v)}\ne0,
\]
按隐函数存在定理，变换 (1)：$x=x(u,v),y=y(u,v)$ 确定了 $u,v$ 作为 $x,y$ 的函数，式 (1) 两边同时对 $x$ 求导，得
\[
  1=x'_u u'_x+x'_v v'_x,\qquad
  0=y'_u u'_x+y'_v v'_x.
\]
由此得
\[
  u'_x=\frac{y'_v}{x'_uy'_v-x'_vy'_u},\qquad
  v'_x=\frac{-y'_u}{x'_uy'_v-x'_vy'_u}. \qquad\text{(6)}
\]
但由已知条件 $\dfrac{\partial\mathbf r}{\partial u}$ 与
$\dfrac{\partial\mathbf r}{\partial v}$ 垂直，有
\[
  x'_ux'_v+y'_uy'_v=0.
\]
如此可得 (5) 中前两式. 同理可证后两式. 证毕.

\noindent\textbf{例 6.5.9}\quad 设有方程
\[
  \frac{x^2}{a^2+u}+\frac{y^2}{b^2+u}+\frac{z^2}{c^2+u}=1, \qquad\text{(1)}
\]
证明
\[
  (\operatorname{grad}u)^2=2A\cdot\operatorname{grad}u, \qquad\text{(2)}
\]
其中 $A=(x,y,z)$.（中国科学技术大学）

\noindent\textbf{分析}\quad 这是一个兼有梯度计算与隐函数求导的问题. 根据向量的数量积公式，(2) 等价于
\[
  u_x'^2+u_y'^2+u_z'^2=2(xu'_x+yu'_y+zu'_z). \qquad\text{(3)}
\]
可见我们的任务在于：由方程 (1) 证明式 (3).

\noindent\textbf{证}\quad 不难验证式 (1) 满足隐函数存在定理的条件，因此由式 (1) 将 $u$ 定义为 $x,y,z$ 的函数，将式 (1) 对 $x$ 求导，得
\[
  \frac{(a^2+u)2x-u'_x x^2}{(a^2+u)^2}
  -
  \frac{y^2u'_x}{(b^2+u)^2}
  -
  \frac{z^2u'_x}{(c^2+u)^2}=0,
\]
即
\[
  \frac{2x}{a^2+u}
  =
  \left[
  \frac{x^2}{(a^2+u)^2}
  +
  \frac{y^2}{(b^2+u)^2}
  +
  \frac{z^2}{(c^2+u)^2}
  \right]u'_x. \qquad\text{(4)}
\]
根据轮换对称性，有
\[
  \frac{2y}{b^2+u}
  =
  \left[
  \frac{x^2}{(a^2+u)^2}
  +
  \frac{y^2}{(b^2+u)^2}
  +
  \frac{z^2}{(c^2+u)^2}
  \right]u'_y, \qquad\text{(5)}
\]
\[
  \frac{2z}{c^2+u}
  =
  \left[
  \frac{x^2}{(a^2+u)^2}
  +
  \frac{y^2}{(b^2+u)^2}
  +
  \frac{z^2}{(c^2+u)^2}
  \right]u'_z. \qquad\text{(6)}
\]
式 (4)，(5)，(6) 平方后相加，在等式两端约去公因子，得
\[
  4=
  \left[
  \frac{x^2}{(a^2+u)^2}
  +
  \frac{y^2}{(b^2+u)^2}
  +
  \frac{z^2}{(c^2+u)^2}
  \right](u_x'^2+u_y'^2+u_z'^2). \qquad\text{(7)}
\]
式 (4)，(5)，(6) 分别乘 $x,y,z$ 后相加，注意式 (1)，有
\[
  2=
  \left[
  \frac{x^2}{(a^2+u)^2}
  +
  \frac{y^2}{(b^2+u)^2}
  +
  \frac{z^2}{(c^2+u)^2}
  \right](xu'_x+yu'_y+zu'_z). \qquad\text{(8)}
\]
将式 (7)，(8) 联立，即得 (3)，从而式 (2) 求证.

\noindent\textbf{例 6.5.10}\quad 假设函数 $f(x,y)$ 在原点 $(0,0)$ 的某邻域 $U$ 内有定义，并满足如下的条件：

1) 在原点，它沿任意方向
$\mathbf l=(\cos\alpha,\sin\alpha)$（$0\le\alpha<2\pi$）的方向导数存在，且
\[
  \frac{\partial f}{\partial l}=A\cos\alpha+B\sin\alpha,
\]
其中 $A,B$ 是两个常数；

2) 存在常数 $M>0$，使得对 $U$ 中任何两点 $(x_1,y_1)$ 与 $(x_2,y_2)$ 成立不等式
\[
  |f(x_1,y_1)-f(x_2,y_2)|\le
  M(|x_1-x_2|+|y_1-y_2|).
\]
试证：$f(x,y)$ 在原点可微，且
\[
  df(0,0)=Adx+Bdy.
\]
（武汉大学）

\noindent\textbf{分析}\quad 当 $\alpha=0$ 时，向量
$\mathbf l=(\cos0,\sin0)=(1,0)$ 代表 $x$ 轴正向单位向量. 因此
\[
  \frac{\partial f(0,0)}{\partial x}
  =
  \left.\frac{\partial f}{\partial l}\right|_{\alpha=0}=A.
\]
同理有 $\dfrac{\partial f(0,0)}{\partial y}=B$. 由此可见，只要证明了 $f$ 在原点可微，则
$df(0,0)=Adx+Bdy$ 明显成立.

为了证明 $f$ 在原点可微，按定义，即要证明
\[
  F(x,y)\equiv
  \frac{f(x,y)-f(0,0)-Ax-By}{\sqrt{x^2+y^2}}
  \to0\quad(\text{当 }x\to0,y\to0\text{ 时}).
\]
令 $x=t\cos\alpha,y=t\sin\alpha$，作变换，记
\[
  \varphi(t,\alpha)=F(t\cos\alpha,t\sin\alpha),
\]
于是问题等价于要证明
\[
  \lim_{t\to0}\varphi(t,\alpha)=0
  \quad(\text{关于 }\alpha\in[0,2\pi]\text{ 一致}).
\]
为此，我们只需证明：

$1^\circ$ $\forall\alpha\in[0,2\pi]$，
$\lim_{t\to0}\varphi(t,\alpha)=0$；

$2^\circ$ $\exists L>0$，使得 $\forall\alpha,\alpha_0\in[0,2\pi]$，
$|\varphi(t,\alpha)-\varphi(t,\alpha_0)|\le L|\alpha-\alpha_0|$.

\noindent\textbf{证}\quad $1^\circ$ 由已知条件，对于
$\mathbf l=(\cos\alpha,\sin\alpha)$（$0\le\alpha\le2\pi$），
\[
  \frac{\partial f}{\partial l}=A\cos\alpha+B\sin\alpha,
\]
因此 $\forall\alpha\in[0,2\pi]$，
\[
\begin{aligned}
  \lim_{t\to0}\varphi(t,\alpha)
  &=
  \lim_{t\to0}\frac1t
  [f(t\cos\alpha,t\sin\alpha)-f(0,0)-At\cos\alpha-Bt\sin\alpha]\\
  &=
  \lim_{t\to0}\left[
  \frac{f(t\cos\alpha,t\sin\alpha)-f(0,0)}{t}
  -(A\cos\alpha+B\sin\alpha)
  \right]=0.
\end{aligned}
\]

$2^\circ$ 根据 $\varphi(t,\alpha)$ 的定义与条件 2)：
\[
\begin{aligned}
|\varphi(t,\alpha)-\varphi(t,\alpha_0)|
&=
\left|
\frac{f(t\cos\alpha,t\sin\alpha)-f(0,0)-At\cos\alpha-Bt\sin\alpha}{t}
\right.\\
&\quad\left.
-
\frac{f(t\cos\alpha_0,t\sin\alpha_0)-f(0,0)-At\cos\alpha_0-Bt\sin\alpha_0}{t}
\right|\\
&\le
\frac1{|t|}
|f(t\cos\alpha,t\sin\alpha)-f(t\cos\alpha_0,t\sin\alpha_0)|\\
&\quad
 +|A||\cos\alpha-\cos\alpha_0|+|B||\sin\alpha-\sin\alpha_0|\\
&\le
M(|\cos\alpha-\cos\alpha_0|+|\sin\alpha-\sin\alpha_0|)
 +|A||\cos\alpha-\cos\alpha_0|+|B||\sin\alpha-\sin\alpha_0|\\
&\le (2M+|A|+|B|)|\alpha-\alpha_0|
=L|\alpha-\alpha_0|,
\end{aligned}
\]
其中 $L=2M+|A|+|B|$ 为常数. $\forall\varepsilon>0$，取
$n>\dfrac{4\pi L}{\varepsilon}$（这时 $\dfrac{2\pi}{n}<\dfrac{\varepsilon}{2L}$），令
$\alpha_k=k\dfrac{2\pi}{n}$（$k=0,1,2,\cdots,n-1$），由 $1^\circ$ 知，对每个 $\alpha_k$，
$\exists\delta_k>0$，使得 $|t|<\delta_k$ 时有
$|\varphi(t,\alpha_k)|<\dfrac\varepsilon2$. 令
$\delta=\min\{\delta_1,\cdots,\delta_n\}$，则当 $|t|<\delta$ 时，
$\forall\alpha\in[0,2\pi]$，必 $\exists k\in\{0,1,\cdots,n-1\}$ 使得
$k\dfrac{2\pi}{n}\le\alpha<(k+1)\dfrac{2\pi}{n}$，从而
$|\alpha-\alpha_k|<\dfrac{2\pi}{n}<\dfrac{\varepsilon}{2L}$. 因此
\[
  |\varphi(t,\alpha)|
  \le |\varphi(t,\alpha)-\varphi(t,\alpha_k)|+|\varphi(t,\alpha_k)|
  \le L|\alpha-\alpha_k|+|\varphi(t,\alpha_k)|
  <L\frac{2\pi}{n}+\frac\varepsilon2<\varepsilon.
\]
这就证明了 $\lim_{t\to0}\varphi(t,\alpha)=0$ 关于
$\alpha\in[0,2\pi]$ 一致. 因而
\[
  \lim_{(x,y)\to(0,0)}F(x,y)=0,
\]
$f$ 在 $(0,0)$ 处可微. 证毕.

\noindent\textbf{\ensuremath{\star}例 6.5.11}\quad 设函数 $f(x,y)$ 在定义域的某点 $(x_0,y_0)$ 上存在非零方向导数，且沿三个不同方向的方向导数相等. 证明：$f(x,y)$ 在 $(x_0,y_0)$ 处不可微.（北京大学）

\noindent\textbf{提示}\quad （反证法）如果 $f(x,y)$ 在该点 $(x_0,y_0)$ 可微，则存在偏导数
\[
  f'_x(x_0,y_0)\overset{\text{记}}{=}A,\quad
  f'_y(x_0,y_0)\overset{\text{记}}{=}B,
\]
且
\[
  df(x_0,y_0)=f'_x(x_0,y_0)dx+f'_y(x_0,y_0)dy=Adx+Bdy.
\]

\noindent\textbf{证 I}\quad 按题设，存在三个不同方向：
$\mathbf e_i=(\cos\theta_i,\sin\theta_i)$（$i=1,2,3$），且
\[
  \frac{\partial f(x_0,y_0)}{\partial e_i}
  =
  A\cos\theta_i+B\sin\theta_i=C\ne0. \qquad\text{(1)}
\]
（$A$ 和 $B$ 不可能同时为 $0$，否则 $f(x,y)$ 在 $(x_0,y_0)$ 处方向导数全为 $0$，与题设矛盾.）故 $A^2+B^2\ne0$，
\[
  A\cos\theta_i+B\sin\theta_i
  =
  \sqrt{A^2+B^2}\left(
  \frac A{\sqrt{A^2+B^2}}\cos\theta_i+
  \frac B{\sqrt{A^2+B^2}}\sin\theta_i
  \right)
  =C\ne0.
\]
记 $\dfrac A{\sqrt{A^2+B^2}}=\cos\theta$，则
$\dfrac B{\sqrt{A^2+B^2}}=\sin\theta$. 代入上式得
\[
  \cos(\theta-\theta_i)=\frac C{\sqrt{A^2+B^2}}\quad(i=1,2,3).
\]
此式表明：应存在向量 $\mathbf e=(\cos\theta,\sin\theta)$，它能充当三个不同向量
$\{\mathbf e_i\}_{i=1,2,3}$ 中每两个 $\mathbf e_i$ 的分角线，但在 $\mathbb R^2$ 里这是不可能的！故上述 $A,B$ 不存在，$f(x,y)$ 在点 $(x_0,y_0)$ 处不可微.

\noindent\textbf{证 II}\quad 三个方向的方向导数相等，即
\[
\begin{cases}
  A\cos\theta_1+B\sin\theta_1-C=0, &(a)\\
  A\cos\theta_2+B\sin\theta_2-C=0, &(b)\\
  A\cos\theta_3+B\sin\theta_3-C=0, &(c)
\end{cases}\qquad\text{(2)}
\]
此式看作 $(A,B,C)$ 的三元一次线性齐次方程组，它只有零解的充要条件是
\[
  \begin{vmatrix}
    \cos\theta_1&\sin\theta_1&1\\
    \cos\theta_2&\sin\theta_2&1\\
    \cos\theta_3&\sin\theta_3&1
  \end{vmatrix}\ne0. \qquad\text{(3)}
\]
因为平面上三个不同向量 $\{\mathbf e_i\}_{i=1,2,3}$ 中一个向量可写成另外两个向量的线性组合，如 $\mathbf e_1$ 可写成 $\mathbf e_2,\mathbf e_3$ 的线性组合，即 $\exists\alpha,\beta\ne0$，使得
\[
  \cos\theta_1=\alpha\cos\theta_2+\beta\cos\theta_3,\quad
  \sin\theta_1=\alpha\sin\theta_2+\beta\sin\theta_3. \qquad\text{(4)}
\]
两式先平方，再相加得
\[
  1=\alpha^2+\beta^2+2\alpha\beta(\cos\theta_2\cos\theta_3+\sin\theta_2\sin\theta_3)
  =
  \alpha^2+\beta^2+2\alpha\beta\cos(\theta_2-\theta_3)\ne(\alpha+\beta)^2
\]
（因为 $\theta_2\ne\theta_3$）. 因此，$\alpha+\beta\ne1$. 式 (3) 的行列式
\[
\begin{vmatrix}
  \cos\theta_1&\sin\theta_1&1\\
  \cos\theta_2&\sin\theta_2&1\\
  \cos\theta_3&\sin\theta_3&1
\end{vmatrix}
\overset{\text{式(4)}}{=}
\begin{vmatrix}
  \alpha\cos\theta_2+\beta\cos\theta_3&
  \alpha\sin\theta_2+\beta\sin\theta_3&1\\
  \cos\theta_2&\sin\theta_2&1\\
  \cos\theta_3&\sin\theta_3&1
\end{vmatrix}
\]
\[
  =
\begin{vmatrix}
  0&0&1-\alpha-\beta\\
  \cos\theta_2&\sin\theta_2&1\\
  \cos\theta_3&\sin\theta_3&1
\end{vmatrix}
  =(1-\alpha-\beta)\sin(\theta_3-\theta_2)\ne0.
\]
此式说明：要式 (3) 成立，则式 (2) 只有零解：
$(A,B,C)=(0,0,0)$，故 $f(x,y)$ 在点 $(x_0,y_0)$ 处不可微.

\noindent\textbf{\ensuremath{\star}例 6.5.12}\quad 设 $f(x,y)$ 在 $P_0\in\mathbb R^2$ 的邻域 $U(P_0)$ 内存在三阶偏导数，并且所有三阶偏导数的绝对值不超过常数 $M>0$. 另设 $P_1,P_2$ 是 $U(P_0)$ 内的两点，且关于 $P_0$ 对称，$P_1$ 与 $P_0$ 的距离为 $l>0$，记 $\mathbf l=\overrightarrow{P_0P_1}$，证明：
\[
  \left|
  \frac{f(P_1)-f(P_2)}{2l}
  -
  \frac{\partial f(P_0)}{\partial l}
  \right|
  \le
  \frac{\sqrt2}{3}Ml^2.
\]
（南开大学）

\noindent\textbf{提示}\quad 取 $P_0$ 作原点，以 $\mathbf l$ 的方向作为数轴的正向，$l$ 作为长度单位，坐标记作 $t$，则 $P_0,P_1,P_2$ 的坐标分别是
$t=0,t=1,t=-1$. 当 $P_1,P_2$ 被选定之后，函数 $f$ 的值只与 $t$ 有关，
$f(x(t),y(t))$ 可看成单变量 $t$ 的函数，记作
\[
  F(t)=f(x(t),y(t)).
\]
记 $P_i$ 的坐标为 $(x_i,y_i)$（$i=0,1,2$），$P_0,P_1$ 连线上点 $P$ 满足：
\[
  P(t)=(x(t),y(t))=(x_0+t(x_1-x_0),y_0+t(y_1-y_0)),
\]
\[
  F(t)=f(x(t),y(t))
  =
  f(x_0+t(x_1-x_0),y_0+t(y_1-y_0)),
\]
\[
  F(0)=f(x_0,y_0),\quad F(1)=f(x_1,y_1),\quad F(-1)=f(x_2,y_2).
\]
原不等式可写为
\[
  \left|\frac{F(1)-F(-1)}2-F'(0)\right|
  \le\frac{\sqrt2}{3}M. \qquad\text{(1)}
\]
利用 Taylor 公式，
\[
  F(1)=F(0)+F'(0)+\frac12F''(0)+\frac16F'''(\xi)\quad(0<\xi<1),
\]
\[
  F(-1)=F(0)-F'(0)+\frac12F''(0)-\frac16F'''(\eta)\quad(-1<\eta<0).
\]
因三阶偏导数的绝对值都不超过常数 $M>0$，上两式相减可得
\[
\left|
\frac{F(1)-F(-1)}2-F'(0)
\right|
=
\left|
\frac12\left(\frac16F'''(\xi)+\frac16F'''(\eta)\right)
\right|
\le\frac{\sqrt2}{3}M. \qquad\text{(2)}
\]
补证式 (2) 里的不等式：$F$ 对 $t$ 的导数实际上都是 $f(x,y)$ 沿 $\mathbf l$ 方向的方向导数，方向导数的界尚不得而知. 为此应用习题 6.5.3 题的方法，先求（三阶）方向导数的表达式. 轴 $\mathbf l$ 上的单位向量记作
$\mathbf e=(\cos\alpha,\cos\beta)$，则
\[
\frac{\partial f(x,y)}{\partial l}
=f'_x\cos\alpha+f'_y\cos\beta
=\left(\frac{\partial}{\partial x}\cos\alpha+\frac{\partial}{\partial y}\cos\beta\right)f.
\]
\[
\frac{\partial^2f(x,y)}{\partial l^2}
=
\left(\frac{\partial}{\partial x}\cos\alpha+\frac{\partial}{\partial y}\cos\beta\right)^2f(x,y)
=f''_{xx}\cos^2\alpha+2f''_{xy}\cos\alpha\cos\beta+f''_{yy}\cos^2\beta.
\]
递推，得
\[
  \frac{\partial^3f(x,y)}{\partial l^3}
  =
  \left(\frac{\partial}{\partial x}\cos\alpha+\frac{\partial}{\partial y}\cos\beta\right)^3f(x,y).
\]
因为所有三阶偏导数的绝对值不超过常数 $M>0$，所以
\[
\begin{aligned}
\left|\frac{\partial^3f(x,y)}{\partial l^3}\right|
&=
\left|
\left(\frac{\partial}{\partial x}\cos\alpha+\frac{\partial}{\partial y}\cos\beta\right)^3f(x,y)
\right|\\
&\le M(\cos^3\alpha+3\cos^2\alpha\cos\beta
  +3\cos\alpha\cos^2\beta+\cos^3\beta)\\
&=M(\cos\alpha+\cos\beta)^3
=M[\sqrt2(\sin45^\circ\cos\alpha+\cos45^\circ\sin\alpha)]^3\\
&=2\sqrt2M[\sin(\alpha+45^\circ)]^3\le2\sqrt2M.
\end{aligned}
\]
因此式 (2) 里：$|F'''(\xi)|\le2\sqrt2M$，$|F'''(\eta)|\le2\sqrt2M$. 故式 (2) 最后的不等式成立：
\[
  \left|\frac12\left(\frac16F'''(\xi)+\frac16F'''(\eta)\right)\right|
  \le\frac{\sqrt2}{3}M.
\]

\subsection*{单元练习 6.5}

\noindent\textbf{6.5.1}\quad 计算函数 $z=\ln(x^2+y^2)$ 在点
$P_0(x_0,y_0)$ 处沿过此点的等位线垂直的方向上的方向导数.
\[
  \left(\pm\frac2{\sqrt{x_0^2+y_0^2}}\right)
\]

\noindent\textbf{提示}\quad （直接用公式.）等位线是原点为中心的同心圆，法线沿径向方向：
\[
  \mathbf n=(\cos\alpha,\cos\beta)
  =
  \pm\left(
  \frac{x}{\sqrt{x^2+y^2}},
  \frac{y}{\sqrt{x^2+y^2}}
  \right).
\]
（外法线方向取“$+$”，内法线方向取“$-$”.）

\noindent\textbf{6.5.2}\quad 计算函数
\[
  z=1-\left(\frac{x^2}{a^2}+\frac{y^2}{b^2}\right)\quad(a,b>0)
\]
在点 $P\left(\dfrac a{\sqrt2},\dfrac b{\sqrt2}\right)$ 处沿曲线
\[
  \frac{x^2}{a^2}+\frac{y^2}{b^2}=1
\]
在此点内法线方向上的导数.

\noindent\textbf{提示}\quad
\[
  z\overset{\text{记}}{=}F(x,y)=1-\left(\frac{x^2}{a^2}+\frac{y^2}{b^2}\right)\quad(a,b>0).
\]
曲线 $F(x,y)=0$ 在点
$P\left(\dfrac a{\sqrt2},\dfrac b{\sqrt2}\right)$ 处内法向量为
\[
  (F'_x,F'_y)\big|_P
  =
  -\left(\frac2{a^2}\cdot\frac a{\sqrt2},
  \frac2{b^2}\cdot\frac b{\sqrt2}\right)
  =
  -\sqrt2\left(\frac1a,\frac1b\right).
\]
$P\left(\dfrac a{\sqrt2},\dfrac b{\sqrt2}\right)$ 点的（内）法向单位向量
\[
  \mathbf n=(\cos\alpha,\cos\beta)
  =
  -\frac1{\sqrt{a^2+b^2}}(b,a),
\]
在点 $P\left(\dfrac a{\sqrt2},\dfrac b{\sqrt2}\right)$ 处：
\[
  (z'_x,z'_y)=\sqrt2\left(\frac1a,\frac1b\right).
\]
于是
\[
  \frac{\partial z}{\partial n}
  =
  z'_x\cos\alpha+z'_y\cos\beta
  =
  \frac{-\sqrt2}{\sqrt{a^2+b^2}}\left(\frac ba+\frac ab\right)
  =
  -\frac{\sqrt2\sqrt{a^2+b^2}}{ab}.
\]

\noindent\textbf{6.5.3}\quad 设 $u=f(x,y,z)$ 为二次可微函数. 若
$\cos\alpha,\cos\beta,\cos\gamma$ 为方向 $l$ 的方向余弦，求
\[
  \frac{\partial^2u}{\partial l^2}
  =
  \frac{\partial}{\partial l}\left(\frac{\partial u}{\partial l}\right).
\]

\noindent\textbf{解}\quad
\[
  \frac{\partial u}{\partial l}
  =
  f'_x\cos\alpha+f'_y\cos\beta+f'_z\cos\gamma
  =
  \left(
  \cos\alpha\frac{\partial}{\partial x}
  +\cos\beta\frac{\partial}{\partial y}
  +\cos\gamma\frac{\partial}{\partial z}
  \right)f,
\]
\[
\begin{aligned}
  \frac{\partial^2u}{\partial l^2}
  &=
  \frac{\partial}{\partial l}\left(\frac{\partial u}{\partial l}\right)
  =
  \left(
  \cos\alpha\frac{\partial}{\partial x}
  +\cos\beta\frac{\partial}{\partial y}
  +\cos\gamma\frac{\partial}{\partial z}
  \right)^2f\\
  &=
  f''_{xx}\cos^2\alpha+f''_{yy}\cos^2\beta+f''_{zz}\cos^2\gamma\\
  &\quad
  +2\left(
  f''_{xy}\cos\alpha\cos\beta+
  f''_{yz}\cos\beta\cos\gamma+
  f''_{zx}\cos\gamma\cos\alpha
  \right).
\end{aligned}
\]

\noindent\textbf{6.5.4}\quad 设 $u=f(x,y,z)$ 为二次可微函数，
$l_1,l_2,l_3$ 为三个互相垂直的方向，证明：

1)
\[
  \left(\frac{\partial u}{\partial l_1}\right)^2+
  \left(\frac{\partial u}{\partial l_2}\right)^2+
  \left(\frac{\partial u}{\partial l_3}\right)^2
  =
  \left(\frac{\partial u}{\partial x}\right)^2+
  \left(\frac{\partial u}{\partial y}\right)^2+
  \left(\frac{\partial u}{\partial z}\right)^2;
\]

2)
\[
  \frac{\partial^2u}{\partial l_1^2}+
  \frac{\partial^2u}{\partial l_2^2}+
  \frac{\partial^2u}{\partial l_3^2}
  =
  \frac{\partial^2u}{\partial x^2}+
  \frac{\partial^2u}{\partial y^2}+
  \frac{\partial^2u}{\partial z^2}.
\]

\noindent\textbf{提示}\quad 现将 $\{l_i\}_{i=1,2,3}$ 看作新直角坐标系，并用
$\mathbf i,\mathbf j,\mathbf k$ 分别表示新坐标轴上的单位向量. 设
\[
  \mathbf e_i=(\cos\alpha_i,\cos\beta_i,\cos\gamma_i)\quad(i=1,2,3),
\]
则
\[
  \frac{\partial u}{\partial l_i}
  =
  f'_x\cos\alpha_i+f'_y\cos\beta_i+f'_z\cos\gamma_i. \qquad\text{(1)}
\]
在原来的坐标系 $Oxyz$ 中，用 $\mathbf i,\mathbf j,\mathbf k$ 表示 $x,y,z$ 三坐标轴上的单位向量，它们在新坐标系里的方向余弦是
\[
  \mathbf i=(\cos\alpha_1,\cos\alpha_2,\cos\alpha_3),\quad
  \mathbf j=(\cos\beta_1,\cos\beta_2,\cos\beta_3),\quad
  \mathbf k=(\cos\gamma_1,\cos\gamma_2,\cos\gamma_3),
\]
作为单位向量：
\[
  \sum_{i=1}^3\cos^2\alpha_i=1,\quad
  \sum_{i=1}^3\cos^2\beta_i=1,\quad
  \sum_{i=1}^3\cos^2\gamma_i=1. \qquad\text{(2)}
\]
因互相垂直，故
\[
  \sum_{i=1}^3\cos\alpha_i\cos\beta_i=0,\quad
  \sum_{i=1}^3\cos\beta_i\cos\gamma_i=0,\quad
  \sum_{i=1}^3\cos\gamma_i\cos\alpha_i=0. \qquad\text{(3)}
\]

\noindent\textbf{再提示}\quad 1) 用式 (1)：
\[
\begin{aligned}
  \sum_{i=1}^3\left(\frac{\partial u}{\partial l_i}\right)^2
  &=
  \sum_{i=1}^3(f'_x\cos\alpha_i+f'_y\cos\beta_i+f'_z\cos\gamma_i)^2\\
  &=
  f_x'^2\sum_{i=1}^3\cos^2\alpha_i+
  f_y'^2\sum_{i=1}^3\cos^2\beta_i+
  f_z'^2\sum_{i=1}^3\cos^2\gamma_i\\
  &\quad+
  2f'_xf'_y\sum_{i=1}^3\cos\alpha_i\cos\beta_i+
  2f'_yf'_z\sum_{i=1}^3\cos\beta_i\cos\gamma_i+
  2f'_zf'_x\sum_{i=1}^3\cos\gamma_i\cos\alpha_i\\
  &\overset{\text{式(2),(3)}}{=}
  f_x'^2+f_y'^2+f_z'^2.
\end{aligned}
\]

2) 利用上题（题 6.5.3）所得的公式，
\[
\begin{aligned}
  \sum_{i=1}^3\left(\frac{\partial^2u}{\partial l_i^2}\right)
  &=
  \sum_{i=1}^3
  \left[
  f''_{xx}\cos^2\alpha_i+
  f''_{yy}\cos^2\beta_i+
  f''_{zz}\cos^2\gamma_i\right.\\
  &\quad\left.
  +2\left(
  f''_{xy}\cos\alpha_i\cos\beta_i+
  f''_{yz}\cos\beta_i\cos\gamma_i+
  f''_{zx}\cos\gamma_i\cos\alpha_i
  \right)\right]\\
  &\quad(\text{利用式 (2)，(3)})\\
  &=f''_{xx}+f''_{yy}+f''_{zz}.
\end{aligned}
\]

\noindent\textbf{6.5.5}\quad 求函数 $u=x+y+z$ 在沿球面
$x^2+y^2+z^2=1$ 上点 $P_0(x_0,y_0,z_0)$ 的外法线方向的方向导数，并问在球面上什么点处此导数
1) 取最大值；2) 取最小值；3) 等于零.

\noindent\textbf{提示}\quad $(u'_x,u'_y,u'_z)=(1,1,1)$. 记
$F=x^2+y^2+z^2-1$，则
\[
  (F'_x,F'_y,F'_z)=(2x,2y,2z),
\]
单位外法向量 $\mathbf n=(x,y,z)$. 所求的方向导数：
\[
  \frac{\partial u}{\partial n}=(1,1,1)\cdot(x,y,z)=x+y+z.
\]
在球面上点 $(x_0,y_0,z_0)$ 处，外法线方向的方向导数：
\[
  \left.\frac{\partial u}{\partial n}\right|_{(x_0,y_0,z_0)}
  =
  x_0+y_0+z_0.
\]

\noindent\textbf{再提示}\quad 设
\[
  L=(x+y+z)-\lambda(x^2+y^2+z^2-1).
\]
令 $L'_x=L'_y=L'_z=0$，解方程得极值可疑点
\[
  (x,y,z)=\left(\pm\frac{\sqrt3}{3},\pm\frac{\sqrt3}{3},\pm\frac{\sqrt3}{3}\right).
\]

\end{document}
